% ============================================================
%  Certifikatshandboken – 1:a upplagan
%  Radioskola.se · Dennis Bengtsson med hjälp av Claude AI
%  Kompilera med: xelatex certifikatshandboken.tex
% ============================================================

\documentclass[a4paper,11pt,twoside]{book}

% ── Teckenkodning & språk ──────────────────────────────────
\usepackage{fontspec}
\usepackage{polyglossia}
\setmainlanguage{swedish}

% ── Typsnitt ──────────────────────────────────────────────
\setmainfont{Source Serif 4}[
  BoldFont = Source Serif 4 SemiBold,
  ItalicFont = Source Serif 4 Italic
]
\newfontfamily\headingfont{Oswald}[
  BoldFont = Oswald Bold,
  UprightFont = Oswald Regular
]
\newfontfamily\monofont{IBM Plex Mono}
\newfontfamily\emojifont{Segoe UI Emoji}
\newcommand{\emoji}[1]{{\emojifont #1}}
\newcommand{\etitle}[2]{{\emojifont\mdseries #1}\space #2}

\usepackage{siunitx}
\sisetup{
  locale = DE,
  output-decimal-marker = {,}
}


% Fallback om Google-typsnitten saknas:
% Byt ut mot: \setmainfont{Georgia} och \newfontfamily\headingfont{Arial Narrow}

% Enhetshjälp (robust i både text och matte)
\newcommand{\ohm}{\ensuremath{\,\Omega}}
\newcommand{\volt}{\ensuremath{\,\mathrm{V}}}
\newcommand{\ampere}{\ensuremath{\,\mathrm{A}}}
\newcommand{\watt}{\ensuremath{\,\mathrm{W}}}
\newcommand{\microhenry}{\ensuremath{\,\mu\mathrm{H}}}
\newcommand{\microfarad}{\ensuremath{\,\mu\mathrm{F}}}
\newcommand{\milliampere}{\ensuremath{\,\mathrm{mA}}}


% ── Sidlayout ─────────────────────────────────────────────
\usepackage[
  a4paper,
  top=22mm, bottom=20mm,
  left=20mm, right=20mm,
  headheight=14pt
]{geometry}

% ── Färger ────────────────────────────────────────────────
\usepackage{xcolor}

% ── Gemensamma färger för kapitel 1.8 ────────────────────
\definecolor{cBlueTitle}{HTML}{1A6EA8}
\definecolor{cBlueSub}{HTML}{4A7FA0}
\definecolor{cPanel}{HTML}{F7FBFF}
\definecolor{cPanelStroke}{HTML}{B0CCE0}
\definecolor{cAxis}{HTML}{888888}
\definecolor{cGridBlue}{HTML}{1A6EA8}
\definecolor{cRed}{HTML}{C0392B}
\definecolor{cBlue}{HTML}{1A6EA8}
\definecolor{cBottom}{HTML}{F0F5FA}
\definecolor{cBottomStroke}{HTML}{C0CCD8}

% Systembokens grundfärger (rör ej)
\definecolor{orange}{HTML}{E67E22}
\definecolor{orangedark}{HTML}{D35400}
\definecolor{orangepale}{HTML}{FEF6EE}
\definecolor{coverbg}{HTML}{07070F}
\definecolor{greydark}{HTML}{1A1A2A}
\definecolor{greybody}{HTML}{2C3440}
\definecolor{greymed}{HTML}{7F8C8D}
\definecolor{greylight}{HTML}{F4F5F6}
\definecolor{greyborder}{HTML}{E0E4E8}
\definecolor{bluedark}{HTML}{1A5276}
\definecolor{bluepale}{HTML}{EAF2F8}
\definecolor{greendark}{HTML}{1E8449}
\definecolor{greenpale}{HTML}{EAFAF1}
\definecolor{reddark}{HTML}{922B21}
\definecolor{redpale}{HTML}{FDEDEC}

% Röda och rosa toner
\definecolor{korall}{RGB}{255, 99, 71}
\definecolor{rosenrod}{RGB}{220, 80, 100}
\definecolor{ljusrosa}{RGB}{255, 182, 193}
\definecolor{cerise}{RGB}{190, 30, 90}
\definecolor{burgundy}{RGB}{128, 0, 32}

% Orange och gula toner
\definecolor{ljusbrun}{RGB}{210, 170, 120}
\definecolor{sandfarg}{RGB}{240, 210, 160}
\definecolor{honung}{RGB}{235, 175, 50}
\definecolor{solgul}{RGB}{255, 215, 0}
\definecolor{ljusgul}{RGB}{255, 245, 150}
\definecolor{ockra}{RGB}{200, 155, 50}
\definecolor{brons}{RGB}{180, 120, 50}

% Gröna toner
\definecolor{ljusgron}{RGB}{180, 220, 150}
\definecolor{mintgron}{RGB}{152, 220, 180}
\definecolor{skogsgron}{RGB}{40, 100, 60}
\definecolor{olivgron}{RGB}{107, 130, 50}
\definecolor{limegron}{RGB}{180, 210, 50}
\definecolor{jadegron}{RGB}{0, 168, 120}
\definecolor{salviegron}{RGB}{140, 170, 130}

% Blå toner
\definecolor{himmelsblaa}{RGB}{135, 195, 235}
\definecolor{ljusblaa}{RGB}{173, 216, 230}
\definecolor{koboltblaa}{RGB}{0, 70, 170}
\definecolor{marineblaa}{RGB}{15, 40, 90}
\definecolor{staalblaa}{RGB}{70, 120, 160}
\definecolor{isblaa}{RGB}{210, 235, 245}
\definecolor{duvblaa}{RGB}{120, 150, 180}
\definecolor{petroleumblaa}{RGB}{0, 100, 120}

% Lila och violetta toner
\definecolor{lavendel}{RGB}{200, 175, 230}
\definecolor{lila}{RGB}{140, 80, 180}
\definecolor{morkviolett}{RGB}{80, 30, 120}
\definecolor{malva}{RGB}{180, 120, 160}
\definecolor{aubergine}{RGB}{90, 30, 70}
\definecolor{ljuslila}{RGB}{220, 200, 240}

% Bruna och beige toner
\definecolor{kakao}{RGB}{110, 65, 40}
\definecolor{cappuccino}{RGB}{175, 130, 95}
\definecolor{beige}{RGB}{235, 220, 195}
\definecolor{kanel}{RGB}{175, 95, 50}
\definecolor{mocka}{RGB}{145, 105, 75}
\definecolor{lerbrun}{RGB}{160, 120, 85}

% Grå och neutrala toner
\definecolor{ljusgra}{RGB}{220, 220, 220}
\definecolor{mellangra}{RGB}{160, 160, 160}
\definecolor{kolgra}{RGB}{80, 80, 80}
\definecolor{perlgra}{RGB}{200, 200, 205}
\definecolor{skiffergra}{RGB}{100, 115, 125}
\definecolor{varmdark}{RGB}{60, 50, 45}

% Metalliska toner
\definecolor{guld}{RGB}{210, 175, 55}
\definecolor{silvergraa}{RGB}{192, 192, 192}
\definecolor{koppar}{RGB}{185, 110, 65}
\definecolor{antracit}{RGB}{55, 60, 65}


% ── Grafik & tabeller ─────────────────────────────────────
\usepackage{graphicx}
\usepackage{booktabs}
\usepackage{tabularx}
\usepackage{array}
\usepackage{longtable}
\usepackage{multirow}
\usepackage{colortbl}
\usepackage{csquotes}

% ── Infoboxar (tcolorbox) ─────────────────────────────────
\usepackage{tcolorbox}
\tcbuselibrary{skins, breakable}

% ── Räkneexempel-box ──────────────────────────────────────
\newtcolorbox{worked}[1][]{
  colback=green!6,
  colframe=green!50!black,
  fonttitle=\bfseries,
  #1
}

% ── Praktisk info-box ─────────────────────────────────────
\newtcolorbox{infobox}[1][]{
  colback=blue!5,
  colframe=blue!40!black,
  fonttitle=\bfseries,
  #1
}

% Orangebox
\newtcolorbox{orangebox}[1][]{
  breakable,
  colback=orangepale,
  colframe=orangedark,
  fonttitle=\headingfont\bfseries\small,
  coltitle=white,
  colbacktitle=orangedark,
  sharp corners,
  #1
}

% Bluebox
\newtcolorbox{bluebox}[1][]{
  breakable,
  colback=bluepale,
  colframe=bluedark,
  fonttitle=\headingfont\bfseries\small,
  coltitle=white,
  colbacktitle=bluedark,
  sharp corners,
  #1
}

% Greenbox
\newtcolorbox{greenbox}[1][]{
  breakable,
  colback=greenpale,
  colframe=greendark,
  fonttitle=\headingfont\bfseries\small,
  coltitle=white,
  colbacktitle=greendark,
  sharp corners,
  #1
}

% Redbox
\newtcolorbox{redbox}[1][]{
  breakable,
  colback=redpale,
  colframe=reddark,
  fonttitle=\headingfont\bfseries\small,
  coltitle=white,
  colbacktitle=reddark,
  sharp corners,
  #1
}

% Greybox
\newtcolorbox{greybox}[1][]{
  breakable,
  colback=greylight,
  colframe=greydark,
  fonttitle=\headingfont\bfseries\small,
  coltitle=white,
  colbacktitle=greydark,
  sharp corners,
  #1
}

% Exempelbox (vänsterborderlinje)
\newtcolorbox{examplebox}[1][]{
  breakable,
  colback=greylight,
  colframe=greydark,
  leftrule=4pt, rightrule=0pt, toprule=0pt, bottomrule=0pt,
  sharp corners,
  fonttitle=\headingfont\bfseries\footnotesize,
  coltitle=white,
  #1
}

% Analogibox
\newtcolorbox{analogybox}[1][]{
  colback=teal!8,
  colframe=teal!60!black,
  fonttitle=\bfseries,
  #1
}

% Formelbox
\newtcolorbox{formulabox}{
  colback=orangepale,
  colframe=orange,
  boxrule=2pt,
  sharp corners,
  halign=center,
  fontupper=\headingfont\bfseries\large\color{greydark},
}

% Storybox (för berättelser och anekdoter)
\newtcolorbox{storybox}[1][]{
  colback=orange!8!white,
  colframe=orange!60!black,
  fonttitle=\bfseries,
  title={#1},
  before upper={\parindent=0pt}
}

% ── Facit-block: fråga (grå) + svar (grön) i en box ──────
\newtcolorbox{facitblock}[1]{
  skin=bicolor,
  breakable,
  sharp corners,
  fonttitle=\headingfont\bfseries\small,
  coltitle=white,
  colbacktitle=greydark,
  colframe=greydark,
  colback=greylight,
  colbacklower=greenpale,
  segmentation style={solid, draw=greydark, line width=0.4pt},
  title={Fråga #1},
}

% ── Sidhuvud/sidfot ───────────────────────────────────────
\usepackage{fancyhdr}
\setlength{\headheight}{16.46pt}
\addtolength{\topmargin}{-2.46pt}
\pagestyle{fancy}
\fancyhf{}
\fancyhead[LO]{\colorbox{orangedark}{\textcolor{white}{\headingfont\footnotesize\MakeUppercase{\leftmark}}}}
\fancyhead[RE]{\colorbox{orangedark}{\textcolor{white}{\headingfont\footnotesize\MakeUppercase{Certifikatshandboken}}}}
\fancyhead[RO]{\headingfont\footnotesize\color{greymed}\MakeUppercase{\rightmark}}
\fancyhead[LE]{\headingfont\footnotesize\color{greymed}\MakeUppercase{Radioamatör Utbildning · 1:a upplagan}}
\renewcommand{\headrulewidth}{0pt}
\fancyfoot[C]{\headingfont\small\thepage}
\fancyfoot[L]{\headingfont\footnotesize\color{greymed}Certifikatshandboken · 1:a upplagan}
\fancyfoot[R]{\headingfont\footnotesize\color{greymed}Radioskola.se}

% ── Rubriker (titlesec) ───────────────────────────────────
\usepackage{titlesec}

\titleformat{\chapter}[display]
  {\headingfont\bfseries\color{greydark}}
  {\color{orange}\footnotesize\MakeUppercase{Kapitel \thechapter}}
  {0pt}
  {\Huge}
  [\vspace{2mm}\textcolor{orange}{\rule{12mm}{3pt}}]

\titleformat{\section}
  {\headingfont\bfseries\large\color{greydark}}
  {\thesection}
  {0.5em}
  {}
  [\vspace{-1mm}\textcolor{greyborder}{\rule{\linewidth}{1.5pt}}]

\titleformat{\subsection}
  {\headingfont\bfseries\normalsize\color{greydark}}
  {\thesubsection}
  {0.5em}
  {}

\titlespacing{\chapter}{0pt}{-10pt}{12pt}
\titlespacing{\section}{0pt}{8pt}{4pt}
\titlespacing{\subsection}{0pt}{6pt}{3pt}

% ── Övrigt ────────────────────────────────────────────────
\usepackage{hyperref}
\hypersetup{
  colorlinks=true,
  linkcolor=orangedark,
  urlcolor=bluedark,
  pdftitle={Certifikatshandboken – 1:a upplagan},
  pdfauthor={Dennis Bengtsson}
}
\usepackage{microtype}
\usepackage{setspace}
\usepackage{enumitem}
\usepackage{amsmath}
\usepackage{amssymb}
\usepackage{lettrine}   % dropcap
\usepackage{tikz}
\usepackage{circuitikz}
\usepackage{pgf}
\usepackage{pifont}    % specialsymboler

% Steg-lista (numrerade med orange cirklar)
\newlist{steplist}{enumerate}{1}
\setlist[steplist]{
  label={\protect\stepnum{\arabic*}},
  leftmargin=*,
  labelindent=0pt,
  itemsep=2pt
}
\newcommand{\stepnum}[1]{%
  \tikz[baseline=(char.base)]{
    \node[circle,fill=orange,inner sep=1.5pt,minimum size=16pt]
      (char) {\textcolor{white}{\headingfont\bfseries\footnotesize #1}};
  }%
}

% Checkboxlista
\newlist{checklist}{itemize}{1}
\setlist[checklist]{label=\textcolor{orange}{›}, leftmargin=*}

% Brödtext
\setlength{\parindent}{0pt}
\setlength{\parskip}{4pt}
\onehalfspacing

% ── Tabellinställningar ───────────────────────────────────
\renewcommand{\arraystretch}{1.3}
\newcolumntype{L}[1]{>{\raggedright\arraybackslash}p{#1}}
\newcolumntype{C}[1]{>{\centering\arraybackslash}p{#1}}

% Tabellstil
\newcommand{\tblhdr}[1]{\cellcolor{greydark}\textcolor{white}{\headingfont\bfseries\footnotesize\MakeUppercase{#1}}}

% Frågeblock – nivåmärken
\newcommand{\qgrund}{\colorbox{greenpale}{\textcolor{greendark}{\headingfont\bfseries\tiny GRUND}}}
\newcommand{\qrakna}{\colorbox{bluepale}{\textcolor{bluedark}{\headingfont\bfseries\tiny RÄKNA}}}
\newcommand{\qprov}{\colorbox{redpale}{\textcolor{reddark}{\headingfont\bfseries\tiny PROV}}}
\newcommand{\qfallan}{\colorbox{ljusbrun}{\textcolor{kakao}{\headingfont\bfseries\tiny FÄLLAN}}}

% Frågeblock-miljö
\newtcolorbox{qblock}[2][]{
  breakable,
  colback=white,
  colframe=greyborder,
  boxrule=0.5pt,
  sharp corners,
  before upper={\textbf{\headingfont\footnotesize\color{orange}#2\quad}},
  #1
}

% Signatur
\newcommand{\signature}[2]{%
  \vspace{4mm}
  \textcolor{greyborder}{\rule{\linewidth}{0.5pt}}
  \vspace{2mm}\\
  \begin{minipage}{0.5\textwidth}
    {\headingfont\bfseries\normalsize #1}\\
    {\footnotesize\itshape\color{greymed} #2}
  \end{minipage}
  \hfill
  \begin{minipage}{0.4\textwidth}
    \raggedleft
    {\headingfont\bfseries\color{orange}Radioskola.se}\\
    {\footnotesize\itshape\color{greymed}Gratis radioamatörutbildning för alla}
  \end{minipage}
}

% ═══════════════════════════════════════════════════════════
%  DOKUMENTET BÖRJAR
% ═══════════════════════════════════════════════════════════
\begin{document}
\let\cleardoublepage\clearpage


% ── FRAMSIDA ──────────────────────────────────────────────
\begin{titlepage}
\pagecolor{coverbg}
\color{white}

% Topband (hela vägen upp, ingen vspace innan)
{\colorbox{orangedark}{\makebox[\linewidth][l]{%
  \hspace{5mm}
  {\monofont\footnotesize\color{white} · · · − − − · · ·}
  \hfill
  {\footnotesize\color{white} Radioskola.se~·~2026}
  \hspace{5mm}
}}}

\vspace{6mm}

% ── ILLUSTRATION ──────────────────────────────────────────
\begin{center}
\begin{tikzpicture}[scale=0.75]

  % Bakgrund
  \fill[coverbg] (-5,-0.5) rectangle (5,7);

  % ── Svensk flagga (övre högra hörnet) ──
  \fill[blue!70!black] (2.8,5.2) rectangle (4.8,6.6);
  \fill[yellow!90!orange] (3.5,5.2) rectangle (3.85,6.6); % vertikalt
  \fill[yellow!90!orange] (2.8,5.8) rectangle (4.8,6.15); % horisontellt
    
% ── Torn (Eiffel-inspirerat, raka ben) ──
  % Raka ben som smalnar uppåt
  \draw[orange!80, line width=2.5pt] (-1.8,0) -- (-0.9,1.8);  % vänster ben, nedre del
  \draw[orange!80, line width=2.5pt] (1.8,0)  -- (0.9,1.8);   % höger ben, nedre del
  \draw[orange!80, line width=2pt]   (-0.9,1.8) -- (-0.3,3.8); % vänster ben, övre del
  \draw[orange!80, line width=2pt]   (0.9,1.8)  -- (0.3,3.8);  % höger ben, övre del

  % Första plattformen
  \draw[orange!70, line width=2pt] (-0.9,1.8) -- (0.9,1.8);

  % Diagonala stag, nedre del
  \draw[orange!50, line width=1pt] (-1.8,0)  -- (-0.0,1.8);
  \draw[orange!50, line width=1pt] (1.8,0)   -- (0.0,1.8);
  \draw[orange!50, line width=1pt] (-1.35,0.9) -- (1.35,0.9);  % horisontellt stag

  % Diagonala stag, övre del
  \draw[orange!50, line width=1pt] (-0.9,1.8) -- (0.3,3.8);
  \draw[orange!50, line width=1pt] (0.9,1.8)  -- (-0.3,3.8);
  \draw[orange!50, line width=1pt] (-0.6,2.6) -- (0.6,2.6);    % horisontellt stag

  % Topp-plattform
  \draw[orange!70, line width=1.5pt] (-0.3,3.8) -- (0.3,3.8);

  % Vertikalt torn ovanpå
  \draw[orange!80, line width=2.5pt] (0,3.8) -- (0,5.0);

  % Antenn-element
  \draw[white, line width=2.5pt] (-0.8,4.4) -- (0.8,4.4);
  \draw[white, line width=2pt]   (-0.55,4.7) -- (0.55,4.7);
  \draw[white, line width=1.5pt] (-0.3,5.0)  -- (0.3,5.0);

  % Toppljus med glow
  \fill[red, opacity=0.15] (0,5.15) circle (10pt);
  \fill[red, opacity=0.25] (0,5.15) circle (7pt);
  \fill[red, opacity=0.45] (0,5.15) circle (5pt);
  \fill[red!80]            (0,5.15) circle (3.5pt);

  % Radiovågor
  \foreach \r/\op in {0.5/0.9, 1.0/0.65, 1.6/0.4, 2.2/0.2}{
    \draw[orange, opacity=\op, line width=1pt]
      (-\r, 4.4) arc(180:0:\r);
  }

\end{tikzpicture}
\end{center}

\vspace{2mm}

% Titel
{\headingfont\tiny\color{orange}\MakeUppercase{1:a upplagan · 2026}}\\[2mm]
{\headingfont\small\color{greymed}\MakeUppercase{Studieguide för radioamatörexamen}}\\[3mm]
{\headingfont\fontsize{42}{44}\selectfont\bfseries\MakeUppercase{Certifikats-}}\\
{\headingfont\fontsize{42}{44}\selectfont\bfseries\textcolor{orange}{\MakeUppercase{handboken}}}

\vspace{4mm}
{\itshape\color{greymed} Allt du behöver för att klara provet – och komma ut på luften.}

\vspace{5mm}

% ── Pills ──
\begin{center}
\foreach \pill in {Elektronik, Radioteknik, Antenner, Vågutbredning, Regler, Trafikmetoder, Säkerhet}{
  \colorbox{orange!25}{%
    \textcolor{orange}{%
      \headingfont\bfseries\small\MakeUppercase{\pill}%
    }%
  }\hspace{2pt}%
}
\end{center}

\vspace{6mm}
\textcolor{orange}{\rule{\linewidth}{2pt}}
\vspace{3mm}

% Info-band
\begin{tabular}{cccc}
  {\headingfont\bfseries\Large\color{orange}10} &
  {\headingfont\bfseries\Large\color{orange}80+} &
  {\headingfont\bfseries\Large\color{orange}Fri} &
  {\headingfont\bfseries\Large\color{orange}SE/SM} \\[2pt]
  {\tiny\color{greymed}\MakeUppercase{Kapitel}} &
  {\tiny\color{greymed}\MakeUppercase{Avsnitt}} &
  {\tiny\color{greymed}\MakeUppercase{Ingen kostnad}} &
  {\tiny\color{greymed}\MakeUppercase{Svensk standard}} \\
\end{tabular}

\vfill

\begin{center}
{\headingfont\footnotesize\color{gray}Radioskola.se · \emph{1:a upplagan} · Dennis Bengtsson med hjälp av Claude AI}
\end{center}

\end{titlepage}
\pagecolor{white}
\color{black}

% ── FÖRORD ────────────────────────────────────────────────
\chapter*{Förord}
\addcontentsline{toc}{chapter}{Förord}
\markboth{Förord}{}

{\itshape Varför den här boken finns – och varför du håller den i handen}

\vspace{4mm}

\lettrine[lines=3, loversize=0.1]{\textcolor{orange}{D}}{en} här boken följer strukturen för det svenska radioamatörcertifikatet, men den är inte skriven som en torr regelsamling. Varje kapitel börjar med \textit{varför} – varför fungerar det så här, vad händer egentligen i kabeln, i luften, i mottagaren? Först när du förstår \textit{varför} blir formlerna och reglerna självklara.

Jag har tänkt på dig som aldrig öppnat en elektronikbok. På dig som kanske är tolv år och nyfiken, eller sjuttiotvå och alltid velat förstå hur radion fungerar. Stegen i den här boken är medvetet små – ingen ska behöva stanna upp och känna sig dummare än när de började. Tvärtom: varje sida ska ge en liten \textit{aha}-känsla som gör att du vill läsa vidare.

Den här boken är skapad i nära samarbete med AI-verktyg. Min roll har varit att se över pedagogiken, välja vad som fungerar, omdefiniera det som inte gör det. Jag har hela tiden frågat mig: fattar en nybörjare det här? Det ansvaret är mitt.

Vi börjar från grunden. Vad är egentligen ström? Vad är spänning? Vad menar man med skindjup? Jaha, ser symbolen för ohm ut så?
Från de frågorna bygger vi – bit för bit, lager för lager – hela vägen upp till antenner, radiovågor och internationell kommunikation. Ingenting tas för givet, och ingenting introduceras innan du har verktygen att förstå det.

\vspace{4mm}

För många har radio varit just det – en inkörsport till ett helt nytt intresse. Elektroniken som från början verkade som ett nödvändigt ont för att klara provet visar sig vara ett djupt och fascinerande område i sig. Kanske slutar det inte med ett certifikat. Kanske börjar du bygga egna antenner, reparera gammalt radiogods, eller fördjupa dig i signalbehandling och mjukvaruradio. Det har hänt förr. Andra vill kunna kommunicera vid katastrofer, när all övrig kommunikation är nedsläckt. Oavsett vad som lockar dig är det här just den bok som tar dig dit.

\vspace{4mm}

Det enda som krävs av dig är nyfikenhet — samma nyfikenhet jag hade när jag skrev den här boken. Jag lärde mig det mesta samtidigt som jag skrev den. Certifikatet är målet, men nyfikenheten är resan — och den slutar inte vid provbordet. Jag hoppas att den här boken inte bara hjälper dig att klara provet, som den gjorde för mig, utan också tänder en gnista som får dig att vilja lära dig mer, bygga mer, och framför allt: kommunicera mer.

\vspace{4mm}

På radioskola.se hittar du övningsfrågor och ett simulerat prov — inspirerat av körkortsteorin, fast för radio.

\vspace{6mm}


{\raggedleft\itshape Dennis Bengtsson\\Radioskola.se, 2026\par}

\chapter*{Innan du börjar}
\addcontentsline{toc}{chapter}{Innan du börjar}
\markboth{Innan du börjar}{}

\begin{orangebox}[title={\etitle{📻}{Vad väntar dig som radioamatör?}}]
Med ett certifikat får du rätt att sända på ett stort antal frekvensband – från kortvåg som studsar runt jorden till VHF och UHF för lokal och satellittrafik. Du kan prata med expeditioner i Antarktis, delta i tävlingar, bygga egna antenner och sändare, och ingå i ett globalt nätverk där tekniken är det gemensamma språket.

Och när mobilnäten och internet går ner vid en katastrof – är det radioamatörerna som fortfarande är på luften.

\textbf{Det finns två certifikatsnivåer:} instegscertifikatet ger dig tillgång till de flesta banden med begränsad effekt, medan HAREC-certifikatet är det fullständiga internationellt erkända certifikatet som den här boken förbereder dig för. Klarar du HAREC klarar du allt.
\end{orangebox}

\begin{bluebox}[title={\emoji{📚} Så här använder du handboken effektivt}]
Du behöver inte mer matematik än grundskolan – addition, multiplikation och lite bråkräkning. Resten lär du dig på vägen. De flesta klarar provet efter fyra till sex veckors studier.

\begin{steplist}
  \item Läs varje kapitel i ordning – grunderna återkommer i alla senare avsnitt
  \item Räkna formelexemplen med penna och papper – det enda sättet att lära sig räkna
  \item Orangefärgade rutor = exakt vad examinatorn frågar om
  \item Gör övningsfrågorna utan att titta på svaren först
  \item Repetera kapitlet om du missade fler än 2 av 5 frågor
  \item Besök Radioskola.se för interaktiva grafer och provsimulatorer
\end{steplist}
\end{bluebox}

\begin{greenbox}[title={\emoji{🌍} Radioamatörer i siffror}]
Det finns över \textbf{tre miljoner licensierade radioamatörer} i världen, varav drygt 10\,000 i Sverige. Under katastrofer som jordbävningar, orkaner och strömavbrott har amatörradionätverk upprepade gånger räddat liv när all annan kommunikation slagits ut. Hobbyn är erkänd av ITU – FN:s organ för telekommunikation.
\end{greenbox}

\begin{greybox}[title={\emoji{⚠️} Viktigt om regelverk}]
Frekvenstillstånd, effektgränser och bandplaner kan ändras. Kontrollera alltid aktuell information på \textbf{pts.se} och \textbf{ssa.se} innan du genomför ditt prov.
\end{greybox}

\signature{Dennis\,}{Grundare, Radioskola.se · 2026}

% ── INNEHÅLLSFÖRTECKNING ──────────────────────────────────
\tableofcontents

% ═══════════════════════════════════════════════════════════
%  KAPITEL 1 – GRUNDLÄGGANDE ELEKTRONIK
% ═══════════════════════════════════════════════════════════
\chapter{Grundläggande Elektronik}

\lettrine[lines=3, loversize=0.1]{\textcolor{orange}{A}}{ll} radioteknik vilar på en och samma grund: förståelsen för hur elektroner rör sig, hur kretsar beter sig och hur vi kan forma elektriska signaler till radiofrekvent energi. Det är inte svårare än snickeri i grunden – det handlar om att lära sig ett nytt material och dess egenskaper.

Du behöver inte ha byggt en enda krets för att förstå det här kapitlet. Det räcker med att du någon gång slangat en trädgård, skruvat på en kran eller undrat varför det bubblar mer när du häller ketchup ur en smal flaska än en bred. Elektronen är inget mystiskt – den är bara ett litet korn i ett enormt vattenfall, och det fallet kan vi lära oss styra.

Det är också viktigt att förstå \emph{varför} vi börjar här. En radioamatör som förstår elektronikens grunder kan felsöka sin egen utrustning, anpassa en antenn, läsa ett kopplingsschema och – kanske viktigast av allt – förstå vad som faktiskt händer inuti lådan. Det ger en frihet som ingen köpt lösning kan ge.

Vi börjar smått: med ström och spänning, de allra grundläggande storheterna. Sedan lägger vi till motstånd och ser hur de tre hänger ihop i Ohms lag – en av de mest användbara formler du kommer att stöta på. Steg för steg bygger vi upp bilden, från likström till växelström, från enkla motstånd till kondensatorer och spolar, tills vi till sist förstår hur ett filter fungerar och varför en resonanskrets kan plocka ut just den station du vill lyssna på.

När du är klar med kapitel~1 kommer du inte bara att kunna räkna på enkla kretsar – du kommer att \emph{se} dem för dig.

\begin{bluebox}[title={Vad du lär dig i kapitel 1}]
\begin{tabular}{@{}ll@{}}
  \textbf{1.1} & Elektrisk ström och spänning \\
  \textbf{1.2} & Resistans och Ohms lag \\
  \textbf{1.3} & Effekt och energi \\
  \textbf{1.4} & Kondensatorer \\
  \textbf{1.5} & Spolar och induktans \\
  \textbf{1.6} & Serie- och parallellkoppling \\
  \textbf{1.7} & Växelström och impedans \\
  \textbf{1.8} & Filter och resonanskretsar \\
\end{tabular}
\end{bluebox}

% ── 1.1 ───────────────────────────────────────────────────
\section{Elektrisk ström och spänning}

\lettrine[lines=2]{\textcolor{orange}{E}}{lektrisk} ström är flödet av elektriska laddningar – elektroner – genom en ledare. Varje gång du tänder en lampa, sänder ett radiomeddelande eller startar din bil rör sig miljarder elektroner genom kopparkablar och kretsar. Att förstå ström och spänning är att förstå den allra mest grundläggande mekanismen bakom all elektronik.

\begin{storybox}[{\emoji{🎙️} En natt i fält – när batteriet bestämmer hur länge du får sända}]
Det är sent en fredagskväll. Du sitter i ett tält någonstans i 
Dalaskogen med en handhållen VHF-radio, ett 7\,Ah-batteri och ett 
uppdrag: hålla kontakten med basläger under de kommande åtta timmarna.

Radiooperatören som tränade dig en gång sa: \emph{``Lär dig ditt 
batteri bättre än du känner dina skor. Det är det enda som håller 
dig i kontakt.''}

Du vet att radion drar 1{,}5\,A när du sänder och 0{,}3\,A när du 
lyssnar. Du planerar att sända ungefär 20\,\% av tiden. Med det du 
lär dig i det här avsnittet räknar du snabbt ut genomsnittlig 
förbrukning och batteritid -- och andas ut. Det räcker med marginal.

Det är ett enkelt räkneexempel. Men ute i skogen, utan ett 
vägguttag i sikte, är det skillnaden mellan att fullfölja uppdraget 
eller inte. Formeln lärde du dig på fem minuter. Förståelsen för 
vad den \emph{betyder} -- den bär du med dig resten av livet.
\end{storybox}

\begin{examplebox}[title={\emoji{💧} Vattenanalogin – ett sätt att tänka}]
Föreställ dig vatten i ett rör. \textbf{Strömmen (I)} är som vattenflödet – hur många liter som passerar en punkt per sekund. \textbf{Spänningen (U)} är som vattentrycket – den kraft som driver vattnet framåt. Utan tryck rör sig inget vatten, och utan spänning flödar ingen elektrisk ström. Ju mer tryck, desto mer flöde – precis som i en elektrisk krets.
\end{examplebox}

% ── 1.1.1 ─────────────────────────────────────────────────
\subsection{Ström – Ampere (A)}

Ström betecknas \textbf{I} och mäts i \textbf{Ampere (A)}, uppkallat efter den franske fysikern André-Marie Ampère. I radioteknik möter vi enorma variationer – från bråkdelar av en mikroampere i en känslig mottagarförstärkare till 20 ampere eller mer i ett HF-slutsteg vid full sändeffekt.

\begin{center}
\begin{tabular}{llll}
\rowcolor{greydark}
\tblhdr{Prefix} & \tblhdr{Symbol} & \tblhdr{Värde} & \tblhdr{Exempel i radioteknik} \\
Mikroampere & µA & 0,000\,001 A & S-meter, svaga ingångssignaler i mottagare \\
\rowcolor{greylight}
Milliampere & mA & 0,001 A & LED-indikator (20 mA), mikrofonförstärkare \\
Ampere & A & 1 A & Handhållen VHF-radio vid mottagning (0,5–2 A) \\
\rowcolor{greylight}
Kiloampere & kA & 1\,000 A & Blixt (20–200 kA) – därav behovet av blixtnedledare \\
\end{tabular}
\end{center}

Det vanligaste misstaget i beräkningar är att glömma konvertera milliampere till ampere. Kom ihåg: \textbf{1 A = 1\,000 mA = 1\,000\,000 µA.} Ska du räkna med Ohms lag måste strömmen alltid vara i ampere – annars blir svaret tusen gånger fel.

\textbf{Snabbomvandling:} 500 mA = 0,5 A \quad | \quad 250 mA = 0,25 A \quad | \quad 50 mA = 0,05 A

\subsubsection*{Övningsfrågor – 1.1.1}

\begin{qblock}{1.}
\qgrund\quad Vilken enhet mäts elektrisk ström i? Vad är beteckningen för ström i formler? Ange också de vanligaste underenheterna och deras relationer till grundenheten.
\end{qblock}

\begin{qblock}{2.}
\qrakna\quad Omvandla följande strömvärden till ampere: a) 250 mA \quad b) 4\,200 µA \quad c) 0,75 kA \quad d) 85 mA. Omvandla dessutom 0,035 A till mA och till µA.
\end{qblock}

% ── 1.1.2 ─────────────────────────────────────────────────
\subsection{Spänning – Volt (V)}

Spänning betecknas \textbf{U} (från tyska \emph{Unterschied} – skillnad) och mäts i \textbf{Volt (V)}. Det är skillnaden i elektrisk potential mellan två punkter. Utan spänningsskillnad händer ingenting – precis som vatten inte rinner om marken är helt plan.

\begin{center}
\begin{tabular}{ll}
\rowcolor{greydark}
\tblhdr{Spänning} & \tblhdr{Källa eller användning} \\
1,5 V & AA/AAA-batteri (ficklampor, enklare handhållna apparater) \\
\rowcolor{greylight}
9 V & Transistorbatteri (äldre handhållna radioapparater) \\
12 V & Bilbatteri, mobila installationer \\
\rowcolor{greylight}
\textbf{13,8 V} & \textbf{Standard för amatörradioutrustning – universell standard inom hobbyn} \\
48 V & Fantommatning för kondensatormikrofoner \\
\rowcolor{greylight}
230 V & Hushållsel i Sverige (växelström – livsfarligt vid felhantering) \\
\end{tabular}
\end{center}

\begin{redbox}[title={\emoji{⚠️} Elsäkerhet}]
Spänningar över \textbf{50 V} kan vara livsfarliga – bland annat nätspänning (230 V AC) och kondensatorer som håller kvar laddning länge efter att apparaten stängts av. Elsäkerhet, jordning och åskskydd behandlas samlat i kapitel~9.
\end{redbox}

\subsubsection*{Övningsfrågor – 1.1.2}

\begin{qblock}{3.}
\qgrund\quad Vilken spänning är standard för amatörradioutrustning, och varför är det \emph{inte} exakt 12 V trots att de flesta mobila installationer körs från ett 12 V-bilbatteri?
\end{qblock}


% ── 1.1.3 ─────────────────────────────────────────────────
\subsection{Likström (DC) och växelström (AC)}

Det finns två fundamentalt olika typer av elektrisk ström. Att veta skillnaden – och att aldrig blanda ihop dem – är avgörande för säker drift av din radiostation.

\begin{examplebox}[title={\emoji{💧} DC och AC i vattenanalogin}]
Likström (DC) är som en flod som alltid rinner åt samma håll – stabilt och förutsägbart. Växelström (AC) är som vattnet i ett tidvattenområde som rytmiskt byter riktning – det skvalpar fram och tillbaka 50 gånger i sekunden. Din radio drivs av DC inuti, men hämtar sin energi från vägguttaget med AC via ett nätaggregat som omvandlar.
\end{examplebox}

\begin{center}
\begin{tabular}{L{2cm}L{1.5cm}L{4.5cm}L{4.5cm}}
\rowcolor{greydark}
\tblhdr{Typ} & \tblhdr{Förk.} & \tblhdr{Beskrivning} & \tblhdr{Användning i radio} \\
Likström & DC & Flödar alltid i samma riktning, från plus till minus & Batterier, nätaggregat (13,8 V DC), all elektronik inuti radion \\
\rowcolor{greylight}
Växelström & AC & Växlar riktning periodiskt. I Sverige 50 Hz = 50 gånger per sekund & Vägguttag (230 V AC, 50 Hz), transformatorer i nätaggregat \\
\end{tabular}
\end{center}

Inom amatörradio är \textbf{13,8 V DC} den universella strömförsörjningsstandarden –
det är den spänning alla stationära transceivrar och de flesta tillbehör är konstruerade
för. Men varför just 13,8 V och inte det mer välbekanta 12 V? Ett 12 V-blybatteri har
i fullt laddat tillstånd faktiskt 12,6–13,2 V, och vid inkopplat laddningsaggregat
stiger spänningen till 13,6–14,4 V. Radioapparater konstrueras för att fungera optimalt
i hela detta spänningsintervall, med 13,8 V som mittenvärde. En radio som matas med
bara 11–12 V levererar märkbart lägre sändeffekt.

Kunskapen om ström och spänning används konkret varje dag i en 
radiostation:

\textbf{Beräkna batteritid:} Batteritid (h) = Kapacitet (Ah) $\div$ Förbrukning (A). Exempel: ett 7 Ah-batteri och en radio som drar 0,5 A ger 14 timmars drift. Blanda sändning och mottagning? Räkna ett viktat medelvärde av strömmen.

\textbf{Dimensionera nätaggregat:} Välj alltid aggregat med god marginal. En 100 W HF-radio drar ungefär 2 A vid mottagning men upp till 22 A vid full sändeffekt. Slutsteget har en verkningsgrad på ca 50–60\,\% – det betyder att radion drar 180–200 W från nätaggregatet för att leverera 100 W RF ut i antennen. 200 W $\div$ 13,8 V $\approx$ 14,5 A, och med övrig elektronik landar totalförbrukningen kring 20–22 A. Välj ett 25–30 A-aggregat, inte ett på 15 A.

\textbf{Kabeltjocklek spelar roll:} Tunn kabel ger spänningsfall. En kabel med 10 m enkel väg (20 m tur och retur) av 2,5 mm² koppar har ungefär 0,14 $\Omega$ DC-resistans. Vid 20 A uppstår ett spänningsfall på $U = R \times I = 0{,}14 \times 20 = 2{,}8$ V – radion får bara 11,0 V i stället för 13,8 V. Använd minst 6 mm² kabel för HF-transceiver med slutsteg, och håll kabellängden så kort som möjligt.

\begin{worked}[title={Räkneexempel 1.1.3-1 – Batteritid under fältövning}]
En handhållen VHF-radio sänder med 1,5 A och tar emot med 0,3 A. Under en fältövning sänder du 20\,\% av tiden och lyssnar 80\,\%. Batteriet är 12 Ah.

\textbf{Genomsnittlig ström:}
\[ I = 0{,}20 \times 1{,}5 + 0{,}80 \times 0{,}3 = 0{,}30 + 0{,}24 = 0{,}54\,\text{A} \]
\textbf{Batteritid:}
\[ t = \frac{12\,\text{Ah}}{0{,}54\,\text{A}} \approx 22\,\text{timmar} \]
\end{worked}

\begin{worked}[title={Räkneexempel 1.1.3-2 – Välja rätt nätaggregat}]
Din HF-transceiver kräver 13,8 V och har maxförbrukning 22 A vid sändning, 2 A vid mottagning. Vilken minsta strömstyrka behöver nätaggregatet klara?

Säkert val: välj \textbf{25–30 A} för att ha marginal och undvika att aggregatet arbetar vid max hela tiden – det förkortar livslängden och ger sämre spänningsstabilitet.
\end{worked}

\subsubsection*{Övningsfrågor – 1.1.3}

\begin{qblock}{4.}
\qgrund\quad Vad är skillnaden mellan likström (DC) och växelström (AC)? Ge ett konkret praktiskt exempel på vardera typ i en typisk radiostation.
\end{qblock}

\begin{qblock}{5.}
\qrakna\quad Du har ett 12 Ah-batteri och en handhållen VHF-radio som förbrukar 1,5 A vid sändning och 0,3 A vid mottagning. Under ett fältpass sänder du 20\,\% och tar emot 80\,\% av tiden. a) Vad är den genomsnittliga strömförbrukningen? b) Hur länge räcker batteriet?
\end{qblock}

\begin{qblock}{6.}
\qprov\quad Vid nödkommunikation ska du driva en HF-transceiver (13,8 V, max 22 A vid sändning / 2 A vid mottagning) från en 60 Ah-batteribank. Du sänder 15 min och tar emot 45 min per timme. a) Genomsnittlig ström? b) Drifttid vid 50\,\% urladdningsgräns? c) Räcker ett 25 A nätaggregat?
\end{qblock}

\begin{orangebox}[title={\emoji{🎯} Viktigt för provet}]
\begin{checklist}
  \item Ström mäts i \textbf{Ampere (A)}, betecknas \textbf{I}
  \item Spänning mäts i \textbf{Volt (V)}, betecknas \textbf{U}
  \item DC flödar alltid åt samma håll – AC byter riktning periodiskt
  \item Amatörradiostandardspänning: \textbf{13,8 V DC}
  \item Konvertera \textbf{alltid} mA till A innan du beräknar (1 A = 1\,000 mA)
  \item Spänningar över 50 V är livsfarliga – se kapitel~9
\end{checklist}
\end{orangebox}

% ── 1.2 ───────────────────────────────────────────────────
\section{Resistans och Ohms lag}

\lettrine[lines=2]{\textcolor{orange}{R}}{esistans} är ett materials förmåga att bromsa elektrisk ström. Alla material har resistans, men skillnaderna är astronomiska: koppars resistivitet är femton miljoner gånger lägre än gummits. Det är denna skillnad som gör det möjligt att leda ström exakt dit vi vill – och blockera den överallt annars.

\begin{storybox}[{\emoji{🔧} Den lysande insikten – om en LED som vägrade lysa}]
Det var en torsdagskväll och Håkan, nybliven radioamatör med ett 
halvfärdigt hemmabygge på bordet, höll ett lysdiodpaket i handen 
och funderade. Han hade anslutit LED:en direkt till 12\,V-matningen. 
Det hade blixtat till, luktat lite konstigt -- och sedan hade dioden 
aldrig lyst igen.

\emph{``Varför?''}, frågade han sin handledare på klubbkvällen.

Svaret var enkelt: en LED vill ha ungefär 2\,V och 20\,mA. Utan en 
serieresistor tar den vad den kan -- och 12\,V är mer än den klarar. 
Handledaren plockade fram ett papper:
\[
R = \frac{U}{I} = \frac{12\,\text{V} - 2\,\text{V}}{0{,}020\,\text{A}} 
= 500\,\Omega
\]
En 470\,$\Omega$-resistor ur lådan, en ny LED -- och indikatorn på 
Håkans radio lyser fortfarande, tio år senare.

Ohms lag räddade inte bara en LED. Den gav Håkan något viktigare: 
förmågan att förstå vad som händer inne i en krets, och att lösa 
problemet nästa gång själv.
\end{storybox}

\begin{examplebox}[title={\emoji{💧} Vattenanalogin – resistans som friktionen i röret}]
Resistans är som friktionen i ett vattenrör. Ett brett, slätt kopparrör (koppartråd) har nästan ingen friktion – vattnet flödar lätt. Ett smalt rör fyllt med grus (kolresistor) har stor friktion – vattnet tvingas fram långsamt. Mer friktion = mer motstånd = lägre strömflöde vid samma tryck (spänning).
\end{examplebox}

Resistans betecknas \textbf{R} och mäts i \textbf{Ohm ($\Omega$)}, uppkallat efter den tyske fysikern Georg Simon Ohm. I praktiken möter vi allt från milliohm i korta kablar till megaohm i isolationsmaterial.

\begin{center}
\begin{tabular}{llll}
\rowcolor{greydark}
\tblhdr{Prefix} & \tblhdr{Symbol} & \tblhdr{Värde} & \tblhdr{Exempel} \\
Milliohm & m$\Omega$ & $10^{-3}\,\Omega$ & Korta kopparledare, kabelskarvkontakter \\
\rowcolor{greylight}
Ohm & $\Omega$ & $1\,\Omega$ & Lågohmiga resistorer, dummylast (50\,$\Omega$) \\
Kilohm & k$\Omega$ & $10^{3}\,\Omega$ & Vanligaste resistorerna i elektronik \\
\rowcolor{greylight}
Megohm & M$\Omega$ & $10^{6}\,\Omega$ & Koaxkabelisolation ($>$1 M$\Omega$), voltmeteringång \\
\end{tabular}
\end{center}

% ── 1.2.1 ─────────────────────────────────────────────────
\subsection{Ohms lag – grundformeln}

Ohms lag beskriver det matematiska sambandet mellan spänning (U), 
ström (I) och resistans (R). Det är den mest fundamentala lagen i 
all elektronik och sitter i ryggmärgen på alla som arbetar med teknik.

Täck det du söker med fingret – det som syns är formeln:

\begin{center}
\begin{tikzpicture}
  \draw[thick] (0,0) -- (3,0) -- (1.5,2.6) -- cycle;
  \draw[thick] (0.75,1.3) -- (2.25,1.3);
  \draw[thick] (1.5,0) -- (1.5,1.3);
  \node at (1.5, 1.95) {\large\bfseries U};
  \node at (0.75, 0.6) {\large\bfseries R};
  \node at (2.25, 0.6) {\large\bfseries I};
  \node at (1.5, 0.6) {\small $\boldsymbol{\times}$};
  \node at (1.5, 1.62) {\small $\boldsymbol{\div}$};
\end{tikzpicture}

\medskip
\textbf{Täck U} $\rightarrow$ $R \times I$ \quad\textbf{·}\quad
\textbf{Täck R} $\rightarrow$ $U \div I$ \quad\textbf{·}\quad
\textbf{Täck I} $\rightarrow$ $U \div R$
\end{center}

\begin{worked}[title={Räkneexempel 1.2.1-1 – LED-serieresistor}]
En LED på din radiofront kräver 2 V och 20 mA. Radion matas med 12 V. Vilken resistans behövs i serie med LED:en?

Spänning över resistorn: $12\,\text{V} - 2\,\text{V} = 10\,\text{V}$ \quad Omvandla: $20\,\text{mA} = 0{,}020\,\text{A}$
\[ R = U \div I = 10\,\text{V} \div 0{,}020\,\text{A} = \mathbf{500\,\Omega} \]
Välj standardvärdet \textbf{470 $\Omega$} eller \textbf{510 $\Omega$} ur E12-serien.
\end{worked}

\begin{worked}[title={Räkneexempel 1.2.1-2 – Spänningsfall i DC-matningskabel}]
En 5 m DC-matningskabel (tur och retur 10 m) med 1,5 mm² tvärsnittsarea har DC-resistansen 0,12 $\Omega$. Under sändning drar radion 20 A. Vilket spänningsfall uppstår?
\[ U = R \times I = 0{,}12\,\Omega \times 20\,\text{A} = \mathbf{2{,}4\,\text{V}} \]
Med 13,8 V från nätaggregatet får radion bara 11,4 V – sändeffekten sjunker märkbart. Lösning: kortare kabel eller större tvärsnittsarea (minst 6 mm² för HF-transceiver).
\end{worked}

\begin{worked}[title={Räkneexempel 1.2.1-3 – Felsökning med multimeter}]
Du mäter 9 V över en okänd resistor och 18 mA ström. Vad är resistansen? Vilken standardresistor är det troligen?

Omvandla: $18\,\text{mA} = 0{,}018\,\text{A}$
\[ R = U \div I = 9\,\text{V} \div 0{,}018\,\text{A} = \mathbf{500\,\Omega} \]
Troligen en \textbf{470 $\Omega$} eller \textbf{510 $\Omega$} resistor (E12-serien) med 5\,\% tolerans. $470\,\Omega \times 1{,}05 = 493\,\Omega$ – stämmer.
\end{worked}

\subsubsection*{Övningsfrågor – 1.2.1}

\begin{qblock}{1.}
\qgrund\quad Ange Ohms lag i alla tre formerna (för U, för I och för R). Vad betecknar varje bokstav och i vilken enhet mäts de? Rita URI-triangeln och förklara hur den används.
\end{qblock}

\begin{qblock}{2.}
\qrakna\quad Beräkna det okända värdet (konvertera enheter om det behövs):
\begin{center}
\begin{tabular}{ll}
a) $U = ?$, $R = 330\,\Omega$, $I = 45\,\text{mA}$ &
b) $I = ?$, $U = 13{,}8\,\text{V}$, $R = 4{,}7\,\text{k}\Omega$ \\[4pt]
c) $R = ?$, $U = 6\,\text{V}$, $I = 12\,\text{mA}$ &
d) $U = ?$, $R = 50\,\Omega$, $I = 2\,\text{A}$
\end{tabular}
\end{center}
\end{qblock}

\begin{qblock}{3.}
\qrakna\quad En antennrotors matningskabel har DC-resistansen 0,8\,$\Omega$ (tur och retur), rotorn drar 3 A. a) Vilket spänningsfall uppstår? b) Nätaggregatet ger 13,8 V – vilken spänning når rotorn? c) Rotorn kräver minst 11 V. Godkänd installation?
\end{qblock}

\begin{qblock}{F.}
\qfallan\quad \textbf{Fällan:} En radioamatör beräknar vilket 
motstånd han behöver för att begränsa strömmen genom en krets 
till 500\,mA vid 12\,V. Han skriver:
\[
R = \frac{U}{I} = \frac{12}{500} = 0{,}024\,\Omega
\]
Han beställer en 0{,}024\,$\Omega$-resistor och undrar varför 
ingen sådan finns i katalogen.

Vad gick fel? Beräkna rätt värde och förklara vilket 
standardvärde ur E12-serien som passar.
\end{qblock}

\medskip
\noindent\rule{\textwidth}{0.4pt}
{\small\itshape Försök själv innan du läser svaret.}
\medskip

\begin{worked}[title={Lösning – Fallgropsfråga 1.2.1\,F}]
Felet är att strömmen aldrig konverterades från milliampere 
till ampere. $500\,\text{mA} = 0{,}5\,\text{A}$.

Rätt beräkning:
\[
R = \frac{U}{I} = \frac{12\,\text{V}}{0{,}5\,\text{A}} = 24\,\Omega
\]
Närmaste standardvärde i E12-serien är \textbf{22\,$\Omega$} 
eller \textbf{27\,$\Omega$}. Välj 27\,$\Omega$ om du vill 
hålla strömmen \emph{under} 500\,mA.

\textbf{Lärdomen:} Kontrollera alltid enheten på strömmen 
innan du beräknar. Ett svar på 0{,}024\,$\Omega$ är en 
tydlig signal att något är fel -- en resistor på under 
1\,$\Omega$ är ovanlig och bör alltid väcka misstankar.
\end{worked}

% ── 1.2.2 ─────────────────────────────────────────────────
\subsection{Material och resistivitet}

Inte alla material leder lika bra. Valet av material avgör om vi har en ledare, en resistor eller en isolator – och i radioteknik är detta val kritiskt för kretskortsspår, antennelement, skärmning och isolering.

\begin{center}
\begin{tabular}{L{2.5cm}L{4cm}L{2.5cm}L{4cm}}
\rowcolor{greydark}
\tblhdr{Kategori} & \tblhdr{Typiska material} & \tblhdr{Resistans} & \tblhdr{Användning i radio} \\
Ledare & Silver, koppar, guld, aluminium & Mycket låg & Kablar, antennelement, kretskortsspår, kontaktytor \\
\rowcolor{greylight}
Halvledare & Kisel, germanium, kol & Kontrollerbar & Transistorer, dioder, resistorer \\
Isolatorer & Plast, gummi, keramik, luft & Mycket hög & Kabelisolering, antennisolatorer, chassimontering \\
\end{tabular}
\end{center}

Silver leder faktiskt bättre än koppar – men är dyrare. Skillnaden i resistans är bara cirka 6\,\%. I praktiken används silver enbart som ytbeläggning på kontakter och i militär- och rymdtillämpningar. Koppar ger det bästa förhållandet mellan kostnad och prestanda och är den universella standarden för alla elektriska ledare.

Koaxkablar och radioapparater konstrueras i dag nästan uteslutande för 50\,$\Omega$. Det är ingen slump – 50\,$\Omega$ är en genomtänkt kompromiss. Lägst dämpning i en luftfylld koaxkabel uppnås vid ungefär 77\,$\Omega$, medan högst effektkapacitet uppnås vid ungefär 30\,$\Omega$. Geometriskt medelvärde: $\sqrt{77 \times 30} \approx 48\,\Omega$, avrundat till 50\,$\Omega$. Det satte sig som världsstandard för amatörradio och professionell RF-teknik under andra världskriget.

Resistansen hos en ledare beror inte bara på vilket material den är gjord av, utan också på dess \textbf{längd} och \textbf{tvärsnittsarea}. En lång, tunn koppartråd har högre resistans än en kort, tjock – precis som ett smalt, långt rör ger mer friktion än ett kort, brett. Det är därför kabeltjocklek spelar roll i en radiostation: dubbel längd ger dubbel resistans, dubbel area ger hälften så hög resistans.

\begin{worked}[title={Räkneexempel 1.2.2-1 – Resistans och geometri}]
En koppartråd har resistansen 2\,$\Omega$. Vad händer om du byter 
till en tråd som är dubbelt så lång men har dubbelt så stor 
tvärsnittsarea?

\textbf{Dubbel längd} $\rightarrow$ dubbel resistans: 
$2\,\Omega \times 2 = 4\,\Omega$\\
\textbf{Dubbel area} $\rightarrow$ hälften så hög resistans: 
$4\,\Omega \div 2 = \mathbf{2\,\Omega}$

Effekterna tar ut varandra – resistansen är oförändrad. En kabel 
som är dubbelt så lång men dubbelt så tjock har alltså samma 
resistans som originalet. Det är därför tjockare kabel krävs när 
kabellängden ökar – annars stiger spänningsfallet.
\end{worked}

\newpage

Ohms lag är inte bara teori – den används konkret varje dag i en 
radiostation:

\textbf{Dimensionera kablar:} Beräkna spänningsfall med $U = R \times I$. Max 0,5 V fall i matningskabel (ca 3,6\,\% av 13,8 V) är en bra tumregel.

\textbf{Välja resistorer:} LED-serien, spänningsdelare och strömgränser beräknas alla med Ohms lag. Välj alltid närmaste standardvärde i E12- eller E24-serien.

\textbf{Felsökning:} Mät spänning och ström, beräkna resistans och jämför med schematvärdet. Nästan noll $\Omega$ = kortslutning. Oändligt $\Omega$ = avbrott.

\textbf{Nätaggregatsstorlek:} Beräknas med hänsyn till slutstegets verkningsgrad – se avsnitt~1.3 för fullständig genomgång.

\begin{redbox}[title={\emoji{⚠️} Det absolut vanligaste misstaget}]
Glöm aldrig att konvertera milliampere till ampere innan beräkningen. 
500 mA = 0,5 A, inte 500 A. Fel enhet ger svar som är tusen gånger 
för stora – och det avslöjas direkt om du stannar upp och frågar dig: 
\emph{är det här rimligt?}

\textbf{Kontrollvana:} Gör alltid en rimlighetscheck. En resistor på 
470 $\Omega$ med 12 V spänning bör ge $I = 12 \div 470 \approx 
25\,\text{mA}$ – ungefär som en LED. Om svaret är 25 A har något 
gått fel.
\end{redbox}

\newpage

\subsubsection*{Övningsfrågor – 1.2.2}

\begin{qblock}{4.}
\qgrund\quad Rangordna dessa material från lägst till högst resistans: 
silver, plast, kol, koppar, gummi, aluminium. Förklara sedan varför 
silver inte används i kablar trots att det leder bättre än koppar.
\end{qblock}

\begin{qblock}{5.}
\qgrund\quad En koppartråd har resistansen 2 $\Omega$. Vad händer med 
resistansen om du a) dubblerar längden? b) dubblerar tvärsnittsarean? 
c) gör båda samtidigt? Motivera utan att räkna.
\end{qblock}

\begin{qblock}{6.}
\qprov\quad En SWR-meter mäter den framåtgående spänningen till 70,7 V 
(RMS) i ett 50 $\Omega$-system. a) Beräkna strömmen genom systemet. 
b) Beräkna effekten ($P = U \times I$). c) En tekniker påstår att 
koaxkabelns resistans (3 dB förlust vid denna effektnivå) är 
acceptabelt. Instämmer du? Motivera med beräkning.
\end{qblock}

\begin{qblock}{7.}
\qprov\quad Varför anges antennimpedansen för de flesta 
amatörradioapparater till 50 $\Omega$ och inte 75 $\Omega$? Beskriv 
det fysikaliska resonemang som leder till kompromissen 50 $\Omega$, 
och nämn vid vilken impedans en luftfylld koaxkabel har lägst dämpning 
respektive högst effektkapacitet.
\end{qblock}

\begin{orangebox}[title={\emoji{🎯} Viktigt för provet}]
\begin{checklist}
  \item Ohms lag: $U = R \times I$ (och de två omskrivna formerna)
  \item Resistans mäts i \textbf{Ohm ($\Omega$)}, betecknas \textbf{R}
  \item URI-triangeln: täck det du söker – det kvarvarande är formeln
  \item Konvertera \textbf{alltid} mA $\rightarrow$ A innan beräkning ($\div$ 1\,000)
  \item Resistans ökar med längd och minskar med större tvärsnittsarea
  \item 50 $\Omega$ är standard-impedans för amatörradio och RF-teknik
  \item Koppar används i kablar pga utmärkt ledningsförmåga till rimlig kostnad
\end{checklist}
\end{orangebox}

\newpage

% ── 1.3 ───────────────────────────────────────────────────
\section{Effekt och energi}

\lettrine[lines=2]{\textcolor{orange}{E}}{ffekt} beskriver hur snabbt energi omvandlas eller förbrukas. Resistans och ström räcker inte – vi behöver veta hur mycket arbete som utförs per sekund. Det är skillnaden mellan en liten LED-indikator och ett 100 W slutsteg: båda drar ström, men slutsteget omvandlar energin enormt mycket snabbare.

\begin{storybox}[{\emoji{📻} Begagnatköpet – när ``100\,W'' inte var 100\,W}]
Annonsens rubrik var tydlig: HF-transceiver, 100\,W, fint skick, 
rimligt pris. Maria, licensierad sedan tre år, hade sparat länge 
och åkte dit en lördag med kontanter i fickan.

Hemma i stationen anslöt hon radion, lade ett samtal och tittade 
på wattmetern. Nålen lade sig runt 60\,W. Inte 100.

Var det fel på radion? Hon räknade. En 100\,W HF-radio med 
60\,\% verkningsgrad behöver ungefär 167\,W in -- vid 13{,}8\,V 
innebär det drygt 12\,A. Hennes nätaggregat var märkt 
\emph{``10\,A max''}.

Det var inte radion. Det var nätaggregatet som inte orkade 
leverera tillräcklig ström -- spänningen sjönk under sändning 
och radion drog ner effekten automatiskt.

Ett nytt 25\,A-aggregat senare visade wattmetern 98\,W. Nära nog.

Effektformeln var inte ett provfrågesvar. Det var verktyget som 
sparade Maria från att skicka tillbaka en fullt fungerande radio.
\end{storybox}

Effekt betecknas \textbf{P} och mäts i \textbf{Watt (W)}, uppkallat efter den skotske uppfinnaren James Watt (1736–1819). Energi betecknas \textbf{E} och mäts i \textbf{Wattimmar (Wh)} eller \textbf{Joule (J)}. Sambandet: $E = P \times t$ – en 100 W sändare som körs i 2 timmar förbrukar 200 Wh energi.

\begin{center}
\begin{tabular}{llll}
\rowcolor{greydark}
\tblhdr{Prefix} & \tblhdr{Symbol} & \tblhdr{Värde} & \tblhdr{Radioexempel} \\
Milliwatt & mW & 0,001 W & QRP-sändare (100 mW), laserpekare (5 mW) \\
\rowcolor{greylight}
Watt & W & 1 W & Handhållen VHF-radio på låg effekt \\
Kilowatt & kW & 1\,000 W & Kraftiga slutsteg – max 1 kW PEP med PTS-tillstånd \\
\end{tabular}
\end{center}

% ── 1.3.1 ─────────────────────────────────────────────────
\subsection{De tre effektformlerna}

Det finns tre sätt att beräkna effekt. Alla ger exakt samma svar, men du väljer den formel som passar de storheter du redan känner till.

\begin{formulabox}
{\large$\mathbf{P = U \times I}$}\\[3mm]
{\small\normalfont\itshape
$P = I^2 \times R$ \quad|\quad $P = U^2 \div R$}
\end{formulabox}

\begin{center}
\begin{tabular}{L{4.5cm}L{3cm}L{4cm}}
\rowcolor{greydark}
\tblhdr{Du känner till…} & \tblhdr{Formel} & \tblhdr{Snabbexempel} \\
Spänning (U) och ström (I) & $P = U \times I$ & $13{,}8 \times 7{,}2 = 99$ W \\
\rowcolor{greylight}
Ström (I) och resistans (R) & $P = I^2 \times R$ & $2^2 \times 50 = 200$ W \\
Spänning (U) och resistans (R) & $P = U^2 \div R$ & $100^2 \div 50 = 200$ W \\
\end{tabular}
\end{center}

\begin{examplebox}[title={\emoji{💡} Välj rätt formel – en enkel tumregel}]
Vet du spänning och ström? Använd $P = U \times I$. Saknar du spänningen men vet ström och resistans? Använd $P = I^2 \times R$. Saknar du strömmen men vet spänning och resistans? Använd $P = U^2 \div R$. Du behöver aldrig gissa – formelns variabler talar om vad du behöver mäta.
\end{examplebox}

\begin{redbox}[title={\emoji{⚠️} De tre vanligaste misstagen}]
\textbf{Glömmer konvertera mA → A.} $P = 12\,\text{V} \times 500\,\text{mA}$ ger inte 6\,000 W – konvertera till 0,5 A och få rätt svar: 6 W.

\textbf{Tror att $I^2 = I \times 2$.} $I^2$ betyder $I \times I$. Om $I = 5$ A är $I^2 = 25$, inte 10.

\textbf{Multiplicerar i stället för att dividera i $U^2/R$.} Formelns sista tecken är divisionstecken: $P = U^2$ \emph{delat med} $R$.
\end{redbox}

\begin{worked}[title={Räkneexempel 1.3.1-1 – HF-transceiver}]
En HF-radio drar 22 A från ett 13,8 V nätaggregat under sändning. Hur stor är totaleffekten?
\[ P = U \times I = 13{,}8\,\text{V} \times 22\,\text{A} = \mathbf{303{,}6\,\text{W} \approx 304\,\text{W}} \]
Notera: detta är totalförbrukning. Ut i antennen går kanske 100 W – resten blir värme i slutsteget.
\end{worked}

\begin{worked}[title={Räkneexempel 1.3.1-2 – Dummylast}]
En 50 $\Omega$ dummylast passeras av 2 A. Hur stor effekt måste den tåla?
\[ I^2 = 2 \times 2 = 4 \quad\rightarrow\quad P = I^2 \times R = 4 \times 50 = \mathbf{200\,\text{W}} \]
200 W kräver kraftig kylning – fläkt eller vätskekylning!
\end{worked}

\begin{worked}[title={Räkneexempel 1.3.1-3 – Antennresistans}]
En antenn med 50 $\Omega$ strålningsresistans matas med 22 V (RF). Hur mycket effekt strålas ut?
\[ U^2 = 22 \times 22 = 484 \quad\rightarrow\quad P = U^2 \div R = 484 \div 50 = \mathbf{9{,}7\,\text{W} \approx 10\,\text{W}} \]
\end{worked}

\subsubsection*{Övningsfrågor – 1.3.1}

\begin{qblock}{1.}
\qgrund\quad Förklara med egna ord vad effekt är och hur den skiljer sig från energi. Vilken enhet mäts effekt i och vad heter den uppkallad efter?
\end{qblock}

\begin{qblock}{2.}
\qgrund\quad Du har tre effektformler att välja mellan. I vilka situationer använder du respektive formel? Ge ett radioexempel för varje.
\end{qblock}

\begin{qblock}{3.}
\qrakna\quad Beräkna effekten. Ange vilken formel du valde och varför:
\begin{center}
\begin{tabular}{ll}
a) $U = 13{,}8$ V, $I = 15$ A &
b) $I = 500$ mA, $R = 8\,\Omega$ \\[4pt]
c) $U = 50$ V, $R = 50\,\Omega$ &
d) $U = 230$ V, $I = 4{,}35$ A
\end{tabular}
\end{center}
\end{qblock}

\begin{qblock}{F.}
\qfallan\quad En radioamatör ska välja dummylast. 
Han mäter strömmen genom sin 50\,$\Omega$-last till 2\,A 
och räknar:
\[
P = I^2 \times R = 2^2 \times 50 = 4 \times 50 = 200\,\text{W}
\]
Han väljer en dummylast märkt \emph{``150\,W''} i butiken --
beräkningen visade ju 200\,W och han antar att radion 
sällan kör på max.

Är beräkningen korrekt? Är valet av dummylast klokt? 
Vad bör han ha tänkt på?
\end{qblock}

\medskip
\noindent\rule{\textwidth}{0.4pt}
{\small\itshape Försök själv innan du läser svaret.}
\medskip

\begin{worked}[title={Lösning – Fallgropsfråga 1.3.1\,F}]
Beräkningen är korrekt: $P = 2^2 \times 50 = 200$\,W.

Valet av dummylast är däremot \textbf{inte} klokt. En komponent 
ska aldrig väljas exakt på -- eller under -- sin effektgräns. 
Välj alltid med 25--50\,\% säkerhetsmarginal för att hantera 
effekttoppar och temperaturstegring utan risk för överhettning.

Rätt val: minst \textbf{250\,W}, helst \textbf{300\,W}.

\textbf{Lärdomen:} Räkna rätt -- men dimensionera med marginal. 
Det gäller dummylaster, nätaggregat, kablar och resistorer. 
En komponent som arbetar vid sin maximala märkeffekt förkortar 
sin livslängd drastiskt och utgör en brandrisk.
\end{worked}

% ── 1.3.2 ─────────────────────────────────────────────────
\subsection{Decibel (dB)}

\lettrine[lines=2]{\textcolor{orange}{N}}{aturen} räknar sällan i jämna steg. Ett stearinljus lyser inte hälften så starkt som två -- det känns nästan lika svagt. En åska på tio kilometers avstånd låter inte tio gånger svagare än på en kilometer -- den låter knappt hörbar. Världen fungerar i \emph{gånger}, inte i steg.

Decibel är måttstocken som följer den logiken.

\begin{storybox}[{\emoji{🌧️} Regndropparna}]
Föreställ dig en helt stilla sjö en tidig morgon. 
Inget vind. Inget ljud. Bara vatten.

Du släpper en liten sten. En ring breder ut sig -- 
fin, rund, precis så stor som du förväntar dig.

Sedan släpper du en sten som är \emph{hundra gånger} tyngre.

Ringen blir inte hundra gånger större. Den blir kanske 
tio gånger större -- och det räcker gott för att kännas 
som en helt annan värld jämfört med den första.

Naturen räknar inte i steg. Den räknar i \emph{gånger}.
\end{storybox}

Precis samma sak händer i radioteknik -- och det skapar ett problem.

En känslig mottagare kan höra signaler på \textbf{0,000\,000\,001 W} (en miljardels watt). Ett kraftigt slutsteg levererar \textbf{1\,000 W}. Det är en skillnad på en miljard gånger.

Försök att skriva det på en axel. Försök att rita det i ett diagram. Det går inte -- antingen blir det lilla osynligt, eller får det stora inte plats.

Decibel löser det genom att byta måttstocken. Istället för att räkna steg räknar vi \emph{gånger}. Och istället för att skriva ut alla nollor skriver vi ett litet, hanterbart tal.

\begin{examplebox}[title={\emoji{💡} dB är alltid ett \emph{förhållande} -- inte ett absolut värde}]
+10\,dB betyder inte ``tio watt''. Det betyder ``\emph{tio gånger mer} än det vi jämför med''.

Fråga alltid: tio gånger mer än \emph{vad}?
\end{examplebox}

\subsubsection*{De tre tumreglerna du faktiskt behöver}

Det finns en hel tabell med dB-värden. Du behöver inte memorera den. 

Det räcker med tre siffror -- och de kan du lära dig på fem minuter:

\begin{center}
\begin{tabular}{lll}
\rowcolor{greydark}
\tblhdr{dB-värde} & \tblhdr{Effektförhållande} & \tblhdr{Kom ihåg det som\ldots} \\
+3 dB  & ×2   & dubbelt \\
\rowcolor{greylight}
+10 dB & ×10  & tio gånger \\
$-$3 dB  & ×0,5 & hälften \\
\end{tabular}
\end{center}

Det är grunden. Allt annat bygger på dessa tre.

\medskip

Och om du undrar varför just +3\,dB råkar bli dubbelt och +10\,dB tio gånger -- det är ingen slump, men du behöver inte förstå matematiken bakom för att använda det. Precis som du inte behöver kunna baka för att uppskatta en kanelbulle.

\subsubsection*{Bygga större värden -- addera steg}

Det fina med dB är att du kan \textbf{addera} dem för att kombinera steg. Sändaren ger 10\,W. Förstärkaren bidrar +13\,dB. Hur mycket effekt ut?

Dela upp: $+13\,\text{dB} = +10\,\text{dB} + 3\,\text{dB}$

Det betyder: tio gånger, sedan dubbelt.

$10\,\text{W} \times 10 \times 2 = \mathbf{200\,\text{W}}$

\medskip

Det fungerar alltid -- för varje gång du lägger till dB multiplicerar du effekten. Och när man multiplicerar tal kan man addera deras ``komprimerade'' version. Du behöver inte veta det -- det räcker att det \emph{funkar}.

\begin{examplebox}[title={\emoji{💡} Hela dB-tabellen -- för den nyfikne}]
Du behöver inte memorera detta. Men om du är nyfiken: varje steg från 1 till 10\,dB multiplicerar effekten med ungefär \textbf{1,26}. Precis som ränta på ränta -- 1,26 multiplicerat med sig självt tio gånger ger ×10. Tre steg ger ungefär ×2. Det är varför +3\,dB och +10\,dB råkar landa på runda tal.

\smallskip
\begin{center}
\begin{tabular}{ccc}
\rowcolor{greydark}
\tblhdr{dB} & \tblhdr{Effektförhållande} & \tblhdr{Ungefär} \\
1  & ×1,26 & lite mer \\
\rowcolor{greylight}
2  & ×1,58 & halvvägs mot dubbelt \\
\textbf{3}  & \textbf{×2,00} & \textbf{dubbelt} \\
\rowcolor{greylight}
4  & ×2,51 & drygt dubbelt \\
5  & ×3,16 & trippelt \\
\rowcolor{greylight}
\textbf{6}  & \textbf{×4,00} & \textbf{fyra gånger} \\
7  & ×5,01 & fem gånger \\
\rowcolor{greylight}
8  & ×6,31 & sex gånger \\
9  & ×7,94 & nära åtta gånger \\
\rowcolor{greylight}
\textbf{10} & \textbf{×10,00} & \textbf{tio gånger} \\
\end{tabular}
\end{center}
\end{examplebox}

\subsubsection*{Räkneexempel}

\begin{worked}[title={Räkneexempel 1.3.2-1 -- Kabelförlust}]
Sändaren ger 100\,W. Kabeln har 3\,dB förlust. Hur mycket når antennen?

$-3\,\text{dB} = \times 0{,}5 \;\rightarrow\; 100\,\text{W} \times 0{,}5 = \mathbf{50\,\text{W}}$

Hälften av effekten försvinner i kabeln som värme. Det är ett starkt skäl att hålla koaxkabeln kort.
\end{worked}

\begin{worked}[title={Räkneexempel 1.3.2-2 -- Förstärkare}]
En förstärkare har +10\,dB förstärkning. In går 5\,W. Hur mycket kommer ut?

$+10\,\text{dB} = \times 10 \;\rightarrow\; 5\,\text{W} \times 10 = \mathbf{50\,\text{W}}$
\end{worked}

\begin{worked}[title={Räkneexempel 1.3.2-3 -- Kombinera steg}]
Sändaren ger 10\,W. Förstärkaren bidrar +13\,dB. Hur mycket effekt ut?

Dela upp $+13\,\text{dB}$ i kända steg:
\[
+13\,\text{dB} = +10\,\text{dB} + 3\,\text{dB} = \times 10 \times 2 = \times 20
\]
\[
10\,\text{W} \times 20 = \mathbf{200\,\text{W}}
\]
\end{worked}

\subsubsection*{Övningsfrågor -- 1.3.2}

\begin{qblock}{4.}
\qrakna\quad Din sändare ger 100\,W. Koaxkabeln har 3\,dB förlust.
a) Hur många watt når antennen?
b) Du byter till en kabel med 1\,dB förlust -- ungefär hur många watt når antennen då? (Ledning: 1\,dB $\approx$ ×1,26 förstärkning, alltså $\div$1,26 förlust.)
c) Hur många dB antennförstärkning kompenserar exakt 3\,dB kabelförlust?
\end{qblock}

\begin{qblock}{5.}
\qprov\quad En radioamatör ökar från 100\,W till 400\,W. Hur många dB är ökningen? Kan han förvänta sig dubbelt så lång räckvidd? Motivera.
\end{qblock}

\begin{qblock}{F.}
\qfallan\quad En radioamatör läser att hans förstärkare har 
``20\,dB förstärkning''. Han matar in 1\,V och förväntar sig 
20\,V ut -- eftersom han vet att $+20\,\text{dB} = \times 100$ 
multiplicerar han med 100 och delar sedan med ingångsnivån:
\[
U_{ut} = 1\,\text{V} \times 100 = 100\,\text{V}
\]
Han mäter spänningen ut och får 10\,V. Är förstärkaren defekt?
\end{qblock}

\medskip
\noindent\rule{\textwidth}{0.4pt}
{\small\itshape Försök själv innan du läser svaret.}
\medskip

\begin{worked}[title={Lösning -- Fallgropsfråga 1.3.2\,F}]
Förstärkaren fungerar felfritt. Felet är att radioamatören 
blandade ihop \textbf{effekt-dB} och \textbf{spännings-dB} --
ett av de vanligaste misstagen i radioteknik.

Sambandet ser ut så här:
\begin{center}
\begin{tabular}{ll}
\rowcolor{greydark}
\tblhdr{Storhet} & \tblhdr{$+20$\,dB betyder} \\
Effekt & $\times 100$ \\
\rowcolor{greylight}
Spänning & $\times 10$ \\
\end{tabular}
\end{center}

Anledningen är egentligen ganska logisk när man tänker efter. 
Om du dubblar spänningen flödar också mer ström -- och mer 
ström med högre spänning ger \emph{fyra gånger} mer effekt, 
inte bara dubbelt. Spänning och ström ökar tillsammans. 
Det betyder att en tiofaldig spänningsökning ger hundrafaldig 
effektökning -- och därför räcker det med halva multiplikatorn 
när vi mäter spänning i dB jämfört med effekt.

Kontrollen: $1\,\text{V} \times 10 = 10\,\text{V}$ -- 
precis vad han mätte. \checkmark

\textbf{Lärdomen:} När någon anger dB -- fråga alltid: 
effekt-dB eller spännings-dB? I de flesta radiosammanhang 
avses effekt, men i förstärkarspecifikationer förekommer 
båda. Kontrollera alltid datablad noggrant.
\end{worked}

% ── 1.3.3 ─────────────────────────────────────────────────
\subsection{Verkningsgrad}

Ingen energiomvandling är perfekt – en del går alltid förlorad som värme. Verkningsgraden ($\eta$, uttalas ''eta'') anger hur stor andel av tillförd energi som blir nyttig effekt.

\begin{formulabox}
{\large$\mathbf{\eta = \dfrac{P_{ut}}{P_{in}} \times 100\,\%}$}\\[3mm]
{\small\normalfont\itshape $P_{ut}$ = nyttig effekt ut \quad|\quad $P_{in}$ = tillförd effekt in}
\end{formulabox}

\begin{worked}[title={Räkneexempel 1.3.3-1 – Slutstegets verkningsgrad}]
Slutsteget drar 300 W från nätaggregatet men levererar 200 W RF. Verkningsgrad?
\[ \eta = \frac{200}{300} \times 100\,\% = \mathbf{66{,}7\,\%} \]
De övriga 100 W blir värme – kylflänsen och fläkten måste klara det!
\end{worked}

Verkningsgraden i ett typiskt HF-slutsteg är 50–65\,\%. Det betyder att radion drar 180–200 W från nätaggregatet för att leverera 100 W RF ut i antennen. Den naiva beräkningen $I = P \div U = 100\,\text{W} \div 13{,}8\,\text{V} \approx 7{,}2\,\text{A}$ stämmer alltså \emph{inte} i praktiken – välj alltid ett 25–30 A aggregat med marginal.

Effektnivåer i Sverige regleras av PTS i två steg. Utan särskilt tillstånd gäller max \textbf{200 W PEP} på de flesta band. Med individuellt tillstånd från PTS är max \textbf{1\,000 W PEP}. PEP (\emph{Peak Envelope Power}) är den högsta momentana effekten i sändarsignalens kuverttopp. Kontrollera alltid \texttt{pts.se} och \texttt{ssa.se} för aktuella villkor.

\subsubsection*{Övningsfrågor – 1.3.3}

\begin{qblock}{6.}
\qrakna\quad En HF-transceiver levererar 100 W RF och drar 22 A från 13,8 V. a) Total ineffekt ($P_{in}$)? b) Verkningsgrad ($\eta$)? c) Hur stor effekt omvandlas till värme? d) Kylflänsen klarar 100 W utan fläkt – behövs fläkt?
\end{qblock}

\begin{qblock}{7.}
\qprov\quad En 470 $\Omega$ resistor är seriekopplad med en LED och ansluten till 12 V. Strömmen är 20 mA. a) Beräkna effekten i resistorn med $P = I^2 \times R$. b) Verifiera med $P = U_R \times I$ genom att först beräkna spänningsfallet. c) Vilken standardeffekt ($\frac{1}{4}$ W, $\frac{1}{2}$ W, 1 W) ska resistorn minst ha?
\end{qblock}

\begin{orangebox}[title={\emoji{🎯} Viktigt för provet}]
\begin{checklist}
  \item Tre effektformler: $P = U \times I$, $P = I^2 \times R$, $P = U^2 \div R$ – välj efter vad du vet
  \item $I^2 = I \times I$ (inte $I \times 2$!), $U^2 = U \times U$ (inte $U \times 2$!)
  \item +3 dB = dubbel effekt \quad|\quad +10 dB = tiofaldig \quad|\quad $-$3 dB = halva
  \item Verkningsgrad: $\eta = (P_{ut} \div P_{in}) \times 100\,\%$
  \item Konvertera \textbf{alltid} mA $\rightarrow$ A och mW $\rightarrow$ W innan beräkning
  \item Max 200 W PEP utan tillstånd – upp till 1\,000 W PEP med individuellt PTS-tillstånd
\end{checklist}
\end{orangebox}

% ── 1.4 ───────────────────────────────────────────────────
\section{Kondensatorer}

\lettrine[lines=2]{\textcolor{orange}{E}}{n} kondensator gör något som inget batteri klarar: den kan ladda upp energi långsamt och lämna ifrån sig allt på ett ögonblick.

\begin{storybox}[{\emoji{⚡} Blixten och kameran}]
Innan mobiltelefoner var blixtljuset i en engångskamera 
ett litet underverk.

Batteriet var svagt -- för svagt för att direkt blända 
ett rum. Men kameran hade en kondensator. Varje sekund 
du väntade pumpade batteriet in lite mer laddning. 
Kondensatorn fylldes, långsamt men säkert.

När du tryckte av -- på en tusendels sekund -- lämnade 
kondensatorn ifrån sig \emph{allt} på en gång. Ett 
bländande blixtsken. Sedan tystnad, och laddningen 
började byggas upp igen.
\end{storybox}

Det gör kondensatorn till en av de mest mångsidiga komponenterna i all elektronik -- och för att förstå hur den fungerar börjar vi med det mest grundläggande: vad den faktiskt lagrar, och hur mycket.

% ── 1.4.1 ─────────────────────────────────────────────────
\subsection{Kapacitans}

En kondensator består av två ledande plattor separerade av ett isolerande material -- \textbf{dielektrikum} (det är bara ett fint ord för ett isolerande material). Den lagrar elektrisk energi i ett elektriskt fält mellan plattorna.

\begin{examplebox}[title={\emoji{💡} Gummihinna-analogin}]
Tänk dig en gummihinna spänd tvärs över ett vattenrör.

\textbf{Likström (DC)} trycker vattnet åt ett håll -- hinnan töjer sig tills den inte ger med sig mer. Sedan stopp. Kondensatorn har laddats upp och strömmen upphör.

\textbf{Växelström (AC)} byter riktning hela tiden -- hinnan trycks fram och tillbaka och vibrerar i takt med signalen. Vattnet rör sig fram och tillbaka på \emph{sin} sida om hinnan, men inget vatten passerar faktiskt förbi den. Ändå överförs rörelsen -- tryckvågen går igenom. På samma sätt förflyttar sig inga elektroner genom kondensatorn, men signalenergin fortplantas ändå vidare. Ju högre frekvens, desto snabbare vibrerar hinnan och desto friare passerar signalen.
\end{examplebox}

Kapacitans betecknas \textbf{C} och mäts i \textbf{Farad (F)}, uppkallat efter den brittiske fysikern Michael Faraday. En farad är enormt stor -- i praktiken arbetar vi nästan alltid med pF, nF och µF.

\begin{center}
\begin{tabular}{llll}
\rowcolor{greydark}
\tblhdr{Prefix} & \tblhdr{Symbol} & \tblhdr{Värde} & \tblhdr{Används till} \\
Pikofarad & pF & $10^{-12}$ F & Mycket små kapacitanser -- precisionsfilter och finavstämning \\
\rowcolor{greylight}
Nanofarad & nF & $10^{-9}$ F & Koppling och filtrering i signalkretsar \\
Mikrofarad & µF & $10^{-6}$ F & Utjämning och laddning -- vanligast i vardagselektronik \\
\rowcolor{greylight}
Millifarad & mF & $10^{-3}$ F & Stora kraftkondensatorer i tunga nätaggregat \\
\end{tabular}
\end{center}

Nu vet vi hur mycket en kondensator kan lagra -- men hur beter den sig när signalen växlar?

\subsubsection*{Övningsfrågor -- 1.4.1}

\begin{qblock}{1.}
\qgrund\quad Förklara med egna ord varför en kondensator blockerar likström men inte växelström. Använd gummihinna-analogin i förklaringen.
\end{qblock}

% ── 1.4.2 ─────────────────────────────────────────────────
\subsection{Kapacitiv reaktans}

En kondensators ``motstånd'' mot växelström kallas \textbf{kapacitiv reaktans} och betecknas $X_C$. Till skillnad från en resistors motstånd -- som är detsamma oavsett frekvens -- beror reaktansen helt på hur snabbt signalen växlar.

\begin{formulabox}
{\large$\mathbf{X_C = \dfrac{1}{2\pi \times f \times C}}$}\\[3mm]
{\small\normalfont\itshape $X_C$ i ohm \quad|\quad $f$ i hertz \quad|\quad $C$ i farad \quad|\quad $2\pi \approx 6{,}28$}
\end{formulabox}

\begin{bluebox}[title={\emoji{🔑} Nyckelsamband -- reaktansen minskar när frekvensen ökar}]
Hög frekvens $\rightarrow$ låg $X_C$ $\rightarrow$ signalen passerar lätt.\\
Låg frekvens $\rightarrow$ hög $X_C$ $\rightarrow$ signalen bromsas.

\medskip
\textbf{Exempel:} En 1\,µF kondensator har $X_C \approx 3\,180\,\Omega$ vid 50\,Hz men bara $0{,}16\,\Omega$ vid 1\,MHz -- en faktor 20\,000 skillnad med samma komponent!
\end{bluebox}

Det betyder att samma kondensator beter sig olika beroende på vad du skickar igenom den. Vid låg frekvens är reaktansen hög -- lite ström passerar, som genom ett stort motstånd. Vid hög frekvens är reaktansen låg -- signalen passerar nästan fritt. Du känner igen logiken från Ohms lag: högre motstånd ger mindre ström.

\begin{worked}[title={Räkneexempel 1.4.2-1 -- Reaktans vid två frekvenser}]
Beräkna $X_C$ för en 100\,pF kondensator vid 7\,MHz och 144\,MHz.

$C = 100\,\text{pF} = 10^{-10}\,\text{F}$

\textbf{Vid 7 MHz} (40\,m-bandet):
\[ X_C = \frac{1}{6{,}28 \times 7 \times 10^6 \times 10^{-10}} = \frac{1}{4{,}396 \times 10^{-4}} \approx \mathbf{227\,\Omega} \]

\textbf{Vid 144 MHz} (2\,m-bandet):
\[ X_C = \frac{1}{6{,}28 \times 144 \times 10^6 \times 10^{-10}} = \frac{1}{9{,}043 \times 10^{-3}} \approx \mathbf{11\,\Omega} \]

Reaktansen sjönk från 227\,$\Omega$ till 11\,$\Omega$ -- kondensatorn släpper igenom 144\,MHz nästan fritt. En och samma komponent beter sig alltså som ett rejält motstånd på ett band och som en nästan öppen ledning på ett annat.
\end{worked}

\subsubsection*{Övningsfrågor -- 1.4.2}

\begin{qblock}{2.}
\qrakna\quad Beräkna $X_C$ för en 100\,pF kondensator vid: a) 7\,MHz \quad b) 144\,MHz. \quad c) Vad händer med reaktansen när frekvensen ökar? Motivera med formeln.
\end{qblock}

\begin{qblock}{3.}
\qrakna\quad En kopplingskondensator på 10\,µF används i audiovägen för att blockera DC men släppa igenom ljud (20\,Hz -- 20\,kHz). a) Beräkna $X_C$ vid 20\,Hz. b) Beräkna $X_C$ vid 1\,000\,Hz. c) Är 10\,µF ett bra val? Motivera med hjälp av reaktansvärdena.
\end{qblock}

% ── Fälla 1.4.2 ───────────────────────────────────────────
\begin{qblock}{F.}
\qfallan\quad En radioamatör beräknar kapacitiv reaktans för
en 100\,µF kondensator vid 50\,Hz. Han skriver:
\[
X_C = \frac{1}{6{,}28 \times 50 \times 100} = \frac{1}{31\,400} \approx 0{,}000032\,\Omega
\]
Han blir förvånad -- en kondensator på 100\,µF borde väl
ha mer motstånd än så vid nätfrekvens? Vad gick fel?
\end{qblock}

\medskip
\noindent\rule{\textwidth}{0.4pt}
{\small\itshape Försök själv innan du läser svaret.}
\medskip

\begin{worked}[title={Lösning -- Fallgropsfråga 1.4.2\,F}]
Kondensatorn är inte fel -- enheten är det. Han stoppade
in µF direkt i formeln i stället för att konvertera till
Farad först.

Rätt omvandling: $100\,\mu\text{F} = 100 \times 10^{-6}\,\text{F} = 10^{-4}\,\text{F}$

Rätt beräkning:
\[
X_C = \frac{1}{6{,}28 \times 50 \times 10^{-4}} = \frac{1}{0{,}0314} \approx \mathbf{31{,}8\,\Omega}
\]
Det är en miljon gånger större än hans svar -- och betydligt
mer rimligt för en kondensator i ett nätaggregatsfilter.

\textbf{Lärdomen:} Formeln $X_C = 1/(2\pi fC)$ kräver att
$C$ anges i Farad. Alltid. Ett svar på $0{,}000032\,\Omega$
borde omedelbart väcka misstankar -- en kondensator med
lägre motstånd än en koppartråd är ett tydligt tecken på
fel enhet.
\end{worked}


% ── 1.4.3 ─────────────────────────────────────────────────
\subsection{Kondensatortyper}

Kondensatorer tillverkas i flera olika utföranden -- varje typ har sina styrkor och passar bäst i vissa situationer. Valet beror på hur mycket kapacitans du behöver, hur noggrann den måste vara och om den sitter i en AC- eller DC-krets.

% ── KERAMISK ──────────────────────────────────────────────
\subsubsection*{Keramisk kondensator}

\begin{storybox}[{\emoji{☕} Alltid på plats}]
Tänk på rökdetektorn i taket. Den sitter där år efter år, 
kostar inte mycket, kräver ingen uppmärksamhet och gör 
jobbet utan att du tänker på den. Det är inte den 
spännande komponenten i rummet -- men den ska vara där.

Den keramiska kondensatorn fyller samma roll i en krets. 
Liten, billig och stabil. Den reagerar snabbt på 
störningar och leder bort dem innan de ställer till problem.
\end{storybox}

Keramiska kondensatorer finns från 1\,pF upp till 10\,µF. 
De saknar polaritet, klarar höga frekvenser utan problem 
och åldras knappt alls. Det är standardvalet när du vill 
stabilisera en matningsspänning eller filtrera bort 
högfrekventa störningar.

% ── FILM ──────────────────────────────────────────────────
\subsubsection*{Filmkondensator}

\begin{storybox}[{\emoji{⚗️} Apotekaren}]
En apotekare väger inte upp ett läkemedel med ögonmåttet. 
Dosen måste stämma på milligrammet -- för lite hjälper 
inte, för mycket kan skada. Vågen används varje gång, 
utan undantag.

Filmkondensatorn är byggd med samma inställning till 
precision. Kapacitansvärdet stämmer med vad det står på 
höljet, förlusterna är låga och egenskaperna håller sig 
stabila över tid.
\end{storybox}

Filmkondensatorer (vanligtvis polyester eller polypropen) 
arbetar i intervallet 100\,pF till 10\,µF. De saknar 
polaritet och är rätt val när noggrannhet väger tyngre 
än pris eller storlek.

% ── ELEKTROLYT ────────────────────────────────────────────
\subsubsection*{Elektrolytkondensator}

\begin{storybox}[{\emoji{🔋} Stora tanken}]
En dieselgenerator har en stor bränsletank -- den kan 
lagra mycket och hålla igång länge. Men tanken har ett 
påfyllningsrör och ett utloppsrör, och de sitter inte 
utbytbara. Fyller du på från fel håll händer ingenting bra.

Elektrolytkondensatorn lagrar mycket laddning -- upp till 
10\,000\,µF -- och används när kretsen behöver en rejäl 
buffert. Men precis som dieseltanken har den en bestämd 
riktning. Sätts den baklänges går det illa.
\end{storybox}

\begin{redbox}[title={\emoji{⚠️} Kritisk polaritet}]
Elektrolytkondensatorer är \textbf{polariserade} -- de har 
en definierad plus- och minussida. Ansluts de baklänges 
startar en kemisk reaktion inuti som ger värme och gas. 
I värsta fall slutar det med en kraftig smäll och 
kemikaliespill.

\medskip
\textbf{Så här känner du igen polerna:}
\begin{itemize}
  \item Grå rand med ${-}$-tecken på höljet = \textbf{minuspolen}
  \item På nya komponenter: \textbf{kortare ben = minus}, längre ben = plus
  \item Mät spänningspolariteten i kretsen med multimeter 
        \emph{innan} du sätter i komponenten
\end{itemize}

\medskip
En snabb koll med multimetern tar tio sekunder. Det är 
värt det.
\end{redbox}

% ── VARIABEL ──────────────────────────────────────────────
\subsubsection*{Variabel kondensator}

\begin{storybox}[{\emoji{🎚️} Dimmern}]
En vanlig strömbrytare är på eller av. En dimmer låter 
dig ställa in precis hur mycket ljus du vill ha.

Den variabla kondensatorn fungerar på samma sätt -- 
kapacitansvärdet kan justeras, antingen med en skruv 
(trimmer) eller med en ratt (vridkondensator). När du 
senare lär dig om kretsar som väljer ut en specifik 
frekvens är det ofta en variabel kondensator som gör 
den finjusteringen möjlig.
\end{storybox}

En trimmer justeras en gång vid inbyggnad och lämnas 
sedan ifred. En vridkondensator sitter på en panel och 
används aktivt. Båda arbetar i området 5--500\,pF.

% ── SNABBREFERENS ─────────────────────────────────────────
\medskip
\noindent\textbf{Snabbreferens -- kondensatortyper}

\begin{center}
\begin{tabular}{L{2.4cm}L{2.2cm}L{4.2cm}L{3cm}}
\rowcolor{greydark}
\tblhdr{Typ} & \tblhdr{Värden} & \tblhdr{Kännetecken} & \tblhdr{Typisk användning} \\
Keramisk & 1\,pF -- 10\,µF & Snabb, stabil, ingen polaritet & Störningsdämpning, avkoppling \\
\rowcolor{greylight}
Film & 100\,pF -- 10\,µF & Noggrann, låga förluster, ingen polaritet & Filter, precisionskretsar \\
Elektrolyt & 1\,µF -- 10\,000\,µF & Stor kapacitans, \textbf{polariserad!} & Buffert, utjämning \\
\rowcolor{greylight}
Variabel (trimmer) & 5 -- 500\,pF & Justerbar med skruv & Finavstämning \\
Variabel (vrid) & 10 -- 500\,pF & Justerbar med ratt & Aktiv inställning \\
\end{tabular}
\end{center}

% ── BYPASS ────────────────────────────────────────────────
\begin{bluebox}[title={\emoji{🛡️} Bypass-kondensatorn}]
En \textbf{bypass-kondensator} kopplas mellan 
matningsspänningen och referenspunkten i kretsen. 
Uppgiften är att leda bort högfrekventa störningar och 
snabba spänningsvariationer innan de når känsliga delar 
av kretsen.

\medskip
Varför kombinerar man ofta en liten \textbf{keramisk 
(100\,nF)} med en stor \textbf{elektrolyt (100\,µF)}? 
Du vet redan svaret från 1.4.2: den keramiska har låg 
reaktans vid höga frekvenser och reagerar snabbt på 
snabba störningar. Elektrolyten har stor kapacitans och 
tar hand om långsammare variationer. De täcker var sin 
del av problemet.

\medskip
Det är därför du nästan alltid ser de två sida vid sida 
på ett kretskort.
\end{bluebox}

Nu vet vi vilka typer som finns och vad som skiljer dem åt. 
Men hur snabbt hinner en kondensator ladda upp -- och vad 
avgör det?

% ── ÖVNINGSFRÅGOR 1.4.3 ───────────────────────────────────
\subsubsection*{Övningsfrågor -- 1.4.3}

\begin{qblock}{4.}
\qgrund\quad Vilken kondensatortyp har polaritet och varför 
är det kritiskt att ansluta den rätt? Beskriv \emph{två} 
sätt att identifiera plus- och minussidan på en ny 
komponent.
\end{qblock}

\begin{qblock}{5.}
\qprov\quad Förklara vad en bypass-kondensator gör i en 
elektronisk krets. Varför placeras ofta en liten keramisk 
(100\,nF) parallellt med en stor elektrolyt (100\,µF)?
Koppla ditt svar till det du lärde dig om kapacitiv 
reaktans i 1.4.2.
\end{qblock}

\begin{qblock}{6.}
\qgrund\quad Du ska välja kondensator till tre olika 
situationer. Vilken typ passar bäst och varför?
\begin{enumerate}[label=\alph*)]
  \item Du vill buffra spänningen i ett nätaggregat så 
        att den håller sig stabil vid varierande belastning.
  \item Du vill dämpa snabba störningar precis intill 
        ett känsligt chip på ett kretskort.
  \item Du vill kunna justera kapacitansen med ett 
        skruvmejselstag vid inbyggnad.
\end{enumerate}
\end{qblock}

% ── FÄLLA 1.4.3 ───────────────────────────────────────────
\begin{qblock}{F.}
\qfallan\quad Lars ska byta en uttjänt elektrolytkondensator 
i sitt nätaggregat. Han hittar en ny med rätt kapacitans och 
spänningshållfasthet i lådan, sätter i den snabbt och slår 
på strömmen. Efter några sekunder luktar det konstigt och 
kondensatorn börjar bli varm. Han stänger av direkt.

Vad hände -- och hur hade det kunnat gå om han inte stängt 
av?
\end{qblock}

\medskip
\noindent\rule{\textwidth}{0.4pt}
{\small\itshape Försök själv innan du läser svaret.}
\medskip

\begin{worked}[title={Lösning -- Fallgropsfråga 1.4.3\,F}]
Lars satte i kondensatorn baklänges. Elektrolytkondensatorer 
är polariserade -- plus och minus måste sitta rätt, annars 
startar en kemisk reaktion inuti som ger värme och gas. 
Hade han inte stängt av: överhettning, kemikaliespill och 
i värsta fall en smäll.

Att han stängde av direkt begränsade skadan.

\medskip
\textbf{Rätt rutin innan montering:}
\begin{itemize}
  \item Grå rand med ${-}$ på höljet = minuspolen
  \item Kortare ben = minus på nya komponenter
  \item Mät spänningspolariteten i kretsen med multimeter 
        \emph{innan} komponenten sätts i -- varje gång
\end{itemize}

\textbf{Lärdomen:} En kondensator som passar i hålet sitter 
inte automatiskt rätt. Tio sekunders kontroll sparar 
komponenten -- och kanske dig själv.
\end{worked}

% ── 1.4.4 ─────────────────────────────────────────────────
\subsection{Tidskonstant}

En kondensator laddas inte omedelbart -- hastigheten beror 
på både kapacitansen och resistansen i kretsen. Sambandet 
beskrivs av \textbf{tidskonstanten} $\tau$ (grekiska 
bokstaven tau).

\begin{formulabox}
{\large$\mathbf{\tau = R \times C}$}\\[3mm]
{\small\normalfont\itshape $\tau$ i sekunder \quad|\quad $R$ i ohm \quad|\quad $C$ i farad}
\end{formulabox}

\begin{storybox}[{\emoji{🚲} Cykeldäcket}]
Pumpa upp ett cykeldäck för hand. De första tagen går 
lätt -- däcket är nästan tomt och luften strömmar in 
utan motstånd. Men ju högre trycket blir, desto segare 
går det. De sista baren kräver dubbelt så mycket kraft 
som de första, och mot slutet kämpar du för varje 
pumptagg.

Kondensatorn laddas på samma sätt. I början flödar 
strömmen fritt -- men ju mer laddning som byggs upp på 
plattorna, desto mer motverkar den ny laddning. Kurvan 
planar ut. I teorin når kondensatorn aldrig exakt full 
laddning -- i praktiken räknar vi den som fulladdat 
efter $5\tau$.
\end{storybox}

Siffrorna i tabellen nedan är inte godtyckliga -- de följer 
av hur uppladdningskurvan alltid ser ut, oavsett vilka 
värden R och C har.

\begin{center}
\begin{tabular}{lll}
\rowcolor{greydark}
\tblhdr{Tid} & \tblhdr{Laddning} & \tblhdr{Urladdning} \\
$1\tau$ & 63\,\% & 37\,\% kvar \\
\rowcolor{greylight}
$2\tau$ & 86\,\% & 14\,\% kvar \\
$3\tau$ & 95\,\% &  5\,\% kvar \\
\rowcolor{greylight}
$5\tau$ & 99\,\% \textbf{(praktiskt fulladdat)} & 1\,\% kvar \textbf{(praktiskt urladdat)} \\
\end{tabular}
\end{center}

\begin{worked}[title={Räkneexempel 1.4.4-1 -- Avkopplingskondensator}]
En 100\,µF kondensator laddas genom en 1\,k$\Omega$ 
resistor. Hur lång tid tar det att ladda den fullständigt?

\[ \tau = R \times C = 1\,000\,\Omega \times 100 \times 10^{-6}\,\text{F} = 0{,}1\,\text{s} \]

Fullt laddad efter $5\tau = 5 \times 0{,}1 = \mathbf{0{,}5\,\text{s}}$.

Det är en halv sekund -- ungefär den fördröjning du märker om du håller 
på/av-knappen och undrar varför ingenting händer direkt.
\end{worked}

\begin{worked}[title={Räkneexempel 1.4.4-2 -- Bypass-kondensator}]
En 100\,nF bypass-kondensator med 50\,$\Omega$ 
källresistans. Hur snabbt reagerar den?

\[ \tau = 50\,\Omega \times 100 \times 10^{-9}\,\text{F} = 5\,\mu\text{s} \]

Tusen gånger snabbare än föregående exempel -- det är 
därför små keramiska kondensatorer reagerar snabbt på 
störningar. En störning som varar en mikrosekund är redan 
borta innan bypass-kondensatorn ens hunnit ladda upp till $1\tau$.
\end{worked}

\begin{bluebox}[title={\emoji{📖} Serie- och parallellkoppling av kondensatorer}]
Kondensatorer följer \textbf{omvända regler} mot resistorer 
vid serie- och parallellkoppling -- kondensatorer i serie 
ger \emph{lägre} kapacitans, parallellt ger \emph{högre}. 
Detta behandlas i avsnitt~1.6 där reglerna jämförs sida 
vid sida.
\end{bluebox}

% ── ÖVNINGSFRÅGOR 1.4.4 ───────────────────────────────────
\subsubsection*{Övningsfrågor -- 1.4.4}

\begin{qblock}{7.}
\qrakna\quad En 47\,µF kondensator laddas genom en 
470\,$\Omega$ resistor. 
a) Beräkna tidskonstanten $\tau$. 
b) Hur lång tid tar det innan kondensatorn är praktiskt 
fulladdat? 
c) Hur stor andel av laddningen finns kvar efter $2\tau$?
\end{qblock}

% ── FÄLLA 1.4.4 ───────────────────────────────────────────
\begin{qblock}{F.}
\qfallan\quad En radioamatör beräknar tidskonstanten för 
en krets med en 470\,$\Omega$ resistor och en 100\,µF 
kondensator. Han skriver:
\[
\tau = R \times C = 470 \times 100 = 47\,000\,\text{s}
\]
Han blir förvirrad -- det kan väl inte ta 47\,000 sekunder 
att ladda en liten kondensator? Han har rätt att vara 
förvirrad. Vad gick fel?
\end{qblock}

\medskip
\noindent\rule{\textwidth}{0.4pt}
{\small\itshape Försök själv innan du läser svaret.}
\medskip

\begin{worked}[title={Lösning -- Fallgropsfråga 1.4.4\,F}]
Samma fälla som i XC-beräkningen -- han glömde konvertera 
µF till Farad.

Rätt omvandling: $100\,\mu\text{F} = 10^{-4}\,\text{F}$

Rätt beräkning:
\[
\tau = 470\,\Omega \times 10^{-4}\,\text{F} = 0{,}047\,\text{s} = \mathbf{47\,\text{ms}}
\]

Fullt laddad efter $5\tau = 5 \times 47\,\text{ms} = 235\,\text{ms}$ -- 
ungefär en fjärdedels sekund. Rimligt.

\textbf{Lärdomen:} Formeln $\tau = R \times C$ kräver $C$ 
i Farad. Alltid. Ett svar på tolv timmar för en liten 
kondensator är ett solklart tecken på fel enhet -- gör 
alltid en rimlighetscheck.
\end{worked}

\begin{orangebox}[title={\emoji{🎯} Viktigt för provet}]
\begin{checklist}
  \item Kapacitans mäts i \textbf{Farad (F)} -- i praktiken pF, nF, µF
  \item Blockerar DC, släpper igenom AC -- mer vid högre frekvens
  \item $X_C = 1/(2\pi fC)$ -- reaktansen \textbf{minskar} vid högre frekvens
  \item Elektrolyter har polaritet -- fel anslutning kan ge explosion
  \item Serie: $C$ minskar \quad|\quad Parallell: $C$ adderas (se avsnitt 1.6)
  \item Tidskonstant: $\tau = R \times C$ -- efter $5\tau$ är kondensatorn praktiskt fulladdat
\end{checklist}
\end{orangebox}

% ── 1.5 ───────────────────────────────────────────────────
\section{Spolar och induktans}

\lettrine[lines=2]{\textcolor{orange}{E}}{n} spole gör precis det som kondensatorn inte gör -- och tvärtom. De är varandras fullständiga motsatser, och när du förstår den ena förstår du egentligen båda.

\begin{storybox}[{\emoji{🌀} Tråden som minns}]
Föreställ dig att du drar en lång tråd runt en penna, varv 
efter varv. Inget märkvärdigt -- det är bara en tråd.

Men när du skickar ström genom den skapas ett magnetfält 
runt varje varv. De slår ihop till ett gemensamt fält längs 
hela spolen. Stänger du av strömmen försöker fältet hålla 
strömmen vid liv lite till -- det tar tid innan det klingar 
av.

Det är det som skiljer spolen från en vanlig ledare.
\end{storybox}

Det är den egenskapen som gör spolen användbar -- för att 
filtrera, bromsa och forma signaler.

% ── 1.5.1 ─────────────────────────────────────────────────
\subsection{Induktans}

En spole (induktor) är en lindad tråd. När ström flödar 
genom den skapas ett magnetfält som lagrar energi -- 
precis som kondensatorn lagrar energi i ett elektriskt 
fält, men med exakt omvända egenskaper mot ström och 
frekvens.

\begin{examplebox}[title={\emoji{💡} Det tunga vattenhjulet}]
Föreställ dig ett tungt vattenhjul i ett rör.

\textbf{Likström (DC)} startar hjulet och det snurrar 
sedan jämnt -- strömmen flödar fritt utan motstånd 
efter uppstarten.

\textbf{Växelström (AC)} vill byta riktning hela tiden, 
men det tunga hjulet motstår varje riktningsbyte. Det 
måste bromsas, stanna och accelerera åt andra hållet. 
Ju högre frekvens -- fler riktningsbyten per sekund -- 
desto mer motstånd bjuder hjulet.
\end{examplebox}

Induktans betecknas \textbf{L} och mäts i \textbf{Henry 
(H)}, uppkallat efter den amerikanske fysikern Joseph 
Henry. En henry är stor -- i praktiken arbetar vi med 
nH, µH och mH.

\begin{center}
\begin{tabular}{llll}
\rowcolor{greydark}
\tblhdr{Prefix} & \tblhdr{Symbol} & \tblhdr{Värde} & \tblhdr{Används till} \\
Nanohenry & nH & $10^{-9}$ H & Mycket små spolar vid höga frekvenser \\
\rowcolor{greylight}
Mikrohenry & µH & $10^{-6}$ H & Vanligast i filter och svängningskretsar \\
Millihenry & mH & $10^{-3}$ H & Audiofilter och större choker \\
\rowcolor{greylight}
Henry & H & $1$ H & Nätaggregat och transformatorer \\
\end{tabular}
\end{center}

Flera faktorer påverkar hur stor induktansen blir:

\begin{center}
\begin{tabular}{ll}
\rowcolor{greydark}
\tblhdr{Förändring} & \tblhdr{Effekt på L} \\
Fler varv & Högre L -- kvadratiskt samband: dubbla varven ger fyra gånger L \\
\rowcolor{greylight}
Järn- eller ferritkärna & Mycket högre L -- 10 till 1\,000 gånger luftkärna \\
Större diameter & Högre L \\
\rowcolor{greylight}
Längre spole (samma antal varv) & Lägre L \\
\end{tabular}
\end{center}

Det kvadratiska sambandet för varv är värt att stanna vid: dubbla varven ger \emph{fyra} gånger induktansen -- inte dubbelt. Tredubbla varven ger nio gånger. Det är en vanlig tankefälla. Järn- och ferritkärnor fungerar som förstärkare för magnetfältet -- de leder det bättre än luft, precis som en ledare leder ström bättre än gummi.

Nu vet vi hur stor en spoles induktans kan vara -- men 
hur beter den sig när signalen växlar?

\subsubsection*{Övningsfrågor -- 1.5.1}

\begin{qblock}{1.}
\qgrund\quad Förklara med vattenhjuls-analogin varför en 
spole bromsar växelström men inte likström. Varför ökar 
motståndet när frekvensen ökar?
\end{qblock}

% ── Fälla 1.5.1 ───────────────────────────────────────────
\begin{qblock}{F.}
\qfallan\quad En radioamatör lindade en spole med 10 varv 
och mätte induktansen till 5\,µH. Han behöver fyra gånger 
så stor induktans. Han lägger till 10 varv till -- totalt 
20 varv -- och är säker på att han nu har 20\,µH. 

Stämmer det?
\end{qblock}

\medskip
\noindent\rule{\textwidth}{0.4pt}
{\small\itshape Försök själv innan du läser svaret.}
\medskip

\begin{worked}[title={Lösning -- Fallgropsfråga 1.5.1\,F}]
Ja -- men av fel anledning. Han trodde att dubbla varven 
ger dubbel induktans. Det stämmer inte.

Sambandet är \textbf{kvadratiskt}: induktansen är 
proportionell mot antalet varv i \emph{kvadrat}.

\begin{center}
\begin{tabular}{lll}
\rowcolor{greydark}
\tblhdr{Varv} & \tblhdr{Faktor} & \tblhdr{Induktans} \\
10 & $1^2 = 1$ & 5\,µH \\
\rowcolor{greylight}
20 & $2^2 = 4$ & 20\,µH \\
30 & $3^2 = 9$ & 45\,µH \\
\end{tabular}
\end{center}

Han fick rätt svar av fel anledning. Nästa gång -- när 
han behöver tredubbla induktansen -- räknar han tre 
gånger fler varv och undrar varför han får nio gånger 
mer induktans istället.

\textbf{Lärdomen:} Dubbla varven ger fyra gånger 
induktansen. Vill du ha exakt dubbel induktans behöver 
du $\sqrt{2} \approx 1{,}41$ gånger fler varv -- 
inte dubbelt.
\end{worked}

% ── 1.5.2 ─────────────────────────────────────────────────
\subsection{Induktiv reaktans}

Spolens ``motstånd'' mot växelström kallas \textbf{induktiv 
reaktans} och betecknas $X_L$. Till skillnad från 
kondensatorns reaktans -- som minskar när frekvensen ökar 
-- ökar spolens reaktans med frekvensen. De är varandras 
motsatser, precis som vattenhjulet och gummihinnan.

\begin{formulabox}
{\large$\mathbf{X_L = 2\pi \times f \times L}$}\\[3mm]
{\small\normalfont\itshape $X_L$ i ohm \quad|\quad $f$ i hertz \quad|\quad $L$ i henry \quad|\quad $2\pi \approx 6{,}28$}
\end{formulabox}

\begin{bluebox}[title={\emoji{🔑} Nyckelsamband -- reaktansen ökar när frekvensen ökar}]
Hög frekvens $\rightarrow$ hög $X_L$ $\rightarrow$ signalen bromsas.\\
Låg frekvens $\rightarrow$ låg $X_L$ $\rightarrow$ signalen passerar lätt.

\medskip
\textbf{Exempel:} En 100\,µH spole har $X_L \approx 
0{,}031\,\Omega$ vid 50\,Hz men hela $90\,500\,\Omega$ 
vid 144\,MHz -- en faktor tre miljoner skillnad med 
samma komponent!
\end{bluebox}

Det är precis tvärtom mot kondensatorn. Du känner igen 
logiken från Ohms lag: högre motstånd ger mindre ström. 
Spolen och kondensatorn kompletterar varandra -- det ena 
bromsar vad det andra släpper igenom. Det kommer vi 
tillbaka till när vi tittar på filter och resonanskretsar.

\begin{worked}[title={Räkneexempel 1.5.2-1 -- Reaktans vid två frekvenser}]
En 100\,µH spole: hur stor är reaktansen vid 50\,Hz 
respektive 144\,MHz?

$L = 100\,\mu\text{H} = 10^{-4}\,\text{H}$

\textbf{Vid 50 Hz:}
\[ X_L = 6{,}28 \times 50 \times 10^{-4} \approx \mathbf{0{,}031\,\Omega} \]

\textbf{Vid 144 MHz:}
\[ X_L = 6{,}28 \times 144 \times 10^6 \times 10^{-4} \approx \mathbf{90\,500\,\Omega} \]

Faktor tre miljoner skillnad -- samma spole är helt 
transparent för nätfrekvens men ett kraftigt hinder 
vid höga frekvenser.
\end{worked}

Nu känner vi reaktansen -- nästa steg är att se hur 
spolen och kondensatorn förhåller sig till varandra, 
sida vid sida.

\subsubsection*{Övningsfrågor -- 1.5.2}

\begin{qblock}{2.}
\qrakna\quad Beräkna $X_L$ för en 10\,µH spole vid: 
a) 3,5\,MHz \quad b) 14\,MHz \quad c) 144\,MHz. \quad 
d) Vad händer med $X_L$ när frekvensen ökar?
\end{qblock}

\begin{qblock}{3.}
\qrakna\quad Du behöver en spole med minst 500\,$\Omega$ 
reaktans vid 7\,MHz. Hur stor induktans krävs?
\end{qblock}

% ── Fälla 1.5.2 ───────────────────────────────────────────
\begin{qblock}{F.}
\qfallan\quad En radioamatör beräknar reaktansen för en 
10\,µH spole vid 14\,MHz. Han skriver:
\[
X_L = 6{,}28 \times 14 \times 10^6 \times 10 = 879\,000\,\Omega
\]
Han blir förvånad -- nästan en megaohm för en liten 
spole? Det borde väl inte stämma. Han har rätt att 
vara förvånad. Vad gick fel?
\end{qblock}

\medskip
\noindent\rule{\textwidth}{0.4pt}
{\small\itshape Försök själv innan du läser svaret.}
\medskip

\begin{worked}[title={Lösning -- Fallgropsfråga 1.5.2\,F}]
Samma klassiska fälla som i kondensatorberäkningarna 
-- han stoppade in µH direkt i formeln i stället för 
att konvertera till Henry.

Rätt omvandling: $10\,\mu\text{H} = 10 \times 10^{-6}\,\text{H} = 10^{-5}\,\text{H}$

Rätt beräkning:
\[ X_L = 6{,}28 \times 14 \times 10^6 \times 10^{-5} \approx \mathbf{879\,\Omega} \]

Det är tusen gånger mindre -- och ett rimligt värde 
för en spole på det bandet.

\textbf{Lärdomen:} Formeln $X_L = 2\pi fL$ kräver 
$L$ i Henry. Alltid. Gör alltid en rimlighetscheck.
\end{worked}

% ── 1.5.3 ─────────────────────────────────────────────────
\subsection{Kondensator vs Spole -- fullständiga motsatser}

Nu när vi känner både kondensatorn och spolen är det 
dags att sätta dem sida vid sida. Tråden som minns och 
gummihinnan som vibrerar -- de är varandras spegelbild i 
nästan allt, och att hålla isär dem är en av de 
viktigaste sakerna att lära sig i grundläggande 
radioteknik.

\begin{center}
\begin{tabular}{L{3.2cm}L{3.5cm}L{3.5cm}}
\rowcolor{greydark}
\tblhdr{Egenskap} & \tblhdr{Kondensator (C)} & \tblhdr{Spole (L)} \\
Likström (DC) & Blockerar & Släpper igenom \\
\rowcolor{greylight}
Växelström (AC) & Släpper igenom & Bromsar \\
Hög frekvens & Låg reaktans (lätt) & Hög reaktans (svårt) \\
\rowcolor{greylight}
Frekvens ökar & Reaktans minskar $\downarrow$ & Reaktans ökar $\uparrow$ \\
Lagrar energi i & Elektriskt fält & Magnetfält \\
\rowcolor{greylight}
Reaktansformel & $X_C = 1/(2\pi fC)$ & $X_L = 2\pi fL$ \\
Tidskonstant & $\tau = R \times C$ & $\tau = L / R$ \\
\end{tabular}
\end{center}

Lär dig den här tabellen. Den dyker upp igen i filter, 
impedans och resonanskretsar -- och varje gång är det 
samma mönster.

Nu när vi vet hur spolen beter sig är det dags att se 
vilka typer som finns -- och vilken som passar till vad.

\subsubsection*{Övningsfrågor -- 1.5.3}

\begin{qblock}{4.}
\qgrund\quad Fyll i tabellen utan att titta i texten:
\begin{center}
\begin{tabular}{L{3.5cm}L{3cm}L{3cm}}
\rowcolor{greydark}
\tblhdr{Egenskap} & \tblhdr{Kondensator} & \tblhdr{Spole} \\
Blockerar DC/AC? & & \\
\rowcolor{greylight}
Reaktans vid hög f? & & \\
Reaktansformel & & \\
\rowcolor{greylight}
Lagrar energi i & & \\
\end{tabular}
\end{center}
\end{qblock}

% ── 1.5.4 ─────────────────────────────────────────────────
\subsection{Spoltyper}

Precis som kondensatorer finns spolar i flera utföranden 
-- valet beror på frekvensområde, hur hög induktans som 
behövs och hur stora förlusterna får vara.

% ── LUFTKÄRNA ─────────────────────────────────────────────
\subsubsection*{Luftkärnespole}

\begin{storybox}[{\emoji{🎯} Inga genvägar}]
En våg som mäter exakt. Inga justeringar, inga 
korrigeringsfaktorer -- det som visas är det som är.

Luftkärnespolen fungerar på samma sätt. Utan kärna 
som påverkar magnetfältet är egenskaperna stabila och 
förutsägbara. Det den ger är det den ger.
\end{storybox}

Luftkärnespolar används där låga förluster och stabil 
induktans väger tyngre än kompakthet -- typiskt i 
precisionsfilter och svängningskretsar.

% ── FERRITKÄRNA ───────────────────────────────────────────
\subsubsection*{Ferritkärnespole (toroid)}

Samma antal varv, samma ström -- men stoppar du in en 
ferritkärna kan induktansen bli hundra gånger högre. 
Kärnan förstärker magnetfältet utan att du behöver 
linda ett enda extra varv.

\begin{figure}[h]
\centering
\begin{tikzpicture}[font=\small]

  % ── Vänster: luftkärna ──────────────────────────────────
  \draw[thick] (0,0) -- (0.6,0);
  \draw[thick, color=red!70!black] 
    (0.6,0) .. controls (0.6,0.5) and (1.0,0.5) .. (1.0,0)
            .. controls (1.0,-0.5) and (1.4,-0.5) .. (1.4,0)
            .. controls (1.4,0.5) and (1.8,0.5) .. (1.8,0)
            .. controls (1.8,-0.5) and (2.2,-0.5) .. (2.2,0)
            .. controls (2.2,0.5) and (2.6,0.5) .. (2.6,0)
            .. controls (2.6,-0.5) and (3.0,-0.5) .. (3.0,0);
  \draw[thick] (3.0,0) -- (3.6,0);
  \node[gray, font=\itshape\scriptsize] at (1.8,0) {luft};
  \node[font=\footnotesize] at (1.8,-1.0) {10 varv};
  \node[font=\large\bfseries, color=red!70!black] at (1.8,-1.5) {5\,µH};
  \node[font=\footnotesize, gray] at (1.8,1.1) {Utan kärna};

  % ── Mittendelare ────────────────────────────────────────
  \draw[dashed, gray!50] (4.2,-2.0) -- (4.2,1.5);
  \filldraw[fill=green!60!black, draw=none, rounded corners=8pt, opacity=0.85] 
    (3.8,0.6) rectangle (4.6,1.4);
  \node[white, font=\bfseries\scriptsize] at (4.2,1.1) {×100};
  \node[white, font=\scriptsize] at (4.2,0.85) {mer};

  % ── Höger: ferritkärna ──────────────────────────────────
  \filldraw[fill=brown!60!black, draw=brown!40!black, rounded corners=3pt, opacity=0.85]
    (4.8,-0.35) rectangle (7.6,0.35);
  \node[white, font=\bfseries\tiny] at (6.2,0) {FERRIT};
  \draw[thick] (4.2,0) -- (4.8,0);
  \draw[thick, color=blue!70!black]
    (4.8,0) .. controls (4.8,0.5) and (5.2,0.5) .. (5.2,0)
            .. controls (5.2,-0.5) and (5.6,-0.5) .. (5.6,0)
            .. controls (5.6,0.5) and (6.0,0.5) .. (6.0,0)
            .. controls (6.0,-0.5) and (6.4,-0.5) .. (6.4,0)
            .. controls (6.4,0.5) and (6.8,0.5) .. (6.8,0)
            .. controls (6.8,-0.5) and (7.2,-0.5) .. (7.2,0)
            .. controls (7.2,0.5) and (7.6,0.5) .. (7.6,0);
  \draw[thick] (7.6,0) -- (8.2,0);
  \node[font=\footnotesize] at (6.2,-1.0) {10 varv};
  \node[font=\large\bfseries, color=blue!70!black] at (6.2,-1.5) {500\,µH};
  \node[font=\footnotesize, gray] at (6.2,1.1) {Med ferritkärna};

  % ── Undertext ───────────────────────────────────────────
  \node[gray, font=\itshape\scriptsize] at (4.1,-2.2) 
    {Samma antal varv — kärnan gör skillnaden};

\end{tikzpicture}
\end{figure}

Ferritspolar lindas ofta på en toroid -- en munkformad 
kärna. Det håller magnetfältet inuti ringen och minskar 
störningar ut mot omgivningen. Används för 
bredbandstransformatorer och som RF-choker.

% ── VARIABEL ──────────────────────────────────────────────
\subsubsection*{Variabel spole}

En variabel spole har en justerbar kärna som kan 
skruvas in och ut. Ju djupare kärnan sitter, desto 
mer förstärks magnetfältet och desto högre blir 
induktansen. Används för finavstämning när ett exakt 
värde behövs vid inbyggnad -- analogt med trimmern 
bland kondensatorerna.

% ── SNABBREFERENS ─────────────────────────────────────────
\medskip
\noindent\textbf{Snabbreferens -- spoltyper}

\begin{center}
\begin{tabular}{L{2.8cm}L{4cm}L{3.5cm}}
\rowcolor{greydark}
\tblhdr{Typ} & \tblhdr{Kännetecken} & \tblhdr{Typisk användning} \\
Luftkärna & Låga förluster, stabil & Precisionsfilter, svängningskretsar \\
\rowcolor{greylight}
Ferritkärna (toroid) & Hög induktans per varv, kompakt & Bredbandstransformator, RF-choker \\
Pulverkärna & Högt Q-värde, låga förluster vid HF & HF-filter, transformatorer \\
\rowcolor{greylight}
Variabel & Justerbar kärna med skruv & Finavstämning vid inbyggnad \\
\end{tabular}
\end{center}

\subsubsection*{Övningsfrågor -- 1.5.4}

\begin{qblock}{5.}
\qgrund\quad Du ska välja spoltyp till tre situationer. 
Vilken passar bäst och varför?
\begin{enumerate}[label=\alph*)]
  \item Du behöver hög induktans i ett litet utrymme 
        och kan acceptera något högre förluster.
  \item Du bygger ett precisionsfilter där stabila 
        egenskaper är viktigare än kompakthet.
  \item Du vill kunna justera induktansen med ett 
        skruvmejselstag vid inbyggnad.
\end{enumerate}
\end{qblock}

% ── 1.5.5 ─────────────────────────────────────────────────
\subsection{Tidskonstant}

Precis som kondensatorn inte laddas omedelbart, når 
strömmen genom en spole inte sitt maxvärde direkt. 
Spolens \textbf{tidskonstant} beskriver hur snabbt 
magnetfältet byggs upp -- och den följer samma mönster 
som du lärde dig i avsnitt 1.4.4, men med omvänd formel.

\begin{formulabox}
{\large$\mathbf{\tau = \dfrac{L}{R}}$}\\[3mm]
{\small\normalfont\itshape $\tau$ i sekunder \quad|\quad $L$ i henry \quad|\quad $R$ i ohm}
\end{formulabox}

Jämför med kondensatorn: $\tau = R \times C$. Notera 
att R och L byter plats -- hög induktans och låg 
resistans ger lång tidskonstant för en spole, precis 
som hög kapacitans och hög resistans ger lång 
tidskonstant för en kondensator.

Tidsförloppet är identiskt: efter $1\tau$ har strömmen 
nått 63\,\% av sitt slutvärde, efter $5\tau$ är den 
praktiskt stabil (99\,\%).

\begin{worked}[title={Räkneexempel 1.5.5-1 -- Spole i krets}]
En 10\,mH spole med 5\,$\Omega$ DC-resistans. 
Hur lång tid tar det innan strömmen stabiliseras?

\[ \tau = \frac{L}{R} = \frac{0{,}010\,\text{H}}{5\,\Omega} = \mathbf{2\,\text{ms}} \]

Strömmen stabiliseras efter $5\tau = 10\,\text{ms}$.
\end{worked}

\begin{orangebox}[title={\emoji{⚡} Reläspolar och spänningstransienter}]
Hög induktans och låg resistans ger lång tidskonstant 
-- och det kan vara farligt. När ett relä bryts snabbt 
försöker magnetfältet hålla strömmen vid liv. Eftersom 
kretsen är bruten kan inte strömmen flöda normalt -- 
istället byggs en hög spänning upp på ett ögonblick. 

Denna \textbf{spänningstopp} (transient) kan förstöra 
känsliga komponenter i kretsen. Lösningen är en 
fridiod parallellt med reläspolen -- den ger 
spänningstransienten en ofarlig väg att försvinna.
\end{orangebox}

\begin{orangebox}[title={\emoji{🎯} Viktigt för provet}]
\begin{checklist}
  \item Induktans mäts i \textbf{Henry (H)} -- i praktiken nH, µH, mH
  \item Släpper igenom DC, bromsar AC -- mer vid högre frekvens
  \item $X_L = 2\pi fL$ -- reaktansen \textbf{ökar} vid högre frekvens (tvärtom mot C!)
  \item Fler varv eller järnkärna $\rightarrow$ högre induktans
  \item Tidskonstant: $\tau = L/R$ -- samma förlopp som $\tau = R \times C$, efter $5\tau$ är strömmen stabil
  \item Spole och kondensator tillsammans bildar en LC-krets -- mer om det senare
  \item C och L är varandras motsatser -- lär dig jämförelsetabellen i 1.5.3
\end{checklist}
\end{orangebox}

\subsubsection*{Övningsfrågor -- 1.5.5}

\begin{qblock}{6.}
\qrakna\quad En 10\,mH spole har 5\,$\Omega$ 
DC-resistans. a) Beräkna tidskonstanten $\tau$. 
b) Hur lång tid tar det innan strömmen är praktiskt 
stabil? c) Vad händer med $\tau$ om du byter till en 
spole med 20\,mH men samma resistans?
\end{qblock}

% ── Fälla 1.5.5 ───────────────────────────────────────────
\begin{qblock}{F.}
\qfallan\quad En radioamatör beräknar tidskonstanten för 
en krets med en 10\,mH spole och en 50\,$\Omega$ resistor. 
Han skriver:
\[
\tau = \frac{R}{L} = \frac{50}{0{,}010} = 5\,000\,\text{s}
\]
Han är säker på att formeln är $\tau = R/L$ -- det är 
ju samma bokstäver som för kondensatorn fast tvärtom. 
Men 5\,000 sekunder för en liten spole? Vad gick fel?
\end{qblock}

\medskip
\noindent\rule{\textwidth}{0.4pt}
{\small\itshape Försök själv innan du läser svaret.}
\medskip

\begin{worked}[title={Lösning -- Fallgropsfråga 1.5.5\,F}]
Han blandade ihop täljaren och nämnaren. Formeln är 
$\tau = L/R$ -- \emph{inte} $R/L$.

Rätt beräkning:
\[ \tau = \frac{L}{R} = \frac{0{,}010\,\text{H}}{50\,\Omega} = \mathbf{0{,}2\,\text{ms}} \]

Strömmen stabiliseras efter $5\tau = 1\,\text{ms}$ -- 
rimligt för en liten spole.

\textbf{Lärdomen:} Det är lätt att komma ihåg att 
$\tau = L/R$ och $\tau = R \times C$ -- men lätt att 
blanda ihop vilket som är täljare och nämnare. En 
minneshake: för spolen ligger L \emph{ovanpå} -- 
$\tau = L/R$. För kondensatorn \emph{multipliceras} 
R och C. Kontrollera alltid att svaret är rimligt.
\end{worked}

% ── 1.6 ───────────────────────────────────────────────────
\section{Serie- och parallellkoppling}

\lettrine[lines=2]{\textcolor{orange}{D}}{et} finns två grundläggande sätt att koppla ihop komponenter: serie och parallell. Valet påverkar hela kretsens beteende -- och fel val kan göra att kretsen fungerar dåligt eller inte alls.

\begin{storybox}[{\emoji{💡} Julgransbelysningen och felsökningen}]
Gamla julgransbelysningar var seriekopplade. En trasig 
lampa -- och hela slingan slocknade. Du fick gå igenom 
lampa för lampa tills du hittade den skyldiga.

Moderna julgransljus är parallellkopplade. En trasig 
lampa påverkar inte de övriga -- de lyser vidare. 
Du märker knappt att en gått sönder.

Samma princip, helt olika konsekvenser. Det är 
skillnaden mellan serie och parallell i ett nötskal.
\end{storybox}

I resten av avsnittet ser du hur kopplingssättet avgör 
hur spänning, ström och resistans fördelas i kretsen.


% ── 1.6.1 ─────────────────────────────────────────────────
\subsection{Seriekoppling}

I en seriekrets flödar samma ström genom varje 
komponent -- det finns bara en väg. Spänningen däremot 
fördelas proportionellt mot resistansen i varje del.

\begin{formulabox}
{\large$\mathbf{R_{tot} = R_1 + R_2 + R_3 + \ldots}$}\\[3mm]
{\small\normalfont\itshape $I = I_1 = I_2 = I_3$ \quad|\quad $U_{tot} = U_1 + U_2 + U_3$}
\end{formulabox}

\begin{bluebox}[title={\emoji{🔑} Nyckelsamband -- serie}]
\textbf{Ström:} Identisk i alla komponenter -- det finns 
bara en väg för elektronen att ta.\\
\textbf{Spänning:} Delas proportionellt mot resistansen 
-- stor R får stor spänningsdel.\\
\textbf{Resistans:} Adderas alltid -- $R_{tot}$ blir 
alltid större än den största enskilda.
\end{bluebox}

\begin{worked}[title={Räkneexempel 1.6.1 -- Spänningsdelare med LED}]
En LED kräver 2\,V och 20\,mA. Matningsspänningen är 
12\,V. Vilket seriemotstånd behövs och vilken 
effektklassning ska det ha?

Spänning över resistorn:
\[ U_R = 12 - 2 = 10\,\text{V} \]

Resistans:
\[ R = \frac{U_R}{I} = \frac{10}{0{,}020} = 500\,\Omega \quad \rightarrow \quad \text{välj }470\,\Omega \]


Effekt i resistorn:
\[ P = U_R \times I = 10 \times 0{,}020 = 0{,}2\,\text{W} \quad \rightarrow \quad \text{välj 0{,}5\,W} \]
\end{worked}

Nu vet vi hur serie fungerar -- men vad händer om vi 
ger strömmen flera vägar att välja på?

\subsubsection*{Övningsfrågor -- 1.6.1}

\begin{qblock}{1.}
\qrakna\quad Tre resistorer -- 100\,$\Omega$, 220\,$\Omega$ 
och 470\,$\Omega$ -- kopplas i serie till 12\,V. 
a) Beräkna $R_{tot}$. 
b) Beräkna strömmen i kretsen. 
c) Beräkna spänningsfallet över varje resistor och 
kontrollera att de summerar till 12\,V.
\end{qblock}

% ── Fälla 1.6.1 ───────────────────────────────────────────
\begin{qblock}{F.}
\qfallan\quad Maria mäter spänningarna i en seriekrets 
med tre resistorer och 12\,V matning. Hon mäter 
3\,V, 4\,V och 4\,V -- totalt 11\,V. Hon drar 
slutsatsen att Kirchhoffs spänningslag inte stämmer 
i den här kretsen.

Har hon rätt?
\end{qblock}

\medskip
\noindent\rule{\textwidth}{0.4pt}
{\small\itshape Försök själv innan du läser svaret.}
\medskip

\begin{worked}[title={Lösning -- Fallgropsfråga 1.6.1\,F}]
Kirchhoffs spänningslag stämmer alltid -- det är 
mätningen som är fel.

KVL säger att spänningsfallen ska summera till 
källspänningen. $3 + 4 + 4 = 11\,\text{V}$, inte 
12\,V. Det betyder att en volt försvinner någonstans 
-- och det gör spänningar aldrig av sig självt.

Troliga orsaker:
\begin{itemize}
  \item Multimetern är felinställd eller har svagt 
        batteri
  \item Hon mäter inte över rätt punkter
  \item Det finns ett spänningsfall hon missat -- 
        t.ex. över en kontakt eller en ledning med 
        oväntat hög resistans
\end{itemize}

\textbf{Lärdomen:} Om spänningsfallen inte summerar 
till källspänningen är mätningen fel -- inte lagen. 
KVL är ett av de säkraste verktygen du har för att 
hitta fel i en krets.
\end{worked}

% ── 1.6.2 ─────────────────────────────────────────────────
\subsection{Parallellkoppling}

I en parallellkrets har varje komponent samma spänning 
-- de sitter alla direkt mellan samma två punkter. 
Strömmen däremot delas -- mer ström flödar genom 
lägre resistans.

\begin{formulabox}
{\large$\mathbf{\dfrac{1}{R_{tot}} = \dfrac{1}{R_1} + \dfrac{1}{R_2} + \dfrac{1}{R_3} + \ldots}$}\\[4mm]
{\small\normalfont\itshape Genväg för två resistorer: $R_{tot} = \dfrac{R_1 \times R_2}{R_1 + R_2}$}
\end{formulabox}

\begin{bluebox}[title={\emoji{🔑} Nyckelsamband -- parallell}]
\textbf{Spänning:} Identisk över alla komponenter.\\
\textbf{Ström:} Delas -- mer ström i lägre resistans.\\
\textbf{Resistans:} Minskar alltid -- $R_{tot}$ blir 
alltid \emph{lägre} än den minsta enskilda.
\end{bluebox}

\begin{worked}[title={Räkneexempel 1.6.2 -- Dummylast av 200\,$\Omega$-resistorer}]
Fyra 200\,$\Omega$ resistorer kopplas parallellt. Vad 
blir $R_{tot}$ och hur fördelas effekten?

\[ R_{tot} = \frac{200}{4} = \mathbf{50\,\Omega} \]

Effekttåligheten fyrdubblas -- varje resistor tar sin 
fjärdedel. Om varje resistor tål 1\,W klarar 
konstruktionen totalt 4\,W.
\end{worked}

\subsubsection*{Övningsfrågor -- 1.6.2}

\begin{qblock}{2.}
\qrakna\quad Samma tre resistorer -- 100\,$\Omega$, 
220\,$\Omega$ och 470\,$\Omega$ -- kopplas nu parallellt 
till 12\,V. 
a) Beräkna $R_{tot}$ med den allmänna formeln. 
b) Beräkna total ström och strömmen i varje gren. 
c) Kontrollera: är $R_{tot}$ lägre än den minsta 
resistorn?
\end{qblock}

\begin{qblock}{3.}
\qrakna\quad Du ska bygga en 50\,$\Omega$ dummylast som 
tål minst 4\,W för testning av QRP-sändare. Du har 
200\,$\Omega$ resistorer som tål 1\,W var. 
a) Hur många behövs och hur kopplas de? 
b) Vilken total effekttålighet får konstruktionen? 
c) Varför räcker inte en enda 50\,$\Omega$ / 1\,W resistor?
\end{qblock}

% ── Fälla 1.6.2 ───────────────────────────────────────────
\begin{qblock}{F.}
\qfallan\quad Erik kopplar 100\,$\Omega$ och 100\,$\Omega$ 
parallellt och räknar:
\[ R_{tot} = 100 + 100 = 200\,\Omega \]
Han tror att parallellkoppling ger samma resultat som 
seriekoppling -- man adderar ju komponenterna i båda 
fallen, eller? Vad gick fel?
\end{qblock}

\medskip
\noindent\rule{\textwidth}{0.4pt}
{\small\itshape Försök själv innan du läser svaret.}
\medskip

\begin{worked}[title={Lösning -- Fallgropsfråga 1.6.2\,F}]
Erik använde serieformeln på en parallellkrets. Det är 
det vanligaste misstaget med kopplingsregler.

Rätt formel för parallell:
\[ R_{tot} = \frac{R_1 \times R_2}{R_1 + R_2} = \frac{100 \times 100}{100 + 100} = \frac{10\,000}{200} = \mathbf{50\,\Omega} \]

Inte 200\,$\Omega$ -- utan hälften av en resistor. 
Det är precis vad regeln säger: parallellkoppling 
ger alltid \emph{lägre} $R_{tot}$ än den minsta 
ingående komponenten. Två lika resistorer parallellt 
ger alltid halva värdet.

\textbf{Lärdomen:} Parallell $\neq$ addera. 
Parallellkoppling minskar alltid resistansen -- 
ju fler resistorer du lägger till, desto lägre 
$R_{tot}$.
\end{worked}

% ── 1.6.3 ─────────────────────────────────────────────────
\subsection{Kirchhoffs lagar}

Reglerna för serie- och parallellkoppling bygger på två
grundläggande lagar formulerade av Gustav Kirchhoff 1845.
Du har redan använt dem i praktiken -- nu får de sina namn.

\begin{bluebox}[title={\emoji{📐} Kirchhoffs lagar}]
\textbf{Strömlagen (KCL) -- Kirchhoffs Current Law:}\\
Summan av alla strömmar \emph{in} i en nod är lika
med summan av alla strömmar \emph{ut}. Ingen ström
försvinner -- det som flödar in måste flöda ut.

\[ I_{\text{in}} = I_{\text{ut}} \quad \Rightarrow \quad I_{\text{tot}} = I_1 + I_2 + I_3 \]

\medskip
\textbf{Spänningslagen (KVL) -- Kirchhoffs Voltage Law:}\\
Summan av alla spänningar i en sluten krets är noll.
All spänning som källan levererar förbrukas över
komponenterna.

\[ U_{\text{källa}} = U_1 + U_2 + U_3 \]
\end{bluebox}

Du känner igen båda: KCL använder du i noder och
parallellgrenar, och KVL varje gång du summerar
spänningsfall i en sluten krets.


% ── Fälla 1.6.3 ───────────────────────────────────────────
\begin{qblock}{F.}
\qfallan\quad Laban ska parallellkoppla två spolar på 
100\,µH vardera. Han minns att kondensatorer i 
parallell adderas, och att spolar är kondensatorernas 
motsatser -- så han räknar att två spolar parallellt 
ger hälften, 50\,µH.

Är hans resonemang rätt?
\end{qblock}

\medskip
\noindent\rule{\textwidth}{0.4pt}
{\small\itshape Försök själv innan du läser svaret.}
\medskip

\begin{worked}[title={Lösning -- Fallgropsfråga 1.6.3\,F}]
Svaret är rätt -- men resonemanget är fel, och det 
kommer leda honom fel nästa gång.

Spolar följer \textbf{samma regler som resistorer} -- 
inte omvända regler som kondensatorer. Parallella 
spolar ger lägre $L_{tot}$, precis som parallella 
resistorer ger lägre $R_{tot}$.

Två lika spolar parallellt:
\[ L_{tot} = \frac{L_1 \times L_2}{L_1 + L_2} = \frac{100 \times 100}{100 + 100} = \mathbf{50\,\mu\text{H}} \]

Han fick rätt svar av fel anledning. Resonemanget 
``spolar är tvärtom mot kondensatorer'' stämmer för 
DC/AC-beteende, men inte för kopplingsregler.

\textbf{Lärdomen:} Kondensatorn är alltid tvärtom 
mot resistorn. Spolen är alltid likadan som 
resistorn. Inte tvärtom mot kondensatorn.
\end{worked}

\subsubsection*{Övningsfrågor -- 1.6.3}

\begin{qblock}{4.}
\qgrund\quad En parallellkrets har tre grenar och matas 
med 12\,V. Du mäter strömmen i två av grenarna: 
120\,mA och 54\,mA. Totalströmmen från källan är 
225\,mA.

a) Vad säger KCL om strömmen i den tredje grenen? 
Beräkna den.

b) Vilken resistans har den tredje grenen?

c) Stämmer totalströmmen med $I = U/R_{tot}$ om du 
beräknar $R_{tot}$ från de tre grenresistanserna?
\end{qblock}


% ── 1.6.4 ─────────────────────────────────────────────────
\subsection{Omvända regler för kondensatorer och spolar}

Resistorer följer de intuitiva reglerna: serie ökar, 
parallell minskar. Kondensatorer och spolar är 
annorlunda -- och att blanda ihop dem är en av de 
vanligaste misstagen på provet.

\begin{center}
\begin{tabular}{L{2.8cm}L{4.2cm}L{4.2cm}}
\rowcolor{greydark}
\tblhdr{Komponent} & \tblhdr{Serie} & \tblhdr{Parallell} \\
Resistorer (R) & $R_{tot}$ \textbf{ökar} -- adderas & $R_{tot}$ \textbf{minskar} \\
\rowcolor{greylight}
Kondensatorer (C) & $C_{tot}$ \textbf{minskar} & $C_{tot}$ \textbf{ökar} -- adderas \\
Spolar (L) & $L_{tot}$ \textbf{ökar} -- adderas & $L_{tot}$ \textbf{minskar} \\
\end{tabular}
\end{center}

Kondensatorn följer \emph{omvända} regler mot 
resistorn -- och precis samma regler som resistorn 
gäller för spolen. Minnesregel: kondensatorn är alltid 
tvärtom, spolen är alltid som resistorn.

\begin{worked}[title={Räkneexempel 1.6.4 -- Kondensatorer i serie och parallellt}]
Du har 47\,µF, 100\,µF och 220\,µF.

\textbf{Parallellt -- kapacitansen adderas:}
\[ C_{tot} = 47 + 100 + 220 = \mathbf{367\,\mu\text{F}} \]

\textbf{47\,µF och 100\,µF i serie:}
\[ \frac{1}{C_{tot}} = \frac{1}{47} + \frac{1}{100} = 0{,}02128 + 0{,}01 = 0{,}03128 \]
\[ C_{tot} = \frac{1}{0{,}03128} \approx \mathbf{32\,\mu\text{F}} \]

Serie gav lägre kapacitans -- tvärtom mot vad 
resistorer gör i serie.
\end{worked}

\subsubsection*{Övningsfrågor -- 1.6.4}

\begin{qblock}{5.}
\qrakna\quad Du har kondensatorerna 47\,µF, 100\,µF 
och 220\,µF. 
a) Vad är $C_{tot}$ om alla tre kopplas parallellt? 
b) Om 47\,µF och 100\,µF kopplas i serie -- vad är 
$C_{tot}$? 
c) Varför är kondensatorer i serie ovanligt i 
praktiska kretsar?
\end{qblock}

\begin{qblock}{6.}
\qprov\quad Du bygger ett LC-filter för 40-metersbandet 
(7\,MHz). Under testerna märker du att resonansfrekvensen 
är för hög -- filtret behöver stämmas ner.

Resonansfrekvensen beror på $L$ och $C$ enligt:
\[ f_0 = \frac{1}{2\pi\sqrt{LC}} \]

Du har tillgång till extra komponenter av samma typ 
som redan sitter i kretsen.

a) Vad måste hända med $L$ och/eller $C$ för att 
$f_0$ ska sjunka? Motivera med formeln.

b) Hur kopplar du extra spolar respektive extra 
kondensatorer för att uppnå det -- serie eller 
parallell? Motivera med kopplingsreglerna från 
avsnitt 1.6.4.

c) En tekniker föreslår att lägga till en extra spole 
parallellt och en extra kondensator parallellt. 
Håller du med? Motivera.
\end{qblock}

% ── Fälla 1.6.4 ───────────────────────────────────────────
\begin{qblock}{F.}
\qfallan\quad Sara ska öka kapacitansen i sitt filter. 
Hon kopplar tre 100\,µF kondensatorer i serie och 
räknar ut att hon får $3 \times 100 = 300$\,µF.

Stämmer det -- och är serie rätt koppling för att 
öka kapacitansen?
\end{qblock}

\medskip
\noindent\rule{\textwidth}{0.4pt}
{\small\itshape Försök själv innan du läser svaret.}
\medskip

\begin{worked}[title={Lösning -- Fallgropsfråga 1.6.4\,F}]
Dubbelt fel. 

\textbf{Fel 1 -- formeln:} Kondensatorer i serie 
adderas \emph{inte}. Rätt formel:
\[ \frac{1}{C_{tot}} = \frac{1}{100} + \frac{1}{100} + \frac{1}{100} = \frac{3}{100} \quad \Rightarrow \quad C_{tot} = \frac{100}{3} \approx \mathbf{33\,\mu\text{F}} \]

\textbf{Fel 2 -- kopplingen:} Serie \emph{minskar} 
kapacitansen. Hon fick en tredjedel av vad hon 
startade med -- raka motsatsen till vad hon ville.

För att öka kapacitansen ska kondensatorer kopplas 
\textbf{parallellt} -- då adderas värdena direkt:
\[ C_{tot} = 100 + 100 + 100 = 300\,\mu\text{F} \]

\textbf{Lärdomen:} Kondensatorer följer omvända 
regler mot resistorer. Serie minskar, parallell ökar. 
Minnesregel: kondensatorn är alltid tvärtom.
\end{worked}

% ── JÄMFÖRELSE ────────────────────────────────────────────
\begin{center}
\noindent\textbf{Sammanfattning -- serie vs parallell}\\[4mm]
\begin{tabular}{L{3cm}L{3.5cm}L{3.5cm}}
\rowcolor{greydark}
\tblhdr{Egenskap} & \tblhdr{Serie} & \tblhdr{Parallell} \\
Ström & Identisk i alla & Delas mellan grenar \\
\rowcolor{greylight}
Spänning & Delas proportionellt & Identisk över alla \\
Resistans (R) & Adderas -- ökar & Minskar \\
\rowcolor{greylight}
Kapacitans (C) & Minskar & Adderas -- ökar \\
Induktans (L) & Adderas -- ökar & Minskar \\
\rowcolor{greylight}
Fel i komponent & Hela kretsen stannar & Övriga fungerar \\
\end{tabular}
\end{center}

\begin{orangebox}[title={\emoji{🎯} Viktigt för provet}]
\begin{checklist}
  \item Serie: Ström identisk, spänning delas, $R$ adderas
  \item Parallell: Spänning identisk, ström delas, $R$ minskar
  \item Kondensatorer: Serie \textbf{minskar} $C$, parallell \textbf{ökar} $C$ -- tvärtom mot $R$!
  \item Spolar: Samma regler som resistorer
  \item Genväg 2 parallella: $R_{tot} = (R_1 \times R_2) / (R_1 + R_2)$
  \item KCL: ström in = ström ut i varje nod
  \item KVL: spänningarna i en sluten krets summerar till noll
\end{checklist}
\end{orangebox}

% ── 1.7 ───────────────────────────────────────────────────
\section{Växelström och impedans}

\lettrine[lines=2]{\textcolor{orange}{I}}{nom} radiotekniken arbetar vi nästan uteslutande med 
växelström -- en ström som ständigt byter riktning. 
Nätströmmen gör det 50 gånger per sekund. RF-signaler 
gör det miljoner eller miljarder gånger per sekund. 
Utan förståelse för hur frekvens, fas och impedans samverkar 
går det inte att förstå varför en antenn fungerar, varför 
en kabel har förluster eller varför sändaren ska matchas 
mot antennen.


% ── 1.7.1 ─────────────────────────────────────────────────
\subsection{Frekvens och period}

Varje växelström upprepar sig i ett mönster -- en
cykel. Hur snabbt den cykeln upprepas är frekvensen,
och hur lång tid en cykel tar är perioden.
I radiotekniken, där signaler nästan alltid är
sinusvågor, upprepas dessa cykler i miljontals per
sekund -- se figur~\ref{fig:sinusvag} för hur en
typisk sinusvåg ser ut.

\begin{storybox}[{\emoji{💡} Hjärtat och metronomen}]
Tänk på en hjärtmonitor. Linjen på skärmen svänger
upp och ner -- en gång per hjärtslag. Om hjärtat slår
60 gånger per minut slår det en gång per sekund.
Frekvensen är 1~Hz, perioden är 1~sekund.

En radioamatör på 14~MHz skickar en signal vars
``hjärtslag'' upprepas 14~000~000 gånger varje sekund.
Perioden -- alltså den tid ett enda svängning tar --
är ungefär 71~nanosekunder. Så fort att inget
mänskligt sinne kan uppfatta det, men elektroniken
lever i den verkligheten.
\end{storybox}

Frekvens och period är varandras inverser --
känner du till det ena beräknar du direkt det andra.
En full cykel går från noll upp till toppvärdet,
ner genom noll till botten och tillbaka -- de orange
punkterna i figuren markerar just dessa nollgenomgångar.

\begin{formulabox}
{\large$\mathbf{f = \dfrac{1}{T}}$\quad|\quad$\mathbf{T = \dfrac{1}{f}}$}\\[3mm]
{\small\normalfont\itshape
$f$ = frekvens i hertz (Hz)\quad|\quad$T$ = period i sekunder (s)
\quad|\quad Exempel: $f = 1\,\text{MHz} \Rightarrow T = 1\,\mu\text{s}$}
\end{formulabox}

% ── Figur: sinusvåg ───────────────────────────────────────
\begin{figure}[htbp]
\centering
\begin{tikzpicture}[>=stealth, font=\small]

  % Inställningar
  \def\A{1.8}        % amplitud (cm)
  \def\T{4.0}        % en period (cm)
  \def\xmax{8.8}     % total x-längd = 2.2 perioder

  % Bakgrundsruta
  \fill[gray!6, rounded corners=4pt]
    (-1.4,-\A-0.7) rectangle (\xmax+1.2,\A+0.9);

  % Axlar
  \draw[->, thick, gray!60!black]
    (-0.4,0) -- (\xmax+0.4,0) node[right] {$t$};
  \draw[->, thick, gray!60!black]
    (0,-\A-0.5) -- (0,\A+0.6) node[above] {$U$};
  \node[gray!60!black] at (-0.2,-0.2) {\footnotesize 0};

  % Streckade referenslinjer vid topp och botten
  \draw[red!70!black, dashed, thin]
    (0,\A) -- ({\T/4},\A);
  \draw[red!70!black, dashed, thin]
    (0,-\A) -- ({3*\T/4},-\A);

  % Sinusvåg -- två hela perioder
  \draw[orange!90!red, very thick, samples=300, domain=0:{2*\T}]
    plot (\x, {\A*sin(360/\T * \x)});

  % Nollgenomgångspunkter
  \foreach \xp in {0, {\T/2}, \T, {3*\T/2}, {2*\T}}
    \fill[orange!80!red] (\xp,0) circle (2.2pt);

  % Period T -- pil under grafen
  \draw[<->, blue!70!black, thick]
    (0,-\A-0.55) -- (\T,-\A-0.55)
    node[midway, below=2pt]
      {$T$ \footnotesize= tid för en hel svängning};
  \draw[blue!70!black, thin] (0,-\A-0.35) -- (0,-\A-0.75);
  \draw[blue!70!black, thin] (\T,-\A-0.35) -- (\T,-\A-0.75);

  % U_topp pil (vänster om y-axeln)
  \draw[<->, red!70!black, thick]
    (-0.9,0) -- (-0.9,\A)
    node[midway, left=2pt] {$U_\text{topp}$};
  \node[red!70!black, align=center, font=\scriptsize] at (-0.9,-0.35)
    {max\\värde};

  % -U_topp etikett
  \node[red!70!black, font=\small, left=6pt] at (0,-\A)
    {$-U_\text{topp}$};

\end{tikzpicture}
\caption{En sinusvåg för en radiosignal. Period $T$ (blå pil) är
tiden för en hel svängning. Toppvärdet $U_\text{topp}$ (röd) är
det högsta värdet -- och $-U_\text{topp}$ det lägsta.
Nollgenomgångarna markeras med orange punkter.}
\label{fig:sinusvag}
\end{figure}

\begin{bluebox}[title={\emoji{🔑} Nyckelsamband -- frekvens och period}]
\textbf{Frekvens} mäts i hertz (Hz) -- antal cykler per sekund.\\
\textbf{Period} mäts i sekunder (s) -- tid för en cykel.\\
\textbf{Samband:} $f \times T = 1$ -- de är varandras inverser.\\
\textbf{Prefix:} kHz = $10^3$~Hz,\; MHz = $10^6$~Hz,\; GHz = $10^9$~Hz.
\end{bluebox}

% ── Frekvenstabell ────────────────────────────────────────
\medskip
\noindent Här är några exempel på frekvenser i vardagen och
inom radio -- lägg märke till hur perioden krymper
när frekvensen stiger:
\medskip

\begin{center}
\begin{tabular}{L{2.4cm}L{2.4cm}L{4.2cm}L{2.2cm}}
\rowcolor{greydark}
\tblhdr{Frekvens} & \tblhdr{Band} & \tblhdr{Radioexempel} & \tblhdr{Period (ca.)} \\
50~Hz         & Nätfrekvens   & Eluttag i Europa (230~V AC) & 20~ms \\
\rowcolor{greylight}
20~Hz--20~kHz & Audio         & Mikrofoner, högtalare, SSB-tal & 50~µs--50~ms \\
3--30~MHz     & HF (kortväg)  & Amatörradio, rundradio, DX & 33--333~ns \\
\rowcolor{greylight}
144--146~MHz  & VHF 2m        & Lokal FM, repeatrar & $\approx$7~ns \\
430--440~MHz  & UHF 70~cm     & Handhållen radio, D-STAR & $\approx$2~ns \\
\rowcolor{greylight}
2,4~GHz       & Mikrovåg      & WiFi, amatörradio (13~cm) & $\approx$0,4~ns \\
\end{tabular}
\end{center}

\begin{worked}[title={Räkneexempel 1.7.1 -- Frekvens och period på 40-metersbandet}]
En amatörradiosignal sänds på 7,050~MHz.

\textbf{a) Vad är perioden?}
\[ T = \frac{1}{f} = \frac{1}{7{,}050 \times 10^6} \approx 142\,\text{ns} \]

\textbf{b) Hur många perioder hinns med på en sekund?}

Det är bara definitionen av frekvens:\\
$7{,}050 \times 10^6 = 7\,050\,000$ perioder per sekund.

\textbf{c) WiFi arbetar på 2,4~GHz. Hur förhåller sig perioden till 7~MHz?}
\[ T_{\text{WiFi}} = \frac{1}{2{,}4 \times 10^9} \approx 0{,}42\,\text{ns} \]
Perioden är ungefär 340 gånger kortare -- signalen
svänger 340 gånger fortare.
\end{worked}

\subsubsection*{Övningsfrågor -- 1.7.1}

\begin{qblock}{1.}
\qrakna\quad En signal har frekvensen 144,500~MHz.

a) Beräkna periodens längd i nanosekunder.\\
b) En annan signal har perioden $T = 200\,\text{ns}$.
Vad är dess frekvens i MHz?\\
c) Hur förhåller sig de två frekvenserna till
varandra -- vilket band tillhör de?
\end{qblock}

\begin{qblock}{F.}
\qfallan\quad Erik hör att ``frekvensen på 40-metersbandet
är 7~MHz''. Han drar slutsatsen att perioden är
$T = 7 \times 10^6$~sekunder -- eftersom period och
frekvens är samma sak fast med olika enhet.

Vad gick fel i hans resonemang?
\end{qblock}

\medskip
\noindent\rule{\textwidth}{0.4pt}
{\small\itshape Försök själv innan du läser svaret.}
\medskip

\begin{worked}[title={Lösning -- Fallgropsfråga 1.7.1\,F}]
Erik har förväxlat frekvens och period -- de är
\emph{inte} samma sak med olika enhet. De är
varandras \textbf{inverser}.

Rätt beräkning:
\[ T = \frac{1}{f} = \frac{1}{7 \times 10^6\,\text{Hz}} \approx 143\,\text{ns} \]

$7 \times 10^6$ sekunder är nästan tre månader --
inte ens nätfrekvensens period (20~ms) i närheten.
Svaret måste hamna i nanosekunder för MHz-signaler.
Enhetskontroll avslöjar felet:
Hz~=~1/s, alltså $T = 1/\text{Hz} = \text{s}$.

\textbf{Lärdomen:} Frekvens och period är inverser.
Hög frekvens = kort period. Låg frekvens = lång period.
\end{worked}


% ── 1.7.2 ─────────────────────────────────────────────────
\subsection{Effektivvärde (RMS)}

Växelströmmen varierar ständigt -- sinusvågen du
känner igen svänger upp, ner, upp, ner. Men hur
mycket \emph{arbete} utför den egentligen? Hur mycket
värme alstras, hur starkt lyser en lampa? Toppvärdet
räcker inte som svar -- vi behöver något mer
representativt.

\begin{storybox}[{\emoji{💡} Snöröjarens genomsnittliga kraft}]
Föreställ dig en snöröjare som arbetar varierande
hårt under dagen -- ibland full kraft, ibland
halvfart, ibland pausar han. Om du ska beräkna
hur mycket snö han röjer totalt bryr du dig inte
om ögonblicksvärdena. Du frågar: \emph{vilken konstant
kraftnivå hade gett samma resultat?}

Det är exakt vad RMS gör med växelström.
RMS -- Root Mean Square -- är det AC-värde
som ger \emph{samma effekt} som ett lika stort DC-värde.
230~V RMS ur vägguttaget värmer en elpatron lika
mycket som 230~V DC skulle göra.
\end{storybox}

För en ren sinusvåg är sambandet mellan toppvärde
och RMS alltid detsamma -- det beror bara på
kurvans form, inte på frekvensen.

\begin{formulabox}
{\large$\mathbf{U_{RMS} = U_{topp} \times 0{,}707}$
\quad|\quad
$\mathbf{U_{topp} = U_{RMS} \times 1{,}414}$}\\[3mm]
{\small\normalfont\itshape
Gäller för sinusvåg.\quad
$0{,}707 = \dfrac{1}{\sqrt{2}}$\quad|\quad
$1{,}414 = \sqrt{2}$\quad|\quad
Exempel: $U_{topp} = 325\,\text{V} \Rightarrow U_{RMS} \approx 230\,\text{V}$}
\end{formulabox}

Varför just 0,707? En sinusvåg tillbringar mest tid
nära noll och minst tid nära toppvärdet -- kurvan
rusar förbi toppen men dröjer kvar i mitten.
Genomsnittet av dess ``arbetsförmåga'' hamnar därför
lägre än toppen, på exakt $1/\sqrt{2} \approx 0{,}707$
av toppvärdet. För andra vågformer gäller andra tal
-- en fyrkantvåg till exempel har RMS = toppvärdet
eftersom den håller full nivå hela halvperioden.
Men för sinusvågen, som dominerar i radioteknik,
är 0,707 alltid svaret.

% ── Figur: sinusvåg med RMS ───────────────────────────────
\begin{figure}[htbp]
\centering
\begin{tikzpicture}[>=stealth, font=\small]

  % Inställningar
  \def\A{1.8}        % amplitud (cm)
  \def\T{4.0}        % en period (cm)
  \def\xmax{8.8}     % total x-längd = 2.2 perioder
  \def\rms{1.2728}   % A * 1/sqrt(2)

  % Bakgrundsruta
  \fill[gray!6, rounded corners=4pt]
    (-1.4,-\A-0.7) rectangle (\xmax+2.6,\A+0.9);

  % Axlar
  \draw[->, thick, gray!60!black]
    (-0.4,0) -- (\xmax+0.4,0) node[right] {$t$};
  \draw[->, thick, gray!60!black]
    (0,-\A-0.5) -- (0,\A+0.6) node[above] {$U$};
  \node[gray!60!black] at (-0.2,-0.2) {\footnotesize 0};

  % Streckade referenslinjer vid topp och botten
  \draw[red!70!black, dashed, thin]
    (0,\A) -- ({\T/4},\A);
  \draw[red!70!black, dashed, thin]
    (0,-\A) -- ({3*\T/4},-\A);

  % RMS-linje
  \draw[green!55!black, dashed, semithick]
    (0,\rms) -- (\xmax,\rms);

  % Sinusvåg -- två hela perioder
  \draw[orange!90!red, very thick, samples=300, domain=0:{2*\T}]
    plot (\x, {\A*sin(360/\T * \x)});

  % Nollgenomgångspunkter
  \foreach \xp in {0, {\T/2}, \T, {3*\T/2}, {2*\T}}
    \fill[orange!80!red] (\xp,0) circle (2.2pt);

  % Period T -- pil under grafen
  \draw[<->, blue!70!black, thick]
    (0,-\A-0.55) -- (\T,-\A-0.55)
    node[midway, below=2pt]
      {$T$ \footnotesize= en hel svängning};
  \draw[blue!70!black, thin] (0,-\A-0.35) -- (0,-\A-0.75);
  \draw[blue!70!black, thin] (\T,-\A-0.35) -- (\T,-\A-0.75);

  % U_topp pil
  \draw[<->, red!70!black, thick]
    (-0.9,0) -- (-0.9,\A)
    node[midway, left=2pt] {$U_\text{topp}$};
  \node[red!70!black, align=center, font=\scriptsize] at (-0.9,-0.35)
    {max\\värde};

  % -U_topp etikett
  \node[red!70!black, font=\small, left=6pt] at (0,-\A)
    {$-U_\text{topp}$};

  % U_RMS etikett höger
  \node[green!50!black, font=\small, right=6pt, align=left]
    at (\xmax,\rms)
    {$U_\text{RMS} = 0{,}707 \times U_\text{topp}$\\
     \scriptsize samma effekt som DC};

\end{tikzpicture}
\caption{Samma sinusvåg som i avsnitt~1.7.1, nu med
effektivvärdet $U_\text{RMS}$ (grön, streckad) tillagt.
RMS-linjen ligger på 70,7\,\% av toppvärdet och
representerar den DC-nivå som ger motsvarande effekt.
Gäller för ren sinusvåg.}
\label{fig:sinusvag_rms}
\end{figure}

\begin{bluebox}[title={\emoji{🔑} Nyckelsamband -- RMS}]
\textbf{RMS} är det ``verkliga'' värdet för effektberäkningar -- inte toppvärdet.\\
\textbf{Toppvärde} är alltid högre än RMS: $U_{topp} = \sqrt{2} \times U_{RMS} \approx 1{,}414 \times U_{RMS}$.\\
\textbf{Multimetrar} visar normalt RMS automatiskt.\\
\textbf{Oscilloskop} visar toppvärde -- dividera med $\sqrt{2}$ för att få RMS.\\
\textbf{Effekt:} Använd \emph{alltid} $P = U_{RMS} \times I_{RMS}$ -- aldrig toppvärden!
\end{bluebox}

\begin{worked}[title={Räkneexempel 1.7.2a -- Nätspänningens toppvärde}]
Vägguttaget ger 230~V RMS.

\textbf{Vilket toppvärde måste isoleringen klara?}
\[ U_{topp} = 230 \times 1{,}414 \approx \mathbf{325\,\text{V}} \]

Isoleringen i nätansluten utrustning måste klara
325~V -- inte 230~V! Felaktig dimensionering kan
leda till isolationsgenomslag.
\end{worked}

\begin{worked}[title={Räkneexempel 1.7.2b -- Toppspänning på koaxkabeln}]
En sändare ger 100~W ut i 50~$\Omega$.
Vad är toppspänningen på koaxkabeln?

Löser ut $U_{RMS}$ ur effektformeln:
\[ P = \frac{U_{RMS}^2}{R}
   \quad\Rightarrow\quad
   U_{RMS} = \sqrt{P \times R} = \sqrt{100 \times 50} = 70{,}7\,\text{V} \]

\[ U_{topp} = 70{,}7 \times 1{,}414 \approx \mathbf{100\,\text{V}} \]

Koaxkabelns spänningshållfasthet och eventuella
kondensatorer i antennsystemet måste dimensioneras
för 100~V toppvärde -- inte för RMS-värdet 70{,}7~V.
Underskattar du toppvärdet riskerar du isolationsgenomslag
vid full sändareffekt.
\end{worked}

\subsubsection*{Övningsfrågor -- 1.7.2}

\begin{qblock}{2.}
\qrakna\quad

a) Varför är RMS-värdet för en sinusvåg lägre än
toppvärdet? Förklara med egna ord -- utan formler.\\
b) Nätspänningen är 230~V RMS. Vad är toppvärdet?\\
c) En sändare matar antennen med 100~V toppvärde.
Vad är RMS-värdet?\\
d) Varför är det viktigt att använda RMS-värden
vid effektberäkningar och inte toppvärden?
\end{qblock}

\begin{qblock}{F.}
\qfallan\quad Sara mäter spänningen på en koaxkabel
med ett oscilloskop och ser ett toppvärde på 141~V.
Hon beräknar effekten i 50~$\Omega$:
\[ P = \frac{U^2}{R} = \frac{141^2}{50} \approx 398\,\text{W} \]
Hon rapporterar att sändaren ger 398~W.
Hur stor är effekten egentligen -- och var gick det fel?
\end{qblock}

\medskip
\noindent\rule{\textwidth}{0.4pt}
{\small\itshape Försök själv innan du läser svaret.}
\medskip

\begin{worked}[title={Lösning -- Fallgropsfråga 1.7.2\,F}]
Sara använde toppvärdet direkt i effektformeln.
Det ger dubbelt för hög effekt.

Oscilloskopet visar toppvärde -- det måste först
konverteras till RMS:
\[ U_{RMS} = \frac{141}{1{,}414} \approx 99{,}7\,\text{V} \approx 100\,\text{V} \]
\[ P = \frac{U_{RMS}^2}{R} = \frac{100^2}{50} = \mathbf{200\,\text{W}} \]

Sändaren ger 200~W -- inte 398~W. Felet uppstår
för att toppvärdet är $\sqrt{2}$ gånger RMS, och
$(\sqrt{2})^2 = 2$ -- effekten överdrivs med
exakt en faktor 2.

\textbf{Lärdomen:} Effektformler kräver alltid
RMS-värden. Oscilloskop visar toppvärde -- konvertera
alltid med $U_{RMS} = U_{topp} / \sqrt{2}$ innan
du beräknar effekt.
\end{worked}


% ── 1.7.3 ─────────────────────────────────────────────────
\subsection{Impedans}

I DC-kretsar räcker resistansen för att beskriva hur
kretsen motstår ström. Men i AC-kretsar med spolar
och kondensatorer tillkommer ytterligare motstånd
som \emph{beror på frekvensen} -- reaktans. Impedansen
samlar ihop allt.

\begin{storybox}[{\emoji{🚣} Roddbåten och de tre krafterna}]
Föreställ dig att du ska ro rakt över en flod från
en brygga till en annan mitt emot.
\textbf{Din rodd -- resistansen R.}
Du ror rakt fram. Vattnet motstånd mot åran är
detsamma oavsett hur fort du ror -- resistansen
bryr sig inte om frekvensen.
\textbf{Flodens sidokrafter -- reaktanserna $X_L$ och $X_C$.}
Men floden driver dig i sidled. Spolen driver dig
nedströms -- och ju snabbare du ror (hög frekvens),
desto hårdare trycker den. Kondensatorn driver dig
uppströms -- men tappar kraft ju snabbare du ror.
De motverkar varandra. Vid rätt roddhastighet tar
de ut varandra helt och du driver inte i sidled alls
-- det är resonans.
\textbf{Var du landar -- impedansen Z.}
Du når inte bryggan mitt emot -- du landar lite
nedströms eller uppströms beroende på vilken
sidokraft som dominerar. Hur långt du faktiskt
färdas från start till mål är inte summan av
roddkraften och sidokraften -- det är diagonalen.
Pythagoras, inte addition. Det är impedansen.
\end{storybox}


\begin{formulabox}
{\large$\mathbf{Z = \sqrt{R^2 + (X_L - X_C)^2}}$}\\[3mm]
{\small\normalfont\itshape
$R$ = resistans (ohm)\quad|\quad$X_L$ = induktiv reaktans\quad|\quad
$X_C$ = kapacitiv reaktans\quad|\quad$Z$ i ohm}\\[3mm]
{\normalsize$X_L = 2\pi f L$\quad|\quad$X_C = \dfrac{1}{2\pi f C}$}
\end{formulabox}

% ── Figur: impedansdiagram ────────────────────────────────
\begin{figure}[htbp]
\centering
\begin{tikzpicture}[>=stealth, font=\small]

  % Koordinater (R=3.6, X=2.4 => Z=sqrt(3.6²+2.4²)=4.33, vinkel=33.7°)
  \def\R{3.6}
  \def\X{2.4}

  % Bakgrundsruta
  \fill[gray!6, rounded corners=4pt]
    (-0.6,-0.5) rectangle (5.8,3.5);

  % Streckade hjälplinjer
  \draw[gray!35, dashed, thin] (0,\X) -- (\R,\X);
  \draw[gray!35, dashed, thin] (\R,0) -- (\R,\X);

  % Axlar
  \draw[->, thick, gray!60!black] (-0.2,0) -- (5.2,0)
    node[right, gray!60!black] {$R$};
  \draw[->, thick, gray!60!black] (0,-0.2) -- (0,3.2)
    node[above, gray!60!black] {$X$};
  \node[gray!60!black] at (-0.25,-0.25) {0};

  % R-vektor (röd)
  \draw[->, red!70!black, very thick, line cap=round]
    (0,0) -- (\R,0);
  \node[red!70!black, below=5pt] at (\R/2, 0) {$R$};

  % X-vektor (blå) -- från (R,0) uppåt
  \draw[->, blue!70!black, very thick, line cap=round]
    (\R,0) -- (\R,\X);
  \node[blue!70!black, right=6pt] at (\R, \X/2)
    {$X$\;{\footnotesize$(X_L - X_C)$}};

  % Z-vektor (grön) -- från origo till (R,X)
  \draw[->, green!50!black, very thick, line cap=round]
    (0,0) -- (\R,\X);
  \node[green!40!black, above=4pt] at (\R/2+0.1, \X/2)
    {\rotatebox{33.7}{$Z$}};

  % Vinkelmarkering θ
  \draw[gray!60, thin] (0.7,0) arc[start angle=0, end angle=33.7, radius=0.7];
  \node[gray!60!black] at (1.05,0.22) {$\theta$};

\end{tikzpicture}
\caption{Impedansdiagram för en induktiv krets ($X_L > X_C$).
$R$ (röd) är resistansen, $X = X_L - X_C$ (blå) är nettoreaktansen
och $Z$ (grön) är den resulterande impedansen.
Vinkeln $\theta = \arctan(X/R)$ -- dess storlek beror på
förhållandet mellan reaktans och resistans.}
\label{fig:impedansdiagram}
\end{figure}

\begin{bluebox}[title={\emoji{🔑} Nyckelsamband -- impedans}]
\textbf{Resistans R:} Frekvensopåverkad -- samma värde vid alla frekvenser.\\
\textbf{Induktiv reaktans $X_L$:} Ökar med frekvensen -- spolen ``bromsar'' mer vid höga frekvenser.\\
\textbf{Kapacitiv reaktans $X_C$:} Minskar med frekvensen -- kondensatorn ``öppnar sig'' vid höga frekvenser.\\
\textbf{Impedans Z:} Alltid $\geq R$ -- reaktanserna kan aldrig ta bort resistansen.
\end{bluebox}

\begin{worked}[title={Räkneexempel 1.7.3 -- Impedans i en RLC-krets}]
En krets har $R = 50\,\Omega$, $X_L = 80\,\Omega$
och $X_C = 30\,\Omega$ vid en viss frekvens.

\textbf{a) Beräkna total impedans Z:}
\[ Z = \sqrt{R^2 + (X_L - X_C)^2} = \sqrt{50^2 + (80-30)^2}
   = \sqrt{2500 + 2500} = \sqrt{5000} \approx \mathbf{70{,}7\,\Omega} \]

\textbf{b) Är kretsen kapacitiv eller induktiv?}

Eftersom $X_L > X_C$ dominerar induktansen --
kretsen är \textbf{induktiv}.

\textbf{c) Vad behöver $X_C$ bli för resonans ($Z = R$)?}

Vid resonans ska $X_L = X_C$, alltså:
\[ X_C = X_L = 80\,\Omega \]
\end{worked}

\subsubsection*{Övningsfrågor -- 1.7.3}

\begin{qblock}{3.}
\qrakna\quad En RF-krets har $R = 75\,\Omega$,
$X_L = 120\,\Omega$ och $X_C = 45\,\Omega$.

a) Beräkna total impedans $Z$.\\
b) Är kretsen kapacitiv eller induktiv? Motivera.\\
c) Om frekvensen sänks -- åt vilket håll rör sig
$X_L$ respektive $X_C$? Vad händer med $Z$?\\
d) Beräkna fasvinkel $\theta = \arctan\!\left(\dfrac{X_L - X_C}{R}\right)$.
\end{qblock}

\begin{qblock}{F.}
\qfallan\quad Lasse ser att $R = 30\,\Omega$,
$X_L = 40\,\Omega$ och $X_C = 40\,\Omega$ i en krets.
Han räknar:
\[ Z = 30 + 40 + 40 = 110\,\Omega \]
Vad gick fel -- och vad är rätt svar?
\end{qblock}

\medskip
\noindent\rule{\textwidth}{0.4pt}
{\small\itshape Försök själv innan du läser svaret.}
\medskip

\begin{worked}[title={Lösning -- Fallgropsfråga 1.7.3\,F}]
Lasse adderade impedanserna rakt av -- det fungerar
inte för AC-kretsar. $X_L$ och $X_C$ är fasförskjutna
och \emph{tar ut varandra}.

Rätt formel:
\[ Z = \sqrt{R^2 + (X_L - X_C)^2} = \sqrt{30^2 + (40-40)^2}
   = \sqrt{900 + 0} = \mathbf{30\,\Omega} \]

Kretsen är i resonans -- $X_L$ och $X_C$ tar ut
varandra helt och impedansen reduceras till bara
resistansen. Det är lägsta möjliga impedans för
denna krets.

\textbf{Lärdomen:} Reaktanser adderas inte som
resistorer. $X_L$ och $X_C$ subtraheras från
varandra -- vid resonans försvinner de ur ekvationen.
\end{worked}


% ── 1.7.4 ─────────────────────────────────────────────────
\subsection{Resonans}

Resonans är det tillstånd då en krets arbetar mest
effektivt på en specifik frekvens. Det är grunden
för hur antenner, filter och oscillatorer fungerar.

{\small\itshape I detta avsnitt avses serie-RLC för
att bygga intuition -- parallellfallet och dess
konsekvenser för SWR och impedansanpassning
behandlas i avsnitt~1.7.6.}

\begin{storybox}[{\emoji{💡} Barnet på gungan}]
Tänk på ett barn som gungar. Om du knuffar i rätt
rytm -- precis när gungan svänger tillbaka mot
dig -- byggs rörelsen upp till maximal amplitud
med minimal ansträngning.

Knuffar du för tidigt eller för sent motarbetar
du rörelsen istället. Det spelar ingen roll hur
mycket kraft du lägger på -- fel frekvens ger
dåligt resultat.

En LC-krets fungerar likadant. Vid resonansfrekvensen
``knuffar'' signalen kretsen i rätt takt. $X_L$
och $X_C$ tar ut varandra, impedansen sjunker till
bara $R$, och kretsen svarar maximalt på signalen.
\end{storybox}

\begin{formulabox}
\textbf{Resonansvillkor:}\quad{\large$\mathbf{X_L = X_C}$\quad$\Rightarrow$\quad$\mathbf{Z = R}$}\\[4mm]
\textbf{Resonansfrekvens:}\quad{\large$\mathbf{f_0 = \dfrac{1}{2\pi\sqrt{LC}}}$}\\[3mm]
{\small\normalfont\itshape
$f_0$ = resonansfrekvens (Hz)\quad|\quad$L$ = induktans (H)\quad|\quad$C$ = kapacitans (F)}
\end{formulabox}

% ── Figur: resonanskurva (serie-RLC) ─────────────────────
\begin{figure}[htbp]
\centering
\begin{tikzpicture}[>=stealth, font=\small]

  % Bakgrundsruta
  \fill[gray!6, rounded corners=4pt]
    (-0.8,-0.5) rectangle (8.4,4.8);

  % Axlar
  \draw[->, thick, gray!60!black]
    (-0.2,0) -- (7.8,0) node[right] {$f$};
  \draw[->, thick, gray!60!black]
    (0,-0.2) -- (0,4.5) node[above] {$|Z|$ ($\Omega$)};
  \node[gray!60!black] at (-0.25,-0.25) {0};

  % R-nivå (streckad röd)
  \draw[red!65!black, dashed, semithick]
    (0,0.33) -- (7.6,0.33);
  \node[red!65!black, left=4pt, font=\small\itshape] at (0,0.33) {$R$};

  % f0-markering (vertikal streckad)
  \draw[gray!50, dashed, thin]
    (3.84,0) -- (3.84,4.3);
  \node[green!45!black, below=5pt, font=\small\itshape] at (3.84,0) {$f_0$};

  % Resonanskurva (log-log skala, Q=12, symmetrisk)
  % Punkter beräknade med Python: f log från 0.25 till 4.0, Z normaliserat
  \draw[green!50!black, very thick, line cap=round, line join=round]
    plot[smooth] coordinates {
      (0.00,3.80) (0.25,3.58) (0.49,3.36) (0.74,3.14)
      (0.99,2.93) (1.23,2.72) (1.48,2.52) (1.73,2.32)
      (1.97,2.13) (2.22,1.95) (2.47,1.77) (2.71,1.60)
      (2.96,1.44) (3.21,1.29) (3.45,1.14) (3.60,1.03)
      (3.70,0.90) (3.76,0.74) (3.80,0.55) (3.84,0.33)
      (3.88,0.55) (3.92,0.74) (3.98,0.90) (4.08,1.03)
      (4.23,1.14) (4.47,1.29) (4.72,1.44) (4.97,1.60)
      (5.21,1.77) (5.46,1.95) (5.71,2.13) (5.95,2.32)
      (6.20,2.52) (6.45,2.72) (6.69,2.93) (6.94,3.14)
      (7.19,3.36) (7.43,3.58) (7.68,3.80)
    };

  % Punkt vid minimum
  \fill[green!50!black] (3.84,0.33) circle (3pt);

  % Z=R etikett
  \node[red!65!black, right=6pt, font=\scriptsize] at (3.84,0.33)
    {$|Z| = R$ \footnotesize(minimum)};

  % Kapacitiv sida
  \node[blue!60!black, font=\scriptsize, align=center] at (1.5,2.8)
    {Kapacitiv\\[2pt]$X_C > X_L$};

  % Induktiv sida
  \node[red!60!black, font=\scriptsize, align=center] at (6.2,2.8)
    {Induktiv\\[2pt]$X_L > X_C$};

\end{tikzpicture}
\caption{Impedansens belopp $|Z|$ som funktion av frekvensen för en
serie-RLC-krets. Vid resonansfrekvensen $f_0$ tar $X_L$ och $X_C$
ut varandra -- $|Z|$ sjunker till sitt minimum $R$.
OBS: för en \emph{parallell} RLC-krets gäller det omvända --
impedansen är \emph{maximal} vid resonans.}
\label{fig:resonanskurva}
\end{figure}

\begin{bluebox}[title={\emoji{🔑} Nyckelsamband -- resonans}]
\textbf{Vid resonans:} $X_L = X_C$ och $Z = R$ (minsta möjliga impedans).\\
\textbf{Under resonans} ($f < f_0$): $X_C > X_L$ -- kretsen är kapacitiv.\\
\textbf{Över resonans} ($f > f_0$): $X_L > X_C$ -- kretsen är induktiv.\\
\textbf{Antenner} konstrueras för att resonera på önskad frekvens -- SWR minimeras.\\
\textbf{Filter} väljer ut frekvenser genom att resonera selektivt.
\end{bluebox}

\begin{worked}[title={Räkneexempel 1.7.4 -- Resonansfrekvens för ett LC-filter}]
Ett LC-filter har $L = 2{,}5\,\mu\text{H}$ och
$C = 100\,\text{pF}$.

\textbf{a) Beräkna resonansfrekvensen:}
\[ f_0 = \frac{1}{2\pi\sqrt{LC}} =
   \frac{1}{2\pi\sqrt{2{,}5 \times 10^{-6} \times 100 \times 10^{-12}}} \]
\[ = \frac{1}{2\pi\sqrt{2{,}5 \times 10^{-16}}} =
   \frac{1}{2\pi \times 1{,}58 \times 10^{-8}} \approx \mathbf{10\,\text{MHz}} \]

\textbf{b) Filtret ska stämmas ner till 7~MHz.
Utan att byta spole -- vad händer med C?}

$f_0$ ska minska, alltså måste $\sqrt{LC}$ öka,
alltså måste $C$ öka. Parallellkoppling av en extra
kondensator ökar den totala kapacitansen -- seriekoppling
hade minskat den.
\end{worked}

\subsubsection*{Övningsfrågor -- 1.7.4}

\begin{qblock}{4.}
\qrakna\quad Ett LC-filter ska resonera på
3,5~MHz (80-metersbandet). Spolen är $L = 5\,\mu\text{H}$.

a) Beräkna nödvändig kapacitans $C$.\\
b) Filtret resonerar nu för högt -- frekvensen
är 4~MHz istället för 3,5~MHz. Utan att byta
spole -- hur justerar du $C$ och i vilken
koppling (serie/parallell) lägger du till mer
kapacitans?\\
c) Vad händer med $X_L$ och $X_C$ när frekvensen
stiger från under till över resonans?
\end{qblock}

\begin{qblock}{F.}
\qfallan\quad Anna har ett LC-filter som resonerar
på 7~MHz. Hon vill höja resonansfrekvensen.
En kamrat säger: ``Lägg till en kondensator parallellt
med den du har -- mer kapacitans ger högre frekvens.''

Stämmer det?
\end{qblock}

\medskip
\noindent\rule{\textwidth}{0.4pt}
{\small\itshape Försök själv innan du läser svaret.}
\medskip

\begin{worked}[title={Lösning -- Fallgropsfråga 1.7.4\,F}]
Fel. Mer kapacitans ger \emph{lägre} resonansfrekvens
-- raka motsatsen.

Formeln visar det direkt:
\[ f_0 = \frac{1}{2\pi\sqrt{LC}} \]

Ökar $C$ så ökar nämnaren, vilket \emph{minskar} $f_0$.

För att \emph{höja} resonansfrekvensen ska $L$ eller
$C$ minskas -- inte ökas. Alternativt: byt ut
kondensatorn mot en med lägre kapacitans.

\textbf{Lärdomen:} Resonansfrekvensen och $LC$-produkten
rör sig i motsatta riktningar. Mer $L$ eller $C$
sänker alltid $f_0$.
\end{worked}

% ── 1.7.5 ─────────────────────────────────────────────────
\subsection{Fasförskjutning -- ELI the ICE man}

I en ren resistiv krets är spänning och ström i fas --
de når sina max- och nollvärden exakt samtidigt.
I kretsar med spolar eller kondensatorer är de
däremot förskjutna i tid mot varandra.
Det kallas fasförskjutning.

\begin{storybox}[{\emoji{💡} Tågen som aldrig kommer samtidigt}]
Föreställ dig en järnvägsstation med två spår.
På det ena spåret anländer spänningens tåg,
på det andra strömmens tåg. Vid ett motstånd
anländer de exakt samtidigt -- alltid i fas.

Vid en spole är spänningens tåg alltid 15 minuter
före strömmens. Spänningen trycker på -- men spolen
``bryr sig inte'' om strömmen förrän lite senare.
Det kallas att spänningen leder strömmen.

Vid en kondensator är det tvärtom: strömmen rusar
in och börjar ladda upp plattorna \emph{innan}
spänningen hunnit byggas upp. Strömmen leder spänningen.

Minnesregeln \textbf{ELI the ICE man} sammanfattar
detta en gång för alla.
\end{storybox}

\begin{bluebox}[title={\emoji{🔑} ELI the ICE man}]
\textbf{ELI:} I en spole (L) leder \textbf{spänningen} \textbf{strömmen}.\\
\phantom{\textbf{ELI:}} Spänning (E) före ström (I) i induktans (L).\\[2mm]
\textbf{ICE:} I en kondensator (C) leder \textbf{strömmen} \textbf{spänningen}.\\
\phantom{\textbf{ICE:}} Ström (I) före spänning (E) i kapacitans (C).\\[3mm]

\begin{center}
\begin{tabular}{L{2.8cm}L{4.0cm}L{4.0cm}}
\rowcolor{greydark}
\tblhdr{Komponent} & \tblhdr{Fasförhållande} & \tblhdr{Vinkel $\phi$} \\
Resistor & I och U i fas & $0°$ \\
\rowcolor{greylight}
Spole (L) & Spänning leder ström & $+90°$ \\
Kondensator (C) & Ström leder spänning & $-90°$ \\
\end{tabular}
\end{center}
\smallskip
{\small\itshape Tecknet på $\phi = \arctan\!\bigl((X_L - X_C)/R\bigr)$:
positivt = induktiv krets, negativt = kapacitiv krets.}
\end{bluebox}

% ── Figur: fasordiagram ──────────────────────────────────
\begin{figure}[htbp]
\centering
\begin{tikzpicture}[>=stealth, font=\small]

  % Bakgrundsruta
  \fill[gray!6, rounded corners=4pt]
    (-0.4,-0.4) rectangle (13.4,5.8);

  % ── Panel 1: Resistor (φ=0°) ─────────────────────────
  \begin{scope}[xshift=1.8cm, yshift=2.8cm]
    \draw[->, gray!45] (-1.2,0) -- (1.8,0) node[right, font=\scriptsize, gray!60] {$R$};
    \draw[->, gray!45] (0,-1.5) -- (0,1.8) node[above, font=\scriptsize, gray!60] {$X$};
    \draw[->, red!70!black, very thick] (0,0) -- (1.4,0);
    \node[red!70!black, above=3pt, font=\small\itshape] at (0.7,0) {$Z$};
    \node[gray!60, font=\scriptsize] at (0.7,-0.25) {$\phi=0°$};
    \node[font=\small\bfseries] at (0,1.95) {Resistor};
    \node[font=\scriptsize] at (0,-1.7) {I och U i fas};
    \node[red!70!black, font=\small\bfseries] at (0,-2.1) {$\phi = 0°$};
  \end{scope}

  % Skiljelinje 1
  \draw[gray!30] (3.8,0.2) -- (3.8,5.6);

  % ── Panel 2: Spole L (φ=+90°) ────────────────────────
  \begin{scope}[xshift=6.4cm, yshift=2.8cm]
    \draw[->, gray!45] (-1.2,0) -- (1.8,0) node[right, font=\scriptsize, gray!60] {$R$};
    \draw[->, gray!45] (0,-1.5) -- (0,1.8) node[above, font=\scriptsize, gray!60] {$X$};
    \draw[->, blue!70!black, very thick] (0,0) -- (0,1.4);
    \node[blue!70!black, right=4pt, font=\small\itshape] at (0,0.7) {$|Z|=X_L$};
    \draw[gray!55, thin] (0.4,0) arc[start angle=0, end angle=90, radius=0.4];
    \node[gray!60, font=\scriptsize\itshape] at (0.32,0.22) {$\phi$};
    \node[blue!70!black, font=\normalsize\bfseries] at (0,2.15) {ELI};
    \node[font=\small\bfseries] at (0,1.72) {Spole (L)};
    \node[font=\scriptsize] at (0,-1.7) {Spänning leder ström};
    \node[blue!70!black, font=\small\bfseries] at (0,-2.1) {$\phi = +90°$};
  \end{scope}

  % Skiljelinje 2
  \draw[gray!30] (9.0,0.2) -- (9.0,5.6);

  % ── Panel 3: Kondensator C (φ=−90°) ──────────────────
  \begin{scope}[xshift=11.6cm, yshift=2.8cm]
    \draw[->, gray!45] (-1.2,0) -- (1.8,0) node[right, font=\scriptsize, gray!60] {$R$};
    \draw[->, gray!45] (0,-1.5) -- (0,1.8) node[above, font=\scriptsize, gray!60] {$X$};
    \draw[->, orange!80!red, very thick] (0,0) -- (0,-1.4);
    \node[orange!80!red, right=4pt, font=\small\itshape] at (0,-0.7) {$|Z|=X_C$};
    \draw[gray!55, thin] (0.4,0) arc[start angle=0, end angle=-90, radius=0.4];
    \node[gray!60, font=\scriptsize\itshape] at (0.32,-0.22) {$\phi$};
    \node[orange!80!red, font=\normalsize\bfseries] at (0,2.15) {ICE};
    \node[font=\small\bfseries] at (0,1.72) {Kondensator (C)};
    \node[font=\scriptsize] at (0,-1.7) {Ström leder spänning};
    \node[orange!80!red, font=\small\bfseries] at (0,-2.1) {$\phi = -90°$};
  \end{scope}

\end{tikzpicture}
\caption{Fasvinkeln $\phi = \arg(Z)$ i $R$--$X$-planet.
Resistorn ger $\phi=0°$, spolen $+90°$ (induktiv, $X>0$, vektor uppåt)
och kondensatorn $-90°$ (kapacitiv, $X<0$, vektor nedåt).
I en blandad krets hamnar $\phi$ mellan dessa extremvärden.}
\label{fig:fasordiagram}
\end{figure}

\begin{worked}[title={Räkneexempel 1.7.5 -- Fasförhållande i en blandad krets}]
En krets har $R = 50\,\Omega$ och $X_L = 50\,\Omega$,
ingen kondensator.

\textbf{Vad är impedansen och fasförskjutningen?}
\[ Z = \sqrt{50^2 + 50^2} = \sqrt{5000} \approx 70{,}7\,\Omega \]

Fasvinkel:
\[ \phi = \arctan\!\left(\frac{X_L}{R}\right) = \arctan\!\left(\frac{50}{50}\right) = \arctan(1) = 45° \]

Spänningen leder strömmen med 45° -- kretsen är
induktiv men inte rent induktiv (det vore 90° om
$R = 0$).
\end{worked}

\subsubsection*{Övningsfrågor -- 1.7.5}

\begin{qblock}{5.}
\qprov\quad Förklara minnesregeln ``ELI the ICE man''
med egna ord.

a) Vad säger den om spolar och kondensatorer?\\
b) Vad händer med fasförskjutningen vid resonans
när $X_L = X_C$?\\
c) Varför är fasförskjutning viktigt att förstå
vid effektberäkningar i AC-kretsar?
\end{qblock}

\begin{qblock}{F.}
\qfallan\quad Björn minns ELI the ICE man och
konstaterar: ``I en kondensator leder ström
spänningen med 90°. Alltså är spänningen 90°
\emph{efter} strömmen -- spänningen är försenad.''

Han applicerar detta på en AC-krets och säger att
det innebär att kondensatorn \emph{bromsar} strömmen
mer vid höga frekvenser.

Vad stämmer och vad stämmer inte?
\end{qblock}

\medskip
\noindent\rule{\textwidth}{0.4pt}
{\small\itshape Försök själv innan du läser svaret.}
\medskip

\begin{worked}[title={Lösning -- Fallgropsfråga 1.7.5\,F}]
Fasförhållandet har han rätt om -- ström leder
spänning i en kondensator med 90°, vilket är
detsamma som att spänningen är 90° försenad. Det stämmer.

Men slutsatsen om bromsning är \textbf{fel -- tvärtom}.

Kondensatorns reaktans är:
\[ X_C = \frac{1}{2\pi f C} \]

$X_C$ \emph{minskar} när frekvensen ökar --
kondensatorn erbjuder \emph{mindre} motstånd vid
höga frekvenser, inte mer. Vid tillräckligt hög
frekvens är $X_C \approx 0$ och kondensatorn
fungerar nästan som en kortslutning.

Det är spolen (induktansen) som bromsar mer vid
höga frekvenser: $X_L = 2\pi f L$ ökar med $f$.

\textbf{Lärdomen:} Fasförskjutning och reaktansens
frekvensberoende är skilda begrepp. Kondensatorn
leder ström men \emph{minskar} sitt motstånd med frekvensen.
\end{worked}


% ── 1.7.6 ─────────────────────────────────────────────────
\subsection{Standard-impedanser och varför just 50\,$\Omega$}

Inom RF-teknik och amatörradio stöter du alltid på
samma impedansvärden -- 50~$\Omega$ för sändare och
antenner, 75~$\Omega$ för TV. Det är inte slumpmässigt.

\begin{storybox}[{\emoji{💡} Kompromissen som blev världsstandard}]
Under andra världskriget behövde militären ett
gemensamt impedansvärde för all RF-utrustning.
Ingenjörerna undersökte koaxkabelns egenskaper:

Vid \emph{lägst dämpning} i en luftfylld koax
är optimal impedans $\approx 77\,\Omega$.
Vid \emph{högst effektkapacitet} (maximal
spänningshållfasthet kombinerat med strömkapacitet)
är optimal impedans $\approx 30\,\Omega$.

Kompromissen -- det geometriska medelvärdet:
\[ Z_{\text{kompromiss}} = \sqrt{77 \times 30} \approx 48\,\Omega \approx \mathbf{50\,\Omega} \]

Valet etablerades tidigt i RF-industrin och
förstärktes av militär och civil standardisering --
och är nu universellt inom amatörradio och RF-teknik
världen över.
\end{storybox}

% ── Figur: kompromissdiagram ─────────────────────────────
\begin{figure}[htbp]
\centering
\begin{tikzpicture}[>=stealth, font=\small]

  % Ytterbakgrund
  \fill[gray!6, rounded corners=6pt]
    (-0.4,-1.8) rectangle (13.8,6.4);

  % Diagrambakgrund
  \fill[white, rounded corners=4pt]
    (0.4,0.6) rectangle (13.2,6.1);
  \draw[gray!25, rounded corners=4pt]
    (0.4,0.6) rectangle (13.2,6.1);

  % Axlar
  % x: Z=20..90 => 0.6..12.8 cm, skala 0.1743 cm/Ω
  % y: relativ fördel 0..1 => 1.0..5.8 cm
  \draw[->, thick, gray!60!black]
    (0.6,1.0) -- (13.2,1.0) node[right, font=\small\itshape] {$Z_0\;(\Omega)$};
  \draw[->, thick, gray!60!black]
    (0.6,1.0) -- (0.6,6.1) node[above, font=\scriptsize, gray!65] {Relativ fördel};

  % X-ticks
  \foreach \xpos/\lbl in {
    0.60/20, 2.34/30, 4.09/40, 5.83/50,
    7.57/60, 9.31/70, 10.53/77, 12.80/90}{
    \draw[gray!45] (\xpos,1.0) -- (\xpos,0.85);
    \node[font=\scriptsize, gray!60] at (\xpos,0.65) {\lbl};
  }

  % Effekttålighetskurva (röd), topp vid ~30 Ω
  \draw[red!65!black, very thick, line cap=round]
    plot[smooth] coordinates {
      (0.60,1.20) (1.20,1.80) (2.00,3.60)
      (2.34,4.80) (2.80,4.60) (3.60,3.60)
      (5.00,2.40) (7.00,1.65) (9.50,1.25)
      (12.80,1.05)
    };

  % Låg dämpningskurva (blå), topp vid ~77 Ω
  \draw[blue!65!black, very thick, line cap=round]
    plot[smooth] coordinates {
      (0.60,1.10) (1.50,1.30) (3.00,1.80)
      (5.00,2.90) (7.00,3.90) (9.00,4.60)
      (10.53,4.85) (11.50,4.60) (12.80,3.80)
    };

  % Vertikala markeringslinjer
  \draw[red!50!black, dashed, semithick]
    (2.34,1.0) -- (2.34,4.80);
  \draw[green!45!black, dashed, semithick]
    (5.83,1.0) -- (5.83,5.80);
  \draw[blue!50!black, dashed, semithick]
    (10.53,1.0) -- (10.53,4.85);

  % Toppunkter
  \fill[red!65!black]  (2.34,4.80) circle (3.5pt);
  \fill[blue!65!black] (10.53,4.85) circle (3.5pt);
  \fill[green!50!black] (5.83,5.65) circle (3.5pt);

  % Etiketter vid toppar
  \node[red!70!black, font=\small, above=2pt] at (2.34,4.80)
    {$\approx 30\,\Omega$};
  \node[blue!70!black, font=\small, above=2pt] at (10.53,4.85)
    {$\approx 77\,\Omega$};
  \node[green!50!black, font=\small\bfseries, above=2pt] at (5.83,5.65)
    {$\approx 50\,\Omega$};

  % Kurvtexter
  \node[red!65!black, font=\scriptsize, align=left] at (1.0,3.20)
    {Max\\effekt-\\tålighet};
  \node[blue!65!black, font=\scriptsize, align=right] at (12.4,3.50)
    {Lägst\\dämpning};

  % Formelruta (grön, rymlig)
  \fill[green!8, rounded corners=5pt]
    (5.20,2.20) rectangle (9.80,4.80);
  \draw[green!40!black, rounded corners=5pt]
    (5.20,2.20) rectangle (9.80,4.80);
  \node[green!40!black, font=\small\bfseries] at (7.50,4.45)
    {Geometrisk kompromiss:};
  \node[green!35!black, font=\normalsize] at (7.50,3.75)
    {$\sqrt{77 \times 30} \approx 48\,\Omega$};
  \node[green!35!black, font=\large\bfseries] at (7.50,3.10)
    {$\approx \mathbf{50\,\Omega}$};
  \node[green!30!black, font=\scriptsize, align=center] at (7.50,2.55)
    {Etablerat i RF-industrin och};
  \node[green!30!black, font=\scriptsize, align=center] at (7.50,2.30)
    {förstärkt genom standardisering};

  % Legend
  \draw[red!65!black, very thick]  (1.0,0.25) -- (2.2,0.25);
  \node[font=\scriptsize] at (3.3,0.25) {Effekttålighet};
  \draw[blue!65!black, very thick]  (5.0,0.25) -- (6.2,0.25);
  \node[font=\scriptsize] at (7.4,0.25) {Låg dämpning};
  \draw[green!50!black, very thick, dashed] (9.0,0.25) -- (10.2,0.25);
  \node[font=\scriptsize] at (11.5,0.25) {50\,$\Omega$ kompromiss};

  % Standardimpedansrutor längst ner
  \fill[blue!8, rounded corners=4pt]
    (0.30,-1.65) rectangle (3.30,-0.15);
  \draw[blue!30, rounded corners=4pt]
    (0.30,-1.65) rectangle (3.30,-0.15);
  \node[blue!70!black, font=\large\bfseries] at (1.80,-0.70) {50\,$\Omega$};
  \node[blue!60!black, font=\scriptsize, align=center] at (1.80,-1.20)
    {Sändare, RF\\amatörradio};

  \fill[purple!8, rounded corners=4pt]
    (3.60,-1.65) rectangle (6.60,-0.15);
  \draw[purple!25, rounded corners=4pt]
    (3.60,-1.65) rectangle (6.60,-0.15);
  \node[purple!70!black, font=\large\bfseries] at (5.10,-0.70) {75\,$\Omega$};
  \node[purple!55!black, font=\scriptsize, align=center] at (5.10,-1.20)
    {TV, kabel-TV\\video};

  \fill[orange!8, rounded corners=4pt]
    (6.90,-1.65) rectangle (9.90,-0.15);
  \draw[orange!30, rounded corners=4pt]
    (6.90,-1.65) rectangle (9.90,-0.15);
  \node[orange!70!black, font=\large\bfseries] at (8.40,-0.70) {300\,$\Omega$};
  \node[orange!55!black, font=\scriptsize, align=center] at (8.40,-1.20)
    {Twin-lead\\foldad dipol};

  \fill[green!8, rounded corners=4pt]
    (10.20,-1.65) rectangle (13.20,-0.15);
  \draw[green!25, rounded corners=4pt]
    (10.20,-1.65) rectangle (13.20,-0.15);
  \node[green!50!black, font=\large\bfseries] at (11.70,-0.70) {600\,$\Omega$};
  \node[green!40!black, font=\scriptsize, align=center] at (11.70,-1.20)
    {Äldre tele\\balanserat audio};

\end{tikzpicture}
\caption{Schematisk/konceptuell kurva över koaxkabelns egenskaper som
funktion av karakteristisk impedans $Z_0$. Lägst dämpning uppnås
nära 77\,$\Omega$, högst effekttålighet nära 30\,$\Omega$.
Det geometriska medelvärdet $\sqrt{77 \times 30} \approx 50\,\Omega$
utgör den praktiska kompromissen. Kurvorna är principiella --
exakt form beror på kabelkonstruktion och dielektrikum.}
\label{fig:impedanskompromiss}
\end{figure}

\begin{center}
\begin{tabular}{L{2.5cm}L{8.0cm}}
\rowcolor{greydark}
\tblhdr{Impedans} & \tblhdr{Användning} \\
50~$\Omega$ & Amatörradio, sändare, koaxkabel, mätutrustning, mobiltelefoni \\
\rowcolor{greylight}
75~$\Omega$ & TV-antenner, kabel-TV, videosystem \\
300~$\Omega$ & Bandkabel (twin-lead) för TV, foldad dipol \\
\rowcolor{greylight}
600~$\Omega$ & Äldre telefonlinjer, balanserade audiosystem \\
\end{tabular}
\end{center}

\begin{bluebox}[title={\emoji{🔑} Nyckelsamband -- impedanser}]
\textbf{Målimpedans i amatörradio är normalt 50~$\Omega$:}
sändare, koax (RG-58, RG-213), antenningång.\\
\textbf{Undvik 75~$\Omega$ TV-kabel} i 50~$\Omega$-sändarsystem om du inte vet vad
du gör -- impedansfelet ger mätbara reflektioner.\\
\textbf{Foldad dipol} ligger ofta kring 200--300~$\Omega$ beroende på konstruktion
och omgivning. Teoretiskt ger 300\,$\Omega \rightarrow 50\,\Omega$ ett
impedansförhållande 6:1; i praktiken används ofta 4:1 balun med
ytterligare anpassning eller tuner.
\end{bluebox}

\subsubsection*{Övningsfrågor -- 1.7.6}

\begin{qblock}{6.}
\qgrund\quad

a) Varför är 50~$\Omega$ ett kompromissvärde --
vad kompromissas det mellan?\\
b) Du hittar en rulle koaxkabel märkt ``75~$\Omega$''
i förrådet. Kan du använda den till din
amatörradiostation istället för 50~$\Omega$-kabel?
Vad händer?\\
c) En foldad dipol har impedansen $\approx 300\,\Omega$.
Du ska mata den med 50~$\Omega$ koax.
Vilket impedansförhållande behöver en balun
kompensera?
\end{qblock}

\begin{qblock}{F.}
\qfallan\quad Erik tittar på en rulle 75~$\Omega$-kabel
och säger: ``Skillnaden mot 50~$\Omega$ är bara
25~$\Omega$ -- det är en liten avvikelse, det
spelar knappast någon roll för SWR.''

Stämmer det?
\end{qblock}

\medskip
\noindent\rule{\textwidth}{0.4pt}
{\small\itshape Försök själv innan du läser svaret.}
\medskip

\begin{worked}[title={Lösning -- Fallgropsfråga 1.7.6\,F}]
Nej -- Erik tänker i differens, men SWR beror på
\emph{kvoten} mellan impedanserna.

I allmänhet gäller:
\[
\Gamma=\frac{Z_L-Z_0}{Z_L+Z_0},
\qquad
\mathrm{SWR}=\frac{1+|\Gamma|}{1-|\Gamma|}
\]

För rent resistiva värden $Z_L=75\,\Omega$, $Z_0=50\,\Omega$:
\[
\Gamma=\frac{75-50}{75+50}=\frac{25}{125}=0{,}2
\]
\[
\mathrm{SWR}=\frac{1+0{,}2}{1-0{,}2}
=\frac{1{,}2}{0{,}8}=1{,}5
\]

SWR~1,5 är mätbart och kan i känsliga system ge
märkbara reflektioner och effektförluster -- även
om det i många praktiska situationer (mottagning,
låg effekt, kort kabel) är acceptabelt.

\textbf{Lärdomen:} Impedansmismatch mäts alltid i
kvot, inte differens. En ``liten'' skillnad i ohm
kan ge ett betydande SWR beroende på referensimpedansen.
\end{worked}



% ── 1.7.7 ─────────────────────────────────────────────────
\subsection{Skineffekt}

Vid likström flödar strömmen jämnt fördelad
genom hela ledarens tvärsnittsarea. Vid växelström
-- särskilt vid höga frekvenser -- trängs strömmen
ut mot ledarens yta. Det kallas skineffekt.

\begin{storybox}[{\emoji{💡} Passagerarna på det överfulla tåget}]
Föreställ dig ett runt tåg (tänk ett långt rör)
fullt med passagerare -- elektroner -- som ska ta
sig från A till B. Vid låg hastighet sprider de
sig jämnt, även i mitten av vagnen.

Men när tåget går allt fortare börjar passagerarna
trängas mot fönstren -- mot ytterväggarna. Mitten
av tåget töms. Till slut är det bara passagerare
längs den yttersta centimetern av väggen.

Det är precis vad som händer med elektroner vid
höga frekvenser. Mitten av ledaren används inte
alls -- bara det tunna ytskiktet, ``skinnet''.
Det gör att ledarens effektiva tvärsnittsarea
minskar och resistansen ökar.
\end{storybox}

\begin{bluebox}[title={\emoji{🔑} Nyckelsamband -- skineffekt}]
\textbf{Skindjupet} $\delta$ beräknas som:
\[ \delta = \sqrt{\frac{\rho}{\pi f \mu}} \]
där $\rho$ = resistivitet\,($\Omega$m), $f$ = frekvens\,(Hz),
$\mu = \mu_0 \approx 4\pi \times 10^{-7}$\,H/m för icke-magnetiska
material som koppar.
Skindjupet minskar med $\sqrt{f}$ -- dubbel frekvens
ger $\sqrt{2}$ gånger tunnare strömskikt.
{\small\itshape (Ekvivalent med den fullständiga formen
$\delta = \sqrt{2\rho/(\omega\mu)}$ där $\omega = 2\pi f$.)}\\[3mm]

\begin{center}
\begin{tabular}{L{3.0cm}L{4.0cm}}
\rowcolor{greydark}
\tblhdr{Frekvens} & \tblhdr{Skindjup (koppar)} \\
1\,MHz   & $\approx 66\,\mu\text{m}$ \\
\rowcolor{greylight}
14\,MHz  & $\approx 18\,\mu\text{m}$ \\
144\,MHz & $\approx 6\,\mu\text{m}$ \\
\rowcolor{greylight}
430\,MHz & $\approx 3\,\mu\text{m}$ \\
\end{tabular}
\end{center}
\smallskip
\textbf{Kopparrör} fungerar lika bra som massiv koppar
för VHF/UHF-antenner -- strömmen går ändå bara i ytan.\\
\textbf{Silverpläterade} kontakter och ledare ger lägre
förlust vid VHF/UHF -- silvrets bättre ledning utnyttjas
just i det tunna ytskiktet.\\
\textbf{Kabelförluster} ökar med frekvensen -- en koax
som fungerar bra på HF kan ha oacceptabla förluster på UHF.
\end{bluebox}

% ── Figur: skindjup vs frekvens ──────────────────────────
\begin{figure}[htbp]
\centering
\begin{tikzpicture}[>=stealth, font=\small]

  % Bakgrundsruta
  \fill[gray!6, rounded corners=4pt]
    (-0.3,-0.4) rectangle (9.8,5.5);

  % Diagrambakgrund
  \fill[white, rounded corners=3pt] (0.5,0.3) rectangle (9.2,5.1);
  \draw[gray!20, rounded corners=3pt] (0.5,0.3) rectangle (9.2,5.1);

  % Rutnätslinjer (korrekt log-placering)
  % Vertikalt vid 1, 10, 100 MHz
  \foreach \xg in {2.625, 4.750, 6.875}{
    \draw[gray!20] (\xg,0.5) -- (\xg,5.1);
  }
  % Horisontellt vid 3, 6, 18, 65 µm (korrekt log-y)
  \draw[gray!20] (0.5,0.905) -- (9.2,0.905);
  \draw[gray!20] (0.5,1.597) -- (9.2,1.597);
  \draw[gray!20] (0.5,2.695) -- (9.2,2.695);
  \draw[gray!20] (0.5,3.977) -- (9.2,3.977);

  % Axlar
  \draw[->, thick, gray!60!black]
    (0.5,0.5) -- (9.4,0.5)
    node[right, font=\small\itshape] {$f$ (MHz, log)};
  \draw[->, thick, gray!60!black]
    (0.5,0.5) -- (0.5,5.3)
    node[above, font=\small\itshape] {$\delta$ ($\mu$m, log)};

  % X-axel ticks och etiketter
  \foreach \xpos/\lbl in {
    0.500/0{,}1, 2.625/1, 4.750/10, 6.875/100, 9.000/1000}{
    \draw[gray!50] (\xpos,0.5) -- (\xpos,0.35);
    \node[font=\scriptsize, gray!60] at (\xpos,0.18) {\lbl};
  }

  % Y-axel ticks och etiketter (korrekt log-placering)
  % delta=2 => y=0.500, 3=>0.905, 6=>1.597, 18=>2.695, 65=>3.977, 200=>5.100
  \foreach \ypos/\lbl in {
    0.500/2, 0.905/3, 1.597/6, 2.695/18, 3.977/65, 5.100/200}{
    \draw[gray!50] (0.5,\ypos) -- (0.35,\ypos);
    \node[font=\scriptsize, gray!60, left=2pt] at (0.5,\ypos) {\lbl};
  }

  % ── Skindjupskurva: rät linje på log-log ─────────────
  \draw[green!50!black, very thick, line cap=round]
    (0.500,5.100) -- (9.000,0.500);

  % ── Markeringspunkter (korrekt log-log position) ──────
  % 1 MHz:   x=2.625, y=3.981
  % 14 MHz:  x=5.061, y=2.663
  % 144 MHz: x=7.212, y=1.499
  % 430 MHz: x=8.221, y=0.952
  \fill[red!65!black] (2.625,3.981) circle (3pt);
  \fill[red!65!black] (5.061,2.663) circle (3pt);
  \fill[red!65!black] (7.212,1.499) circle (3pt);
  \fill[red!65!black] (8.221,0.952) circle (3pt);

  % Etiketter vid markeringar
  \node[red!65!black, font=\scriptsize, right=4pt] at (2.625,3.981)
    {1\,MHz · 65\,$\mu$m};
  \node[red!65!black, font=\scriptsize, right=4pt] at (5.061,2.663)
    {14\,MHz · 18\,$\mu$m};
  \node[red!65!black, font=\scriptsize, right=4pt] at (7.212,1.499)
    {144\,MHz · 6\,$\mu$m};
  \node[red!65!black, font=\scriptsize, right=4pt] at (8.221,0.952)
    {430\,MHz · 3\,$\mu$m};

  % δ ∝ 1/√f etikett
  \node[green!40!black, font=\small\itshape] at (3.2,2.0)
    {$\delta \propto 1/\!\sqrt{f}$};

\end{tikzpicture}
\caption{Skindjup $\delta$ som funktion av frekvens $f$ för koppar,
beräknat med $\delta = \sqrt{\rho/(\pi f \mu)}$.
$\delta$ minskar med $\sqrt{f}$ -- vid högre radiofrekvenser
används bara ett allt tunnare ytskikt av ledaren.
På log-log skala syns beroendet som en rät linje.}
\label{fig:skineffekt}
\end{figure}

\begin{worked}[title={Räkneexempel 1.7.7 -- Skindjup i koppar vid 14\,MHz}]
Koppar har resistivitet $\rho = 1{,}68 \times 10^{-8}\,\Omega\text{m}$
och permeabilitet $\mu = \mu_0 = 4\pi \times 10^{-7}\,\text{H/m}$.

\[ \delta = \sqrt{\frac{\rho}{\pi f \mu}}
   = \sqrt{\frac{1{,}68 \times 10^{-8}}{\pi \times 14 \times 10^6
     \times 4\pi \times 10^{-7}}} \]

Nämnaren: $\pi \times 14 \times 10^6 \times 4\pi \times 10^{-7}
= 4\pi^2 \times 1{,}4 \approx 55{,}3$

\[ \delta = \sqrt{\frac{1{,}68 \times 10^{-8}}{55{,}3}}
   = \sqrt{3{,}04 \times 10^{-10}} \approx \mathbf{17{,}4\,\mu\text{m} \approx 18\,\mu\text{m}} \]

En kopparledare som är tjockare än $\sim$0,1\,mm har
vid 14\,MHz sin ström koncentrerad till
ett ytskikt tunnare än ett mänskligt hårstrå.
\end{worked}

\subsubsection*{Övningsfrågor -- 1.7.7}

\begin{qblock}{7.}
\qgrund\quad

a) Varför kan ett kopparrör ersätta massiv koppar
i en UHF-antenn utan att prestanda försämras?\\
b) Varför är silverpläterade kontakter motiverade
i en VHF/UHF-station men knappast vid DC-tillämpningar?\\
c) Du har 30~m RG-58-kabel till din
HF-station (14~MHz) och funderar på att även
använda den till UHF-trafik (430~MHz).
Vad är problemet?
\end{qblock}

\begin{qblock}{F.}
\qfallan\quad Magnus resonerar: ``Skineffekten
ökar ledarens resistans vid höga frekvenser.
Alltså borde en tjockare ledare alltid ge
lägre förlust -- mer koppar är alltid bättre.''

Stämmer det vid VHF/UHF?
\end{qblock}

\medskip
\noindent\rule{\textwidth}{0.4pt}
{\small\itshape Försök själv innan du läser svaret.}
\medskip

\begin{worked}[title={Lösning -- Fallgropsfråga 1.7.7\,F}]
Inte alltid. Vid VHF/UHF är skindjupet bara några
mikrometer. En 2\,mm-ledare och en 10\,mm-ledare
har strömmen fördelad i exakt samma ytskikt --
inget av det extra kopparet i mitten bidrar till
strömtransporten.

Det som spelar roll vid VHF/UHF är ledarens
\textbf{ytarea} -- omkrets gånger längd -- inte
tvärsnittsarean. En tjockare ledare hjälper
därför bara för att en större diameter ger
en större omkrets och därmed mer yta för strömmen
att flöda i. \textbf{Ytans beskaffenhet} spelar
också stor roll: silverpläterat ger lägre förlust
än ren koppar eftersom silver leder bättre i
just det tunna ytskikt som faktiskt används.

\textbf{Lärdomen:} Vid höga frekvenser är yta
viktigare än volym. Kopparrör är lika bra som
massiv koppar om ytarean är densamma.
\end{worked}


% ── 1.7.8 ─────────────────────────────────────────────────
\subsection{SWR -- Standing Wave Ratio}

SWR mäter hur väl sändarens impedans matchar 
antennens impedans. En dålig matchning innebär 
att en del av effekten reflekteras tillbaka mot 
sändaren -- bortslösad effekt som dessutom kan 
skada slutsteget.

\begin{storybox}[{\emoji{💡} Ekot i bergsdalen}]
Föreställ dig att du ropar in i en bergsdal. 
Om dalen absorberar all din röst -- perfekt 
anpassning -- hörs inget eko. Allt roper vidare 
i rätt riktning.

Men om bergväggen är blank och hård studsar 
rösten tillbaka mot dig. Du hör ekot -- och 
den energi som ekot bär kan inte längre göra 
nytta i rätt riktning.

En antenn med fel impedans är som bergväggen. 
Sändareffekten studsar tillbaka längs koaxkabeln. 
SWR-metern mäter hur starkt detta eko är.
\end{storybox}

% ── SWR-diagram (TikZ) ────────────────────────────────────
% Fysikaliskt korrekt envelope: V+·sqrt(1+Γ²+2Γcos(2kx))
% Γ=0.5 → SWR=3:1  |  antinod=82.5 px  |  nod=27.5 px (≠0)
% Canvas 16×9.6 cm  (SVG 800×480, skala 0.02 cm/px, y-speglad)
\begin{figure}[h]
\centering
\begin{tikzpicture}[x=1cm, y=1cm, font=\sffamily\small]

  % ── Färger ────────────────────────────────────────────
  \colorlet{panelblue}{blue!4!white}
  \colorlet{panelred}{red!3!white}
  \colorlet{strokeblue}{blue!25!white}
  \colorlet{strokered}{red!25!white}
  \colorlet{coaxblue}{blue!15!white}
  \colorlet{coaxbluedark}{blue!35!white}
  \colorlet{coaxred}{red!15!white}
  \colorlet{coaxreddark}{red!30!white}
  \colorlet{wavegreen}{green!55!black}
  \colorlet{wavered}{red!75!black}
  \colorlet{nodeorange}{orange!70!black}
  \colorlet{titleblue}{blue!65!black}
  \colorlet{titlered}{red!65!black}
  \colorlet{subtitleblue}{blue!50!black}
  \colorlet{subtitlered}{red!40!black}
  \colorlet{infoblue}{blue!7!white}
  \colorlet{inforef}{red!5!white}
  \colorlet{badgered}{red!8!white}
  \colorlet{badgegrn}{green!8!white}
  \colorlet{bottombar}{blue!5!white}
  \colorlet{formblue}{blue!60!black}

  % ── Panelramar ────────────────────────────────────────
  \filldraw[fill=panelblue, draw=strokeblue, line width=0.4pt, rounded corners=3pt]
    (0.4,0.9) rectangle (7.6,9.2);
  \filldraw[fill=panelred, draw=strokered, line width=0.4pt, rounded corners=3pt]
    (8.4,0.9) rectangle (15.6,9.2);

  % ── Bottenmeny ────────────────────────────────────────
  \filldraw[fill=bottombar, draw=strokeblue!60, line width=0.4pt, rounded corners=2pt]
    (0.4,0.16) rectangle (15.6,0.72);
  \node[anchor=center, text=formblue, font=\sffamily\fontsize{5.5}{6.5}\selectfont]
    at (8.0,0.44)
    {$\mathrm{SWR}=\tfrac{1{+}|\Gamma|}{1{-}|\Gamma|}$
     \enspace$|$\enspace
     $\Gamma=\tfrac{Z_L-Z_0}{Z_L+Z_0}$
     \enspace$|$\enspace
     $P_\text{ref}=|\Gamma|^2\times P_\text{in}$
     \enspace$|$\enspace
     $\mathrm{SWR}\approx Z_\text{hög}/Z_\text{låg}$
     \enspace\textit{\small(förenklad, resistiv last)}};

  % ════════════════════════════════════════════════════
  % VÄNSTER PANEL — SWR 1:1
  % ════════════════════════════════════════════════════

  \node[anchor=center, text=titleblue,
        font=\sffamily\bfseries\fontsize{8}{9}\selectfont]
    at (4.0,8.6) {SWR 1:1 — Perfekt matchning};
  \node[anchor=center, text=subtitleblue,
        font=\sffamily\itshape\fontsize{6.5}{7}\selectfont]
    at (4.0,8.26) {Ingen reflektion — all energi når antennen};

  % Framåtvåg-pil
  \draw[-{stealth}, color=wavegreen, line width=0.9pt]
    (1.9,7.74) -- (5.5,7.74);
  \node[anchor=west, text=wavegreen,
        font=\sffamily\bfseries\fontsize{6}{7}\selectfont]
    at (2.1,7.9) {Framåtvåg (100\,W)};

  % Koaxkabel vänster
  \filldraw[fill=coaxblue, draw=coaxbluedark,
            line width=0.3pt, rounded corners=2pt]
    (0.9,5.14) rectangle (7.1,5.7);
  \filldraw[fill=coaxbluedark, rounded corners=1pt]
    (0.9,5.34) rectangle (7.1,5.5);

  % TX vänster
  \filldraw[fill=blue!10!white, draw=titleblue,
            line width=0.4pt, rounded corners=1pt]
    (0.6,5.22) rectangle (1.0,5.62);
  \node[text=titleblue, font=\sffamily\bfseries\fontsize{5}{6}\selectfont]
    at (0.8,5.42) {TX};
  \node[anchor=south, text=subtitleblue,
        font=\sffamily\fontsize{5}{6}\selectfont]
    at (0.8,5.62) {Sändare};

  % Antenn vänster
  \draw[titleblue, line width=0.4pt] (7.04,5.42) -- (7.22,5.42);
  \draw[titleblue, line width=0.5pt] (7.22,5.14) -- (7.22,5.7);
  \draw[titleblue, line width=0.4pt] (7.08,5.04) -- (7.36,5.04);
  \draw[titleblue, line width=0.3pt] (7.12,4.9)  -- (7.32,4.9);
  \node[anchor=south, text=subtitleblue,
        font=\sffamily\fontsize{5}{6}\selectfont]
    at (7.22,5.7) {Antenn};

  % Sinusvåg SWR 1:1 — konstant amplitud
  \draw[wavegreen, line width=0.65pt, line cap=round]
    (1.1,5.42)--(1.174,5.671)--(1.248,5.913)--(1.323,6.136)--(1.397,6.332)--
    (1.471,6.494)--(1.546,6.615)--(1.62,6.691)--(1.694,6.72)--(1.768,6.699)--
    (1.843,6.631)--(1.917,6.516)--(1.991,6.361)--(2.065,6.17)--(2.139,5.95)--
    (2.214,5.711)--(2.288,5.461)--(2.362,5.209)--(2.437,4.965)--(2.511,4.739)--
    (2.585,4.538)--(2.659,4.37)--(2.734,4.242)--(2.808,4.158)--(2.882,4.121)--
    (2.956,4.134)--(3.03,4.195)--(3.105,4.302)--(3.179,4.452)--(3.253,4.637)--
    (3.328,4.853)--(3.402,5.089)--(3.476,5.338)--(3.55,5.591)--(3.624,5.836)--
    (3.699,6.066)--(3.773,6.272)--(3.847,6.446)--(3.921,6.581)--(3.996,6.672)--
    (4.07,6.716)--(4.144,6.711)--(4.219,6.658)--(4.293,6.558)--(4.367,6.415)--
    (4.441,6.235)--(4.516,6.024)--(4.59,5.79)--(4.664,5.542)--(4.738,5.29)--
    (4.812,5.043)--(4.887,4.809)--(4.961,4.599)--(5.035,4.42)--(5.109,4.278)--
    (5.184,4.18)--(5.258,4.128)--(5.332,4.125)--(5.406,4.17)--(5.481,4.263)--
    (5.555,4.399)--(5.629,4.574)--(5.704,4.78)--(5.778,5.011)--(5.852,5.257)--
    (5.926,5.509)--(6.0,5.758)--(6.075,5.994)--(6.149,6.209)--(6.223,6.393)--
    (6.298,6.542)--(6.372,6.647)--(6.446,6.707)--(6.52,6.718)--(6.595,6.68)--
    (6.669,6.595)--(6.743,6.466)--(6.817,6.297)--(6.891,6.095)--(6.966,5.868)--
    (7.04,5.623);

  % Amplitud-annotation
  \draw[wavegreen, line width=0.25pt, dashed, opacity=0.5]
    (2.3,6.72) -- (2.3,5.42);
  \node[anchor=west, text=wavegreen,
        font=\sffamily\fontsize{5}{6}\selectfont, opacity=0.8]
    at (2.35,6.07) {amp.\ konstant};

  % Badges vänster
  \filldraw[fill=badgered, draw=wavered,
            line width=0.3pt, rounded corners=2pt]
    (0.8,3.18) rectangle (3.8,3.9);
  \node[text=wavered, font=\sffamily\bfseries\fontsize{6}{7}\selectfont]
    at (2.3,3.62) {0\,\% reflekteras \checkmark};
  \node[text=subtitlered, font=\sffamily\itshape\fontsize{5}{6}\selectfont]
    at (2.3,3.3) {(förlustfri kabel)};

  \filldraw[fill=badgegrn, draw=wavegreen,
            line width=0.3pt, rounded corners=2pt]
    (4.0,3.18) rectangle (7.0,3.9);
  \node[text=wavegreen, font=\sffamily\bfseries\fontsize{6}{7}\selectfont]
    at (5.5,3.62) {100\,\% når antennen};
  \node[text=subtitleblue, font=\sffamily\itshape\fontsize{5}{6}\selectfont]
    at (5.5,3.3) {(vid lasten)};

  % Info-ruta vänster
  \filldraw[fill=infoblue, draw=strokeblue,
            line width=0.3pt, rounded corners=2pt]
    (0.8,2.12) rectangle (6.9,3.0);
  \node[text=titleblue, font=\sffamily\bfseries\fontsize{6.5}{7.5}\selectfont]
    at (3.85,2.7) {Reflekterad effekt: 0\,W};
  \node[text=subtitleblue, font=\sffamily\itshape\fontsize{5.5}{6.5}\selectfont]
    at (3.85,2.34) {Dalen absorberar all röst — inget eko};

  % ════════════════════════════════════════════════════
  % HÖGER PANEL — SWR 3:1
  % ════════════════════════════════════════════════════

  \node[anchor=center, text=titlered,
        font=\sffamily\bfseries\fontsize{8}{9}\selectfont]
    at (12.0,8.72) {SWR 3:1 — Dålig matchning};
  \node[anchor=center, text=subtitlered,
        font=\sffamily\itshape\fontsize{6.5}{7}\selectfont]
    at (12.0,8.42) {25\,\% reflekteras tillbaka mot sändaren};
  \node[anchor=center, text=titlered,
        font=\sffamily\fontsize{5.5}{6.5}\selectfont]
    at (12.0,8.14)
    {$|\Gamma|=0{,}5\;\Rightarrow\;
      \mathrm{SWR}=(1{+}0{,}5)/(1{-}0{,}5)=3{:}1$};

  % Pilar höger
  \draw[-{stealth}, color=wavegreen, line width=0.9pt]
    (9.5,7.76) -- (11.8,7.76);
  \node[anchor=west, text=wavegreen,
        font=\sffamily\bfseries\fontsize{6}{7}\selectfont]
    at (9.6,7.9) {Framåtvåg (100\,W)};
  \draw[-{stealth}, color=wavered, line width=0.9pt]
    (11.8,7.44) -- (9.5,7.44);
  \node[anchor=west, text=wavered,
        font=\sffamily\bfseries\fontsize{6}{7}\selectfont]
    at (9.6,7.3) {Reflektion (25\,W)};

  % Koaxkabel höger
  \filldraw[fill=coaxred, draw=coaxreddark,
            line width=0.3pt, rounded corners=2pt]
    (8.9,5.14) rectangle (15.1,5.7);
  \filldraw[fill=coaxreddark, rounded corners=1pt]
    (8.9,5.34) rectangle (15.1,5.5);

  % TX höger
  \filldraw[fill=red!10!white, draw=titlered,
            line width=0.4pt, rounded corners=1pt]
    (8.6,5.22) rectangle (9.0,5.62);
  \node[text=titlered, font=\sffamily\bfseries\fontsize{5}{6}\selectfont]
    at (8.8,5.42) {TX};
  \node[anchor=south, text=subtitlered,
        font=\sffamily\fontsize{5}{6}\selectfont]
    at (8.8,5.62) {Sändare};

  % Antenn höger
  \draw[titlered, line width=0.4pt] (15.04,5.42) -- (15.22,5.42);
  \draw[titlered, line width=0.5pt] (15.22,5.14) -- (15.22,5.7);
  \draw[titlered, line width=0.4pt] (15.08,5.04) -- (15.36,5.04);
  \draw[titlered, line width=0.3pt] (15.12,4.9)  -- (15.32,4.9);
  \node[anchor=south, text=subtitlered,
        font=\sffamily\fontsize{5}{6}\selectfont]
    at (15.22,5.7) {Antenn};

  % ── Stående våg-envelope (SWR 3:1) ───────────────────
  % Fyllning
  \fill[wavered, opacity=0.10]
    (9.1,7.07)--(9.174,7.042)--(9.248,6.959)--(9.323,6.826)--(9.397,6.65)--
    (9.471,6.444)--(9.545,6.23)--(9.62,6.048)--(9.694,5.97)--(9.768,6.048)--
    (9.842,6.23)--(9.917,6.444)--(9.991,6.65)--(10.065,6.826)--(10.139,6.959)--
    (10.214,7.042)--(10.288,7.07)--(10.362,7.042)--(10.437,6.959)--(10.511,6.826)--
    (10.585,6.65)--(10.659,6.444)--(10.733,6.23)--(10.808,6.048)--(10.882,5.97)--
    (10.956,6.048)--(11.03,6.23)--(11.105,6.444)--(11.179,6.65)--(11.253,6.826)--
    (11.328,6.959)--(11.402,7.042)--(11.476,7.07)--(11.55,7.042)--(11.625,6.959)--
    (11.699,6.826)--(11.773,6.65)--(11.847,6.444)--(11.922,6.23)--(11.996,6.048)--
    (12.07,5.97)--(12.144,6.048)--(12.218,6.23)--(12.293,6.444)--(12.367,6.65)--
    (12.441,6.826)--(12.515,6.959)--(12.59,7.042)--(12.664,7.07)--(12.738,7.042)--
    (12.812,6.959)--(12.887,6.826)--(12.961,6.65)--(13.035,6.444)--(13.11,6.23)--
    (13.184,6.048)--(13.258,5.97)--(13.332,6.048)--(13.407,6.23)--(13.481,6.444)--
    (13.555,6.65)--(13.629,6.826)--(13.703,6.959)--(13.778,7.042)--(13.852,7.07)--
    (13.926,7.042)--(14.0,6.959)--(14.075,6.826)--(14.149,6.65)--(14.223,6.444)--
    (14.297,6.23)--(14.372,6.048)--(14.446,5.97)--(14.52,6.048)--(14.595,6.23)--
    (14.669,6.444)--(14.743,6.65)--(14.817,6.826)--(14.892,6.959)--(14.966,7.042)--
    (15.04,7.07)--(15.04,3.77)--(14.966,3.798)--(14.892,3.881)--(14.817,4.014)--
    (14.743,4.19)--(14.669,4.396)--(14.595,4.61)--(14.52,4.792)--(14.446,4.87)--
    (14.372,4.792)--(14.297,4.61)--(14.223,4.396)--(14.149,4.19)--(14.075,4.014)--
    (14.0,3.881)--(13.926,3.798)--(13.852,3.77)--(13.778,3.798)--(13.703,3.881)--
    (13.629,4.014)--(13.555,4.19)--(13.481,4.396)--(13.407,4.61)--(13.332,4.792)--
    (13.258,4.87)--(13.184,4.792)--(13.11,4.61)--(13.035,4.396)--(12.961,4.19)--
    (12.887,4.014)--(12.812,3.881)--(12.738,3.798)--(12.664,3.77)--(12.59,3.798)--
    (12.515,3.881)--(12.441,4.014)--(12.367,4.19)--(12.293,4.396)--(12.218,4.61)--
    (12.144,4.792)--(12.07,4.87)--(11.996,4.792)--(11.922,4.61)--(11.847,4.396)--
    (11.773,4.19)--(11.699,4.014)--(11.625,3.881)--(11.55,3.798)--(11.476,3.77)--
    (11.402,3.798)--(11.328,3.881)--(11.253,4.014)--(11.179,4.19)--(11.105,4.396)--
    (11.03,4.61)--(10.956,4.792)--(10.882,4.87)--(10.808,4.792)--(10.733,4.61)--
    (10.659,4.396)--(10.585,4.19)--(10.511,4.014)--(10.437,3.881)--(10.362,3.798)--
    (10.288,3.77)--(10.214,3.798)--(10.139,3.881)--(10.065,4.014)--(9.991,4.19)--
    (9.917,4.396)--(9.842,4.61)--(9.768,4.792)--(9.694,4.87)--(9.62,4.792)--
    (9.545,4.61)--(9.471,4.396)--(9.397,4.19)--(9.323,4.014)--(9.248,3.881)--
    (9.174,3.798)--(9.1,3.77)--cycle;

  % Övre envelope
  \draw[wavered, line width=0.6pt, line cap=round, line join=round]
    (9.1,7.07)--(9.174,7.042)--(9.248,6.959)--(9.323,6.826)--(9.397,6.65)--
    (9.471,6.444)--(9.545,6.23)--(9.62,6.048)--(9.694,5.97)--(9.768,6.048)--
    (9.842,6.23)--(9.917,6.444)--(9.991,6.65)--(10.065,6.826)--(10.139,6.959)--
    (10.214,7.042)--(10.288,7.07)--(10.362,7.042)--(10.437,6.959)--(10.511,6.826)--
    (10.585,6.65)--(10.659,6.444)--(10.733,6.23)--(10.808,6.048)--(10.882,5.97)--
    (10.956,6.048)--(11.03,6.23)--(11.105,6.444)--(11.179,6.65)--(11.253,6.826)--
    (11.328,6.959)--(11.402,7.042)--(11.476,7.07)--(11.55,7.042)--(11.625,6.959)--
    (11.699,6.826)--(11.773,6.65)--(11.847,6.444)--(11.922,6.23)--(11.996,6.048)--
    (12.07,5.97)--(12.144,6.048)--(12.218,6.23)--(12.293,6.444)--(12.367,6.65)--
    (12.441,6.826)--(12.515,6.959)--(12.59,7.042)--(12.664,7.07)--(12.738,7.042)--
    (12.812,6.959)--(12.887,6.826)--(12.961,6.65)--(13.035,6.444)--(13.11,6.23)--
    (13.184,6.048)--(13.258,5.97)--(13.332,6.048)--(13.407,6.23)--(13.481,6.444)--
    (13.555,6.65)--(13.629,6.826)--(13.703,6.959)--(13.778,7.042)--(13.852,7.07)--
    (13.926,7.042)--(14.0,6.959)--(14.075,6.826)--(14.149,6.65)--(14.223,6.444)--
    (14.297,6.23)--(14.372,6.048)--(14.446,5.97)--(14.52,6.048)--(14.595,6.23)--
    (14.669,6.444)--(14.743,6.65)--(14.817,6.826)--(14.892,6.959)--(14.966,7.042)--
    (15.04,7.07);

  % Undre envelope
  \draw[wavered, line width=0.6pt, line cap=round, line join=round]
    (9.1,3.77)--(9.174,3.798)--(9.248,3.881)--(9.323,4.014)--(9.397,4.19)--
    (9.471,4.396)--(9.545,4.61)--(9.62,4.792)--(9.694,4.87)--(9.768,4.792)--
    (9.842,4.61)--(9.917,4.396)--(9.991,4.19)--(10.065,4.014)--(10.139,3.881)--
    (10.214,3.798)--(10.288,3.77)--(10.362,3.798)--(10.437,3.881)--(10.511,4.014)--
    (10.585,4.19)--(10.659,4.396)--(10.733,4.61)--(10.808,4.792)--(10.882,4.87)--
    (10.956,4.792)--(11.03,4.61)--(11.105,4.396)--(11.179,4.19)--(11.253,4.014)--
    (11.328,3.881)--(11.402,3.798)--(11.476,3.77)--(11.55,3.798)--(11.625,3.881)--
    (11.699,4.014)--(11.773,4.19)--(11.847,4.396)--(11.922,4.61)--(11.996,4.792)--
    (12.07,4.87)--(12.144,4.792)--(12.218,4.61)--(12.293,4.396)--(12.367,4.19)--
    (12.441,4.014)--(12.515,3.881)--(12.59,3.798)--(12.664,3.77)--(12.738,3.798)--
    (12.812,3.881)--(12.887,4.014)--(12.961,4.19)--(13.035,4.396)--(13.11,4.61)--
    (13.184,4.792)--(13.258,4.87)--(13.332,4.792)--(13.407,4.61)--(13.481,4.396)--
    (13.555,4.19)--(13.629,4.014)--(13.703,3.881)--(13.778,3.798)--(13.852,3.77)--
    (13.926,3.798)--(14.0,3.881)--(14.075,4.014)--(14.149,4.19)--(14.223,4.396)--
    (14.297,4.61)--(14.372,4.792)--(14.446,4.87)--(14.52,4.792)--(14.595,4.61)--
    (14.669,4.396)--(14.743,4.19)--(14.817,4.014)--(14.892,3.881)--(14.966,3.798)--
    (15.04,3.77);

  % ── Nod-markeringar ───────────────────────────────────
  \foreach \nx in {9.7, 10.88, 13.26}{
    \draw[nodeorange, line width=0.3pt, dashed] (\nx,5.98) -- (\nx,4.86);
    \node[anchor=south, text=nodeorange,
          font=\sffamily\fontsize{5}{6}\selectfont] at (\nx,6.22) {nod};
  }

  % ── Antinod-markeringar ───────────────────────────────
  \foreach \ax in {10.28, 12.66}{
    \draw[wavered, line width=0.3pt, opacity=0.6] (\ax,7.07) -- (\ax,6.75);
    \node[anchor=south, text=wavered,
          font=\sffamily\fontsize{5}{6}\selectfont] at (\ax,6.4) {antinod};
  }

  % ── Amplitud-parenteser ───────────────────────────────
  % Max (röd)
  \draw[wavered, line width=0.3pt]
    (15.12,7.07) -- (15.28,7.07) -- (15.28,3.77) -- (15.12,3.77);
  \node[anchor=west, text=wavered,
        font=\sffamily\fontsize{4.5}{5}\selectfont]
    at (15.3,5.42) {max};

  % Min (orange)
  \draw[nodeorange, line width=0.3pt]
    (15.12,5.97) -- (15.44,5.97) -- (15.44,4.87) -- (15.12,4.87);
  \node[anchor=west, text=nodeorange,
        font=\sffamily\fontsize{4.5}{5}\selectfont]
    at (15.46,5.5) {min};
  \node[anchor=west, text=nodeorange,
        font=\sffamily\fontsize{4}{4.5}\selectfont]
    at (15.46,5.28) {$\neq 0!$};

  % ── Badges höger ──────────────────────────────────────
  \filldraw[fill=badgered, draw=wavered,
            line width=0.3pt, rounded corners=2pt]
    (8.64,2.68) rectangle (11.64,3.4);
  \node[text=wavered, font=\sffamily\bfseries\fontsize{6}{7}\selectfont]
    at (10.14,3.12) {25\,\% reflekteras \texttimes};
  \node[text=subtitlered, font=\sffamily\itshape\fontsize{5}{6}\selectfont]
    at (10.14,2.82) {(förlustfri kabel)};

  \filldraw[fill=badgegrn, draw=wavegreen,
            line width=0.3pt, rounded corners=2pt]
    (11.88,2.68) rectangle (14.84,3.4);
  \node[text=wavegreen, font=\sffamily\bfseries\fontsize{6}{7}\selectfont]
    at (13.36,3.12) {75\,\% når antennen};
  \node[text=subtitleblue, font=\sffamily\itshape\fontsize{5}{6}\selectfont]
    at (13.36,2.82) {(vid lasten)};

  % Info-ruta höger
  \filldraw[fill=inforef, draw=strokered,
            line width=0.3pt, rounded corners=2pt]
    (8.64,1.18) rectangle (15.24,2.5);
  \node[text=titlered, font=\sffamily\bfseries\fontsize{6.5}{7.5}\selectfont]
    at (11.94,2.2) {Reflekterad effekt: 25\,W (av 100\,W)};
  \node[text=subtitlered, font=\sffamily\itshape\fontsize{5.5}{6.5}\selectfont]
    at (11.94,1.88) {Bergsväggen studsar tillbaka rösten — du hör ekot};
  \node[text=black!65, font=\sffamily\fontsize{5}{6}\selectfont]
    at (11.94,1.6)
    {Nod $=$ min.\ amplitud $\cdot$ Antinod $=$ max.\ amplitud};
  \node[text=nodeorange, font=\sffamily\itshape\fontsize{5}{6}\selectfont]
    at (11.94,1.32)
    {Vid SWR 3:1 når noden aldrig noll — vågorna tar inte ut varandra helt};

\end{tikzpicture}
\caption{Stående vågor längs koaxkabeln. \textit{Vänster:} SWR~1:1, ingen 
reflektion, konstant amplitud. \textit{Höger:} SWR~3:1, $|\Gamma|=0{,}5$, 
enveloppens amplitud varierar mellan antinod ($A_\text{max}$) och nod 
($A_\text{min}\neq 0$) — vågorna tar inte ut varandra helt.}
\label{fig:swr-envelope}
\end{figure}

\begin{formulabox}
\textbf{Reflektionsfaktor} -- hur stor del av signalen som reflekteras:\\[2mm]
{\large$\mathbf{\Gamma = \dfrac{Z_L - Z_0}{Z_L + Z_0}}$}
\quad{\small\itshape ($\Gamma = 0$ = ingen reflektion\quad|\quad$\Gamma = 1$ = total reflektion)}\\[3mm]
\textbf{SWR från reflektionsfaktor:}\quad
{\large$\mathbf{SWR = \dfrac{1 + |\Gamma|}{1 - |\Gamma|}}$}\\[3mm]
\textbf{Reflekterad effekt:}\quad$\mathbf{P_{ref} = |\Gamma|^2 \times P_{in}}$\\[3mm]
\hrule\vspace{2mm}
{\small\normalfont\itshape
Förenkling vid ren resistiv obalans:\quad$SWR \approx \dfrac{Z_{hög}}{Z_{låg}}$\\[1mm]
$Z_0$ = systemimpedans (50~$\Omega$)\quad|\quad$Z_L$ = lastimpedans (antennens impedans)}
\end{formulabox}

\begin{worked}[title={Räkneexempel -- $\Gamma$ och SWR från impedanser}]
En antenn mäts till 150~$\Omega$ i ett 50~$\Omega$-system.

\textbf{Steg 1 -- Beräkna reflektionsfaktorn:}
\[\Gamma = \frac{Z_L - Z_0}{Z_L + Z_0} = \frac{150 - 50}{150 + 50} = \frac{100}{200} = 0{,}5\]
\textbf{Steg 2 -- Beräkna SWR:}
\[\mathrm{SWR} = \frac{1 + |\Gamma|}{1 - |\Gamma|} = \frac{1 + 0{,}5}{1 - 0{,}5} = \frac{1{,}5}{0{,}5} = \mathbf{3{:}1}\]
\textbf{Steg 3 -- Hur mycket effekt reflekteras?}
\[P_{ref} = |\Gamma|^2 \times P_{in} = 0{,}5^2 \times 100 = 0{,}25 \times 100 = \mathbf{25\,\text{W}}\]
Stämmer med tabellen: SWR 3:1 $\rightarrow$ 25~\% reflekteras.~\checkmark
\end{worked}

\begin{center}
\begin{tabular}{L{2.2cm}L{2.8cm}L{4.5cm}}
\rowcolor{greydark}
\tblhdr{SWR} & \tblhdr{Reflekterad effekt} & \tblhdr{Bedömning} \\
1{,}0:1 & 0~\% & Perfekt -- alla sändare klarar detta \\
\rowcolor{greylight}
1{,}5:1 & 4~\% & Mycket bra -- acceptabelt \\
2{,}0:1 & 11~\% & Bra -- de flesta sändare klarar \\
\rowcolor{greylight}
3{,}0:1 & 25~\% & Högt -- kräver ATU \\
>3{,}0:1 & >25~\% & Riskabelt -- kan skada slutsteget \\
\end{tabular}
\end{center}

\begin{worked}[title={Räkneexempel 1.7.8 -- Reflekterad effekt}]
En sändare ger 100~W ut i ett 50~$\Omega$-system. 
SWR-metern visar 2:1.

\textbf{a) Hur stor procentandel reflekteras?}

Tabell: SWR 2:1 $\rightarrow$ 11~\% reflekteras.

\textbf{b) Hur många watt reflekteras?}
\[ P_{ref} = 0{,}11 \times 100 = \mathbf{11\,\text{W}} \]

\textbf{c) Hur mycket når antennen?}
\[ P_{ant} = 100 - 11 = \mathbf{89\,\text{W}} \]

SWR 2:1 är acceptabelt för de flesta sändare -- 
förlusten är liten och inget slutsteg skadas.
\end{worked}

\subsubsection*{Övningsfrågor -- 1.7.8}

\begin{qblock}{8.}
\qrakna\quad En sändare ger 100~W ut. SWR-metern 
visar 3:1.

a) Hur stor andel av effekten reflekteras?\\
b) Hur mycket watt når faktiskt antennen?\\
c) Din sändare är specad för maximalt SWR 3:1. 
SWR plötsligt hoppar till 5:1 -- vad gör du 
och varför?
\end{qblock}

\begin{qblock}{F.}
\qfallan\quad Peter mäter SWR 1,5:1 på sin antenn 
och är nöjd. Men sedan byter han till en 30~m 
längre koaxkabel med höga förluster. Nu visar 
SWR-metern 1,2:1. Han jublar -- ``bättre SWR 
än förut!''

Har antennen förbättrats?
\end{qblock}

\medskip
\noindent\rule{\textwidth}{0.4pt}
{\small\itshape Försök själv innan du läser svaret.}
\medskip

\begin{facitblock}{8}
\qrakna\quad Sändare 100~W, SWR 3:1.
\tcblower
\textbf{a) Reflekterad andel:}

Ur tabellen: SWR 3:1 $\rightarrow$ \textbf{25~\%} reflekteras.

\textbf{b) Watt till antennen:}
\[ P_{ref} = 0{,}25 \times 100 = 25\,\text{W} \qquad
   P_{ant} = 100 - 25 = \mathbf{75\,\text{W}} \]

\textbf{c) SWR hoppar till 5:1 -- vad gör du?}

Sänk effekten eller stäng av omedelbart. SWR 5:1 innebär att över 
40~\% av effekten studsar tillbaka mot slutsteget -- samma princip 
som bergsdalen fast värre. De flesta moderna sändare har inbyggt 
SWR-skydd som reducerar effekten automatiskt, men lita inte blint 
på det.

Felsök sedan lugnt: lös koaxkontakt eller ett brott i kabeln är 
vanligaste orsaken till ett plötsligt SWR-hopp.
\end{facitblock}

\begin{worked}[title={Lösning -- Fallgropsfråga 1.7.8\,F}]
Nej -- antennen har inte förbättrats. SWR-metern 
ljuger.

En förlustig kabel dämpar \emph{både} utsänd och 
reflekterad signal. Den reflekterade signalen dämpas 
dubbelt (ut och tillbaka) innan den når SWR-metern. 
Resultatet är ett falskt lågt SWR-tal.

Det verkliga SWR vid antennen är fortfarande 1,5:1 
eller sämre -- kabeln döljer bara problemet.

\textbf{Konsekvens:} Kabelförluster ger lägre 
synbart SWR, men förvärrar det totala systemet. 
En lång, dålig kabel kan se ut att ``förbättra'' 
SWR medan den i verkligheten slösar bort effekten.

\textbf{Lärdomen:} SWR ska mätas vid \emph{antennen}, 
inte vid sändaren om kabelförlusterna är höga. 
Lågt SWR vid sändaren med hög kabelförlust är 
inte ett framgångstecken.
\end{worked}


% ── Sammanfattning 1.7 ────────────────────────────────────
\subsubsection*{Sammanfattning -- Kapitel 1.7}

\begin{orangebox}[title={\emoji{🎯} Viktigt för provet -- Kapitel 1.7}]

\textbf{1.7.1--1.7.2\quad Frekvens, period och RMS}
\begin{checklist}
  \item $f = 1/T$ -- frekvens och period är varandras inverser. 50~Hz = 20~ms per cykel
  \item RMS är det \textbf{verkliga} effektvärdet: $U_{RMS} = U_{topp} \times 0{,}707$
  \item Vänd på det: $U_{topp} = U_{RMS} \times 1{,}414$ -- nätspänning 230~V RMS = 325~V i topp
  \item Använd \textbf{alltid} RMS i $P = U \times I$ -- toppvärden ger fel svar
\end{checklist}

\textbf{1.7.3--1.7.5\quad Impedans, resonans och fasförskjutning}
\begin{checklist}
  \item $Z = \sqrt{R^2 + (X_L - X_C)^2}$ -- impedans är motståndet mot växelström
  \item Vid resonans (serie-RLC, given spänning): $X_L = X_C \Rightarrow Z = R$ -- minsta impedans, maximal ström
  \item \textbf{ELI the ICE man} -- spole: spänning före ström (ELI), kondensator: ström före spänning (ICE)
  \item Ren spole/kondensator ger $\pm 90^\circ$ fasförskjutning -- blandade kretsar ger mellanvinklar
\end{checklist}

\textbf{1.7.6--1.7.7\quad 50~$\Omega$, standarder och skineffekt}
\begin{checklist}
  \item 50~$\Omega$ är kompromissen mellan högst effekttålighet ($\approx$30~$\Omega$) och lägst dämpning ($\approx$77~$\Omega$)
  \item Undvik att blanda 50/75~$\Omega$ i sändarsystem utan matchning -- impedansfelet ger reflektioner
  \item Skineffekt: vid höga frekvenser kryper strömmen ut mot ledarytan -- mer koppar i mitten hjälper inte
  \item Större diameter kan ändå hjälpa via ökad ytarea -- $\delta \propto 1/\sqrt{f}$\quad(skintjocklek minskar med frekvensen)
\end{checklist}

\textbf{1.7.8\quad SWR -- Standing Wave Ratio}
\begin{checklist}
  \item $\Gamma = \dfrac{Z_L - Z_0}{Z_L + Z_0}$ \quad|\quad $\mathrm{SWR} = \dfrac{1+|\Gamma|}{1-|\Gamma|}$ \quad|\quad $P_{ref} = |\Gamma|^2 \times P_{in}$
  \item SWR 1:1 = 0~\% reflekterat\quad|\quad SWR 2:1 = 11~\%\quad|\quad SWR 3:1 = 25~\% \quad{\small\itshape(förlustfri linje, vid lasten)}
  \item SWR $>$ 3:1 -- sänk effekten omedelbart, risk för skyddsingrepp eller skada på slutsteget
  \item Lång förlustig kabel kan visa \textbf{lägre} SWR vid sändaren än vid antennen -- kabeln ljuger
  \item Stående vågor: nod = minamplitud (ej noll vid SWR $\neq \infty$), antinod = maxamplitud
\end{checklist}

\end{orangebox}

% ── 1.8 ───────────────────────────────────────────────────
\section{Filter och resonanskretsar}

Filter är kretsar som släpper igenom vissa frekvenser och blockerar
andra. I amatörradio används filter överallt -- för att välja rätt
band, dämpa övertoner och skydda mottagaren från störningar.

\begin{storybox}[{\emoji{💡} Vattensilar i olika storlekar}]
Föreställ dig en hink med vatten och sand. En grov sil släpper igenom
stora partiklar men håller kvar de fina -- det är ett \emph{högpassfilter}.
En fin sil släpper bara igenom det allra minsta -- ett \emph{lågpassfilter}.

Ett bandpassfilter är som en sil med hål av en exakt storlek -- bara
partiklar av rätt storlek passerar. Resten stoppas.

I radio är ''partiklarna'' frekvenser -- och precis som större
partiklar är svårare att sila bort motsvarar högre frekvenser
större partiklar i analogin. Filtret bestämmer vilka frekvenser
som får passera vidare i kretsen.
\end{storybox}

% ── 1.8.1 ─────────────────────────────────────────────────
\subsection{Filtertyper}

Det finns fyra grundläggande filtertyper. Varje typ definieras av
vilket frekvensområde den \emph{passerar} respektive \emph{blockerar}.

\begin{center}
\begin{tabular}{L{2.5cm}L{3cm}L{3cm}L{3cm}}
\rowcolor{greydark}
\tblhdr{Typ} & \tblhdr{Passerar} & \tblhdr{Blockeras} & \tblhdr{Radioexempel} \\
Lågpass (LP) & Låga frekvenser & Höga & Övertonsdämpning efter slutsteg \\
\rowcolor{greylight}
Högpass (HP) & Höga frekvenser & Låga & Blockera LW/AM-störning i mottagare \\
Bandpass (BP) & Ett frekvensband & Utanför bandet & Välj 20m-bandet ur spektrum \\
\rowcolor{greylight}
Bandstopp (BS) & Allt utom ett band & Ett smalt band & Eliminera specifik störsignal \\
\end{tabular}
\end{center}

% ── Diagram 1: De fyra filtertyperna ─────────────────────
% Placeras direkt efter filtertypstabellen i 1.8.1,
% innan stycket om spole/kondensator.

\begin{figure}[h]
\centering
\begin{tikzpicture}[x=1cm, y=1cm, font=\sffamily\small]

  % ── Färger ──────────────────────────────────────────
  \definecolor{panelblue}{RGB}{247,251,255}
  \definecolor{panelred}{RGB}{255,248,247}
  \definecolor{strokeblue}{RGB}{176,204,224}
  \definecolor{strokered}{RGB}{224,176,176}
  \definecolor{passgreen}{RGB}{26,158,74}
  \definecolor{stopred}{RGB}{192,57,43}
  \definecolor{titleblue}{RGB}{26,110,168}
  \definecolor{titlered}{RGB}{192,57,43}
  \definecolor{subtitleblue}{RGB}{74,127,160}
  \definecolor{subtitlered}{RGB}{160,74,74}
  \definecolor{axiscol}{RGB}{136,136,136}
  \definecolor{badgegrnlight}{RGB}{232,248,238}
  \definecolor{badgeredlight}{RGB}{253,236,234}
  \definecolor{bottombar}{RGB}{240,245,250}
  \definecolor{formblue}{RGB}{42,90,128}

  % ══════════════════════════════════════════════════════
  % PANELRAMAR
  % ══════════════════════════════════════════════════════
  \filldraw[fill=panelblue, draw=strokeblue, line width=0.4pt, rounded corners=2pt]
    (0.4,5.4) rectangle (7.9,10.0);
  \filldraw[fill=panelblue, draw=strokeblue, line width=0.4pt, rounded corners=2pt]
    (8.1,5.4) rectangle (15.6,10.0);
  \filldraw[fill=panelred, draw=strokered, line width=0.4pt, rounded corners=2pt]
    (0.4,0.5) rectangle (7.9,5.1);
  \filldraw[fill=panelred, draw=strokered, line width=0.4pt, rounded corners=2pt]
    (8.1,0.5) rectangle (15.6,5.1);

  % ══════════════════════════════════════════════════════
  % LP — LÅGPASSFILTER
  % ══════════════════════════════════════════════════════
  \node[anchor=center, text=titleblue,
        font=\sffamily\bfseries\fontsize{8}{9}\selectfont]
    at (4.15,9.7) {Lågpassfilter (LP)};
  \node[anchor=center, text=subtitleblue,
        font=\sffamily\itshape\fontsize{6}{7}\selectfont]
    at (4.15,9.42) {Passerar låga frekvenser, blockerar höga};

  % Axlar LP
  \draw[axiscol, line width=0.5pt, ->] (1.1,6.7) -- (7.3,6.7);
  \draw[axiscol, line width=0.5pt, ->] (1.1,6.6) -- (1.1,8.9);
  \node[anchor=north, font=\sffamily\fontsize{5.5}{6}\selectfont, text=axiscol]
    at (4.2,6.55) {Frekvens (log) $\rightarrow$};
  \node[anchor=center, font=\sffamily\fontsize{5.5}{6}\selectfont, text=axiscol,
        rotate=90] at (0.72,7.75) {Signalnivå (dB)};
  \node[anchor=east, font=\sffamily\fontsize{5}{5.5}\selectfont, text=axiscol]
    at (1.05,8.64) {0};
  \node[anchor=east, font=\sffamily\fontsize{5}{5.5}\selectfont, text=axiscol]
    at (1.05,6.72) {$-\infty$};

  % fc linje LP
  \draw[titleblue, line width=0.5pt, dashed, opacity=0.6]
    (3.9,8.85) -- (3.9,6.6);
  \node[anchor=north, text=titleblue,
        font=\sffamily\fontsize{5.5}{6}\selectfont]
    at (3.9,6.55) {$f_c$};

  % Passband-skuggning LP
  \fill[passgreen, opacity=0.08] (1.1,6.7) rectangle (3.9,8.8);
  % Stoppband-skuggning LP
  \fill[stopred, opacity=0.06] (3.9,6.7) rectangle (7.2,8.8);

  % LP kurva
  \draw[titleblue, line width=1.4pt, line cap=round, line join=round]
    (1.1,8.64) -- (3.4,8.64)
    .. controls (3.65,8.64) and (3.82,8.56) .. (3.9,8.48)
    .. controls (4.0,8.28) and (4.2,7.7) .. (4.6,7.44)
    .. controls (5.0,7.18) and (5.6,6.94) .. (6.4,6.82)
    -- (7.2,6.78);

  % -3dB punt LP
  \filldraw[fill=titleblue, draw=titleblue] (3.9,8.48) circle (0.08);
  \node[anchor=west, text=titleblue,
        font=\sffamily\fontsize{5}{5.5}\selectfont]
    at (4.0,8.52) {$-3\,\text{dB}$};

  % Badges LP
  \filldraw[fill=badgegrnlight, draw=passgreen, line width=0.3pt, rounded corners=1.5pt]
    (1.3,7.0) rectangle (3.2,7.38);
  \node[anchor=center, text=passgreen,
        font=\sffamily\bfseries\fontsize{5.5}{6}\selectfont]
    at (2.25,7.19) {Passband};

  \filldraw[fill=badgeredlight, draw=stopred, line width=0.3pt, rounded corners=1.5pt]
    (4.8,7.0) rectangle (6.9,7.38);
  \node[anchor=center, text=stopred,
        font=\sffamily\bfseries\fontsize{5.5}{6}\selectfont]
    at (5.85,7.19) {Stoppband};

  % ══════════════════════════════════════════════════════
  % HP — HÖGPASSFILTER
  % ══════════════════════════════════════════════════════
  \node[anchor=center, text=titleblue,
        font=\sffamily\bfseries\fontsize{8}{9}\selectfont]
    at (11.85,9.7) {Högpassfilter (HP)};
  \node[anchor=center, text=subtitleblue,
        font=\sffamily\itshape\fontsize{6}{7}\selectfont]
    at (11.85,9.42) {Passerar höga frekvenser, blockerar låga};

  % Axlar HP
  \draw[axiscol, line width=0.5pt, ->] (8.8,6.7) -- (15.0,6.7);
  \draw[axiscol, line width=0.5pt, ->] (8.8,6.6) -- (8.8,8.9);
  \node[anchor=north, font=\sffamily\fontsize{5.5}{6}\selectfont, text=axiscol]
    at (11.9,6.55) {Frekvens (log) $\rightarrow$};
  \node[anchor=center, font=\sffamily\fontsize{5.5}{6}\selectfont, text=axiscol,
        rotate=90] at (8.42,7.75) {Signalnivå (dB)};
  \node[anchor=east, font=\sffamily\fontsize{5}{5.5}\selectfont, text=axiscol]
    at (8.75,8.64) {0};
  \node[anchor=east, font=\sffamily\fontsize{5}{5.5}\selectfont, text=axiscol]
    at (8.75,6.72) {$-\infty$};

  % fc linje HP
  \draw[titleblue, line width=0.5pt, dashed, opacity=0.6]
    (11.6,8.85) -- (11.6,6.6);
  \node[anchor=north, text=titleblue,
        font=\sffamily\fontsize{5.5}{6}\selectfont]
    at (11.6,6.55) {$f_c$};

  % Stoppband-skuggning HP
  \fill[stopred, opacity=0.06] (8.8,6.7) rectangle (11.6,8.8);
  % Passband-skuggning HP
  \fill[passgreen, opacity=0.08] (11.6,6.7) rectangle (14.9,8.8);

  % HP kurva
  \draw[titleblue, line width=1.4pt, line cap=round, line join=round]
    (8.8,6.78) -- (9.6,6.8)
    .. controls (10.2,6.84) and (10.8,7.18) .. (11.2,7.44)
    .. controls (11.4,7.64) and (11.55,8.28) .. (11.6,8.48)
    .. controls (11.65,8.56) and (11.8,8.64) .. (12.2,8.66)
    -- (14.9,8.66);

  % -3dB punkt HP
  \filldraw[fill=titleblue, draw=titleblue] (11.6,8.48) circle (0.08);
  \node[anchor=west, text=titleblue,
        font=\sffamily\fontsize{5}{5.5}\selectfont]
    at (11.7,8.52) {$-3\,\text{dB}$};

  % Badges HP
  \filldraw[fill=badgeredlight, draw=stopred, line width=0.3pt, rounded corners=1.5pt]
    (8.95,7.0) rectangle (11.0,7.38);
  \node[anchor=center, text=stopred,
        font=\sffamily\bfseries\fontsize{5.5}{6}\selectfont]
    at (9.975,7.19) {Stoppband};

  \filldraw[fill=badgegrnlight, draw=passgreen, line width=0.3pt, rounded corners=1.5pt]
    (12.4,7.0) rectangle (14.5,7.38);
  \node[anchor=center, text=passgreen,
        font=\sffamily\bfseries\fontsize{5.5}{6}\selectfont]
    at (13.45,7.19) {Passband};

  % ══════════════════════════════════════════════════════
  % BP — BANDPASSFILTER
  % ══════════════════════════════════════════════════════
  \node[anchor=center, text=titlered,
        font=\sffamily\bfseries\fontsize{8}{9}\selectfont]
    at (4.15,4.8) {Bandpassfilter (BP)};
  \node[anchor=center, text=subtitlered,
        font=\sffamily\itshape\fontsize{6}{7}\selectfont]
    at (4.15,4.52) {Passerar ett frekvensband, blockerar utanför};

  % Axlar BP
  \draw[axiscol, line width=0.5pt, ->] (1.1,1.8) -- (7.3,1.8);
  \draw[axiscol, line width=0.5pt, ->] (1.1,1.7) -- (1.1,4.1);
  \node[anchor=north, font=\sffamily\fontsize{5.5}{6}\selectfont, text=axiscol]
    at (4.2,1.65) {Frekvens (log) $\rightarrow$};
  \node[anchor=center, font=\sffamily\fontsize{5.5}{6}\selectfont, text=axiscol,
        rotate=90] at (0.72,2.9) {Signalnivå (dB)};
  \node[anchor=east, font=\sffamily\fontsize{5}{5.5}\selectfont, text=axiscol]
    at (1.05,3.88) {0};
  \node[anchor=east, font=\sffamily\fontsize{5}{5.5}\selectfont, text=axiscol]
    at (1.05,1.85) {$-\infty$};

  % f1, f0, f2 BP
  \draw[titlered, line width=0.5pt, dashed, opacity=0.6]
    (2.9,4.05) -- (2.9,1.7);
  \draw[axiscol, line width=0.4pt, dashed, opacity=0.4]
    (4.14,4.05) -- (4.14,1.7);
  \draw[titlered, line width=0.5pt, dashed, opacity=0.6]
    (5.4,4.05) -- (5.4,1.7);
  \node[anchor=north, text=titlered,
        font=\sffamily\fontsize{5.5}{6}\selectfont] at (2.9,1.65) {$f_1$};
  \node[anchor=north, text=axiscol,
        font=\sffamily\fontsize{5.5}{6}\selectfont] at (4.14,1.65) {$f_0$};
  \node[anchor=north, text=titlered,
        font=\sffamily\fontsize{5.5}{6}\selectfont] at (5.4,1.65) {$f_2$};

  % Skuggningar BP
  \fill[stopred, opacity=0.06]  (1.1,1.8) rectangle (2.9,4.0);
  \fill[passgreen, opacity=0.08] (2.9,1.8) rectangle (5.4,4.0);
  \fill[stopred, opacity=0.06]  (5.4,1.8) rectangle (7.2,4.0);

  % BP kurva
  \draw[titlered, line width=1.4pt, line cap=round, line join=round]
    (1.1,1.84) -- (2.0,1.86)
    .. controls (2.4,1.88) and (2.7,2.2) .. (2.9,2.6)
    .. controls (3.1,3.0) and (3.4,3.7) .. (3.7,3.82)
    .. controls (3.9,3.9) and (4.0,3.92) .. (4.14,3.92)
    .. controls (4.28,3.92) and (4.4,3.9) .. (4.6,3.82)
    .. controls (4.9,3.7) and (5.2,3.0) .. (5.36,2.6)
    .. controls (5.52,2.2) and (5.9,1.88) .. (6.2,1.86)
    -- (7.2,1.84);

  % -3dB linje BP
  \draw[titlered, line width=0.5pt, dashed, opacity=0.5]
    (2.9,3.64) -- (5.4,3.64);
  \filldraw[fill=titlered, draw=titlered] (2.9,3.64) circle (0.07);
  \filldraw[fill=titlered, draw=titlered] (5.4,3.64) circle (0.07);
  \node[anchor=west, text=titlered,
        font=\sffamily\fontsize{5}{5.5}\selectfont]
    at (5.45,3.64) {$-3\,\text{dB}$};

  % Badge BP
  \filldraw[fill=badgegrnlight, draw=passgreen, line width=0.3pt, rounded corners=1.5pt]
    (3.1,2.1) rectangle (5.2,2.48);
  \node[anchor=center, text=passgreen,
        font=\sffamily\bfseries\fontsize{5.5}{6}\selectfont]
    at (4.15,2.29) {Passband (BW)};

  % ══════════════════════════════════════════════════════
  % BS — BANDSTOPPFILTER
  % ══════════════════════════════════════════════════════
  \node[anchor=center, text=titlered,
        font=\sffamily\bfseries\fontsize{8}{9}\selectfont]
    at (11.85,4.8) {Bandstoppfilter (BS)};
  \node[anchor=center, text=subtitlered,
        font=\sffamily\itshape\fontsize{6}{7}\selectfont]
    at (11.85,4.52) {Blockerar ett smalt band, passerar allt annat};

  % Axlar BS
  \draw[axiscol, line width=0.5pt, ->] (8.8,1.8) -- (15.0,1.8);
  \draw[axiscol, line width=0.5pt, ->] (8.8,1.7) -- (8.8,4.1);
  \node[anchor=north, font=\sffamily\fontsize{5.5}{6}\selectfont, text=axiscol]
    at (11.9,1.65) {Frekvens (log) $\rightarrow$};
  \node[anchor=center, font=\sffamily\fontsize{5.5}{6}\selectfont, text=axiscol,
        rotate=90] at (8.42,2.9) {Signalnivå (dB)};
  \node[anchor=east, font=\sffamily\fontsize{5}{5.5}\selectfont, text=axiscol]
    at (8.75,3.88) {0};
  \node[anchor=east, font=\sffamily\fontsize{5}{5.5}\selectfont, text=axiscol]
    at (8.75,1.85) {$-\infty$};

  % f1, f0, f2 BS
  \draw[titlered, line width=0.5pt, dashed, opacity=0.6]
    (10.6,4.05) -- (10.6,1.7);
  \draw[axiscol, line width=0.4pt, dashed, opacity=0.4]
    (11.84,4.05) -- (11.84,1.7);
  \draw[titlered, line width=0.5pt, dashed, opacity=0.6]
    (13.1,4.05) -- (13.1,1.7);
  \node[anchor=north, text=titlered,
        font=\sffamily\fontsize{5.5}{6}\selectfont] at (10.6,1.65) {$f_1$};
  \node[anchor=north, text=axiscol,
        font=\sffamily\fontsize{5.5}{6}\selectfont] at (11.84,1.65) {$f_0$};
  \node[anchor=north, text=titlered,
        font=\sffamily\fontsize{5.5}{6}\selectfont] at (13.1,1.65) {$f_2$};

  % Skuggningar BS
  \fill[passgreen, opacity=0.08] (8.8,1.8)  rectangle (10.6,4.0);
  \fill[stopred,   opacity=0.06] (10.6,1.8) rectangle (13.1,4.0);
  \fill[passgreen, opacity=0.08] (13.1,1.8) rectangle (14.9,4.0);

  % BS kurva
  \draw[titlered, line width=1.4pt, line cap=round, line join=round]
    (8.8,3.84) -- (9.8,3.84)
    .. controls (10.2,3.84) and (10.4,3.68) .. (10.6,3.5)
    .. controls (10.8,3.28) and (11.0,2.56) .. (11.2,2.24)
    .. controls (11.5,1.92) and (11.7,1.86) .. (11.84,1.86)
    .. controls (11.98,1.86) and (12.2,1.92) .. (12.48,2.24)
    .. controls (12.7,2.56) and (12.9,3.28) .. (13.08,3.5)
    .. controls (13.3,3.68) and (13.5,3.84) .. (13.88,3.84)
    -- (14.9,3.84);

  % -3dB linje BS
  \draw[titlered, line width=0.5pt, dashed, opacity=0.5]
    (10.6,3.64) -- (13.1,3.64);
  \filldraw[fill=titlered, draw=titlered] (10.6,3.64) circle (0.07);
  \filldraw[fill=titlered, draw=titlered] (13.1,3.64) circle (0.07);
  \node[anchor=west, text=titlered,
        font=\sffamily\fontsize{5}{5.5}\selectfont]
    at (13.15,3.64) {$-3\,\text{dB}$};

  % Badge BS
  \filldraw[fill=badgeredlight, draw=stopred, line width=0.3pt, rounded corners=1.5pt]
    (10.8,2.1) rectangle (12.9,2.48);
  \node[anchor=center, text=stopred,
        font=\sffamily\bfseries\fontsize{5.5}{6}\selectfont]
    at (11.85,2.29) {Stoppband (BW)};

  % ══════════════════════════════════════════════════════
  % BOTTENMENY
  % ══════════════════════════════════════════════════════
  \filldraw[fill=bottombar, draw=strokeblue!60, line width=0.3pt, rounded corners=1.5pt]
    (0.4,0.06) rectangle (15.6,0.42);
  \node[anchor=center, text=formblue,
        font=\sffamily\fontsize{5.5}{6}\selectfont]
    at (8.0,0.24) {LP: spole i serie
      \quad$|$\quad HP: kondensator i serie
      \quad$|$\quad BP/BS: kombinationer av $L$ och $C$
      \quad$|$\quad $-3\,\text{dB}$ = gränsfrekvens = gräns mellan passband och stoppband};

\end{tikzpicture}
\caption{De fyra grundläggande filtertyperna och deras överföringskurvor.
Grönt = passband, rött = stoppband. $f_c$ = gränsfrekvens vid $-3\,\text{dB}$
(LP/HP). $f_1$, $f_2$ = bandgränser vid $-3\,\text{dB}$ (BP/BS),
$f_0$ = centerfrekvens.}
\label{fig:filtertyper}
\end{figure}

Nyckeln till att förstå filter är att spolar och kondensatorer
reagerar olika på frekvens. \textbf{Spolen} har låg impedans vid
låga frekvenser och hög impedans vid höga -- den \emph{släpper igenom}
lågfrekventa signaler och \emph{blockerar} högfrekventa. Placerad i
serie ger den ett lågpassfilter. \textbf{Kondensatorn} är tvärtom --
hög impedans vid låga frekvenser, låg vid höga. Placerad i serie ger
den ett högpassfilter, placerad parallellt till jord leder den bort
högfrekventa störningar. Kombinerar man spolar och kondensatorer i
olika kopplingar får man bandpass- och bandstoppfilter.

\medskip
\noindent\textbf{Gränsfrekvens -- var slutar filtret?}
\medskip

Inget filter stänger av som en ljusbrytare. Istället finns en
\emph{gränsfrekvens} $f_c$ där signalen börjat dämpas märkbart --
per definition där effekten halverats, vilket motsvarar
$-3\;\text{dB}$. Under $f_c$ passerar signalen nästan opåverkad,
över $f_c$ dämpas den alltmer.

\begin{formulabox}
\textbf{Gränsfrekvens (RC-lågpass):}\\[2mm]
{\large$\mathbf{f_c = \dfrac{1}{2\pi RC}}$}\\[3mm]
{\small\normalfont\itshape
$f_c$ = gränsfrekvens (Hz)\quad|\quad
$R$ = resistans ($\Omega$)\quad|\quad
$C$ = kapacitans (F)}
\end{formulabox}

% ── Diagram 2: RC lågpass — exakt kurva vs Bode-asymptot ─
% Placeras direkt efter gränsfrekvens-formelboxen i 1.8.1,
% innan stycket "Filterordning -- hur brant stänger filtret?"

\begin{figure}[h]
\centering
\begin{tikzpicture}[x=1cm, y=1cm, font=\sffamily\small]

  % ── Färger ──────────────────────────────────────────
  \definecolor{panelblue}{RGB}{247,251,255}
  \definecolor{strokeblue}{RGB}{176,204,224}
  \definecolor{passgreen}{RGB}{26,158,74}
  \definecolor{stopred}{RGB}{192,57,43}
  \definecolor{titleblue}{RGB}{26,110,168}
  \definecolor{subtitleblue}{RGB}{74,127,160}
  \definecolor{axiscol}{RGB}{136,136,136}
  \definecolor{asymcol}{RGB}{230,126,34}
  \definecolor{badgegrnlight}{RGB}{232,248,238}
  \definecolor{badgeredlight}{RGB}{253,236,234}
  \definecolor{bottombar}{RGB}{240,245,250}
  \definecolor{formblue}{RGB}{42,90,128}

  % ── Panel ───────────────────────────────────────────
  \filldraw[fill=panelblue, draw=strokeblue, line width=0.4pt, rounded corners=2pt]
    (0.4,1.3) rectangle (15.6,8.4);

  % ── Titlar ──────────────────────────────────────────
  \node[anchor=center, text=titleblue,
        font=\sffamily\bfseries\fontsize{9}{10}\selectfont]
    at (8.0,8.1) {Lågpassfilter — 1:a ordningen (RC)};
  \node[anchor=center, text=subtitleblue,
        font=\sffamily\itshape\fontsize{6.5}{7}\selectfont]
    at (8.0,7.76) {Exakt kurva vs.\ Bode-asymptot (6\,dB/oktav)};

  % ── Axlar ───────────────────────────────────────────
  \draw[axiscol, line width=0.6pt, ->]
    (1.8,2.96) -- (14.8,2.96);
  \draw[axiscol, line width=0.6pt, ->]
    (1.8,2.86) -- (1.8,7.2);
  \node[anchor=north, text=axiscol,
        font=\sffamily\fontsize{5.5}{6}\selectfont]
    at (8.3,2.8) {Frekvens (log) $\rightarrow$};
  \node[anchor=center, text=axiscol,
        font=\sffamily\fontsize{5.5}{6}\selectfont, rotate=90]
    at (1.22,5.0) {Signalnivå (dB)};

  % ── Y-ticks ─────────────────────────────────────────
  \foreach \yy/\lbl in {6.8/0, 6.32/$-3$, 5.84/$-6$, 4.88/$-12$, 3.92/$-18$}{
    \draw[axiscol, line width=0.4pt] (1.7,\yy) -- (1.9,\yy);
    \node[anchor=east, text=axiscol,
          font=\sffamily\fontsize{5}{5.5}\selectfont]
      at (1.65,\yy) {\lbl\,dB};
  }

  % ── X-ticks ─────────────────────────────────────────
  \draw[axiscol, line width=0.4pt] (6.8,2.88) -- (6.8,3.04);
  \node[anchor=north, text=titleblue,
        font=\sffamily\bfseries\fontsize{6}{6.5}\selectfont]
    at (6.8,2.8) {$f_c$};

  \draw[axiscol, line width=0.4pt] (10.4,2.88) -- (10.4,3.04);
  \node[anchor=north, text=asymcol,
        font=\sffamily\bfseries\fontsize{6}{6.5}\selectfont]
    at (10.4,2.8) {$2{\cdot}f_c$};

  \draw[axiscol, line width=0.4pt] (14.0,2.88) -- (14.0,3.04);
  \node[anchor=north, text=axiscol,
        font=\sffamily\fontsize{5.5}{6}\selectfont]
    at (14.0,2.8) {$4{\cdot}f_c$};

  % ── Skuggningar ─────────────────────────────────────
  \fill[passgreen, opacity=0.07] (1.8,2.96) rectangle (6.8,7.04);
  \fill[stopred,   opacity=0.05] (6.8,2.96) rectangle (14.6,7.04);

  % ── Vertikala referenslinjer ─────────────────────────
  \draw[titleblue, line width=0.5pt, dashed, opacity=0.45]
    (6.8,2.96) -- (6.8,7.1);
  \draw[asymcol, line width=0.5pt, dashed, opacity=0.45]
    (10.4,2.96) -- (10.4,7.1);
  \draw[axiscol, line width=0.4pt, dashed, opacity=0.3]
    (14.0,2.96) -- (14.0,7.1);

  % 0 dB horisontell referens
  \draw[axiscol, line width=0.4pt, dashed, opacity=0.3]
    (1.8,6.8) -- (14.6,6.8);
  % -3 dB horisontell referens
  \draw[axiscol, line width=0.3pt, dashed, opacity=0.25]
    (1.8,6.32) -- (14.6,6.32);
  % -6 dB horisontell referens
  \draw[axiscol, line width=0.3pt, dashed, opacity=0.25]
    (1.8,5.84) -- (14.6,5.84);

  % ── Bode-asymptot ────────────────────────────────────
  % Platt 0 dB till fc, sedan -6 dB/oktav
  % fc=(6.8,6.8), 2fc=(10.4,5.84), 4fc=(14.0,4.88)
  \draw[asymcol, line width=1.2pt, dashed, opacity=0.85]
    (1.8,6.8) -- (6.8,6.8);
  \draw[asymcol, line width=1.2pt, dashed, opacity=0.85]
    (6.8,6.8) -- (14.6,4.72);

  % Asymptot-etikett
  \filldraw[fill=white, fill opacity=0.9, draw=none]
    (11.6,4.28) rectangle (14.3,5.0);
  \node[anchor=west, text=asymcol,
        font=\sffamily\bfseries\fontsize{5.5}{6}\selectfont]
    at (11.65,4.82) {Bode-asymptot};
  \node[anchor=west, text=asymcol,
        font=\sffamily\fontsize{5}{5.5}\selectfont]
    at (11.65,4.62) {$-6$\,dB/oktav};
  \node[anchor=west, text=asymcol,
        font=\sffamily\itshape\fontsize{4.5}{5}\selectfont]
    at (11.65,4.44) {(approximation)};

  % ── Exakt kurva ──────────────────────────────────────
  \draw[titleblue, line width=1.6pt, line cap=round, line join=round]
    (1.8,6.8)
    .. controls (2.6,6.8) and (3.0,6.76) .. (3.2,6.64)
    .. controls (4.6,6.52) and (5.8,6.42) .. (6.8,6.32)
    .. controls (8.0,6.18) and (9.2,5.94) .. (10.4,5.68)
    .. controls (11.8,5.4) and (13.0,5.12) .. (14.6,4.72);

  % Kurv-etikett
  \filldraw[fill=white, fill opacity=0.9, draw=none]
    (2.1,5.22) rectangle (4.5,5.76);
  \node[anchor=west, text=titleblue,
        font=\sffamily\bfseries\fontsize{5.5}{6}\selectfont]
    at (2.16,5.6) {Exakt kurva};
  \node[anchor=west, text=subtitleblue,
        font=\sffamily\fontsize{4.5}{5}\selectfont]
    at (2.16,5.36) {$1/\!\sqrt{1+(f/f_c)^2}$};

  % ── Markerade punkter ────────────────────────────────

  % Horisontell linje vid -3 dB till fc
  \draw[titleblue, line width=0.6pt, dashed, opacity=0.55]
    (1.8,6.32) -- (6.8,6.32);

  % Punkt fc exakt: (6.8, 6.32) = -3.0 dB
  \filldraw[fill=titleblue, draw=white, line width=0.5pt]
    (6.8,6.32) circle (0.12);
  \filldraw[fill=white, fill opacity=0.92, draw=none]
    (7.0,6.14) rectangle (10.3,6.62);
  \node[anchor=west, text=titleblue,
        font=\sffamily\bfseries\fontsize{5.5}{6}\selectfont]
    at (7.05,6.5) {$-3{,}0$\,dB vid $f_c$ (exakt)};
  \node[anchor=west, text=subtitleblue,
        font=\sffamily\itshape\fontsize{5}{5.5}\selectfont]
    at (7.05,6.26) {= halverad effekt};

  % Horisontell linje vid -7 dB till 2fc
  \draw[titleblue, line width=0.6pt, dashed, opacity=0.55]
    (1.8,5.68) -- (10.4,5.68);

  % Punkt 2fc exakt: (10.4, 5.68) = -7.0 dB
  \filldraw[fill=titleblue, draw=white, line width=0.5pt]
    (10.4,5.68) circle (0.12);
  \filldraw[fill=white, fill opacity=0.92, draw=none]
    (10.6,5.52) rectangle (14.5,5.82);
  \node[anchor=west, text=titleblue,
        font=\sffamily\bfseries\fontsize{5.5}{6}\selectfont]
    at (10.65,5.7) {$-7{,}0$\,dB vid $2{\cdot}f_c$ (exakt)};

  % Punkt 2fc asymptot: (10.4, 5.84) = -6.0 dB
  \draw[asymcol, line width=1.2pt] (10.4,5.84) circle (0.1);
  \filldraw[fill=white, fill opacity=0.92, draw=none]
    (10.6,5.84) rectangle (14.1,6.18);
  \node[anchor=west, text=asymcol,
        font=\sffamily\bfseries\fontsize{5.5}{6}\selectfont]
    at (10.65,6.1) {$-6{,}0$\,dB (asymptot)};
  \node[anchor=west, text=asymcol,
        font=\sffamily\itshape\fontsize{5}{5.5}\selectfont]
    at (10.65,5.9) {Bode vid $2{\cdot}f_c$};

  % Pil ~1 dB skillnad
  \draw[black!60, line width=0.7pt, ->]
    (10.46,5.82) -- (10.46,5.7);
  \filldraw[fill=white, fill opacity=0.88, draw=none]
    (10.5,5.77) rectangle (10.96,5.84);
  \node[anchor=west, text=black!60,
        font=\sffamily\fontsize{4.5}{5}\selectfont]
    at (10.5,5.8) {$\approx$1\,dB};

  % ── Badges ──────────────────────────────────────────
  \filldraw[fill=badgegrnlight, draw=passgreen, line width=0.3pt, rounded corners=1.5pt]
    (2.1,3.26) rectangle (3.8,3.7);
  \node[anchor=center, text=passgreen,
        font=\sffamily\bfseries\fontsize{5.5}{6}\selectfont]
    at (2.95,3.48) {Passband};

  \filldraw[fill=badgeredlight, draw=stopred, line width=0.3pt, rounded corners=1.5pt]
    (11.8,3.26) rectangle (13.5,3.7);
  \node[anchor=center, text=stopred,
        font=\sffamily\bfseries\fontsize{5.5}{6}\selectfont]
    at (12.65,3.48) {Stoppband};

  % ── Bottenmeny ───────────────────────────────────────
  \filldraw[fill=bottombar, draw=strokeblue!60, line width=0.3pt, rounded corners=1.5pt]
    (0.4,1.32) rectangle (15.6,1.98);
  \node[anchor=center, text=formblue,
        font=\sffamily\bfseries\fontsize{5.5}{6}\selectfont]
    at (8.0,1.76) {Bode-asymptot: 0\,dB till $f_c$ — sedan $-6$\,dB per oktav (fördubblad frekvens)};
  \node[anchor=center, text=subtitleblue,
        font=\sffamily\itshape\fontsize{5}{5.5}\selectfont]
    at (8.0,1.5) {Vid $2{\cdot}f_c$: asymptot $= -6{,}0$\,dB
      $\;|\;$ exakt $= -7{,}0$\,dB
      $\;|\;$ skillnad $\approx 1$\,dB
      $\;|\;$ approximationen förbättras längre från $f_c$};

\end{tikzpicture}
\caption{Frekvensrespons för ett RC-lågpassfilter av första ordningen.
Blå kurva = exakt svar $|H(f)| = 1/\!\sqrt{1+(f/f_c)^2}$.
Orange streckad = Bode-asymptot (0\,dB till $f_c$, sedan $-6$\,dB/oktav).
Vid $2f_c$ är skillnaden $\approx 1$\,dB -- asymptoten är en god approximation
men stämmer exakt bara långt ovanför $f_c$.}
\label{fig:rc-lp-bode}
\end{figure}

\medskip
\noindent\textbf{Filterordning -- hur brant stänger filtret?}
\medskip

Ett enkelt RC-filter (första ordningens filter) dämpar med
\textbf{6~dB per oktav} ovanför $f_c$ -- varje gång frekvensen
fördubblas dämpas signalen ytterligare 6~dB (halverad spänning).
Detta är en asymptotisk approximation som gäller väl ovanför $f_c$,
men exakt vid $f_c$ är dämpningen bara $-3\;\text{dB}$.

Lägger man till fler LC-steg -- andra, tredje ordningens filter --
blir lutningen brantare: 12, 18~dB per oktav och så vidare.
Pi-filtret efter ett slutsteg är typiskt tredje ordningen och ger
18~dB/oktav, vilket räcker för att dämpa övertoner till laglig nivå.

\begin{storybox}[{\emoji{💡} Trappan nedför}]
Tänk dig en trappa där varje steg är 6~dB nedåt. Ett enkelt filter
har ett steg -- det går långsamt utför. Tredje ordningens filter
har tre steg -- brantare, snabbare nedåt.

Vill du stoppa övertoner effektivt behöver du tillräckligt många
steg så att signalen är tillräckligt dämpad vid dubbla frekvensen.
\end{storybox}

% ── Diagram 3: Filterordning 1/2/3 ───────────────────────
% Placeras direkt efter storybox "Trappan nedför" i 1.8.1,
% innan \subsubsection*{Övningsfrågor -- 1.8.1}

\begin{figure}[h]
\centering
\begin{tikzpicture}[x=1cm, y=1cm, font=\sffamily\small]

  % ── Färger ──────────────────────────────────────────
  \definecolor{panelblue}{RGB}{247,251,255}
  \definecolor{strokeblue}{RGB}{176,204,224}
  \definecolor{passgreen}{RGB}{26,158,74}
  \definecolor{stopred}{RGB}{192,57,43}
  \definecolor{titleblue}{RGB}{26,110,168}
  \definecolor{subtitleblue}{RGB}{74,127,160}
  \definecolor{axiscol}{RGB}{136,136,136}
  \definecolor{curve1}{RGB}{26,110,168}
  \definecolor{curve2}{RGB}{192,57,43}
  \definecolor{curve3}{RGB}{26,158,74}
  \definecolor{piorange}{RGB}{230,126,34}
  \definecolor{piorangelight}{RGB}{255,243,224}
  \definecolor{piorangedark}{RGB}{176,90,16}
  \definecolor{badgegrnlight}{RGB}{232,248,238}
  \definecolor{badgeredlight}{RGB}{253,236,234}
  \definecolor{bottombar}{RGB}{240,245,250}
  \definecolor{formblue}{RGB}{42,90,128}

  % ── Panel ───────────────────────────────────────────
  \filldraw[fill=panelblue, draw=strokeblue, line width=0.4pt, rounded corners=2pt]
    (0.4,1.3) rectangle (15.6,8.4);

  % ── Titlar ──────────────────────────────────────────
  \node[anchor=center, text=titleblue,
        font=\sffamily\bfseries\fontsize{9}{10}\selectfont]
    at (8.0,8.1) {Filterordning — hur brant stänger filtret?};
  \node[anchor=center, text=subtitleblue,
        font=\sffamily\itshape\fontsize{6.5}{7}\selectfont]
    at (8.0,7.76) {Butterworth lågpass: 1:a, 2:a och 3:e ordningen};

  % ── Axlar ───────────────────────────────────────────
  \draw[axiscol, line width=0.6pt, ->]
    (1.8,2.96) -- (14.8,2.96);
  \draw[axiscol, line width=0.6pt, ->]
    (1.8,2.86) -- (1.8,7.2);
  \node[anchor=north, text=axiscol,
        font=\sffamily\fontsize{5.5}{6}\selectfont]
    at (8.3,2.8) {Frekvens (log) $\rightarrow$};
  \node[anchor=center, text=axiscol,
        font=\sffamily\fontsize{5.5}{6}\selectfont, rotate=90]
    at (1.22,5.0) {Signalnivå (dB)};

  % ── Y-ticks ─────────────────────────────────────────
  \foreach \yy/\lbl in {6.8/0, 6.32/$-3$, 5.84/$-6$, 4.88/$-12$, 3.92/$-18$, 2.96/$-24$}{
    \draw[axiscol, line width=0.4pt] (1.7,\yy) -- (1.9,\yy);
    \node[anchor=east, text=axiscol,
          font=\sffamily\fontsize{5}{5.5}\selectfont]
      at (1.65,\yy) {\lbl\,dB};
  }

  % ── X-ticks ─────────────────────────────────────────
  \draw[axiscol, line width=0.4pt] (6.8,2.88) -- (6.8,3.04);
  \node[anchor=north, text=titleblue,
        font=\sffamily\bfseries\fontsize{6}{6.5}\selectfont]
    at (6.8,2.8) {$f_c$};

  \draw[axiscol, line width=0.4pt] (10.4,2.88) -- (10.4,3.04);
  \node[anchor=north, text=axiscol,
        font=\sffamily\fontsize{5.5}{6}\selectfont]
    at (10.4,2.8) {$2{\cdot}f_c$};

  \draw[axiscol, line width=0.4pt] (14.0,2.88) -- (14.0,3.04);
  \node[anchor=north, text=axiscol,
        font=\sffamily\fontsize{5.5}{6}\selectfont]
    at (14.0,2.8) {$4{\cdot}f_c$};

  % ── Referenslinjer ──────────────────────────────────
  \draw[titleblue, line width=0.5pt, dashed, opacity=0.35]
    (6.8,2.96) -- (6.8,7.1);
  \draw[axiscol, line width=0.4pt, dashed, opacity=0.28]
    (10.4,2.96) -- (10.4,7.1);
  \draw[axiscol, line width=0.4pt, dashed, opacity=0.28]
    (14.0,2.96) -- (14.0,7.1);
  \draw[axiscol, line width=0.4pt, dashed, opacity=0.28]
    (1.8,6.8) -- (14.6,6.8);
  \draw[axiscol, line width=0.3pt, dashed, opacity=0.2]
    (1.8,6.32) -- (14.6,6.32);

  % ── Passband-skuggning ──────────────────────────────
  \fill[passgreen, opacity=0.05] (1.8,2.96) rectangle (6.8,7.04);

  % ── Kurva 1: 1:a ordningen (blå) ────────────────────
  \draw[curve1, line width=1.5pt, line cap=round, line join=round, opacity=0.95]
    (1.954,6.7) -- (2.543,6.677) -- (3.132,6.649) -- (3.72,6.615) --
    (4.309,6.575) -- (4.897,6.527) -- (5.486,6.472) -- (6.075,6.409) --
    (6.663,6.336) -- (7.252,6.255) -- (7.84,6.165) -- (8.429,6.067) --
    (9.018,5.96) -- (9.606,5.846) -- (10.195,5.725) -- (10.783,5.598) --
    (11.372,5.466) -- (11.961,5.33) -- (12.549,5.19) --
    (13.138,5.046) -- (13.726,4.9);

  % ── Kurva 2: 2:a ordningen (röd) ────────────────────
  \draw[curve2, line width=1.5pt, line cap=round, line join=round, opacity=0.95]
    (1.954,6.784) -- (2.543,6.774) -- (3.132,6.76) -- (3.72,6.738) --
    (4.309,6.705) -- (4.897,6.656) -- (5.486,6.585) -- (6.075,6.486) --
    (6.663,6.354) -- (7.252,6.187) -- (7.84,5.986) -- (8.429,5.754) --
    (9.018,5.498) -- (9.606,5.223) -- (10.195,4.934) -- (10.783,4.637) --
    (11.372,4.333) -- (11.961,4.025) -- (12.549,3.715) --
    (13.138,3.403) -- (13.726,3.09);

  % ── Kurva 3: 3:e ordningen (grön) ───────────────────
  \draw[curve3, line width=1.5pt, line cap=round, line join=round, opacity=0.95]
    (1.954,6.797) -- (2.543,6.795) -- (3.132,6.79) -- (3.72,6.78) --
    (4.309,6.762) -- (4.897,6.727) -- (5.486,6.662) -- (6.075,6.55) --
    (6.663,6.371) -- (7.252,6.114) -- (7.84,5.782) -- (8.429,5.394) --
    (9.018,4.968) -- (9.606,4.521) -- (10.195,4.061) --
    (10.783,3.595) -- (11.372,3.126);

  % ── Gemensam -3 dB punkt vid fc ─────────────────────
  \draw[axiscol, line width=0.5pt, dashed, opacity=0.4]
    (1.8,6.32) -- (6.8,6.32);
  \filldraw[fill=white, draw=axiscol, line width=0.8pt]
    (6.8,6.32) circle (0.12);
  \filldraw[fill=white, fill opacity=0.92, draw=none]
    (6.96,6.18) rectangle (10.26,6.5);
  \node[anchor=west, text=axiscol,
        font=\sffamily\bfseries\fontsize{5.5}{6}\selectfont]
    at (7.0,6.34) {$-3$\,dB (alla tre vid $f_c$)};

  % ── Kontrollpunkter vid 2fc ──────────────────────────
  \filldraw[fill=curve1, draw=curve1] (10.4,5.725) circle (0.08);
  \filldraw[fill=curve2, draw=curve2] (10.4,4.934) circle (0.08);
  \filldraw[fill=curve3, draw=curve3] (10.4,4.061) circle (0.08);

  % ── Lutningsetiketter ────────────────────────────────
  % 1:a ordningen
  \filldraw[fill=white, fill opacity=0.92, draw=none]
    (11.96,3.92) rectangle (14.46,4.6);
  \draw[curve1, line width=1.5pt] (11.7,4.36) -- (11.96,4.36);
  \node[anchor=west, text=curve1,
        font=\sffamily\bfseries\fontsize{5.5}{6}\selectfont]
    at (12.0,4.48) {1:a ordningen};
  \node[anchor=west, text=curve1,
        font=\sffamily\fontsize{5}{5.5}\selectfont]
    at (12.0,4.24) {$-6$\,dB/oktav};

  % 2:a ordningen
  \filldraw[fill=white, fill opacity=0.92, draw=none]
    (10.6,3.28) rectangle (13.1,3.96);
  \draw[curve2, line width=1.5pt] (10.34,3.72) -- (10.6,3.72);
  \node[anchor=west, text=curve2,
        font=\sffamily\bfseries\fontsize{5.5}{6}\selectfont]
    at (10.64,3.84) {2:a ordningen};
  \node[anchor=west, text=curve2,
        font=\sffamily\fontsize{5}{5.5}\selectfont]
    at (10.64,3.6) {$-12$\,dB/oktav};

  % 3:e ordningen
  \filldraw[fill=white, fill opacity=0.92, draw=none]
    (8.6,3.66) rectangle (11.1,4.34);
  \draw[curve3, line width=1.5pt] (8.34,4.1) -- (8.6,4.1);
  \node[anchor=west, text=curve3,
        font=\sffamily\bfseries\fontsize{5.5}{6}\selectfont]
    at (8.64,4.22) {3:e ordningen};
  \node[anchor=west, text=curve3,
        font=\sffamily\fontsize{5}{5.5}\selectfont]
    at (8.64,3.98) {$-18$\,dB/oktav};

  % ── Pi-filter annotation ─────────────────────────────
  \filldraw[fill=piorangelight, draw=piorange, line width=0.5pt, rounded corners=2pt]
    (1.92,4.84) rectangle (5.76,5.84);
  \node[anchor=west, text=piorange,
        font=\sffamily\bfseries\fontsize{5.5}{6}\selectfont]
    at (2.0,5.68) {Pi-filter = 3:e ordningen};
  \node[anchor=west, text=piorangedark,
        font=\sffamily\fontsize{5}{5.5}\selectfont]
    at (2.0,5.46) {18\,dB/oktav --- kan uppfylla krav};
  \node[anchor=west, text=piorangedark,
        font=\sffamily\fontsize{5}{5.5}\selectfont]
    at (2.0,5.26) {på övertonsdämpning (t.ex.\ $<\!-40$\,dB)};
  \node[anchor=west, text=piorangedark,
        font=\sffamily\fontsize{5}{5.5}\selectfont]
    at (2.0,5.06) {vid rätt dimensionering};
  % Pil mot grön kurva
  \draw[->, piorange, line width=0.6pt, dashed, opacity=0.7]
    (4.0,5.84) -- (7.1,6.32);

  % ── Badges ──────────────────────────────────────────
  \filldraw[fill=badgegrnlight, draw=passgreen, line width=0.3pt, rounded corners=1.5pt]
    (2.1,3.4) rectangle (3.8,3.84);
  \node[anchor=center, text=passgreen,
        font=\sffamily\bfseries\fontsize{5.5}{6}\selectfont]
    at (2.95,3.62) {Passband};

  \filldraw[fill=badgeredlight, draw=stopred, line width=0.3pt, rounded corners=1.5pt]
    (11.8,3.4) rectangle (13.5,3.84);
  \node[anchor=center, text=stopred,
        font=\sffamily\bfseries\fontsize{5.5}{6}\selectfont]
    at (12.65,3.62) {Stoppband};

  % ── Bottenmeny ───────────────────────────────────────
  \filldraw[fill=bottombar, draw=strokeblue!60, line width=0.3pt, rounded corners=1.5pt]
    (0.4,1.32) rectangle (15.6,1.98);
  \node[anchor=center, text=formblue,
        font=\sffamily\bfseries\fontsize{5.5}{6}\selectfont]
    at (8.0,1.76) {Varje ordning adderar $-6$\,dB/oktav:
      1:a\,=\,6 $|$ 2:a\,=\,12 $|$ 3:e\,=\,18\,dB/oktav};
  \node[anchor=center, text=subtitleblue,
        font=\sffamily\itshape\fontsize{5}{5.5}\selectfont]
    at (8.0,1.5) {Alla tre passerar $-3$\,dB exakt vid $f_c$
      $\;|\;$ Butterworth: $|H|\!=\!1/\!\sqrt{1+(f/f_c)^{2n}}$,\;
      $n = 1, 2, 3$};

\end{tikzpicture}
\caption{Jämförelse av Butterworth lågpassfilter för 1:a, 2:a och 3:e ordningen.
Alla tre passerar $-3$\,dB vid $f_c$. Varje tilläggsordning adderar
$-6$\,dB/oktav i stoppbandet. Pi-filtret (C--L--C) är ett 3:e ordningens
filter med $-18$\,dB/oktav.}
\label{fig:filterordning}
\end{figure}

\subsubsection*{Övningsfrågor -- 1.8.1}

\begin{qblock}{1.}
\qgrund\quad Namnge de fyra filtertyperna och ge ett konkret
radioexempel för varje.
\end{qblock}

\begin{qblock}{2.}
\qrakna\quad Ett RC-lågpassfilter har $R = 1\;\text{k}\Omega$ och
$C = 1\;\text{nF}$.

a) Beräkna gränsfrekvensen $f_c$.\\
b) Uppskatta dämpningen vid $2f_c$ med 6~dB/oktav-tumregeln.
Vad är den exakta dämpningen och varför skiljer det sig något?
\end{qblock}

% ── 1.8.2 ─────────────────────────────────────────────────
\subsection{Resonansfrekvens}

En spole och en kondensator kopplade tillsammans bildar en
\emph{resonanskrets}. Vid resonansfrekvensen $f_0$ balanserar
$X_L$ och $X_C$ varandra exakt -- spolens och kondensatorns
reaktanser tar ut varandra och kretsen svänger maximalt vid
just den frekvensen.

\begin{storybox}[{\emoji{💡} Impedans och koppling avgör tillsammans}]
Att en krets har hög eller låg impedans vid $f_0$ är bara halva
bilden -- var i kretsen den sitter avgör om det blir bandpass
eller bandstopp:

\begin{itemize}
  \item Serie-resonator \emph{i serieväg} $\rightarrow$
    låg impedans vid $f_0$ $\rightarrow$ \textbf{bandpass}
  \item Serie-resonator \emph{i shuntväg} (till jord)
    $\rightarrow$ låg impedans kortsluter $f_0$ $\rightarrow$
    \textbf{bandstopp}
  \item Parallell-resonator \emph{i serieväg} $\rightarrow$
    hög impedans blockerar $f_0$ $\rightarrow$ \textbf{bandstopp}
  \item Parallell-resonator \emph{i shuntväg} (till jord)
    $\rightarrow$ hög impedans låter $f_0$ passera $\rightarrow$
    \textbf{bandpass}
\end{itemize}

I en antenntuner kombineras båda typerna för att matcha impedansen
exakt vid önskad frekvens.
\end{storybox}

\begin{formulabox}
\textbf{Resonansfrekvens:}\\[2mm]
{\large$\mathbf{f_0 = \dfrac{1}{2\pi\sqrt{L \times C}}}$}\\[3mm]
{\small\normalfont\itshape
$f_0$ = resonansfrekvens (Hz)\quad|\quad
$L$ = induktans (H)\quad|\quad
$C$ = kapacitans (F)}\\[2mm]
\hrule\vspace{2mm}
{\small\normalfont\itshape
Vid resonans: $X_L = X_C \;\Rightarrow\;$
seriekrets: $Z = R$ (minimum)\quad|\quad
parallellkrets: $Z$ = maximum}
\end{formulabox}

\begin{worked}[title={Räkneexempel 1.8.2 -- Resonansfrekvens}]
Beräkna $f_0$ för $L = 10\;\mu\text{H}$, $C = 100\;\text{pF}$.

\textbf{Steg 1 -- Sätt in i formeln:}
\[
f_0 = \frac{1}{2\pi\sqrt{L \times C}}
    = \frac{1}{2\pi\sqrt{10 \times 10^{-6}
      \times 100 \times 10^{-12}}}
\]

\textbf{Steg 2 -- Beräkna under rottecknet:}
\[
L \times C = 10^{-5} \times 10^{-10} = 10^{-15}
\quad\Rightarrow\quad
\sqrt{10^{-15}} \approx 3{,}162 \times 10^{-8}
\]

\textbf{Steg 3 -- Slutresultat:}
\[
f_0 = \frac{1}{2\pi \times 3{,}162 \times 10^{-8}}
    \approx \mathbf{5{,}03\;\text{MHz}} \checkmark
\]
\end{worked}

\noindent\textbf{Seriekrets vs parallellkrets vid resonans}
\medskip

Samma $f_0$-formel gäller för både serie- och parallellresonanskretsar,
men beteendet vid resonans är motsatt -- och det är viktigt att hålla
isär dem.

I en \textbf{seriekrets} tar $X_L$ och $X_C$ ut varandra helt vid $f_0$
och kvar är bara resistansen $R$. Impedansen når sitt \emph{minimum}
och strömmen sitt \emph{maximum}. Det är som en öppen dörr -- signalen
vid $f_0$ passerar utan motstånd.

I en \textbf{parallellkrets} är det tvärtom. Vid $f_0$ cirkulerar
energin runt i kretsen utan att vilja lämna den -- impedansen når
sitt \emph{maximum} och strömmen från källan sitt \emph{minimum}.
Det är som en stängd dörr -- signalen vid $f_0$ blockeras från att
passera vidare.

\begin{storybox}[{\emoji{💡} Öppen eller stängd dörr}]
Serie\-resonans = öppen dörr vid $f_0$: signalen passerar lätt,
allt annat bromsas. Används som bandpassfilter i serie med signalvägen.

Parallell\-resonans = stängd dörr vid $f_0$: signalen blockeras,
allt annat passerar. Används som bandstopp parallellt med signalvägen.

I en antenntuner kombineras båda typerna för att matcha impedansen
exakt vid önskad frekvens.
\end{storybox}

% ── Diagram 4: Serie- och parallellresonans ───────────────
% Placeras direkt efter storyboxen "Impedans och koppling
% avgör tillsammans" i 1.8.2, innan
% \subsubsection*{Övningsfrågor -- 1.8.2}

\begin{figure}[h]
\centering
\begin{tikzpicture}[x=1cm, y=1cm, font=\sffamily\small]

  % ── Färger ──────────────────────────────────────────
  \definecolor{panelblue}{RGB}{247,251,255}
  \definecolor{panelred}{RGB}{255,248,247}
  \definecolor{strokeblue}{RGB}{176,204,224}
  \definecolor{strokered}{RGB}{224,176,176}
  \definecolor{passgreen}{RGB}{26,158,74}
  \definecolor{stopred}{RGB}{192,57,43}
  \definecolor{titleblue}{RGB}{26,110,168}
  \definecolor{titlered}{RGB}{192,57,43}
  \definecolor{subtitleblue}{RGB}{74,127,160}
  \definecolor{subtitlered}{RGB}{160,74,74}
  \definecolor{axiscol}{RGB}{136,136,136}
  \definecolor{badgegrnlight}{RGB}{232,248,238}
  \definecolor{badgeredlight}{RGB}{253,236,234}
  \definecolor{bottombar}{RGB}{240,245,250}
  \definecolor{formblue}{RGB}{42,90,128}

  % ══════════════════════════════════════════════════════
  % PANELRAMAR
  % ══════════════════════════════════════════════════════
  \filldraw[fill=panelblue, draw=strokeblue, line width=0.4pt, rounded corners=2pt]
    (0.4,1.2) rectangle (7.8,8.0);
  \filldraw[fill=panelred, draw=strokered, line width=0.4pt, rounded corners=2pt]
    (8.2,1.2) rectangle (15.6,8.0);

  % ══════════════════════════════════════════════════════
  % TITLAR
  % ══════════════════════════════════════════════════════
  \node[anchor=center, text=titleblue,
        font=\sffamily\bfseries\fontsize{8}{9}\selectfont]
    at (4.1,7.72) {Seriekrets vid resonans};
  \node[anchor=center, text=subtitleblue,
        font=\sffamily\itshape\fontsize{6}{6.5}\selectfont]
    at (4.1,7.44) {Impedans $Z$ når minimum vid $f_0$};

  \node[anchor=center, text=titlered,
        font=\sffamily\bfseries\fontsize{8}{9}\selectfont]
    at (11.9,7.72) {Parallellkrets vid resonans};
  \node[anchor=center, text=subtitlered,
        font=\sffamily\itshape\fontsize{6}{6.5}\selectfont]
    at (11.9,7.44) {Impedans $Z$ når maximum vid $f_0$};

  % ══════════════════════════════════════════════════════
  % SERIE — AXLAR
  % ══════════════════════════════════════════════════════
  \draw[axiscol, line width=0.6pt, ->] (1.2,2.6) -- (7.5,2.6);
  \draw[axiscol, line width=0.6pt, ->] (1.2,2.5) -- (1.2,7.0);
  \node[anchor=north, text=axiscol,
        font=\sffamily\fontsize{5.5}{6}\selectfont]
    at (4.3,2.44) {Frekvens $\rightarrow$};
  \node[anchor=center, text=axiscol,
        font=\sffamily\fontsize{5.5}{6}\selectfont, rotate=90]
    at (0.72,4.7) {Impedans $Z$ ($\Omega$)};

  % Y-labels serie
  \node[anchor=east, text=axiscol,
        font=\sffamily\fontsize{5}{5.5}\selectfont]
    at (1.15,6.8) {hög};
  \node[anchor=east, text=titleblue,
        font=\sffamily\bfseries\fontsize{5}{5.5}\selectfont]
    at (1.15,3.52) {$Z_{min}\!=\!R$};
  \node[anchor=east, text=axiscol,
        font=\sffamily\fontsize{5}{5.5}\selectfont]
    at (1.15,2.64) {låg};

  % f0 linje serie
  \draw[titleblue, line width=0.5pt, dashed, opacity=0.45]
    (4.1,2.6) -- (4.1,7.1);
  \node[anchor=north, text=titleblue,
        font=\sffamily\bfseries\fontsize{6}{6.5}\selectfont]
    at (4.1,2.44) {$f_0$};

  % Zmin referenslinje
  \draw[titleblue, line width=0.4pt, dashed, opacity=0.28]
    (1.2,3.52) -- (7.4,3.52);

  % Passband-skuggning serie
  \fill[titleblue, opacity=0.04] (3.2,2.6) rectangle (5.0,7.0);

  % ── Serie kurva (U-form) ─────────────────────────────
  \draw[titleblue, line width=1.6pt, line cap=round, line join=round]
    (1.3,6.3)
    .. controls (1.6,6.24) and (2.2,5.8) .. (2.8,4.9)
    .. controls (3.2,4.3) and (3.7,3.64) .. (4.1,3.52)
    .. controls (4.5,3.64) and (5.0,4.3) .. (5.4,4.9)
    .. controls (6.0,5.8) and (6.6,6.24) .. (7.1,6.3);

  % Minimum punkt
  \filldraw[fill=titleblue, draw=white, line width=0.5pt]
    (4.1,3.52) circle (0.12);
  \filldraw[fill=white, fill opacity=0.92, draw=none]
    (4.24,3.38) rectangle (7.1,3.7);
  \node[anchor=west, text=titleblue,
        font=\sffamily\bfseries\fontsize{5.5}{6}\selectfont]
    at (4.28,3.54) {$Z = R$ (minimum)};

  % BW markering serie
  \draw[titleblue, line width=0.6pt] (3.2,3.2) -- (5.0,3.2);
  \draw[titleblue, line width=0.6pt] (3.2,3.14) -- (3.2,3.26);
  \draw[titleblue, line width=0.6pt] (5.0,3.14) -- (5.0,3.26);
  \node[anchor=north, text=titleblue,
        font=\sffamily\fontsize{5}{5.5}\selectfont]
    at (4.1,3.12) {BW ($-3$\,dB)};

  % Topologinot serie
  \filldraw[fill=white, fill opacity=0.88, draw=none]
    (1.36,5.76) rectangle (6.88,6.12);
  \node[anchor=center, text=subtitleblue,
        font=\sffamily\itshape\fontsize{5}{5.5}\selectfont]
    at (4.1,5.94) {(som serieelement i signalvägen)};

  % Badges serie
  \filldraw[fill=badgegrnlight, draw=passgreen, line width=0.3pt, rounded corners=1.5pt]
    (1.36,5.16) rectangle (2.8,5.56);
  \node[anchor=center, text=passgreen,
        font=\sffamily\bfseries\fontsize{5.5}{6}\selectfont]
    at (2.08,5.36) {Stoppband};

  \filldraw[fill=badgegrnlight, draw=passgreen, line width=0.3pt, rounded corners=1.5pt]
    (5.46,5.16) rectangle (6.9,5.56);
  \node[anchor=center, text=passgreen,
        font=\sffamily\bfseries\fontsize{5.5}{6}\selectfont]
    at (6.18,5.36) {Stoppband};

  \filldraw[fill=badgeredlight, draw=stopred, line width=0.3pt, rounded corners=1.5pt]
    (3.3,4.46) rectangle (4.9,4.86);
  \node[anchor=center, text=stopred,
        font=\sffamily\bfseries\fontsize{5.5}{6}\selectfont]
    at (4.1,4.66) {Passband};

  % ══════════════════════════════════════════════════════
  % PARALLELL — AXLAR
  % ══════════════════════════════════════════════════════
  \draw[axiscol, line width=0.6pt, ->] (9.0,2.6) -- (15.3,2.6);
  \draw[axiscol, line width=0.6pt, ->] (9.0,2.5) -- (9.0,7.0);
  \node[anchor=north, text=axiscol,
        font=\sffamily\fontsize{5.5}{6}\selectfont]
    at (12.1,2.44) {Frekvens $\rightarrow$};
  \node[anchor=center, text=axiscol,
        font=\sffamily\fontsize{5.5}{6}\selectfont, rotate=90]
    at (8.52,4.7) {Impedans $Z$ ($\Omega$)};

  % Y-labels parallell
  \node[anchor=east, text=titlered,
        font=\sffamily\bfseries\fontsize{5}{5.5}\selectfont]
    at (8.95,6.48) {$Z_{max}$};
  \node[anchor=east, text=axiscol,
        font=\sffamily\fontsize{5}{5.5}\selectfont]
    at (8.95,2.64) {låg};

  % f0 linje parallell
  \draw[titlered, line width=0.5pt, dashed, opacity=0.45]
    (11.9,2.6) -- (11.9,7.1);
  \node[anchor=north, text=titlered,
        font=\sffamily\bfseries\fontsize{6}{6.5}\selectfont]
    at (11.9,2.44) {$f_0$};

  % Zmax referenslinje
  \draw[titlered, line width=0.4pt, dashed, opacity=0.28]
    (9.0,6.48) -- (15.2,6.48);

  % Passband-skuggning parallell
  \fill[titlered, opacity=0.04] (11.0,2.6) rectangle (12.8,7.0);

  % ── Parallell kurva (inverterad U) ──────────────────
  \draw[titlered, line width=1.6pt, line cap=round, line join=round]
    (9.1,2.9)
    .. controls (9.4,2.96) and (10.0,3.3) .. (10.6,4.0)
    .. controls (11.0,4.5) and (11.5,6.0) .. (11.9,6.48)
    .. controls (12.3,6.0) and (12.8,4.5) .. (13.2,4.0)
    .. controls (13.8,3.3) and (14.4,2.96) .. (14.9,2.9);

  % Maximum punkt
  \filldraw[fill=titlered, draw=white, line width=0.5pt]
    (11.9,6.48) circle (0.12);
  \filldraw[fill=white, fill opacity=0.92, draw=none]
    (12.04,6.34) rectangle (15.3,6.66);
  \node[anchor=west, text=titlered,
        font=\sffamily\bfseries\fontsize{5.5}{6}\selectfont]
    at (12.08,6.5) {$Z = $ maximum ($\gg R$)};

  % BW markering parallell
  \draw[titlered, line width=0.6pt] (11.0,5.84) -- (12.8,5.84);
  \draw[titlered, line width=0.6pt] (11.0,5.78) -- (11.0,5.9);
  \draw[titlered, line width=0.6pt] (12.8,5.78) -- (12.8,5.9);
  \node[anchor=south, text=titlered,
        font=\sffamily\fontsize{5}{5.5}\selectfont]
    at (11.9,5.86) {BW ($-3$\,dB)};

  % Topologinot parallell
  \filldraw[fill=white, fill opacity=0.88, draw=none]
    (9.16,4.6) rectangle (14.88,4.96);
  \node[anchor=center, text=subtitlered,
        font=\sffamily\itshape\fontsize{5}{5.5}\selectfont]
    at (11.9,4.78) {(som shuntelement till jord)};

  % Badges parallell
  \filldraw[fill=badgeredlight, draw=stopred, line width=0.3pt, rounded corners=1.5pt]
    (9.16,3.96) rectangle (10.6,4.36);
  \node[anchor=center, text=stopred,
        font=\sffamily\bfseries\fontsize{5.5}{6}\selectfont]
    at (9.88,4.16) {Passband};

  \filldraw[fill=badgeredlight, draw=stopred, line width=0.3pt, rounded corners=1.5pt]
    (13.2,3.96) rectangle (14.64,4.36);
  \node[anchor=center, text=stopred,
        font=\sffamily\bfseries\fontsize{5.5}{6}\selectfont]
    at (13.92,4.16) {Passband};

  \filldraw[fill=badgegrnlight, draw=passgreen, line width=0.3pt, rounded corners=1.5pt]
    (11.04,3.96) rectangle (12.74,4.36);
  \node[anchor=center, text=passgreen,
        font=\sffamily\bfseries\fontsize{5.5}{6}\selectfont]
    at (11.89,4.16) {Stoppband};

  % ══════════════════════════════════════════════════════
  % BOTTENMENY
  % ══════════════════════════════════════════════════════
  \filldraw[fill=bottombar, draw=strokeblue!60, line width=0.3pt, rounded corners=1.5pt]
    (0.4,1.22) rectangle (15.6,1.88);
  \node[anchor=center, text=formblue,
        font=\sffamily\bfseries\fontsize{5.5}{6}\selectfont]
    at (8.0,1.68) {Samma formel $f_0 = 1/(2\pi\sqrt{LC})$ ---
      men motsatt impedansbeteende vid resonans};
  \node[anchor=center, text=subtitleblue,
        font=\sffamily\itshape\fontsize{5}{5.5}\selectfont]
    at (8.0,1.42) {Serie: $Z = $ minimum ($= R$) vid $f_0$
      $\;|\;$ Parallell: $Z = $ maximum ($\gg R$) vid $f_0$
      $\;|\;$ $BW = f_0/Q$ (hög-$Q$-approximation)};

\end{tikzpicture}
\caption{Impedans som funktion av frekvens för serie- och parallellresonanskrets.
Båda har samma resonansfrekvens $f_0 = 1/(2\pi\sqrt{LC})$ men motsatt
impedansbeteende. Notera att pass/stoppband-karaktären även beror på
var i kretsen resonatorn är kopplad.}
\label{fig:resonanskurvor}
\end{figure}

\subsubsection*{Övningsfrågor -- 1.8.2}

\begin{qblock}{3.}
\qrakna\quad Beräkna $f_0$ för:\\
a) $L = 10\;\mu\text{H}$, $C = 100\;\text{pF}$\\
b) $L = 1\;\mu\text{H}$, $C = 50\;\text{pF}$\\
c) $L = 500\;\text{nH}$, $C = 20\;\text{pF}$
\end{qblock}

% ── 1.8.3 ─────────────────────────────────────────────────
\subsection{Q-värde}

Q-värdet beskriver hur \emph{selektivt} ett filter eller en
resonanskrets är -- hur smalt och spetsigt resonanstoppen är
i förhållande till centerfrekvensen. Ett högre Q ger smalare
bandbredd och skarpare urval. Ett lägre Q ger bredare bandbredd
och rundare topp.

\begin{storybox}[{\emoji{💡} Stämgaffeln och bomullsbollen}]
En stämgaffel har ett mycket högt Q -- den ringer länge på exakt
rätt ton och reagerar knappt på andra frekvenser. Energin
bevaras länge i svängningen och förluster är minimala.

En bomullsboll dämpar alla ljud ungefär lika -- lågt Q, energin
försvinner snabbt och ingenting svänger länge på någon speciell
frekvens.

I radio vill du ha högt Q när du behöver välja ut en smal
frekvens ur ett tätt spektrum -- till exempel i ett mottagarfilter.
Lägre Q passar när du vill dämpa ett brett frekvensområde jämnt,
som i ett lågpassfilter efter ett slutsteg.
\end{storybox}

\begin{formulabox}
\textbf{Q-värde (resonant krets):}\\[2mm]
{\large$\mathbf{Q = \dfrac{f_0}{BW}}$}\\[3mm]
{\small\normalfont\itshape
$f_0$ = centerfrekvens (Hz)\quad|\quad
$BW$ = bandbredd vid $-3\;\text{dB}$ (Hz)}\\[2mm]
\hrule\vspace{2mm}
{\small\normalfont\itshape
Formeln avser resonanta kretsar (bandpass/bandstopp) med
$-3\;\text{dB}$-bandbredd runt $f_0$ -- ej generella lågpassfilter.\\
Högre Q = smalare bandbredd\quad|\quad
Lägre Q = bredare bandbredd}
\end{formulabox}

\begin{worked}[title={Räkneexempel 1.8.3 -- Q-värde}]
En resonanskrets har $f_0 = 7{,}1\;\text{MHz}$ och en
$-3\;\text{dB}$-bandbredd på $10\;\text{kHz}$.

\[
Q = \frac{f_0}{BW}
  = \frac{7{,}1 \times 10^6}{10 \times 10^3}
  = \frac{7100}{10}
  = \mathbf{710}
\]

Q = 710 innebär en mycket smal och selektiv resonanskrets --
bandbredden är bara $1/710$ av centerfrekvensen. Typiskt för
en högkvalitativ LC-krets med låga förluster.
\end{worked}

% Kräver:
% \usepackage{tikz}
% \usetikzlibrary{arrows.meta,calc}
% \usepackage{xcolor}

\begin{figure}[htbp]
\centering
\begin{tikzpicture}[x=0.02cm,y=-0.02cm] % SVG-lik koordinat (y nedåt)

% --- Färger ---
\definecolor{cBlueTitle}{HTML}{1A6EA8}
\definecolor{cBlueSub}{HTML}{4A7FA0}
\definecolor{cPanel}{HTML}{F7FBFF}
\definecolor{cPanelStroke}{HTML}{B0CCE0}
\definecolor{cAxis}{HTML}{888888}
\definecolor{cGridBlue}{HTML}{1A6EA8}
\definecolor{cRed}{HTML}{C0392B}
\definecolor{cBlue}{HTML}{1A6EA8}
\definecolor{cBottom}{HTML}{F0F5FA}
\definecolor{cBottomStroke}{HTML}{C0CCD8}

% --- Bakgrund ---
\fill[white] (0,0) rectangle (800,420);

% Panel
\filldraw[fill=cPanel,draw=cPanelStroke,line width=1.5,rounded corners=8]
  (20,20) rectangle (780,360);

% Titel
\node[font=\bfseries\fontsize{14}{16}\selectfont,text=cBlueTitle] at (400,46)
  {Q-värde -- selektivitet vid resonans};
\node[font=\itshape\fontsize{10.5}{12}\selectfont,text=cBlueSub] at (400,63)
  {Hög Q = smal spets, låg Q = bred rund topp -- samma $f_0$};

% Axlar
\draw[cAxis,line width=1.3] (80,295)--(745,295);
\fill[cAxis] (745,291)--(753,295)--(745,299)--cycle;
\draw[cAxis,line width=1.3] (80,88)--(80,300);
\fill[cAxis] (76,88)--(80,80)--(84,88)--cycle;

% "Frekvens ->" flyttad neråt (y=334) så f0-etiketten inte döljs
\node[font=\fontsize{10}{11}\selectfont,text=black!65] at (415,334) {Frekvens $\rightarrow$};
\node[font=\fontsize{10}{11}\selectfont,text=black!65,rotate=90] at (30,196) {Överföring $|H(f)|$ (dB)};

% Y-ticks
\foreach \yy/\lab in {100/{0 dB},124/{-3 dB},180/{-10 dB},260/{-20 dB}}{
  \draw[cAxis] (75,\yy)--(85,\yy);
  \node[font=\fontsize{9}{10}\selectfont,text=black!65,anchor=east] at (72,\yy+4) {\lab};
}
\node[font=\fontsize{7.5}{8.5}\selectfont,text=black!45,anchor=east] at (72,138) {(BW mäts här)};

% X-ticks
\draw[cAxis,line width=1.2] (400,290)--(400,300);
\node[font=\bfseries\fontsize{10.5}{11}\selectfont,text=cBlueTitle] at (400,322) {$f_0$};

\draw[black!35] (200,291)--(200,299);
\node[font=\fontsize{9}{10}\selectfont,text=black!45] at (200,313) {$0{,}5\cdot f_0$};

\draw[black!35] (600,291)--(600,299);
\node[font=\fontsize{9}{10}\selectfont,text=black!45] at (600,313) {$1{,}5\cdot f_0$};

% Referenslinjer
\draw[cAxis,line width=0.5,dash pattern=on 2pt off 5pt,opacity=.3] (80,100)--(745,100);
\draw[cGridBlue,line width=0.6,dash pattern=on 3pt off 4pt,opacity=.3] (80,124)--(745,124);
\draw[cGridBlue,line width=0.8,dash pattern=on 4pt off 3pt,opacity=.35] (400,88)--(400,300);

% Kurva lag Q (rod)
\draw[cRed,line width=2.5,line cap=round,line join=round,opacity=.9]
plot coordinates {
(220,276.7) (223.2,274.8) (226.4,273.0) (229.6,271.1) (232.8,269.1) (236,267.2)
(239.2,265.3) (242.4,263.3) (245.6,261.3) (248.8,259.4) (252,257.3) (255.2,255.3)
(258.4,253.3) (261.6,251.2) (264.8,249.1) (268,246.9) (271.2,244.8) (274.4,242.6)
(277.6,240.3) (280.8,238.1) (284,235.8) (287.2,233.4) (290.4,231.0) (293.6,228.6)
(296.8,226.1) (300,223.6) (303.2,221.0) (306.4,218.3) (309.6,215.6) (312.8,212.8)
(316,209.9) (319.2,206.9) (322.4,203.9) (325.6,200.8) (328.8,197.5) (332,194.2)
(335.2,190.7) (338.4,187.2) (341.6,183.5) (344.8,179.6) (348,175.6) (351.2,171.4)
(354.4,167.1) (357.6,162.5) (360.8,157.7) (364,152.8) (367.2,147.6) (370.4,142.2)
(373.6,136.6) (376.8,130.8) (380,125.0) (383.2,119.2) (386.4,113.6) (389.6,108.5)
(392.8,104.3) (396,101.4) (399.2,100.1) (400,100.0) (400.8,100.1) (404.0,101.3)
(407.2,104.2) (410.4,108.1) (413.6,112.8) (416.8,118.0) (420.0,123.3) (423.2,128.5)
(426.4,133.7) (429.6,138.7) (432.8,143.4) (436.0,148.0) (439.2,152.3) (442.4,156.4)
(445.6,160.3) (448.8,164.1) (452.0,167.6) (455.2,171.0) (458.4,174.3) (461.6,177.3)
(464.8,180.3) (468.0,183.1) (471.2,185.8) (474.4,188.4) (477.6,191.0) (480.8,193.4)
(484.0,195.7) (487.2,197.9) (490.4,200.1) (493.6,202.2) (496.8,204.2) (500.0,206.2)
(503.2,208.1) (506.4,209.9) (509.6,211.7) (512.8,213.5) (516.0,215.1) (519.2,216.8)
(522.4,218.4) (525.6,219.9) (528.8,221.5) (532.0,222.9) (535.2,224.4) (538.4,225.8)
(541.6,227.2) (544.8,228.5) (548.0,229.8) (551.2,231.1) (554.4,232.4) (557.6,233.6)
(560.8,234.8) (564.0,236.0) (567.2,237.2) (570.4,238.3) (573.6,239.4) (576.8,240.5)
(580.0,241.6)
};

% Kurva hog Q (bla)
\draw[cBlue,line width=2.5,line cap=round,line join=round,opacity=.9]
plot coordinates {
(327.2,293.6) (330.4,290.1) (333.6,286.5) (336.8,282.7) (340.0,278.7) (343.2,274.6)
(346.4,270.3) (349.6,265.7) (352.8,260.8) (356.0,255.7) (359.2,250.2) (362.4,244.3)
(365.6,237.9) (368.8,230.9) (372.0,223.3) (375.2,214.8) (378.4,205.3) (381.6,194.6)
(384.8,182.1) (388.0,167.3) (391.2,149.6) (394.4,128.5) (397.6,107.2) (399.2,100.9)
(400.0,100.0) (400.8,100.9) (402.4,107.2) (405.6,128.0) (408.8,148.4) (412.0,165.5)
(415.2,179.7) (418.4,191.6) (421.6,201.8) (424.8,210.7) (428.0,218.6) (431.2,225.6)
(434.4,232.0) (437.6,237.8) (440.8,243.1) (444.0,248.1) (447.2,252.7) (450.4,256.9)
(453.6,260.9) (456.8,264.7) (460.0,268.3) (463.2,271.6) (466.4,274.8) (469.6,277.8)
(472.8,280.7) (476.0,283.5) (479.2,286.1) (482.4,288.7) (485.6,291.1) (488.8,293.4)
(490.4,294.6)
};

% 0 dB punkt
\filldraw[fill=white,draw=black!55,line width=1.8] (400,100) circle (5.5);

\filldraw[fill=white,fill opacity=.92,draw=none,rounded corners=2] (408,88) rectangle (526,106);
\node[anchor=west,font=\bfseries\fontsize{9}{10}\selectfont,text=black!65] at (414,100) {0 dB vid $f_0$ (båda)};

% BW-markeringar
\draw[cRed,line width=1.5] (380,124)--(420,124);
\draw[cRed,line width=1.5] (380,119)--(380,129);
\draw[cRed,line width=1.5] (420,119)--(420,129);

\draw[cBlue,line width=1.5] (395,124)--(405,124);
\draw[cBlue,line width=1.5] (395,119)--(395,129);
\draw[cBlue,line width=1.5] (405,119)--(405,129);

\draw[cRed,line width=0.8] (380,131)--(380,140)--(420,140)--(420,131);

% BW bred-etikett flyttad åt vänster bort från kurvan
\filldraw[fill=white,fill opacity=.9,draw=none,rounded corners=2] (170,141) rectangle (278,155);
\node[anchor=west,font=\bfseries\fontsize{8.5}{9.5}\selectfont,text=cRed] at (176,152) {BW bred (Q$\approx$10)};
\draw[cRed,line width=0.8,dash pattern=on 2pt off 2pt,opacity=.7] (278,148)--(380,138);

% BW smal-etikett flyttad åt vänster bort från kurvan
\draw[cBlue,line width=0.8] (395,117)--(395,108)--(405,108)--(405,117);
\filldraw[fill=white,fill opacity=.9,draw=none,rounded corners=2] (170,104) rectangle (278,118);
\node[anchor=west,font=\bfseries\fontsize{8.5}{9.5}\selectfont,text=cBlue] at (176,114) {BW smal (Q$\approx$40)};
\draw[cBlue,line width=0.8,dash pattern=on 2pt off 2pt,opacity=.7] (278,111)--(395,108);

% Legend vanster (lag Q) -- utokad ruta
\filldraw[fill=white,fill opacity=.92,draw=cRed,line width=0.8,rounded corners=4]
  (86,158) rectangle (254,214);
\draw[cRed,line width=2.5] (94,174)--(112,174);
\node[anchor=west,font=\bfseries\fontsize{9.5}{10.5}\selectfont,text=cRed] at (116,175) {Låg Q $\approx 10$};
\node[anchor=west,font=\fontsize{8.5}{9.5}\selectfont,text=cRed] at (94,190) {BW = $f_0/Q$ (bred)};
\node[anchor=west,font=\itshape\fontsize{8.5}{9.5}\selectfont,text=red!60!black] at (94,204) {Bomullsboll-analogin};

% Legend hoger (hog Q) -- utokad ruta
\filldraw[fill=white,fill opacity=.92,draw=cBlue,line width=0.8,rounded corners=4]
  (540,128) rectangle (718,182);
\draw[cBlue,line width=2.5] (548,146)--(566,146);
\node[anchor=west,font=\bfseries\fontsize{9.5}{10.5}\selectfont,text=cBlue] at (570,147) {Hög Q $\approx 40$};
\node[anchor=west,font=\fontsize{8.5}{9.5}\selectfont,text=cBlue] at (548,162) {BW = $f_0/Q$ (smal)};
\node[anchor=west,font=\itshape\fontsize{8.5}{9.5}\selectfont,text=blue!60!black] at (548,176) {Stämgaffel-analogin};

% Bottenrad
\filldraw[fill=cBottom,draw=cBottomStroke,line width=1,rounded corners=6]
  (20,368) rectangle (780,408);
\node[font=\bfseries\fontsize{9.5}{10.5}\selectfont,text=cBlueTitle] at (400,384)
  {$Q = f_0 / BW$ -- gäller resonanta kretsar med $-3$ dB-bandbredd runt $f_0$};
\node[font=\itshape\fontsize{9}{10}\selectfont,text=cBlueSub] at (400,399)
  {Båda kurvor normaliserade till 0 dB vid $f_0$ \;|\; Högre Q = smalare bandbredd = skarpare urval};

\end{tikzpicture}
\caption{Q-värdets betydelse för selektivitet i resonanskretsar.}
\label{fig:q-selektivitet}
\end{figure}

\noindent\textbf{Q-värde i praktiken}
\medskip

Q-värdet beror på komponenternas förluster. En luftspole har
betydligt lägre förluster än en spole med järnkärna och når
därmed högre Q. En keramisk resonator eller kvartskristall kan
nå Q-värden på tusentals eller mer -- långt över vad en enkel
LC-krets klarar.

Det finns inget universellt "bra" eller "dåligt" Q -- värdet
måste tolkas i förhållande till tillämpningen. Ett IF-filter i
en mottagare behöver högt Q för att skilja på kanaler tätt
intill varandra. Ett lågpassfilter efter ett slutsteg behöver
inte högt Q -- där är dämpningen vid övertonsfrekvenserna det
viktiga, inte bandbredden kring en resonansfrekvens.

\subsubsection*{Övningsfrågor -- 1.8.3}

\begin{qblock}{4.}
\qrakna\quad En resonanskrets i en mottagare har
$f_0 = 14{,}2\;\text{MHz}$ och bandbredd $BW = 20\;\text{kHz}$.
Beräkna Q-värdet. Vad säger det om kretsens selektivitet?
\end{qblock}

\begin{qblock}{F.}
\qfallan\quad Lisa bygger ett lågpassfilter för att dämpa övertoner
från sitt slutsteg på 40m (7~MHz). Hon mäter att
$-3\;\text{dB}$-punkten ligger på 8~MHz.
``Bra'', tänker hon, ``övertonen på 14~MHz dämpas ordentligt.''

Hennes kompis Erik säger: ``Du måste kontrollera dämpningen vid
14~MHz, inte bara vid $-3\;\text{dB}$-punkten.''

Vem har rätt och varför?
\end{qblock}

\medskip
\noindent\rule{\textwidth}{0.4pt}
{\small\itshape Försök själv innan du läser svaret.}
\medskip

\begin{facitblock}{4}
\qrakna\quad $f_0 = 14{,}2\;\text{MHz}$, $BW = 20\;\text{kHz}$.
\tcblower
\[
Q = \frac{f_0}{BW}
  = \frac{14{,}2 \times 10^6}{20 \times 10^3}
  = \mathbf{710}
\]
Q = 710 innebär en mycket smal och selektiv resonanskrets.
Bandbredden är bara $1/710$ av centerfrekvensen -- kretsen
väljer ut ett smalt fönster runt 14,2~MHz och undertrycker
allt annat effektivt.
\end{facitblock}

\begin{facitblock}{F}
\qfallan\quad Vem har rätt -- Lisa eller Erik?
\tcblower
\textbf{Erik har rätt.}

$-3\;\text{dB}$-punkten säger bara var filtret \emph{börjar}
dämpa, inte hur mycket det dämpar vid en specifik frekvens.
Ett lågpassfilter med $-3\;\text{dB}$ vid 8~MHz ger kanske
bara $-10$ till $-15\;\text{dB}$ dämpning vid 14~MHz,
beroende på filterordning.

För att uppfylla sändningsregler (typiskt $-40\;\text{dB}$
eller bättre på övertoner) måste man mäta den faktiska
dämpningen vid övertonsfrekvenserna -- inte bara kontrollera
$-3\;\text{dB}$-punkten.

\textbf{Lärdomen:} $-3\;\text{dB}$-punkten beskriver filtrets
\emph{bandbredd}, inte dess förmåga att undertrycka oönskade
frekvenser. Dessa är olika saker.
\end{facitblock}

% ── 1.8.4 ─────────────────────────────────────────────────
\subsection{Pi-filtret}

Pi-filtret (C--L--C) är standardlösningen efter ett effektslutssteg
i sändaren. Namnet kommer från att schemat liknar det grekiska
bokstaven $\Pi$ -- en spole i serie med en kondensator till jord på
var sida.

\begin{storybox}[{\emoji{💡} Tre funktioner i ett}]
Pi-filtret gör tre saker samtidigt:
\begin{enumerate}
  \item \textbf{Dämpar övertoner} -- slutsteget genererar harmoniska
    övertoner (2$f$, 3$f$, 4$f$\ldots) som stör andra. LP-karaktären
    hos C--L--C dämpar allt ovanför grundfrekvensen kraftigt.
  \item \textbf{Impedansmatchning} -- transformerar slutstegets
    utgångsimpedans (ofta hundratals $\Omega$) ned till
    koaxkabelns 50~$\Omega$.
  \item \textbf{Blockerar DC} -- kondensatorerna hindrar likspänning
    från att nå antennen.
\end{enumerate}
\end{storybox}

\begin{figure}[htbp]
\centering

\begin{circuitikz}[american, voltage dir=old, scale=1.1, transform shape]

% ── Signalflödespil över HELA filtret ───────────────────
\draw[color=cBlueTitle, ->, line width=1.2pt]
  (-0.6,0.7) -- (7.6,0.7)
  node[midway, above, font=\footnotesize, text=cBlueTitle]
  {signalflöde};

% ── C1 ned till jord ────────────────────────────────────
\draw (0,0)
  to[C, l=$C_1$, a={\small 100\,pF}] (0,-2.5)
  node[ground]{};

% ── Ingångspil ──────────────────────────────────────────
\draw (0,0)
  to[short, -o] (-0.8,0);

% ── Toppledare + L i serie ──────────────────────────────
\draw (0,0)
  to[short] (2.0,0)
  to[L, l=$L$, a={\small 10\,$\mu$H}] (5.0,0)
  to[short] (7.0,0);

% ── C2 ned till jord ────────────────────────────────────
\draw (7.0,0)
  to[C, l=$C_2$, a={\small 100\,pF}] (7.0,-2.5)
  node[ground]{};

% ── Utgångspil ──────────────────────────────────────────
\draw (7.0,0)
  to[short, -o] (7.8,0);

% ── Etiketter ingång ────────────────────────────────────
\node[anchor=east, font=\small] at (-0.95,0.25)
  {Fr\r{a}n slutsteg};
\node[anchor=east, font=\small, text=cRed] at (-0.95,-0.28)
  {${\sim}500\;\Omega$};

% ── Etiketter utgång ────────────────────────────────────
\node[anchor=west, font=\small] at (7.85,0.25)
  {Till antenn};
\node[anchor=west, font=\small, text=cBlueTitle] at (7.85,-0.28)
  {$50\;\Omega$};

\end{circuitikz}

\medskip

\begin{tcolorbox}[
  colback=cBottom,
  colframe=cBottomStroke,
  boxrule=1pt,
  arc=6pt,
  left=8pt, right=8pt, top=5pt, bottom=5pt,
  width=\linewidth
]
{\bfseries\small\color{cBlueTitle}%
  (1)~D\"{a}mpar \"{o}vertoner (LP)\quad
  (2)~Impedansmatchning ${\sim}500\;\Omega \rightarrow 50\;\Omega$\quad
  (3)~Blockerar DC}\\[3pt]
{\itshape\small\color{cBlueSub}%
  $f_0 = \dfrac{1}{2\pi\sqrt{LC}}$\quad
  Dimensioneras f\"{o}r r\"{a}tt band --
  80m-filter fungerar inte p\r{a} 10m-bandet.}
\end{tcolorbox}

\caption{Pi-filtrets kopplingsschema (C--L--C). $C_1$ och $C_2$ till jord
bildar $\Pi$-formen med $L$ i serie.}
\label{fig:pi-filter}
\end{figure}

\begin{redbox}[title={\emoji{⚠️} Fallgrop -- fel filter på fel band}]
Om du använder ett Pi-filter dimensionerat för 80m (3,5~MHz)
på 10m-bandet (28~MHz) händer två saker:
resonansfrekvensen är fel -- filtret dämpar nu \emph{i} bandet
istället för \emph{ovanför} det -- och impedansmatchningen
fungerar inte, SWR stiger kraftigt.

Varje band kräver ett filter dimensionerat för rätt frekvensområde.
\end{redbox}

\noindent\textbf{Q-värde i praktiken}
\medskip

Q-värdet beror på komponenternas förluster. En luftspole har
betydligt lägre förluster än en spole med järnkärna och når
därmed högre Q. En keramisk resonator eller kvartskristall kan
nå Q-värden på tusentals eller mer -- långt över vad en enkel
LC-krets klarar.

Det finns inget universellt ``bra'' eller ``dåligt'' Q -- värdet
måste tolkas i förhållande till tillämpningen. Ett IF-filter i
en mottagare behöver högt Q för att skilja på kanaler tätt
intill varandra. Ett lågpassfilter efter ett slutsteg behöver
inte högt Q -- där är dämpningen vid övertonsfrekvenserna det
viktiga, inte bandbredden kring en resonansfrekvens.

% ── Övningsfrågor ─────────────────────────────────────────
\subsubsection*{Övningsfrågor -- 1.8.4}

\begin{qblock}{5.}
\qprov\quad Förklara Pi-filtrets tre funktioner efter ett
effektslutssteg. Vad händer om man använder ett 80m-filter
på 10m-bandet?
\end{qblock}

\begin{qblock}{6.}
\qrakna\quad Ett Pi-filter har $L = 10\;\mu\text{H}$ och
$C_1 = C_2 = 100\;\text{pF}$.
\begin{enumerate}[label=\alph*)]
  \item Beräkna resonansfrekvensen $f_0$.
  \item Vilket amatörband passar detta filter bäst för?
  \item Vad händer med dämpningen av en överton på $2f_0$
        om Q-värdet är högt respektive lågt?
\end{enumerate}
\end{qblock}


% ── Sammanfattning 1.8 ────────────────────────────────────
\subsubsection*{Sammanfattning -- Kapitel 1.8}

\begin{orangebox}[title={\emoji{🎯} Viktigt för provet -- Kapitel 1.8}]

\textbf{Filtertyper}
\begin{checklist}
  \item \textbf{LP:} passerar låga, blockerar höga -- övertonsfilter efter slutsteg
  \item \textbf{HP:} passerar höga, blockerar låga -- skyddar mot LW/AM-störning
  \item \textbf{BP:} passerar ett band -- väljer ut ett amatörband ur spektrum
  \item \textbf{BS:} blockerar ett smalt band -- eliminerar specifik störsignal
  \item Spole i serie = LP\quad|\quad Kondensator i serie = HP\quad|\quad
        Kombinationer = BP/BS
\end{checklist}

\textbf{Resonans och Q-värde}
\begin{checklist}
  \item $f_0 = \dfrac{1}{2\pi\sqrt{LC}}$ \quad -- resonansfrekvensen där $X_L = X_C$
  \item Seriekrets vid resonans: $Z = R$ (minimum)\quad|\quad
        Parallellkrets: $Z$ = maximum
  \item $Q = f_0/BW$ \quad -- högt Q = smalt filter, lågt Q = brett filter
  \item $-3\;\text{dB}$-punkten = bandbredden, \textbf{inte} dämpningen vid övertoner
\end{checklist}

\textbf{Pi-filter (C--L--C)}
\begin{checklist}
  \item Tre funktioner: (1)~dämpar övertoner \quad
                        (2)~impedansmatchning \quad
                        (3)~blockerar DC
  \item Måste dimensioneras för rätt band -- 80m-filter fungerar inte på 10m
  \item Standard efter effektslutssteg i alla moderna sändare
\end{checklist}

\end{orangebox}

% ══════════════════════════════════════════════════════════════
%  1.9  HALVLEDARE OCH TRANSISTORN
% ══════════════════════════════════════════════════════════════
\section{Halvledare och transistorn}

\lettrine[lines=2]{\textcolor{orange}{H}}{ittills} har vi arbetat med passiva komponenter -- resistorer,
kondensatorer och spolar. De kan bromsa, lagra och fördröja, men
de kan aldrig \emph{förstärka}. En svag signal in ger alltid en
ännu svagare signal ut.

Transistorn förändrar allt. Med en liten signal kan den styra
en stor -- och det är den egenskapen som gör modern elektronik
möjlig. Utan transistorn finns ingen radio, ingen dator, ingen
mobiltelefon. Det är inte en överdrift att säga att transistorn
är 1900-talets viktigaste uppfinning.

Men för att förstå transistorn måste vi börja ett steg tidigare:
med halvledarmaterialet som gör den möjlig.


\usetikzlibrary{arrows.meta, positioning, shapes.geometric, calc}


% Konvertering av transistor-SVG till TikZ
% SVG viewBox: -450 -230 900 430  (px, y nedåt)
% Skalning: 1 px = 0.026 cm  →  divisor ≈ 38.5
% y-spegling: y_tikz = -y_svg * 0.026
% x-offset: x_tikz = x_svg * 0.026  (origo är redan centrerat i SVG)

% Transistordiagram – NPN-förstärkare
% Kräver: \usepackage{tikz}  och  \usetikzlibrary{arrows.meta}

\begin{tikzpicture}[
  font=\sffamily,
  scale=0.026,              % 1 SVG-px = 0.026 cm
  yscale=-1,                % spegla y-axeln (SVG har y nedåt)
  shift={(450,230)},        % flytta origo till SVG-bildens övre vänstra hörn
                            % (SVG viewBox börjar vid -450,-230)
]

% ───────────────────────────────────────────────
% BAKGRUND (vit, ingen ram)
% ───────────────────────────────────────────────
\fill[white] (-450,-230) rectangle (450,200);

% ───────────────────────────────────────────────
% TRE BLOCK
% ───────────────────────────────────────────────

% IN-block (blå)
\draw[rounded corners=5pt, fill=blue!10, draw=blue!50, line width=1.8pt]
  (-370,-75) rectangle (-210,75);

% Transistor-block (orange)
\draw[rounded corners=5pt, fill=orange!12, draw=orange!80, line width=1.8pt]
  (-80,-75) rectangle (80,75);

% OUT-block (grön)
\draw[rounded corners=5pt, fill=green!10, draw=green!60!black, line width=1.8pt]
  (210,-75) rectangle (370,75);

% ───────────────────────────────────────────────
% ETIKETTER: IN
% ───────────────────────────────────────────────
\node[font=\small\sffamily\bfseries] at (-290,-15) {Liten signal};
\node[font=\footnotesize\sffamily]   at (-290,  7) {t.ex. 1\,mA};

% ───────────────────────────────────────────────
% ETIKETTER: OUT
% ───────────────────────────────────────────────
\node[font=\small\sffamily\bfseries] at (290,-15) {Stor signal};
\node[font=\footnotesize\sffamily]   at (290,  7) {t.ex. 100\,mA};

% ───────────────────────────────────────────────
% TRANSISTORETIKETT (överst i orange ruta)
% ───────────────────────────────────────────────
\node[font=\small\sffamily\bfseries] at (0,-50) {Transistor};

% ───────────────────────────────────────────────
% NPN TRANSISTORSYMBOL
% ───────────────────────────────────────────────

% Vertikal linje (body)
\draw[line width=2pt] (10,-32) -- (10,32);

% Bas
\draw[line width=2pt] (-27,0) -- (10,0);

% Kollektor
\draw[line width=2pt] (10,-16) -- (47,-46);

% Emitter med pil
\draw[-{Stealth[length=6pt,width=5pt]}, line width=2pt]
  (10,16) -- (47,46);

% Bas-etikett
\node[font=\scriptsize\sffamily\bfseries] at (-46,0) {Bas};

% ───────────────────────────────────────────────
% K OCH E MED FÖRKLARINGAR
% ───────────────────────────────────────────────
\node[anchor=west, font=\scriptsize\sffamily] at (60,-52)
  {\textbf{K} {\color{black!70}– Kollektor (in-port för ström)}};

\node[anchor=west, font=\scriptsize\sffamily] at (60,62)
  {\textbf{E} {\color{black!70}– Emitter (ut-port för ström)}};

% ───────────────────────────────────────────────
% NPN-ETIKETT
% ───────────────────────────────────────────────
\node[font=\scriptsize\sffamily\bfseries] at (0,68) {NPN};

% ───────────────────────────────────────────────
% STYRSIGNAL: blå pil
% ───────────────────────────────────────────────
\draw[-{Stealth[length=7pt,width=6pt]}, line width=1.8pt, blue!60]
  (-208,0) -- (-88,0);
\node[font=\footnotesize\sffamily] at (-148,-14) {Styrsignal};

% ───────────────────────────────────────────────
% UTSIGNAL: grön tjock pil
% ───────────────────────────────────────────────
\draw[-{Stealth[length=10pt,width=9pt]}, line width=4.5pt, green!60!black]
  (83,0) -- (205,0);
\node[font=\footnotesize\sffamily] at (147,-16) {Utsignal};

% ───────────────────────────────────────────────
% MATNINGSSPÄNNING UPPIFRÅN
% ───────────────────────────────────────────────
\draw[-{Stealth[length=6pt]}, line width=1.3pt, black!70]
  (0,-75) -- (0,-158);
\node[font=\footnotesize\sffamily] at (0,-175)
  {$V_{\!CC}$ – matningsspänning};

% ───────────────────────────────────────────────
% SKYDDSJORD NEDÅT
% ───────────────────────────────────────────────
\draw[line width=1.3pt, black!70] (0,75) -- (0,118);

% Skyddsjordssymbol: tre avtagande linjer
\draw[line width=2.2pt, black!70] (-26,118) -- (26,118);
\draw[line width=2.2pt, black!70] (-17,131) -- (17,131);
\draw[line width=2.2pt, black!70] ( -8,144) -- ( 8,144);

\node[font=\footnotesize\sffamily] at (0,163) {Skyddsjord (GND)};

% ───────────────────────────────────────────────
% BILDTEXT
% ───────────────────────────────────────────────
\node[font=\footnotesize\sffamily\itshape, text=black!70, align=center]
  at (0,200)
  {En liten styrsignal vid basen kontrollerar ett mycket större flöde.};

\end{tikzpicture}

% ──────────────────────────────────────────────────────────────
\subsection{Halvledarmaterial -- vad skiljer en ledare från en isolator?}

Alla material kan delas in i tre grupper beroende på hur lätt
elektroner rör sig genom dem:

\begin{bluebox}[title={\emoji{🔑} Ledare, halvledare och isolatorer}]
\begin{center}
\begin{tabular}{llll}
\rowcolor{greydark}
\tblhdr{Typ} & \tblhdr{Resistivitet} & \tblhdr{Exempel} & \tblhdr{Användning} \\
Ledare & Mycket låg & Koppar, silver, guld & Kablar, ledningar \\
\rowcolor{greylight}
Halvledare & Mellan & Kisel, germanium & Transistorer, dioder \\
Isolator & Mycket hög & Glas, plast, luft & Isolation, kretskort \\
\end{tabular}
\end{center}

\smallskip
Det speciella med halvledare är att deras ledningsförmåga kan
\textbf{styras} -- genom temperatur, ljus eller, viktigast av
allt, genom att tillsätta små mängder föroreningar (dopning).
\end{bluebox}

\begin{storybox}[{\emoji{💡} Kiselatomens fyra händer}]
En kiselatom har fyra ytterelektroner -- fyra ``händer'' som den
gärna håller i med grannarna. I rent kisel bildar atomerna ett
perfekt kristallgitter där alla händer är upptagna. Inga fria
elektroner, inget strömflöde -- kisel i sin renaste form leder
knappt alls.

Nu tillsätter vi en liten mängd fosfor. Fosforatomen har
\emph{fem} ytterelektroner -- en hand för mycket. Den femte
elektronen har inget att hålla i och rör sig fritt genom
kristallen. Vi har skapat ett \textbf{n-type} halvledarmaterial
(n för negativ -- fria elektroner är negativt laddade).

Tillsätter vi istället bor, som bara har tre ytterelektroner,
uppstår ett ``hål'' -- en plats där en elektron saknas. Hålet
beter sig som en positiv laddning som rör sig genom materialet.
Det kallar vi \textbf{p-type} halvledarmaterial.
\end{storybox}

\begin{bluebox}[title={\emoji{🔑} n-type och p-type kisel}]
\textbf{n-type} (negative): Dopat med fosfor, arsenik eller
antimon. Fria \textbf{elektroner} är laddningsbärare.

\medskip
\textbf{p-type} (positive): Dopat med bor, aluminium eller
gallium. Fria \textbf{hål} är laddningsbärare.

\medskip
{\small\itshape Varken n-type eller p-type är elektriskt laddade
som helhet -- de är neutrala. Det är bara \emph{typen} av
laddningsbärare som skiljer dem åt.}
\end{bluebox}

\subsubsection*{Övningsfrågor -- 1.9.1}

\begin{qblock}{1.}
\qgrund\quad Vad menas med att ett material är en halvledare?
Ge ett exempel på ett halvledarmaterial och förklara kortfattat
skillnaden mellan n-type och p-type.
\end{qblock}

% ──────────────────────────────────────────────────────────────
\subsection{Dioden -- strömmen får bara gå en väg}

När vi sätter ett p-type och ett n-type material mot varandra
uppstår något remarkabelt vid gränssnittet -- ett \textbf{p-n-övergång}
som låter ström flöda i en riktning men blockerar den i den andra.
Det är en diod.

\begin{storybox}[{\emoji{🚪} Envägsdörren}]
Föreställ dig en dörr som bara går att öppna inifrån. Vill du
gå ut -- inga problem, du trycker och dörren öppnar. Vill du
komma in utifrån -- dörren står emot, oavsett hur hårt du trycker.

En diod är elektronikens envägsdörr. Ström i en riktning
(framåtriktning) -- dörren öppnar, strömmen flödar. Ström i
andra riktningen (backriktning) -- dörren håller emot,
ingen ström passerar.
\end{storybox}

\begin{bluebox}[title={\emoji{🔑} Dioden -- grundegenskaper}]
\textbf{Anod} (A): p-type sidan. Positiv spänning här = framåtspänning.\\
\textbf{Katod} (K): n-type sidan. Markeras med ett streck på symbolen.

\medskip
\textbf{Framåtspänning} ($V_f$): Den spänning som krävs för att
``öppna'' dioden. Typiskt:
\begin{itemize}
  \item Kiseldiod: ca \textbf{0,6--0,7\,V}
  \item Germaniumdiod: ca \textbf{0,2--0,3\,V}
  \item LED: ca \textbf{1,8--3,5\,V} (beror på färg)
\end{itemize}

\textbf{Spärrspänning} ($V_R$): Maximalt tillåten backspänning
innan dioden slår igenom (destruktivt). Typiskt 50--1000\,V
beroende på typ.

\medskip
\textbf{Användningsområden:}
\begin{itemize}
  \item Likriktning (AC till DC i nätaggregat)
  \item Skyddsdiod (förhindrar felvänd polning)
  \item AM-detektor i mottagare (envelop-detektor)
  \item Zenerdiod -- spänningsstabilisering
  \item LED -- ljusemission
\end{itemize}
\end{bluebox}

% ── Figur: Diodkurva ─────────────────────────────────────────
\begin{figure}[htbp]
\centering
\begin{tikzpicture}[font=\small, >=stealth]

  \fill[gray!6, rounded corners=4pt] (-0.4,-0.6) rectangle (10.4,6.2);

  % ── Panel 1: Diodsymbol ──────────────────────────────────
  \begin{scope}[xshift=1.8cm, yshift=3.5cm]
    \node[font=\small\bfseries] at (0,1.8) {Diodsymbol};
    % Anod--katod linje
    \draw[thick] (-1.2,0) -- (-0.4,0);
    \draw[thick] (0.4,0)  -- (1.2,0);
    % Triangel (pil)
    \fill[gray!70] (-0.4,-0.5) -- (-0.4,0.5) -- (0.4,0) -- cycle;
    % Katodstreck
    \draw[thick] (0.4,-0.5) -- (0.4,0.5);
    % Etiketter
    \node[font=\scriptsize] at (-1.0,-0.3) {A};
    \node[font=\scriptsize] at (1.0,-0.3)  {K};
    \node[font=\scriptsize, gray] at (-0.6,0.7) {p};
    \node[font=\scriptsize, gray] at (0.6,0.7)  {n};
  \end{scope}

  % Skiljelinje
  \draw[gray!30] (4.0,0.2) -- (4.0,6.0);

  % ── Panel 2: I-V-kurva ───────────────────────────────────
  \begin{scope}[xshift=7.0cm, yshift=3.0cm]
    \node[font=\small\bfseries] at (0,2.8) {Ström--spänningskurva};

    % Axlar
    \draw[->, gray!60] (-2.8,0) -- (2.5,0) node[right,font=\scriptsize,gray]{$U$\,(V)};
    \draw[->, gray!60] (0,-1.8) -- (0,2.5) node[above,font=\scriptsize,gray]{$I$};

    % Framåtkurva (exponentiell form, förenklad)
    \draw[red!70!black, very thick, domain=0.0:1.0, samples=60]
      plot (\x, {2.2*(exp(2.8*\x - 1.9) - 0.15)});

    % Spärrkurva (nästan noll, liten läckström)
    \draw[red!70!black, very thick] (-2.5,-0.08) -- (0,0);

    % Genomslagsknä (stiliserat)
    \draw[red!70!black, very thick] (-2.5,-0.08) -- (-2.5,-1.6);
    \draw[red!70!black, dashed] (-2.5,-1.6) -- (-2.5,-1.9);

    % Markeringar
    \draw[gray!50, dashed] (0.68,0) node[below,font=\scriptsize]{$0{,}6$\,V} -- (0.68,0.5);
    \node[font=\scriptsize, red!60!black] at (1.3,1.8) {Framåt};
    \node[font=\scriptsize, gray!70] at (-1.5,-0.4) {Spärr};
    \node[font=\scriptsize, gray!70] at (-2.1,-1.5) {Genomslag};

  \end{scope}

\end{tikzpicture}
\caption{Diodsymbol och I-V-karakteristik. I framåtriktningen
börjar strömmen flöda vid ca 0,6\,V (kisel). I backriktningen
flödar nästan ingen ström förrän genomslagsspänningen nås --
vilket förstör dioden om ingen strömgränsning finns.}
\label{fig:diod_iv}
\end{figure}

\subsubsection*{Övningsfrågor -- 1.9.2}

\begin{qblock}{2.}
\qgrund\quad Förklara med egna ord varför en diod bara leder
ström i en riktning. Vad kallas de två anslutningarna på en diod?
\end{qblock}

\begin{qblock}{3.}
\qrakna\quad En kiseldiod är seriekopplad med ett motstånd på
$R = 470\,\Omega$ och en spänningskälla på $U = 5\,\text{V}$.
Dioden är framåtspänd. Beräkna strömmen genom kretsen.
\emph{Tips: Kiseldioden har en framåtspänning på 0,7\,V.}
\end{qblock}

\medskip
\noindent\rule{\textwidth}{0.4pt}
{\small\itshape Försök själv innan du läser svaret.}
\medskip

\begin{worked}[title={Lösning -- Räkneexempel 1.9.2}]
Spänningen över motståndet är källspänningen minus diodens
framåtspänning:
\[ U_R = 5{,}0 - 0{,}7 = 4{,}3\,\text{V} \]

Strömmen genom kretsen (Ohms lag):
\[ I = \frac{U_R}{R} = \frac{4{,}3}{470} \approx 9{,}1\,\text{mA} \]
\end{worked}

% ──────────────────────────────────────────────────────────────
\subsection{Transistorn -- den aktiva komponenten}

En transistor är i grunden två p-n-övergångar byggda ovanpå
varandra. Det ger oss tre skikt och tre anslutningar -- och
möjligheten att låta en liten ström styra en stor.

\begin{storybox}[{\emoji{🚰} Vattenkranen}]
En vattenkran kontrollerar ett stort vattenflöde med ett litet
finger. Du behöver inte vara stark för att öppna kranen --
fingrets lilla rörelse styr ett kraftigt vattenflöde.

En transistor är elektronikens vattenkran. En liten \emph{styrsignal}
(fingret) kontrollerar ett stort \emph{lastflöde} (vattnet).
Det är förstärkning: litet in, stort ut -- men med exakt samma
form på signalen.
\end{storybox}

\begin{bluebox}[title={\emoji{🔑} Bipolär transistor (BJT) -- grundbegrepp}]
En NPN-transistor har tre skikt: n--p--n, med tre anslutningar:

\begin{itemize}
  \item \textbf{Bas (B):} Styringången -- den lilla strömmen
  \item \textbf{Kollektor (C):} Den stora strömmen flödar in här
  \item \textbf{Emitter (E):} Den sammansatta strömmen lämnar här
\end{itemize}

\medskip
\textbf{Grundrelationen:}
\[ I_C = \beta \cdot I_B \qquad \text{och} \qquad I_E = I_C + I_B \]

där $\beta$ (strömförstärkning, även kallad $h_{FE}$) typiskt
är 50--500 för vanliga transistorer.

\medskip
\textbf{Exempel:} $\beta = 100$, $I_B = 1\,\text{mA}$
$\Rightarrow$ $I_C = 100\,\text{mA}$

En basström på 1\,mA styr en kollektorström på 100\,mA.
\end{bluebox}

% ── Figur: NPN-transistor ────────────────────────────────────
\begin{figure}[htbp]
\centering
\begin{tikzpicture}[font=\small, >=stealth]

  \fill[gray!6, rounded corners=4pt] (-0.4,-0.4) rectangle (13.4,5.8);

  % ── Panel 1: Struktur ────────────────────────────────────
  \begin{scope}[xshift=2.0cm, yshift=0.8cm]
    \node[font=\small\bfseries] at (0,4.6) {NPN-struktur};

    % n-emitter
    \fill[blue!15] (-1.2,0) rectangle (1.2,1.2);
    \node[font=\scriptsize\bfseries] at (0,0.6) {n (Emitter)};

    % p-bas
    \fill[red!15] (-1.2,1.2) rectangle (1.2,2.2);
    \node[font=\scriptsize\bfseries] at (0,1.7) {p (Bas)};

    % n-kollektor
    \fill[blue!15] (-1.2,2.2) rectangle (1.2,3.8);
    \node[font=\scriptsize\bfseries] at (0,3.0) {n (Kollektor)};

    % Anslutningar
    \draw[thick, ->] (-2.0,0.6)  node[left,font=\scriptsize] {E} -- (-1.2,0.6);
    \draw[thick, ->] (-2.0,1.7)  node[left,font=\scriptsize] {B} -- (-1.2,1.7);
    \draw[thick, ->] (1.2,3.0)   -- (2.0,3.0) node[right,font=\scriptsize] {C};
  \end{scope}

  % Skiljelinje
  \draw[gray!30] (5.0,0.2) -- (5.0,5.6);

  % ── Panel 2: Kretsymbol ──────────────────────────────────
  \begin{scope}[xshift=7.8cm, yshift=2.8cm]
    \node[font=\small\bfseries] at (0,2.5) {Kretsymbol (NPN)};

    % Baslinje
    \draw[thick] (-0.8,-1.2) -- (-0.8,1.2);
    % Bas
    \draw[thick] (-1.5,0) -- (-0.8,0) node[left=8pt,font=\scriptsize] {B};
    % Kollektor
    \draw[thick] (-0.8,0.6) -- (0.8,1.8) node[right,font=\scriptsize] {C};
    % Emitter med pil
    \draw[thick] (-0.8,-0.6) -- (0.8,-1.8);
    \draw[->, thick] (0.3,-1.2) -- (0.8,-1.8) node[right,font=\scriptsize] {E};

    % Strömpilar
    \draw[->, blue!60, dashed] (-1.8,0) node[left,font=\scriptsize,blue!60]{$I_B$ (liten)} -- (-1.2,0);
    \draw[->, red!60, dashed]  (0.8,2.2) node[above,font=\scriptsize,red!60]{$I_C$ (stor)} -- (0.8,1.9);
  \end{scope}

  % ── Panel 3: Förstärkningsformel ─────────────────────────
  \begin{scope}[xshift=11.5cm, yshift=2.8cm]
    \node[font=\small\bfseries] at (0,2.5) {Förstärkning};
    \node[font=\normalsize] at (0,1.5) {$I_C = \beta \cdot I_B$};
    \node[font=\normalsize] at (0,0.5) {$I_E = I_C + I_B$};
    \draw[gray!30] (-1.3,0) -- (1.3,0);
    \node[font=\scriptsize, gray] at (0,-0.4) {Typiskt $\beta$:};
    \node[font=\scriptsize] at (0,-0.9) {50 -- 500};
  \end{scope}

\end{tikzpicture}
\caption{NPN-transistorns struktur, kretsymbol och grundrelationer.
Basströmmen $I_B$ styr kollektorströmmen $I_C$ med förstärkningen
$\beta$. Emitterströmmen är summan av de två.}
\label{fig:npn_transistor}
\end{figure}

\subsubsection*{Övningsfrågor -- 1.9.3}

\begin{qblock}{4.}
\qgrund\quad Namnge de tre anslutningarna på en bipolär transistor
och förklara kortfattat vad var och en gör.
\end{qblock}

\begin{qblock}{5.}
\qrakna\quad En NPN-transistor har $\beta = 150$.
Basströmmen är $I_B = 0{,}5\,\text{mA}$.
Beräkna kollektorströmmen $I_C$ och emitterströmmen $I_E$.
\end{qblock}

\medskip
\noindent\rule{\textwidth}{0.4pt}
{\small\itshape Försök själv innan du läser svaret.}
\medskip

\begin{worked}[title={Lösning -- Räkneexempel 1.9.3}]
\[ I_C = \beta \cdot I_B = 150 \times 0{,}5\,\text{mA}
   = \textbf{75\,mA} \]
\[ I_E = I_C + I_B = 75 + 0{,}5 = \textbf{75{,}5\,mA} \]

Notera att $I_B$ är så liten i förhållande till $I_C$ att
$I_E \approx I_C$ i de flesta praktiska beräkningar.
\end{worked}

% ──────────────────────────────────────────────────────────────
\subsection{Transistorn som förstärkare}

När transistorn arbetar i sitt linjära område -- varken fullt
av eller fullt på -- kan den förstärka en växelspänningssignal.
Det är detta som används i radiomottagarens förstärkarsteg,
mikrofonförstärkare och slutsteg.

\begin{bluebox}[title={\emoji{🔑} Transistorförstärkaren -- arbetspunkten}]
För att transistorn ska kunna förstärka en växelsignal (t.ex. ett
ljud eller en radiosignal) måste den \textbf{förspännas} med
likström till en lämplig \textbf{arbetspunkt} (Q-point) mitt i
det linjära området.

\medskip
Utan förspänning arbetar transistorn bara som en switch --
den är antingen av eller på. Med rätt förspänning kan
växelsignalen ``rida'' ovanpå likströmmen och förstärkas
utan distortion.

\medskip
\textbf{Spänningsförstärkning:}
\[ A_v = \frac{U_{\text{ut}}}{U_{\text{in}}} = -\frac{R_C}{r_e} \]

där $r_e \approx 26\,\text{mV} / I_C$ är transistorns inre
emitterresistans vid rumstemperatur.

\medskip
{\small\itshape Minustecknet visar att signalen inverteras --
när ingången går upp går utgången ned. Det kallas fasförskjutning
på 180°.}
\end{bluebox}

% ── Figur: Förstärkarkoppling ────────────────────────────────
\begin{figure}[htbp]
\centering
\begin{tikzpicture}[font=\small, >=stealth]

  \fill[gray!6, rounded corners=4pt] (-0.4,-0.6) rectangle (10.4,6.4);

  \begin{scope}[xshift=5.0cm, yshift=2.8cm]

    % +Vcc
    \node[font=\small] at (1.8,3.2) {$+V_{CC}$};
    \draw[thick] (1.8,3.0) -- (1.8,2.2);

    % Rc
    \draw[thick] (1.8,2.2) -- (1.8,1.6);
    \draw[thick] (1.5,1.6) rectangle (2.1,0.8);
    \node[font=\scriptsize] at (2.6,1.2) {$R_C$};
    \draw[thick] (1.8,0.8) -- (1.8,0.4); % till kollektor

    % Transistor (förenklad)
    \draw[thick] (1.1,-0.8) -- (1.1,0.8); % baslinje
    \draw[thick] (1.1,0.4)  -- (1.8,0.4); % kollektor
    \draw[thick] (1.1,-0.4) -- (1.8,-1.0); % emitter
    \draw[->, thick] (1.4,-0.7) -- (1.8,-1.0); % emitterpil

% Bas
    \draw[thick] (-0.5,0) -- (1.1,0);          % baslinje fram till transistorn
    \node[font=\scriptsize] at (-1.1,0) {$u_{\text{in}}$};
    \draw[->, thin] (-0.8,0.25) -- (-0.8,-0.25); % kort inpil bredvid etiketten

    % Rb (förspänning)
    \draw[thick] (-0.5,3.0) -- (-0.5,0.6);
    \draw[thick] (-0.8,0.6) rectangle (-0.2,-0.2);
    \node[font=\scriptsize] at (-1.1,0.2) {$R_B$};
    \draw[thick] (-0.5,-0.2) -- (-0.5,0);
    \draw[thick] (-0.5,3.0) -- (1.8,3.0);

    % Emittermotstånd
    \draw[thick] (1.8,-1.0) -- (1.8,-1.6);
    \draw[thick] (1.5,-1.6) rectangle (2.1,-2.4);
    \node[font=\scriptsize] at (2.6,-2.0) {$R_E$};
    \draw[thick] (1.8,-2.4) -- (1.8,-3.0);

    % GND
    \draw[thick] (1.1,-3.0) -- (2.5,-3.0);
    \draw[thick] (1.3,-3.2) -- (2.3,-3.2);
    \draw[thick] (1.5,-3.4) -- (2.1,-3.4);

    % Utgång
    \draw[thick] (1.8,0.4) -- (3.5,0.4);
    \node[font=\scriptsize] at (3.9,0.4) {$u_{\text{ut}}$};
    \draw[->, thin] (3.8,0.7) -- (3.8,0.1);

  \end{scope}

\end{tikzpicture}
\caption{Emitterkopplad förstärkare -- den vanligaste grundkopplingen.
$R_B$ förspänner transistorn till arbetspunkten. $R_C$ omvandlar
kollektorströmsvariationen till en spänningssignal. $R_E$ stabiliserar
arbetspunkten mot temperaturvariationer.}
\label{fig:forstarkare}
\end{figure}

\subsubsection*{Övningsfrågor -- 1.9.4}

\begin{qblock}{6.}
\qprov\quad Vad menas med en transistors ``arbetspunkt''
och varför måste transistorn förspännas för att fungera
som förstärkare?
\end{qblock}

\begin{qblock}{F.}
\qfallan\quad Anna bygger en förstärkare och glömmer
att koppla in förspänningsmotståndet $R_B$.
Hon matar in en 1\,kHz sinussignal och mäter utgången --
men ser bara en starkt distorderad signal som liknar
en fyrkantsvåg.

Förklara vad som händer och varför utgångssignalen
ser ut som den gör.
\end{qblock}

\medskip
\noindent\rule{\textwidth}{0.4pt}
{\small\itshape Försök själv innan du läser svaret.}
\medskip

\begin{worked}[title={Lösning -- Fallgropsfråga 1.9.4\,F}]
Utan $R_B$ får transistorns bas ingen förspänning -- basen
ligger på 0\,V (jord). Transistorn är i \textbf{spärrläge}
(helt av) vid vila.

När insignalen går positiv öppnar transistorn och leder.
När insignalen går negativ stänger den. Transistorn beter sig
som en \textbf{switch} istället för en förstärkare -- den klipper
bort den negativa halvcykeln helt och hållet.

Resultatet är en halvvågslikriktad signal som liknar en
fyrkantsvåg -- inte en förstärkt sinusvåg.

\textbf{Lärdomen:} Förspänning är inte valfritt. Utan rätt
arbetspunkt kan transistorn inte förstärka -- den kan bara
switcha.
\end{worked}

% ──────────────────────────────────────────────────────────────
\subsection{Transistorn som switch}

I digital elektronik och styrkretsars används transistorn nästan
uteslutande som en switch -- antingen fullt av (spärrläge) eller
fullt på (mättnadsläge). Det är snabbt, effektivt och enkelt
att styra digitalt.

\begin{bluebox}[title={\emoji{🔑} Transistorn som switch -- två lägen}]
\textbf{Spärrläge (OFF):}
$I_B = 0$ $\Rightarrow$ $I_C \approx 0$ $\Rightarrow$
Transistorn leder inte. Hela $V_{CC}$ faller över transistorn.

\medskip
\textbf{Mättnadsläge (ON):}
$I_B$ tillräckligt stor $\Rightarrow$ transistorn är fullt
öppen. Spänningsfallet över transistorn är mycket litet
(typiskt 0,1--0,3\,V). Nästan all $V_{CC}$ faller över lasten.

\medskip
\textbf{Regel för mättnad:}
\[ I_B \geq \frac{I_{C(\text{sat})}}{\beta} \]
Basströmmen måste vara tillräckligt stor för att ``öppna kranen
helt''. I praktiken räknar man med en säkerhetsfaktor 5--10.

\medskip
\textbf{Användning i radio:}
\begin{itemize}
  \item T/R-relädrivare (kopplar TX/RX vid sändning)
  \item PTT-krets
  \item Frekvensdelare i PLL
  \item Morsenyckling (CW)
\end{itemize}
\end{bluebox}

\subsubsection*{Övningsfrågor -- 1.9.5}

\begin{qblock}{7.}
\qrakna\quad En NPN-transistor ska driva ett relä med
$R_{\text{relä}} = 100\,\Omega$ från $V_{CC} = 12\,\text{V}$.
Transistorns $\beta = 100$.

a) Beräkna kollektorströmmen $I_C$ när transistorn är
   i mättnad.\\
b) Hur stor basström $I_B$ behövs minst för att säkerställa
   mättnad?\\
c) Om baskretsen drivs från en mikroprocessor-utgång via ett
   motstånd $R_B$, och processorn ger 3,3\,V -- vilket värde
   ska $R_B$ ha? (Använd $V_{BE} = 0{,}7\,\text{V}$ och
   en säkerhetsfaktor på 5.)
\end{qblock}

\medskip
\noindent\rule{\textwidth}{0.4pt}
{\small\itshape Försök själv innan du läser svaret.}
\medskip

\begin{worked}[title={Lösning -- Räkneexempel 1.9.5}]
a) Kollektorströmmen i mättnad (Ohms lag, spänningsfall
   över transistorn försummas):
\[ I_C \approx \frac{V_{CC}}{R_{\text{relä}}}
   = \frac{12}{100} = \textbf{120\,mA} \]

b) Minsta basström för mättnad:
\[ I_{B(\text{min})} = \frac{I_C}{\beta}
   = \frac{120\,\text{mA}}{100} = 1{,}2\,\text{mA} \]

c) Med säkerhetsfaktor 5: $I_B = 5 \times 1{,}2 = 6\,\text{mA}$

Spänning över $R_B$:
$U_{R_B} = 3{,}3 - V_{BE} = 3{,}3 - 0{,}7 = 2{,}6\,\text{V}$

\[ R_B = \frac{U_{R_B}}{I_B}
   = \frac{2{,}6}{6 \times 10^{-3}}
   \approx 433\,\Omega \]

Närmaste standardvärde: \textbf{390\,Ω} eller \textbf{470\,Ω}.
\end{worked}

% ──────────────────────────────────────────────────────────────
\subsection{Transistortyper -- BJT och FET}

Hittills har vi pratat om den bipolära transistorn (BJT).
Det finns en annan viktig familj: fälteffekttransistorn (FET),
som styr strömmen med \emph{spänning} istället för ström.

\begin{bluebox}[title={\emoji{🔑} BJT vs FET -- en jämförelse}]
\begin{center}
\begin{tabular}{lll}
\rowcolor{greydark}
\tblhdr{Egenskap} & \tblhdr{BJT} & \tblhdr{FET (MOSFET)} \\
Styrs av & Basström ($I_B$) & Gattspänning ($V_{GS}$) \\
\rowcolor{greylight}
Ingångsimpedans & Låg (k$\Omega$) & Mycket hög (G$\Omega$) \\
Strömförstärkning & $\beta$ (50--500) & Transkonduktans $g_m$ \\
\rowcolor{greylight}
Brus & Något högre & Lägre vid låga frekvenser \\
Switchhastighet & Snabb & Mycket snabb \\
\rowcolor{greylight}
Typisk användning & RF-förstärkare, slutsteg & Digitalt, switchar, LNA \\
\end{tabular}
\end{center}

\medskip
\textbf{Anslutningar:}\\
BJT: Bas (B) -- Kollektor (C) -- Emitter (E)\\
FET: Gatt (G) -- Drain (D) -- Source (S)

\medskip
{\small\itshape I moderna transceivrar används ofta MOSFET-transistorer
i slutsteget på grund av deras höga effektivitet och robusthet.
BJT dominerar fortfarande i linjära RF-förstärkare och
lågbrusapplikationer.}
\end{bluebox}

\subsubsection*{Övningsfrågor -- 1.9.6}

\begin{qblock}{8.}
\qprov\quad Vad är den viktigaste skillnaden mellan hur
en BJT och en MOSFET styrs? Vilken typ har högst
ingångsimpedans och varför är det en fördel?
\end{qblock}

\begin{qblock}{F.}
\qfallan\quad Erik väljer en BJT-transistor med $\beta = 200$
till sin förstärkare och resonerar: ``Ju högre $\beta$, desto
bättre förstärkning och desto stabilare krets.''

Vad är problematiskt med detta resonemang?
\end{qblock}

\medskip
\noindent\rule{\textwidth}{0.4pt}
{\small\itshape Försök själv innan du läser svaret.}
\medskip

\begin{worked}[title={Lösning -- Fallgropsfråga 1.9.6\,F}]
Hög $\beta$ ger \emph{inte} automatiskt en stabilare krets --
tvärtom kan det orsaka problem:

\textbf{1. $\beta$ varierar kraftigt} mellan exemplar av samma
transistortyp (ofta faktorn 3--5) och med temperatur. En krets
dimensionerad för $\beta = 200$ kan fungera dåligt med en annan
transistor av samma typ som råkar ha $\beta = 80$.

\textbf{2. Återkoppling, inte $\beta$, ger stabilitet.}
Väldesignade förstärkare använder negativ återkoppling
(t.ex. $R_E$ utan avkopplingskondensator) för att göra
förstärkningen \emph{oberoende} av transistorns exakta $\beta$.
En krets som förlitar sig på ett specifikt högt $\beta$-värde
är i praktiken instabil.

\textbf{3. Hög $\beta$ kan ge oscillationsrisk} i RF-förstärkare
om kretsen inte är ordentligt stabiliserad.

\textbf{Slutsats:} Välj transistor utifrån frekvensegenskaper,
effekt och brusfaktor -- inte maximalt $\beta$.
\end{worked}

% ── KAPITEL 1 AVSLUTNING ──────────────────────────────────
\section*{Grunden är lagd}
\addcontentsline{toc}{section}{Grunden är lagd}

Du har nu arbetat dig igenom det som många anser vara det tyngsta kapitlet i hela handboken. Kanske kändes det tungt på vägen – det ska det göra. Det betyder att du faktiskt tänkte. Ge dig själv en välförtjänt paus – du har precis byggt det fundament som resten av utbildningen vilar på.

Tänk på det så här: en radioamatör utan elektronikkunskap är som en bilmekaniker som aldrig lärt sig vad en motor är. Nu vet du vad som händer inne i apparaterna, inte bara hur du vrider på knopparna.

\noindent
Men kom ihåg – formlerna är bara \textit{kartan}. Det verkliga lärandet sitter i förståelsen bakom dem: varför motståndet bromsar, varför kondensatorn lagrar och släpper, varför spolen gör precis tvärtom. Den förståelsen har du nu.

\medskip
\noindent
Och här är något värt att stanna upp vid: det du just lärt dig är inte \enquote{lite enkel teori för att klara ett prov}. Det är samma grund som utbildade elektriker och elektronikingenjörer har med sig från sina första år. Du har arbetat dig igenom reaktans, impedans, resonans och filter -- begrepp som många aldrig ens stöter på. Det är faktiskt ganska imponerande.

\medskip
\noindent
Det roliga är att det svåraste nu ligger bakom dig – det som kommer härnäst är där elektroniken börjar leva.

\begin{bluebox}[title={}]
\textbf{Du har fundamentet. Nu bygger vi vidare.}

\smallskip
I nästa kapitel kliver vi in i radiotekniken – och du kommer att märka att allt du just lärt dig plötsligt dyker upp igen, fast i nya sammanhang. Så är det med kunskap: den förräntar sig.
\end{bluebox}

% ═══════════════════════════════════════════════════════════
%  FACIT – KAPITEL 1
% ═══════════════════════════════════════════════════════════
\chapter*{Facit – Kapitel 1}
\addcontentsline{toc}{chapter}{Facit – Kapitel 1}
\markboth{Facit – Kapitel 1}{}

% ── Facit 1.1 ─────────────────────────────────────────────
\subsection*{1.1 – Elektrisk ström och spänning}

\begin{facitblock}{1}
\qgrund\quad Vilken enhet mäts elektrisk ström i? Vad är beteckningen 
för ström i formler? Ange också de vanligaste underenheterna och deras 
relationer till grundenheten.
\tcblower
Elektrisk ström mäts i \textbf{Ampere (A)}, betecknas \textbf{I}. 
Underenheter: µA = $10^{-6}$ A, mA = $10^{-3}$ A, kA = 1\,000 A. 
\textbf{1 A = 1\,000 mA = 1\,000\,000 µA.}
\end{facitblock}

\begin{facitblock}{2}
\qrakna\quad Omvandla följande strömvärden till ampere: a) 250 mA 
\quad b) 4\,200 µA \quad c) 0,75 kA \quad d) 85 mA. Omvandla 
dessutom 0,035 A till mA och till µA.
\tcblower
a) 250 mA = \textbf{0,25 A} \quad b) 4\,200 µA = \textbf{0,0042 A}\\
c) 0,75 kA = \textbf{750 A} \quad d) 85 mA = \textbf{0,085 A}\\
0,035 A = \textbf{35 mA} = \textbf{35\,000 µA}
\end{facitblock}

\begin{facitblock}{3}
\qgrund\quad Vilken spänning är standard för amatörradioutrustning, 
och varför är det \emph{inte} exakt 12 V trots att de flesta mobila 
installationer körs från ett 12 V-bilbatteri?
\tcblower
Standard är \textbf{13,8 V DC}. Ett fulladdat 12 V-blybatteri är 
12,6–13,2 V och med laddare 13,6–14,4 V. Radion konstrueras för 
hela intervallet med 13,8 V som mittenvärde. En radio som bara matas med 11–12 V levererar märkbart lägre sändeffekt.
\end{facitblock}

\begin{facitblock}{4}
\qgrund\quad Vad är skillnaden mellan likström (DC) och växelström 
(AC)? Ge ett konkret praktiskt exempel på vardera typ i en typisk 
radiostation.
\tcblower
DC flödar alltid åt samma håll. AC byter riktning periodiskt 
(50 Hz i Sverige).\\
DC-exempel: Nätaggregatet levererar 13,8 V DC till HF-transceivern.\\
AC-exempel: Vägguttaget ger 230 V AC som nätaggregatet omvandlar 
till DC.
\end{facitblock}

\begin{facitblock}{5}
\qrakna\quad Du har ett 12 Ah-batteri och en handhållen VHF-radio 
som förbrukar 1,5 A vid sändning och 0,3 A vid mottagning. Under 
ett fältpass sänder du 20\,\% och tar emot 80\,\% av tiden. 
a) Vad är den genomsnittliga strömförbrukningen? 
b) Hur länge räcker batteriet?
\tcblower
a) $I = 0{,}20 \times 1{,}5 + 0{,}80 \times 0{,}3 = 0{,}30 + 0{,}24 
= \mathbf{0{,}54\ A}$\\
b) $t = 12\ \text{Ah} \div 0{,}54\ \text{A} \approx 
\mathbf{22\ \text{timmar}}$
\end{facitblock}

\begin{facitblock}{6}
\qprov\quad Vid nödkommunikation ska du driva en HF-transceiver 
(13,8 V, max 22 A vid sändning / 2 A vid mottagning) från en 
60 Ah-batteribank. Du sänder 15 min och tar emot 45 min per timme. 
a) Genomsnittlig ström? b) Drifttid vid 50\,\% urladdningsgräns? 
c) Räcker ett 25 A nätaggregat?
\tcblower
a) $I = 0{,}25 \times 22 + 0{,}75 \times 2 = 5{,}5 + 1{,}5 = 
\mathbf{7{,}0\ A}$\\
b) Användbar kapacitet: $60 \times 0{,}50 = 30$ Ah.\ \ 
Drifttid: $30 \div 7{,}0 \approx \mathbf{4{,}3\ \text{timmar}}$\\
c) Ja – max 22 A $<$ 25 A, men marginalen är liten (3 A). 
Aggregatet arbetar nära max under sändning och kan bli varmt. 
Ett 30 A-aggregat rekommenderas för längre drift.
\end{facitblock}

% ── Facit 1.2 ─────────────────────────────────────────────
\subsection*{1.2 – Resistans och Ohms lag}

\begin{facitblock}{1}
\qgrund\quad Ange Ohms lag i alla tre formerna (för U, för I och för R). 
Vad betecknar varje bokstav och i vilken enhet mäts de? Rita 
URI-triangeln och förklara hur den används.
\tcblower
$U = R \times I$ \quad | \quad $I = U \div R$ \quad | \quad $R = U \div I$\\
U = spänning (V), R = resistans ($\Omega$), I = ström (A).\\
URI-triangeln: täck det du söker – resten av triangeln visar formeln. 
Täck U $\rightarrow$ multiplicera. Täck R eller I $\rightarrow$ dividera.
\end{facitblock}

\begin{facitblock}{2}
\qrakna\quad Beräkna det okända värdet (konvertera enheter om det behövs):
\begin{center}
\begin{tabular}{ll}
a) $U = ?$, $R = 330\,\Omega$, $I = 45\,\text{mA}$ &
b) $I = ?$, $U = 13{,}8\,\text{V}$, $R = 4{,}7\,\text{k}\Omega$ \\[4pt]
c) $R = ?$, $U = 6\,\text{V}$, $I = 12\,\text{mA}$ &
d) $U = ?$, $R = 50\,\Omega$, $I = 2\,\text{A}$
\end{tabular}
\end{center}
\tcblower
a) $U = 330 \times 0{,}045 = \mathbf{14{,}85\ \text{V}}$\\
b) $I = 13{,}8 \div 4\,700 \approx \mathbf{2{,}94\ \text{mA}}$\\
c) $R = 6 \div 0{,}012 = \mathbf{500\ \Omega}$\\
d) $U = 50 \times 2 = \mathbf{100\ \text{V}}$
\end{facitblock}

\begin{facitblock}{3}
\qrakna\quad En antennrotors matningskabel har DC-resistansen 
0,8\,$\Omega$ (tur och retur), rotorn drar 3 A. 
a) Vilket spänningsfall uppstår? b) Nätaggregatet ger 13,8 V – 
vilken spänning når rotorn? c) Rotorn kräver minst 11 V. 
Godkänd installation?
\tcblower
a) $U_{fall} = R \times I = 0{,}8 \times 3 = \mathbf{2{,}4\ \text{V}}$\\
b) $U_{rotor} = 13{,}8 - 2{,}4 = \mathbf{11{,}4\ \text{V}}$\\
c) Ja – 11,4 V $>$ 11 V, men marginalen är bara 0,4 V. Vid sjunkande 
batterispänning eller åldrad kabel kan gränsen underskridas. 
Rekommendation: kortare eller tjockare kabel.
\end{facitblock}

\newpage

\begin{facitblock}{4}
\qgrund\quad Rangordna dessa material från lägst till högst resistans: 
silver, plast, kol, koppar, gummi, aluminium. Förklara sedan varför 
silver inte används i kablar trots att det leder bättre än koppar.
\tcblower
Lägst $\rightarrow$ högst: Silver, Koppar, Aluminium, Kol, Gummi, Plast.\\
Silver leder ca 6\,\% bättre än koppar men är avsevärt dyrare. 
Skillnaden motiverar inte kostnaden i kablar. Koppar ger bäst 
pris/prestanda och är universell standard. Silver används däremot som 
ytbeläggning på kontakter och i militär-/rymdtillämpningar.
\end{facitblock}

\begin{facitblock}{5}
\qgrund\quad En koppartråd har resistansen 2 $\Omega$. Vad händer med 
resistansen om du a) dubblerar längden? b) dubblerar tvärsnittsarean? 
c) gör båda samtidigt? Motivera utan att räkna.
\tcblower
a) Dubbel längd $\rightarrow$ dubbel resistans: \textbf{4\,$\Omega$}. 
Fler atomer i strömvägen bromsar elektronerna mer.\\
b) Dubbel tvärsnittsarea $\rightarrow$ hälften så hög resistans: 
\textbf{1\,$\Omega$}. Bredare väg ger elektronerna fler parallella 
banor att röra sig i.\\
c) Båda samtidigt tar ut varandra – resistansen är oförändrad: 
\textbf{2\,$\Omega$}.
\end{facitblock}

\begin{facitblock}{6}
\qprov\quad En SWR-meter mäter den framåtgående spänningen till 70,7 V 
(RMS) i ett 50 $\Omega$-system. a) Beräkna strömmen genom systemet. 
b) Beräkna effekten ($P = U \times I$). c) En tekniker påstår att 
koaxkabelns resistans (3 dB förlust vid denna effektnivå) är 
acceptabelt. Instämmer du? Motivera med beräkning.
\tcblower
a) $I = U \div R = 70{,}7 \div 50 = \mathbf{1{,}414\ \text{A}}$\\
b) $P = U \times I = 70{,}7 \times 1{,}414 \approx \mathbf{100\ \text{W}}$\\
c) Nej – $-3$ dB innebär att \textbf{halva effekten} försvinner som 
värme i kabeln. Av 100 W når bara 50 W antennen. I praktiken eftersträvas 
max 1 dB förlust ($\approx$ 80 W når antennen). Lösning: byt till kabel 
med lägre förlust (t.ex.\ Aircell 7, RG-213 eller Ecoflex 10) eller 
förkorta kabellängden.
\end{facitblock}

\newpage

\begin{facitblock}{7}
\qprov\quad Varför anges antennimpedansen för de flesta 
amatörradioapparater till 50 $\Omega$ och inte 75 $\Omega$? Beskriv 
det fysikaliska resonemang som leder till kompromissen 50 $\Omega$, 
och nämn vid vilken impedans en luftfylld koaxkabel har lägst dämpning 
respektive högst effektkapacitet.
\tcblower
I en luftfylld koaxkabel gäller:\\
Lägst dämpning vid ca \textbf{77\,$\Omega$}. Högst effektkapacitet vid 
ca \textbf{30\,$\Omega$}.\\[4pt]
Kompromissen beräknas som det geometriska medelvärdet:
\[ \sqrt{77 \times 30} \approx 48\ \Omega \rightarrow 
\text{avrundat till \textbf{50\,$\Omega$}} \]
Standarden etablerades under andra världskriget och används nu 
universellt för amatörradio, mobiltelefoni och mätutrustning.\\[4pt]
TV-system väljer 75\,$\Omega$ eftersom de prioriterar lägsta möjliga 
dämpning vid låga effektnivåer – effektkapaciteten är mindre viktig.
\end{facitblock}

% ── Facit 1.3 ─────────────────────────────────────────────
\subsection*{1.3 – Effekt och energi}

\begin{facitblock}{1}
\qgrund\quad Förklara med egna ord vad effekt är och hur den skiljer sig från energi. Vilken enhet mäts effekt i och vad heter den uppkallad efter?
\tcblower
Effekt är takten – hur snabbt energi omvandlas per sekund. Energi är den totala mängden arbete som utförts. Sambandet: $E = P \times t$.

Effekt mäts i \textbf{Watt (W)}, uppkallat efter James Watt (1736–1819). Exempel: en 100 W sändare som körs i 2 timmar förbrukar 200 Wh energi.
\end{facitblock}

\begin{facitblock}{2}
\qgrund\quad Du har tre effektformler att välja mellan. I vilka situationer använder du respektive formel? Ge ett radioexempel för varje.
\tcblower
$P = U \times I$ – när du vet spänning och ström.\\
Exempel: 13,8 V $\times$ 22 A $\approx$ 304 W (nätaggregat till transceiver).\\[4pt]
$P = I^2 \times R$ – när du vet ström och resistans.\\
Exempel: 2 A genom 50 $\Omega$ dummylast $\rightarrow$ $4 \times 50 = 200$ W.\\[4pt]
$P = U^2 \div R$ – när du vet spänning och resistans.\\
Exempel: 70,7 V RMS i 50 $\Omega$ $\rightarrow$ $5\,000 \div 50 = 100$ W.
\end{facitblock}

\begin{facitblock}{3}
\qrakna\quad Beräkna effekten. Ange vilken formel du valde och varför:
\begin{center}
\begin{tabular}{ll}
a) $U = 13{,}8$ V, $I = 15$ A &
b) $I = 500$ mA, $R = 8\,\Omega$ \\[4pt]
c) $U = 50$ V, $R = 50\,\Omega$ &
d) $U = 230$ V, $I = 4{,}35$ A
\end{tabular}
\end{center}
\tcblower
a) Vet U och I $\rightarrow$ $P = U \times I = 13{,}8 \times 15 = \mathbf{207\,\text{W}}$\\
b) Vet I och R $\rightarrow$ $P = I^2 \times R$. Konvertera: 500 mA = 0,5 A.
$P = 0{,}5^2 \times 8 = 0{,}25 \times 8 = \mathbf{2\,\text{W}}$\\
c) Vet U och R $\rightarrow$ $P = U^2 \div R = 50^2 \div 50 = 2\,500 \div 50 = \mathbf{50\,\text{W}}$\\
d) Vet U och I $\rightarrow$ $P = U \times I = 230 \times 4{,}35 \approx \mathbf{1\,000\,\text{W} = 1\,\text{kW}}$
\end{facitblock}

\begin{facitblock}{4}
\qrakna\quad Din sändare ger 100 W. Koaxkabeln har 3 dB förlust. a) Hur många watt når antennen? b) Du byter till en kabel med 1 dB förlust – hur många watt når antennen då? c) Hur många dB antennförstärkning kompenserar 3 dB kabelförlust?
\tcblower
a) $-3\,\text{dB} = \times 0{,}5 \rightarrow \mathbf{50\,\text{W}}$ når antennen.\\
b) $-1\,\text{dB} \approx \times 0{,}79 \rightarrow 100 \times 0{,}79 \approx \mathbf{79\,\text{W}}$ når antennen.\\
Skillnad: $79 - 50 = \mathbf{29\,\text{W}}$ mer med den bättre kabeln – nästan 60\,\% mer utstrålad effekt.\\
c) \textbf{+3 dB} antennförstärkning. dB-värden adderas: $-3 + 3 = 0\,\text{dB}$ nettoresultat.
\end{facitblock}

\begin{facitblock}{5}
\qprov\quad En radioamatör ökar från 100 W till 400 W. Hur många dB är ökningen? Kan han förvänta sig dubbelt så lång räckvidd? Motivera.
\tcblower
$400\,\text{W} \div 100\,\text{W} = 4 \rightarrow \mathbf{+6\,\text{dB}}$ (fyrdubbel effekt = +6 dB).\\[4pt]
I teorin: ja – i fri rymd krävs just fyrdubbel effekt för att fördubbla räckvidden.\\
I praktiken: sällan. Räckvidden beror på terräng, antennförstärkning och utbredningsförhållanden. Motstationen märker skillnaden som ungefär ett S-steg – märkbart men inte dramatiskt. En bättre antenn ger ofta mer räckvidd per investerad krona.
\end{facitblock}

\begin{facitblock}{6}
\qrakna\quad En HF-transceiver levererar 100 W RF och drar 22 A från 13,8 V. a) Total ineffekt? b) Verkningsgrad? c) Effekt till värme? d) Behövs fläkt om kylflänsen klarar 100 W?
\tcblower
a) $P_{in} = U \times I = 13{,}8 \times 22 \approx \mathbf{303{,}6\,\text{W}}$\\
b) $\eta = \dfrac{100}{303{,}6} \times 100\,\% \approx \mathbf{32{,}9\,\%}$\\
c) $P_{värme} = 303{,}6 - 100 \approx \mathbf{203{,}6\,\text{W}}$\\
d) Ja – 203,6 W överstiger kylflänens 100 W med god marginal. Utan fläkt riskerar slutsteget att överhettas.
\end{facitblock}

\begin{facitblock}{7}
\qprov\quad En 470 $\Omega$ resistor seriekopplad med en LED, 12 V, 20 mA. a) Effekt med $P = I^2 \times R$. b) Verifiera med $P = U_R \times I$. c) Vilken standardeffekt?
\tcblower
a) Konvertera: 20 mA = 0,02 A.
$P = 0{,}02^2 \times 470 = 0{,}0004 \times 470 = \mathbf{0{,}188\,\text{W}}$\\[4pt]
b) $U_R = I \times R = 0{,}02 \times 470 = 9{,}4\,\text{V}$\\
$P = 9{,}4 \times 0{,}02 = \mathbf{0{,}188\,\text{W}}$ \checkmark Stämmer exakt.\\[4pt]
c) 0,188 W $<$ $\frac{1}{4}$ W (0,25 W) – tekniskt räcker $\frac{1}{4}$ W, men marginalen är bara 25\,\%. En \textbf{$\frac{1}{2}$ W} resistor rekommenderas för god säkerhetsmarginal och längre livslängd.
\end{facitblock}

% ── Facit 1.4 ─────────────────────────────────────────────
\subsection*{1.4 -- Kondensatorer}

\begin{facitblock}{1}
\qgrund\quad Förklara med egna ord varför en kondensator blockerar likström men inte växelström. Använd gummihinna-analogin i förklaringen.
\tcblower
Tänk dig en gummihinna spänd tvärs över ett vattenrör. Likström (DC) trycker vattnet åt ett håll -- hinnan töjer sig tills den inte ger med sig mer. Sedan stopp: inget vatten rinner längre. Kondensatorn har laddats upp och strömmen upphör.

Växelström (AC) byter riktning hela tiden -- hinnan trycks fram och tillbaka och vibrerar i takt med signalen. Energin passerar igenom trots att inga elektroner faktiskt tar sig förbi hinnan. Ju högre frekvens, desto snabbare vibrerar hinnan och desto friare passerar signalen.
\end{facitblock}

\begin{facitblock}{2}
\qrakna\quad Beräkna $X_C$ för en 100\,pF kondensator vid: a) 7\,MHz \quad b) 144\,MHz. \quad c) Vad händer med reaktansen när frekvensen ökar?
\tcblower
$C = 100\,\text{pF} = 10^{-10}\,\text{F}$

a) Vid 7\,MHz:
\[ X_C = \frac{1}{6{,}28 \times 7 \times 10^6 \times 10^{-10}} = \frac{1}{4{,}396 \times 10^{-4}} \approx \mathbf{227\,\Omega} \]

b) Vid 144\,MHz:
\[ X_C = \frac{1}{6{,}28 \times 144 \times 10^6 \times 10^{-10}} = \frac{1}{9{,}043 \times 10^{-3}} \approx \mathbf{11\,\Omega} \]

c) Reaktansen \textbf{minskar} när frekvensen ökar. I formeln $X_C = 1/(2\pi fC)$ står frekvensen i nämnaren -- större $f$ ger mindre $X_C$. Vid 144\,MHz är reaktansen bara 11\,$\Omega$ jämfört med 227\,$\Omega$ vid 7\,MHz -- kondensatorn släpper igenom högfrekventa signaler mycket lättare.
\end{facitblock}

\begin{facitblock}{3}
\qrakna\quad En kopplingskondensator på 10\,µF i audiovägen. a) $X_C$ vid 20\,Hz. b) $X_C$ vid 1\,000\,Hz. c) Är 10\,µF ett bra val?
\tcblower
$C = 10\,\mu\text{F} = 10^{-5}\,\text{F}$

a) Vid 20\,Hz:
\[ X_C = \frac{1}{6{,}28 \times 20 \times 10^{-5}} = \frac{1}{1{,}256 \times 10^{-3}} \approx \mathbf{796\,\Omega} \]

b) Vid 1\,000\,Hz:
\[ X_C = \frac{1}{6{,}28 \times 1\,000 \times 10^{-5}} = \frac{1}{6{,}28 \times 10^{-2}} \approx \mathbf{15{,}9\,\Omega} \]

c) 10\,µF fungerar bra för mellanregister och diskant -- vid 1\,kHz är reaktansen bara 16\,$\Omega$, vilket är försumbart. Däremot är 796\,$\Omega$ vid 20\,Hz problematiskt -- den höga reaktansen dämpar basen märkbart. För fullständig audiorespons ner till 20\,Hz bör man välja 47--100\,µF. I radioteknik, där talfrekvenser (300--3\,000\,Hz) dominerar, är 10\,µF dock helt tillräckligt.
\end{facitblock}

\begin{facitblock}{4}
\qgrund\quad Vilken kondensatortyp har polaritet och varför är det kritiskt att ansluta den rätt? Hur känner man igen plus- och minussidan?
\tcblower
\textbf{Elektrolytkondensatorer} är polariserade -- de har en definierad plus- och minussida. Fel anslutning kan leda till överhettning, läckage av kemikalier och i värsta fall en kraftig smäll.

Identifiering:
\begin{itemize}
  \item Grå rand med minustecken (${-}$) på höljet markerar minuspolen.
  \item På nya komponenter med ben är det \textbf{kortare benet minus} och det längre benet plus.
\end{itemize}
\end{facitblock}

\begin{facitblock}{5}
\qprov\quad Förklara vad en bypass-kondensator gör. Varför kombinerar man ofta 100\,nF keramisk + 100\,µF elektrolyt?
\tcblower
En bypass-kondensator leder bort högfrekventa störningar och snabba spänningsvariationer innan de når känsliga delar av kretsen. På så sätt hålls matningsspänningen stabil.

De två typerna kompletterar varandra:
\begin{itemize}
  \item Den \textbf{lilla keramiska (100\,nF)} har låg reaktans vid höga frekvenser -- du vet från 1.4.2 att reaktansen minskar när frekvensen ökar. Den reagerar snabbt på snabba störningar.
  \item Den \textbf{stora elektrolyten (100\,µF)} har hög kapacitans och tar hand om långsammare variationer och plötsliga belastningar.
\end{itemize}
Tillsammans täcker de var sin del av problemet.
\end{facitblock}

\begin{facitblock}{6}
\qgrund\quad Du ska välja kondensator till tre olika situationer. Vilken typ passar bäst och varför?
\tcblower
a) \textbf{Elektrolytkondensator.} Den har hög kapacitans -- upp till 10\,000\,µF -- och kan lagra tillräckligt med laddning för att jämna ut spänningsfall vid varierande belastning. Kom ihåg att ansluta den rätt: polariteten är kritisk.

b) \textbf{Keramisk kondensator.} Den är liten, snabb och har låg reaktans vid höga frekvenser. Placeras direkt intill det känsliga chippet så att störningarna leds bort innan de hinner orsaka problem.

c) \textbf{Variabel kondensator (trimmer).} Den justeras med ett skruvmejselstag och låses sedan i det inställda värdet. Rätt verktyg när kapacitansen behöver finjusteras vid inbyggnad.
\end{facitblock}

\begin{facitblock}{7}
\qrakna\quad En 47\,µF kondensator laddas genom en 470\,$\Omega$ resistor. a) Beräkna $\tau$. b) Hur lång tid tills praktiskt fulladdat? c) Hur stor andel finns kvar efter $2\tau$?
\tcblower
a)
\[ \tau = R \times C = 470\,\Omega \times 47 \times 10^{-6}\,\text{F} \approx \mathbf{22\,\text{ms}} \]

b) Praktiskt fulladdat efter $5\tau$:
\[ 5 \times 22\,\text{ms} = \mathbf{110\,\text{ms}} \]

c) Efter $2\tau$ är kondensatorn laddad till 86\,\% -- det betyder att \textbf{14\,\%} av laddningen återstår vid urladdning, eller omvänt att 14\,\% saknas vid uppladdning.
\end{facitblock}

% ── Facit 1.5 ─────────────────────────────────────────────
\subsection*{1.5 -- Spolar och induktans}

\begin{facitblock}{1}
\qgrund\quad Förklara med vattenhjuls-analogin varför en 
spole bromsar växelström men inte likström. Varför ökar 
motståndet när frekvensen ökar?
\tcblower
Tänk dig ett tungt vattenhjul i ett rör. Likström driver 
hjulet åt ett håll -- det startar trögt men snurrar sedan 
fritt med jämn hastighet. Ingen broms efter uppstarten: 
strömmen flödar fritt.

Växelström vill ständigt byta riktning på vattenflödet. 
Det tunga hjulet motstår varje riktningsbyte -- det måste 
bromsas, stanna och accelerera åt andra hållet. Ju högre 
frekvens (fler riktningsbyten per sekund), desto mer 
motstånd bjuder hjulet. Det är därför spolens reaktans 
ökar med frekvensen.
\end{facitblock}

\begin{facitblock}{2}
\qrakna\quad Beräkna $X_L$ för en 10\,µH spole vid: 
a) 3,5\,MHz \quad b) 14\,MHz \quad c) 144\,MHz. \quad 
d) Vad händer med $X_L$ när frekvensen ökar?
\tcblower
$L = 10\,\mu\text{H} = 10^{-5}\,\text{H}$

a) Vid 3,5\,MHz:
\[ X_L = 6{,}28 \times 3{,}5 \times 10^6 \times 10^{-5} \approx \mathbf{220\,\Omega} \]

b) Vid 14\,MHz:
\[ X_L = 6{,}28 \times 14 \times 10^6 \times 10^{-5} \approx \mathbf{880\,\Omega} \]

c) Vid 144\,MHz:
\[ X_L = 6{,}28 \times 144 \times 10^6 \times 10^{-5} \approx \mathbf{9\,043\,\Omega} \]

d) $X_L$ ökar proportionellt med frekvensen -- dubbel 
frekvens ger dubbel reaktans. Vid 144\,MHz är reaktansen 
över 9\,k$\Omega$ jämfört med 220\,$\Omega$ vid 3,5\,MHz.
\end{facitblock}

\begin{facitblock}{3}
\qrakna\quad Du behöver en spole med minst 500\,$\Omega$ 
reaktans vid 7\,MHz. Hur stor induktans krävs?
\tcblower
Lös ut $L$ ur formeln $X_L = 2\pi f L$:

\[ L = \frac{X_L}{2\pi f} = \frac{500}{6{,}28 \times 7 \times 10^6} \approx \mathbf{11{,}4\,\mu\text{H}} \]

Välj en standardkomponent på \textbf{12\,µH} eller högre 
för god marginal.
\end{facitblock}

\begin{facitblock}{4}
\qgrund\quad Fyll i tabellen utan att titta i texten.
\tcblower
\begin{center}
\begin{tabular}{L{3.5cm}L{3cm}L{3cm}}
\rowcolor{greydark}
\tblhdr{Egenskap} & \tblhdr{Kondensator} & \tblhdr{Spole} \\
Blockerar DC/AC? & Blockerar DC, släpper AC & Släpper DC, bromsar AC \\
\rowcolor{greylight}
Reaktans vid hög f? & Låg (minskar) & Hög (ökar) \\
Reaktansformel & $X_C = 1/(2\pi fC)$ & $X_L = 2\pi fL$ \\
\rowcolor{greylight}
Lagrar energi i & Elektriskt fält & Magnetfält \\
\end{tabular}
\end{center}
\end{facitblock}

\begin{facitblock}{5}
\qgrund\quad Du ska välja spoltyp till tre situationer. 
Vilken passar bäst och varför?
\tcblower
a) \textbf{Ferritkärnespole (toroid).} Hög induktans per 
varv i ett kompakt utförande. Förlusterna är något högre 
än luftkärna men acceptabla när utrymmet är begränsat.

b) \textbf{Luftkärnespole.} Låga förluster och stabila 
egenskaper över tid. Inget kärnmaterial som påverkar 
induktansen vid temperaturändringar.

c) \textbf{Variabel spole.} Kärnan skruvas in eller ut 
för att justera induktansen -- analogt med trimmern bland 
kondensatorerna.
\end{facitblock}

\begin{facitblock}{6}
\qrakna\quad En 10\,mH spole har 5\,$\Omega$ DC-resistans. 
a) Beräkna $\tau$. b) Hur lång tid tills strömmen är 
praktiskt stabil? c) Vad händer med $\tau$ om $L = 20$\,mH?
\tcblower
a)
\[ \tau = \frac{L}{R} = \frac{0{,}010\,\text{H}}{5\,\Omega} = \mathbf{2\,\text{ms}} \]

b) Praktiskt stabil efter $5\tau = 5 \times 2\,\text{ms} = \mathbf{10\,\text{ms}}$

c) Med $L = 20$\,mH:
\[ \tau = \frac{0{,}020}{5} = 4\,\text{ms} \quad \Rightarrow \quad 5\tau = 20\,\text{ms} \]
Dubbel induktans ger dubbel tidskonstant.
\end{facitblock}

% ── Facit 1.6 ─────────────────────────────────────────────
\subsection*{1.6 -- Serie- och parallellkoppling}

\begin{facitblock}{1}
\qrakna\quad Tre resistorer -- 100\,$\Omega$, 220\,$\Omega$ 
och 470\,$\Omega$ -- kopplas i serie till 12\,V.
a) Beräkna $R_{tot}$.
b) Beräkna strömmen i kretsen.
c) Beräkna spänningsfallet över varje resistor och 
kontrollera att de summerar till 12\,V.
\tcblower
\textbf{a) Total resistans (serie -- adderas):}
\[ R_{tot} = 100 + 220 + 470 = \mathbf{790\,\Omega} \]

\textbf{b) Ström (identisk i hela kretsen):}
\[ I = \frac{U}{R_{tot}} = \frac{12}{790} \approx \mathbf{15{,}2\,\text{mA}} \]

\textbf{c) Spänningsfall över varje resistor:}
\[ U_{100} = 0{,}0152 \times 100 = 1{,}52\,\text{V} \]
\[ U_{220} = 0{,}0152 \times 220 = 3{,}34\,\text{V} \]
\[ U_{470} = 0{,}0152 \times 470 = 7{,}14\,\text{V} \]

Kontroll: $1{,}52 + 3{,}34 + 7{,}14 = 12{,}0\,\text{V}$ ✓

KVL håller -- spänningsfallen summerar exakt till 
källspänningen.
\end{facitblock}


\begin{facitblock}{2}
\qrakna\quad Tre resistorer -- 100\,$\Omega$, 220\,$\Omega$ 
och 470\,$\Omega$ -- kopplas i serie till 12\,V.
a) Beräkna $R_{tot}$.
b) Beräkna strömmen i kretsen.
c) Beräkna spänningsfallet över varje resistor och 
kontrollera att de summerar till 12\,V.
\tcblower
\textbf{a) Total resistans (serie -- adderas):}
\[ R_{tot} = 100 + 220 + 470 = \mathbf{790\,\Omega} \]

\textbf{b) Ström (identisk i hela kretsen):}
\[ I = \frac{U}{R_{tot}} = \frac{12}{790} \approx \mathbf{15{,}2\,\text{mA}} \]

\textbf{c) Spänningsfall över varje resistor:}
\[ U_{100} = 0{,}0152 \times 100 = 1{,}52\,\text{V} \]
\[ U_{220} = 0{,}0152 \times 220 = 3{,}34\,\text{V} \]
\[ U_{470} = 0{,}0152 \times 470 = 7{,}14\,\text{V} \]

Kontroll: $1{,}52 + 3{,}34 + 7{,}14 = 12{,}0\,\text{V}$ ✓

KVL håller -- spänningsfallen summerar exakt till 
källspänningen.
\end{facitblock}

\begin{facitblock}{3}
\qrakna\quad Samma tre resistorer kopplas nu parallellt 
till 12\,V.
a) Beräkna $R_{tot}$ med den allmänna formeln.
b) Beräkna total ström och strömmen i varje gren.
c) Kontrollera: är $R_{tot}$ lägre än den minsta resistorn?
\tcblower
\textbf{a) Total resistans (parallell):}
\[ \frac{1}{R_{tot}} = \frac{1}{100} + \frac{1}{220} + \frac{1}{470} 
   = 0{,}0100 + 0{,}00455 + 0{,}00213 = 0{,}01668 \]
\[ R_{tot} = \frac{1}{0{,}01668} \approx \mathbf{59{,}9\,\Omega} \]

\textbf{b) Ström i varje gren (12\,V överallt):}
\[ I_{100} = \frac{12}{100} = 120\,\text{mA} \]
\[ I_{220} = \frac{12}{220} \approx 54{,}5\,\text{mA} \]
\[ I_{470} = \frac{12}{470} \approx 25{,}5\,\text{mA} \]
\[ I_{tot} = 120 + 54{,}5 + 25{,}5 = \mathbf{200\,\text{mA}} \]

Verifiering: $I_{tot} = 12 / 59{,}9 \approx 200\,\text{mA}$ ✓

\textbf{c)} Ja -- $R_{tot} = 59{,}9\,\Omega$ är lägre än 
den minsta enskilda resistorn (100\,$\Omega$). 
Parallellkoppling ger alltid lägre totalresistans än 
den minsta ingående komponenten. ✓
\end{facitblock}

\begin{facitblock}{4}
\qgrund\quad En parallellkrets har tre grenar och matas 
med 12\,V. Du mäter strömmen i två av grenarna: 
120\,mA och 54\,mA. Totalströmmen från källan är 
225\,mA.
\tcblower
\textbf{a) Strömmen i den tredje grenen (KCL):}

KCL säger att summan av grenströmmarna ska vara lika 
med totalströmmen:
\[ I_3 = I_{tot} - I_1 - I_2 = 225 - 120 - 54 = \mathbf{51\,\text{mA}} \]

\textbf{b) Resistansen i den tredje grenen:}

Spänningen är 12\,V över alla grenar:
\[ R_3 = \frac{U}{I_3} = \frac{12}{0{,}051} \approx \mathbf{235\,\Omega} \]

\textbf{c) Verifiering med $R_{tot}$:}

Grenresistanserna:
\[ R_1 = \frac{12}{0{,}120} = 100\,\Omega \qquad R_2 = \frac{12}{0{,}054} \approx 222\,\Omega \qquad R_3 \approx 235\,\Omega \]

\[ \frac{1}{R_{tot}} = \frac{1}{100} + \frac{1}{222} + \frac{1}{235} = 0{,}01000 + 0{,}00450 + 0{,}00426 = 0{,}01876 \]
\[ R_{tot} \approx 53{,}3\,\Omega \]

\[ I_{tot} = \frac{U}{R_{tot}} = \frac{12}{53{,}3} \approx 225\,\text{mA} \quad \checkmark \]

KCL och Ohms lag ger samma svar -- lagen stämmer.
\end{facitblock}

\begin{facitblock}{5}
\qrakna\quad Du har kondensatorerna 47\,µF, 100\,µF 
och 220\,µF.
a) Vad är $C_{tot}$ om alla tre kopplas parallellt?
b) Om 47\,µF och 100\,µF kopplas i serie -- vad är $C_{tot}$?
c) Varför är kondensatorer i serie ovanligt i 
praktiska kretsar?
\tcblower
\textbf{a) Parallellt -- kapacitansen adderas:}
\[ C_{tot} = 47 + 100 + 220 = \mathbf{367\,\mu\text{F}} \]

\textbf{b) 47\,µF och 100\,µF i serie:}
\[ \frac{1}{C_{tot}} = \frac{1}{47} + \frac{1}{100} 
   = 0{,}02128 + 0{,}01 = 0{,}03128 \]
\[ C_{tot} = \frac{1}{0{,}03128} \approx \mathbf{32\,\mu\text{F}} \]

\textbf{c) Varför ovanligt?}
Seriekoppling \emph{minskar} kapacitansen -- raka 
motsatsen till vad man nästan alltid vill uppnå. 
Det är billigare och enklare att köpa rätt 
kondensator direkt.

Undantaget: serie används för att \emph{höja 
spänningshållfastheten} -- om varje kondensator 
tål 16\,V klarar två i serie 32\,V. Spänningen 
fördelas mellan dem, precis som resistorer i serie 
delar på spänningen.
\end{facitblock}

\begin{facitblock}{6}
\qprov\quad Du bygger ett LC-filter för 40-metersbandet 
(7\,MHz). Resonansfrekvensen är för hög och måste sänkas.
\tcblower
\textbf{a) Vad måste hända med $L$ och/eller $C$?}

Formeln visar att $f_0$ minskar när nämnaren $\sqrt{LC}$ 
ökar -- det vill säga när $L$ eller $C$ (eller båda) 
ökar:
\[ f_0 = \frac{1}{2\pi\sqrt{LC}} \quad \Rightarrow \quad \text{större } L \text{ eller } C \Rightarrow \text{lägre } f_0 \]

\textbf{b) Hur kopplar man extra komponenter?}

\textbf{Spolen:} Spolar följer samma regler som 
resistorer. För att öka $L_{tot}$ kopplas extra 
spolar i \textbf{serie} -- induktanserna adderas:
\[ L_{tot} = L_1 + L_2 \]

\textbf{Kondensatorn:} Kondensatorer följer omvända 
regler mot resistorer. För att öka $C_{tot}$ kopplas 
extra kondensatorer \textbf{parallellt} -- 
kapacitanserna adderas:
\[ C_{tot} = C_1 + C_2 \]

\textbf{c) Teknikerns förslag -- spole parallellt och 
kondensator parallellt?}

Halvrätt. Kondensator parallellt är korrekt -- det 
ökar $C_{tot}$ och sänker $f_0$.

Men spole parallellt är fel. Parallella spolar ger 
\emph{lägre} $L_{tot}$ (samma formel som parallella 
resistorer) -- det höjer $f_0$ i stället för att 
sänka den. Teknikern blandade ihop parallellregeln 
för spolar med parallellregeln för kondensatorer.

\textbf{Rätt svar:} Spole i serie, kondensator 
parallellt.
\end{facitblock}


%% ─────────────────────────────────────────────────────────
%% FACIT 1.7
%% ─────────────────────────────────────────────────────────
\subsection*{1.7 -- Växelström och impedans}

% ── Facit 1 ───────────────────────────────────────────────
\begin{facitblock}{1}
\qrakna\quad Frekvens och period på 144,500~MHz.
\tcblower
\textbf{a) Periodens längd:}
\[ T = \frac{1}{f} = \frac{1}{144{,}5 \times 10^6} \approx \mathbf{6{,}9\,\text{ns}} \]

\textbf{b) Frekvensen för $T = 200\,\text{ns}$:}
\[ f = \frac{1}{T} = \frac{1}{200 \times 10^{-9}} = \mathbf{5\,\text{MHz}} \]

\textbf{c) Vilket band?}

$144{,}5$~MHz är VHF, 2-metersbandet (amatörradio).\\
$5$~MHz är HF (kortvåg) -- nära 60-metersbandet.\\
Frekvensförhållande: $144{,}5 / 5 = 28{,}9$ -- den 
första signalen är nästan 29 gånger högre i frekvens.
\end{facitblock}


% ── Facit 2 ───────────────────────────────────────────────
\begin{facitblock}{2}
\qrakna\quad RMS och toppvärden.
\tcblower
\textbf{a) Varför är RMS lägre än toppvärdet?}

En sinusvåg tillbringar mest tid nära noll och minst
tid nära toppvärdet -- kurvan rusar förbi toppen men
dröjer kvar i mitten. Det genomsnittliga arbetet
är därför lägre än vad toppvärdet antyder:
$0{,}707 \times U_{topp}$, inte $1{,}0 \times U_{topp}$.

\textbf{b) Toppvärde från 230~V RMS:}
\[ U_{topp} = 230 \times 1{,}414 \approx \mathbf{325\,\text{V}} \]

\textbf{c) RMS från 100~V toppvärde:}
\[ U_{RMS} = 100 \times 0{,}707 = \mathbf{70{,}7\,\text{V}} \]

\textbf{d) Varför RMS vid effektberäkningar?}

Effektformeln $P = U \times I$ är definierad för
konstanta (DC) värden. RMS är det AC-värde som ger
\emph{samma effekt} som ett lika stort DC-värde.
Om toppvärde används istället överskattas effekten
med en faktor 2 -- eftersom $(\sqrt{2})^2 = 2$.
\end{facitblock}


\begin{facitblock}{3}
\qrakna\quad $R = 75\,\Omega$, $X_L = 120\,\Omega$, $X_C = 45\,\Omega$.
\tcblower
\textbf{a) Total impedans:}
\[ Z = \sqrt{75^2 + (120 - 45)^2} = \sqrt{5625 + 5625} = \sqrt{11250} \approx \mathbf{106\,\Omega} \]

\textbf{b) Kapacitiv eller induktiv?}

$X_L = 120\,\Omega > X_C = 45\,\Omega$ --
kretsen är \textbf{induktiv}. Spänningen leder strömmen.

\textbf{c) Vad händer om frekvensen sänks?}

$X_L = 2\pi f L$ minskar (induktans minskar med $f$).\\
$X_C = 1/(2\pi f C)$ ökar (kapacitans ökar när $f$ minskar).\\

Differensen $(X_L - X_C)$ minskar -- kretsen rör sig
mot resonans, och $Z$ sjunker mot $R = 75\,\Omega$.
Sänks frekvensen ytterligare passerar kretsen resonans
och blir kapacitiv.

\textbf{d) Fasvinkel:}
\[ \theta = \arctan\!\left(\frac{120 - 45}{75}\right) = \arctan(1{,}0) = \mathbf{45°} \]

Kretsen är induktiv med fasvinkel 45° -- spänningen
leder strömmen med 45°.
\end{facitblock}

% ── Facit 4 ───────────────────────────────────────────────
\begin{facitblock}{4}
\qrakna\quad LC-filter för 3,5~MHz, $L = 5\,\mu\text{H}$.
\tcblower
\textbf{a) Nödvändig kapacitans:}

Löser ut $C$ ur resonansformeln:
\[ f_0 = \frac{1}{2\pi\sqrt{LC}} \quad\Rightarrow\quad C = \frac{1}{(2\pi f_0)^2 L} \]
\[ C = \frac{1}{(2\pi \times 3{,}5 \times 10^6)^2 \times 5 \times 10^{-6}} \]
\[ C = \frac{1}{(2{,}199 \times 10^7)^2 \times 5 \times 10^{-6}} \approx \frac{1}{2{,}415 \times 10^9} \approx \mathbf{414\,\text{pF}} \]

\textbf{b) Resonansen är för hög -- justera $C$:}

$f_0$ är 4~MHz, ska vara 3,5~MHz -- den ska sänkas.\\
Lösning: öka $C$. Lägg till kapacitans \textbf{parallellt} 
(parallella kondensatorer adderas).

\textbf{c) $X_L$ och $X_C$ vid stigande frekvens:}

$X_L = 2\pi f L$ -- \textbf{ökar} med frekvensen.\\
$X_C = 1/(2\pi f C)$ -- \textbf{minskar} med frekvensen.\\

Under resonans: $X_C > X_L$ (kapacitiv karaktär).\\
Vid resonans: $X_C = X_L$.\\
Över resonans: $X_L > X_C$ (induktiv karaktär).
\end{facitblock}


% ── Facit 5 ───────────────────────────────────────────────
\begin{facitblock}{5}
\qprov\quad ELI the ICE man.
\tcblower
\textbf{a) ELI the ICE man:}

\textbf{ELI} -- i en \textbf{L}induktans leder 
\textbf{E}lspänningen \textbf{I}nduktionsströmmen med 90°.\\
\textbf{ICE} -- i en \textbf{C}ondensator leder 
\textbf{I}nduktionsströmmen \textbf{E}lspänningen med 90°.

Fysisk förklaring: en spole \emph{motverkar} 
strömförändringar (inducerar en motspänning). 
Spänningen måste byggas upp \emph{innan} strömmen 
kan flöda -- spänning före ström.

En kondensator \emph{motverkar} spänningsförändringar 
(strömmen flödar in för att ladda upp plattorna 
\emph{innan} spänningen hinner byggas upp) -- ström 
före spänning.

\textbf{b) Vid resonans ($X_L = X_C$):}

Fasförskjutningarna tar ut varandra -- ström och 
spänning är i fas (0°). Kretsen beter sig som en 
ren resistans.

\textbf{c) Varför viktigt vid effektberäkningar?}

Skenbar effekt $S = U \times I$ -- men verklig effekt 
$P = U \times I \times \cos\phi$, där $\phi$ är 
fasförskjutningen. Vid 90° fasförskjutning är $\cos 90° = 0$ 
och ingen \emph{verklig} effekt levereras trots att 
ström och spänning är stora. Fasförskjutning gör att 
ström och spänning inte ``samarbetar''.
\end{facitblock}


% ── Facit 6 ───────────────────────────────────────────────
\begin{facitblock}{6}
\qgrund\quad Standard-impedanser.
\tcblower
\textbf{a) Varför 50~$\Omega$ är en kompromiss:}

I en luftfylld koaxkabel uppnås lägst dämpning vid 
$\approx 77\,\Omega$ och högst effektkapacitet vid 
$\approx 30\,\Omega$. Kompromissen är det geometriska 
medelvärdet: $\sqrt{77 \times 30} \approx 48\,\Omega$, 
avrundat till 50~$\Omega$.

\textbf{b) Kan 75~$\Omega$-kabel användas?}

Tekniskt går det, men det ger impedansobalans mot 
50~$\Omega$-systemet. Reflektioner uppstår vid varje 
anslutning (sändare, antenn), SWR ökar och en del 
effekt förloras. Korta kablar ger liten skillnad; 
långa kablar ger märkbar prestandaförsämring. 
Rekommendationen är att alltid använda 50~$\Omega$-kabel 
i amatörradiosystem.

\textbf{c) Impedansförhållande för foldad dipol:}
\[ \frac{Z_{antenn}}{Z_{koax}} = \frac{300}{50} = \mathbf{6:1} \]

En 1:6-balun behövs, alternativt en 1:4-balun 
kombinerad med impedanstransformation i antennsystemet.
\end{facitblock}


% ── Facit 7 ───────────────────────────────────────────────
\begin{facitblock}{7}
\qgrund\quad Skineffekt.
\tcblower
\textbf{a) Kopparrör vs massiv koppar vid UHF:}

Vid 430~MHz är skindjupet i koppar bara $\approx 3\,\mu\text{m}$. 
Strömmen flödar \emph{enbart} i ett tunt ytskikt -- 
mitten av ledaren är oanvänd. Eftersom röret och 
den massiva staven har identisk yta (ytteryta) är 
deras RF-resistans praktiskt taget lika. Kopparröret 
är lättare och billigare men elektiskt likvärdigt.

\textbf{b) Silverpläterade kontakter vid VHF/UHF:}

Silver har bättre ledningsförmåga än koppar. 
Vid VHF/UHF flödar strömmen bara i ett $\mu$m-tunt 
ytskikt -- silvrets överlägsna ledningsförmåga 
utnyttjas fullt ut just där strömmen är. Vid DC 
däremot flödar strömmen genom hela ledarens 
tvärsnittsarea; silverplätans tjocklek är helt 
negligerbar jämfört med ledarens yta.

\textbf{c) RG-58 på 430~MHz:}

RG-58 är specificerad för HF och lägre VHF. Vid 
430~MHz är förlusterna per meter väsentligt högre 
-- en 30~m kabel kan ge 10+ dB förlust, vilket 
innebär att mer än 90~\% av sändareffekten försvinner 
i kabeln. En kabel med lägre förlust (t.ex. 
Aircell~7 eller LMR-400) behövs för UHF.
\end{facitblock}


% ── Facit 8 ───────────────────────────────────────────────
\begin{facitblock}{8}
\qrakna\quad Sändare 100~W, SWR 3:1.
\tcblower
\textbf{a) Reflekterad andel:}

Ur tabellen: SWR 3:1 $\rightarrow$ \textbf{25~\%} reflekteras.

\textbf{b) Watt till antennen:}
\[ P_{ref} = 0{,}25 \times 100 = 25\,\text{W} \]
\[ P_{ant} = 100 - 25 = \mathbf{75\,\text{W}} \]

\textbf{c) SWR hoppar till 5:1 -- vad gör du?}

Sänk effekten eller stäng av omedelbart. SWR 5:1 innebär att över 40~% av effekten studsar tillbaka mot slutsteget -- samma princip som bergsdalen fast värre. De flesta moderna sändare har inbyggt SWR-skydd som reducerar effekten automatiskt, men lita inte blint på det.
Felsök sedan lugnt: lös koaxkontakt eller ett brott i kabeln är vanligaste orsaken till ett plötsligt SWR-hopp.
\end{facitblock}

\subsection*{1.8 – Filter och resonanskretsar}

\begin{facitblock}{1}
\qgrund\quad Namnge de fyra filtertyperna med ett radioexempel för varje.
\tcblower
\textbf{Lågpass (LP):} Passerar låga frekvenser, blockerar höga.\\
\textit{Exempel: Övertonsfilter efter slutsteg -- dämpar 2f, 3f osv.}

\textbf{Högpass (HP):} Passerar höga frekvenser, blockerar låga.\\
\textit{Exempel: Blockerar LW/AM-störning i mottagarens ingång.}

\textbf{Bandpass (BP):} Passerar ett frekvensband, blockerar utanför.\\
\textit{Exempel: Väljer ut 20m-bandet (14~MHz) ur hela spektrum.}

\textbf{Bandstopp (BS):} Blockerar ett smalt band, passerar allt annat.\\
\textit{Exempel: Eliminerar en specifik störsignal utan att påverka resten.}
\end{facitblock}

\begin{facitblock}{2}
\qrakna\quad RC-lågpassfilter, $R=1\;\text{k}\Omega$, $C=1\;\text{nF}$.
\tcblower
\textbf{a) Gränsfrekvens:}
\[
f_c = \frac{1}{2\pi RC}
    = \frac{1}{2\pi \times 1000 \times 1\times10^{-9}}
    \approx \mathbf{159\;\text{kHz}}
\]

\textbf{b) Dämpning vid $2f_c$:}

\textit{Tumregel (asymptotisk):} $-3\;\text{dB}$ vid $f_c$, plus
$-6\;\text{dB}$ per oktav ger $-3 - 6 = \mathbf{-9\;\text{dB}}$.

\textit{Exakt värde:} För första ordningens RC-lågpass gäller
$|H| = 1/\sqrt{1+(f/f_c)^2}$. Vid $f = 2f_c$:
\[
|H| = \frac{1}{\sqrt{1+4}} = \frac{1}{\sqrt{5}} \approx 0{,}447
\quad\Rightarrow\quad
20\log(0{,}447) \approx \mathbf{-7{,}0\;\text{dB}}
\]

Skillnaden ($-9$ mot $-7$~dB) beror på att 6~dB/oktav är en
asymptotisk approximation -- den är exakt bara långt ovanför $f_c$.
Nära $f_c$ är lutningen ännu inte helt etablerad.
\end{facitblock}

\begin{facitblock}{3}
\qrakna\quad Beräkna $f_0$ för tre LC-kretsar.
\tcblower
\textbf{a)} $L=10\;\mu\text{H}$, $C=100\;\text{pF}$:
\[
f_0 = \frac{1}{2\pi\sqrt{10\times10^{-6}\times100\times10^{-12}}}
    \approx \mathbf{5{,}03\;\text{MHz}}
\]

\textbf{b)} $L=1\;\mu\text{H}$, $C=50\;\text{pF}$:
\[
f_0 = \frac{1}{2\pi\sqrt{1\times10^{-6}\times50\times10^{-12}}}
    \approx \mathbf{22{,}5\;\text{MHz}}
\]

\textbf{c)} $L=500\;\text{nH}$, $C=20\;\text{pF}$:
\[
f_0 = \frac{1}{2\pi\sqrt{500\times10^{-9}\times20\times10^{-12}}}
    \approx \mathbf{50{,}3\;\text{MHz}}
\]
\end{facitblock}

\begin{facitblock}{4}
\qrakna\quad $f_0 = 14{,}2\;\text{MHz}$, $BW = 20\;\text{kHz}$.
\tcblower
\[
Q = \frac{f_0}{BW}
  = \frac{14{,}2\times10^6}{20\times10^3}
  = \mathbf{710}
\]

Q = 710 är mycket högt -- kretsen är extremt selektiv och väljer ut
ett smalt fönster runt 14,2~MHz. Utmärkt för en mottagares
ingångskrets på 20m-bandet.
\end{facitblock}

\begin{facitblock}{5}
\qprov\quad Pi-filtrets funktioner och fel band.
\tcblower
\textbf{Pi-filtrets tre funktioner:}

\textbf{1. Dämpar övertoner} -- slutsteget genererar harmoniska
övertoner (2$f$, 3$f$\ldots) som stör andra stationer och bryter
mot sändningsregler. LP-karaktären hos C--L--C dämpar allt ovanför
grundfrekvensen kraftigt.

\textbf{2. Impedansmatchning} -- transformerar slutstegets
utgångsimpedans (ofta hundratals $\Omega$) ned till koaxkabelns
50~$\Omega$. Utan matchning stiger SWR och effekten når inte antennen.

\textbf{3. Blockerar DC} -- kondensatorerna hindrar likspänning från
slutsteget att nå antennen.

\medskip
\textbf{80m-filter på 10m-bandet:}

Filtret är dimensionerat för resonans runt 3,5~MHz. På 28~MHz ligger
man långt ovanför resonansen -- filtret dämpar nu \emph{i} bandet
istället för ovanför det, och impedansmatchningen är helt fel.
Resultatet är kraftigt förhöjt SWR och utebliven övertonsdämpning.
\end{facitblock}

\begin{facitblock}{6}
\qrakna\quad Pi-filter -- resonansfrekvens och bandval.
\tcblower

\textbf{a) Beräkna $f_0$:}

\[
f_0 = \frac{1}{2\pi\sqrt{LC}}
    = \frac{1}{2\pi\sqrt{10\times10^{-6} \times 100\times10^{-12}}}
    = \frac{1}{2\pi\sqrt{10^{-15}}}
    = \frac{1}{2\pi \times 3{,}162\times10^{-8}}
    \approx \mathbf{5{,}03\;\text{MHz}}
\]

\textbf{b) Vilket amatörband passar?}

5,03~MHz ligger mellan 80m-bandet (3,5--3,8~MHz) och 40m-bandet
(7,0--7,2~MHz). Närmast passar \textbf{40m-bandet} -- filtret
skulle behöva justeras något (minska $L$ eller $C$) för att
resonera exakt vid 7~MHz. Som tumregel: $f_0$ bör ligga strax
ovanför bandets överkant så att dämpningen är tillräcklig redan
vid den lägsta övertonen.

\textbf{c) Högt Q vs lågt Q vid $2f_0 \approx 10\;\text{MHz}$:}

\textbf{Högt Q} ger en smalare och brantare kurva -- dämpningen
ökar snabbt utanför $f_0$. En överton på $2f_0$ hamnar långt
utanför passband och dämpas kraftigt. Bra för övertonsfiltrering.

\textbf{Lågt Q} ger en bredare och flackare kurva -- dämpningen
ökar långsamt. Samma överton på $2f_0$ dämpas betydligt mindre.
Sämre för övertonsfiltrering, men mer okänsligt för
komponenttoleranser.

\medskip
\textit{Slutsats:} Högt Q är önskvärt i ett Pi-filter efter ett
slutsteg -- det ger skarp dämpning av övertoner samtidigt som
grundfrekvensen passerar med låga förluster.
\end{facitblock}



% ╔══════════════════════════════════════════════════════════════╗
% ║  KAPITEL 2 – RADIOTEKNIK                                     ║
% ║  Certifikatshandboken · Radioskola.se                        ║
% ╚══════════════════════════════════════════════════════════════╝

\chapter{Radioteknik}

\lettrine[lines=2]{\textcolor{orange}{R}}{adiovågor} och ljus är egentligen samma sak -- elektromagnetisk strålning. Skillnaden är bara hur snabbt de svänger. Det du lär dig i det här kapitlet är kärnan i all radioteknik -- från en enkel walkie-talkie till ett rymdsondssystem 20 miljarder kilometer bort.

Kapitel 1 gav dig verktygen -- resistorer, kondensatorer, spolar och resonans. Nu ska vi använda dem. Det finns egentligen bara fyra frågor att besvara: vad \emph{är} en radiovåg, hur packar vi information i den, hur skapar vi den, och hur plockar vi ut den ur luften igen?

% ═══════════════════════════════════════════════════════════════
\section{Vad är en radiovåg?}
% ═══════════════════════════════════════════════════════════════

\subsection{Elektromagnetiska vågor -- ljus och radio är släktingar}

Ljus, värme, röntgenstrålning och radiovågor är egentligen samma
sak -- elektromagnetisk strålning. Skillnaden är bara hur snabbt
de \emph{svänger}.

\begin{storybox}[{\emoji{🌊} Dammet och stenen}]
Föreställ dig en stilla sjö. Du kastar i en sten och ser ringar
breda ut sig från nedslagsplatsen. Ringarna är inte vatten som
rör sig utåt -- vattnet gungar bara upp och ned på stället.
Det som rör sig utåt är en \emph{störning}, ett mönster av rörelse.

En elektromagnetisk våg fungerar på samma sätt, men istället för
vatten är det elektriska och magnetiska fält som gungar upp och ned.
Och istället för att röra sig med några meter per sekund rör sig
dessa störningar med ljusets hastighet -- 300\,000 kilometer
\emph{per sekund}.

Det finns ingen skillnad i hur radiovågor och ljus ``fungerar''.
En radioamatörs sändare och en ficklampa gör exakt samma sak --
de skapar elektromagnetiska störningar som sprider sig utåt.
Skillnaden är bara hur fort de svänger.
\end{storybox}

En elektromagnetisk våg består av två fält som alltid följs åt:
ett elektriskt fält och ett magnetiskt fält, vinkelräta mot varandra
och mot rörelseriktningen. De driver varandra -- det elektriska
fältet skapar det magnetiska, som skapar ett nytt elektriskt, och
så vidare ut genom rymden.

\begin{bluebox}[title={\emoji{🔑} Elektromagnetiskt spektrum}]
Alla elektromagnetiska vågor rör sig med samma hastighet i vakuum:
\[ c = 3 \times 10^8 \text{ m/s} \approx 300\,000 \text{ km/s} \]

Det som skiljer dem åt är \textbf{frekvensen} -- hur många gånger
per sekund fälten svänger:

\begin{center}
\begin{tabular}{L{3.5cm}L{3.5cm}L{4cm}}
\rowcolor{greydark}
\tblhdr{Typ} & \tblhdr{Frekvens} & \tblhdr{Våglängd} \\
Radiovågor & 3\,Hz -- 300\,GHz & 1\,mm -- 100\,000\,km \\
\rowcolor{greylight}
Mikrovågor & 300\,MHz -- 300\,GHz & 1\,mm -- 1\,m \\
Infrarött & 300\,GHz -- 430\,THz & 0,7 -- 1000\,µm \\
\rowcolor{greylight}
Synligt ljus & 430 -- 770\,THz & 390 -- 700\,nm \\
Röntgen & $>10^{16}$\,Hz & $<30$\,nm \\
\end{tabular}
\end{center}
\smallskip
{\small\itshape Radioamatörer arbetar med frekvenser från
ca 1,8\,MHz (160\,m-bandet) upp till 250\,GHz.}
\end{bluebox}

% ── Figur: elektromagnetiskt spektrum ───────────────────────
\begin{figure}[htbp]
\centering
\begin{tikzpicture}[font=\small]
  \fill[gray!6, rounded corners=4pt] (-0.4,-1.2) rectangle (13.4,3.2);

  % Spektrumbar
  \foreach \x/\col/\lbl in {
    0/violet!40/Röntgen,
    1.5/blue!40/UV,
    2.8/cyan!50/Ljus,
    4.0/green!40/IR,
    5.5/yellow!50/Mikro,
    7.5/orange!40/UHF,
    9.0/red!30/VHF,
    10.5/red!20/HF,
    12.0/red!10/LF}{
    \fill[\col] (\x,1.2) rectangle (\x+1.4,2.2);
    \node[font=\scriptsize, rotate=45, anchor=west] at (\x+0.7,2.25) {\lbl};
  }

  % Pil
  \draw[->, thick, gray!70] (0,1.7) -- (13.2,1.7);

  % Frekvensaxel
  \foreach \x/\f in {0/GHz,2.8/THz,5.5/300 GHz,7.5/3 GHz,9.0/300 MHz,10.5/30 MHz,12.0/3 MHz}{
    \draw[gray!50] (\x+0.7,1.2) -- (\x+0.7,1.0);
    \node[gray!60, font=\scriptsize] at (\x+0.7,0.75) {\f};
  }

  % Amatörradio markering
  \draw[red!70!black, very thick, ->] (9.8,0.2) -- (9.8,1.1);
  \node[red!70!black, font=\scriptsize\bfseries] at (9.8,-0.1) {Amatörradio};
  \node[red!70!black, font=\scriptsize] at (9.8,-0.5) {HF/VHF/UHF};

\end{tikzpicture}
\caption{Det elektromagnetiska spektrumet. Radiovågor och synligt
ljus är samma typ av strålning -- bara frekvensen skiljer.
Amatörradio spänner över ett enormt område, från LF upp till
mikrovågsområdet.}
\label{fig:em_spektrum}
\end{figure}

% ────────────────────────────────────────────────────────────
\subsection{Frekvens och våglängd -- sambandet $\lambda = c/f$}

Varje radiovåg har en frekvens -- hur många gånger per sekund
fältet svänger -- och en våglängd -- hur långt en full svängning
sträcker sig i rummet. Dessa två är alltid kopplade till varandra
via ljushastigheten.

\begin{storybox}[{\emoji{🚂} Tåg som passerar}]
Föreställ dig att du står vid ett järnvägsspår och räknar vagnar.
Tåget rör sig med konstant hastighet. Om vagnarna är långa
passerar färre vagnar per minut -- frekvensen är låg, våglängden
är lång. Korta vagnar ger fler per minut -- hög frekvens, kort
våglängd.

Radiovågor fungerar likadant. Ljushastigheten är konstant (tågets
hastighet), frekvensen är antalet svängningar per sekund (vagnar
per minut), och våglängden är längden på varje svängning (vagnlängden).
\end{storybox}

\begin{bluebox}[title={\emoji{🔑} Sambandet frekvens--våglängd}]
\[ \lambda = \frac{c}{f} \]
där:
\begin{itemize}
  \item $\lambda$ = våglängd (meter)
  \item $c$ = ljushastigheten $\approx 3 \times 10^8$ m/s
  \item $f$ = frekvens (Hz)
\end{itemize}

\textbf{Praktisk tumregel:} $\lambda \approx \dfrac{300}{f_{\text{MHz}}}$ meter

\medskip
\begin{center}
\begin{tabular}{L{2.5cm}L{2.5cm}L{3cm}L{3cm}}
\rowcolor{greydark}
\tblhdr{Band} & \tblhdr{Frekvens} & \tblhdr{Våglängd} & \tblhdr{Antennlängd*} \\
80\,m & 3,5\,MHz & 85,7\,m & 42,9\,m \\
\rowcolor{greylight}
40\,m & 7\,MHz & 42,9\,m & 21,4\,m \\
20\,m & 14\,MHz & 21,4\,m & 10,7\,m \\
\rowcolor{greylight}
2\,m & 144\,MHz & 2,08\,m & 1,04\,m \\
70\,cm & 432\,MHz & 0,69\,m & 34,5\,cm \\
\end{tabular}
\end{center}
{\small\itshape *Halvvågsdipol -- mer om detta i kapitel 3 (Antenner).}
\end{bluebox}

\begin{worked}[title={Räkneexempel 2.1.2 -- Våglängd på 14\,MHz}]
En radioamatör vill sända på 14,2\,MHz (20\,m-bandet).
Vad är våglängden?

\[ \lambda = \frac{c}{f} = \frac{3 \times 10^8}{14{,}2 \times 10^6}
   = \frac{300}{14{,}2} \approx 21{,}1 \text{ m} \]

Bandets namn ``20\,m'' är alltså en ungefärlig avrundning av
den faktiska våglängden vid bandets mitt.
\end{worked}

% ────────────────────────────────────────────────────────────
\subsection{Frekvensspektrum och amatörband}

Radioamatörer får inte sända på vilka frekvenser som helst --
det internationella radioreglementet (ITU) delar upp spektrumet
mellan olika tjänster. Amatörradio har tilldelats ett antal band
spridda över hela spektrumet, från lågfrekvensområdet (LF) upp
till mikrovågor.

\begin{bluebox}[title={\emoji{🔑} Viktigaste amatörbanden}]
\begin{center}
\begin{tabular}{L{1.8cm}L{3.0cm}L{2.5cm}L{4.0cm}}
\rowcolor{greydark}
\tblhdr{Band} & \tblhdr{Frekvens} & \tblhdr{Område} & \tblhdr{Typisk användning} \\
160\,m & 1,810--2,000\,MHz & MF & Nattrafik, DX \\
\rowcolor{greylight}
80\,m & 3,500--3,800\,MHz & HF & Lokal/regional trafik \\
40\,m & 7,000--7,200\,MHz & HF & DX, dagstrafik \\
\rowcolor{greylight}
20\,m & 14,000--14,350\,MHz & HF & Världsomspännande DX \\
15\,m & 21,000--21,450\,MHz & HF & DX vid högt solflöde \\
\rowcolor{greylight}
10\,m & 28,000--29,700\,MHz & HF & DX, lokal FM \\
2\,m & 144--146\,MHz & VHF & Lokal trafik, repeater \\
\rowcolor{greylight}
70\,cm & 430--440\,MHz & UHF & Lokal trafik, satellit \\
\end{tabular}
\end{center}
{\small\itshape DX = långdistansförbindelser. Mer om vågutbredning i kapitel 4.}
\end{bluebox}

\subsubsection*{Övningsfrågor -- 2.1}

\begin{qblock}{1.}
\qprov\quad En radioamatör sänder på 7,05\,MHz.

a) Beräkna våglängden.\\
b) Vilket band tillhör frekvensen?\\
c) Hur lång är en halvvågsdipol för detta band (ungefär)?
\end{qblock}

\begin{qblock}{F.}
\qfallan\quad Kalle säger: ``Radiovågor rör sig snabbare än ljus
eftersom de har längre våglängd och lägre frekvens.''

Vad är fel med detta påstående?
\end{qblock}

\medskip
\noindent\rule{\textwidth}{0.4pt}
{\small\itshape Försök själv innan du läser svaret.}
\medskip

\begin{worked}[title={Lösning -- Fallgropsfråga 2.1\,F}]
Påståendet är helt fel. Alla elektromagnetiska vågor -- oavsett
frekvens eller våglängd -- rör sig med exakt samma hastighet i
vakuum: $c \approx 3 \times 10^8$ m/s.

Längre våglängd och lägre frekvens är bara två sidor av samma mynt
($\lambda = c/f$). Hastigheten förändras inte -- bara hur tätt
svängningarna sitter.

\textbf{Obs:} I ett medium (t.ex. glas eller koaxialkabel) kan
hastigheten vara lägre, men det beror på mediet, inte på frekvensen.
\end{worked}

% ═══════════════════════════════════════════════════════════════
\section{Modulering -- att packa information i en våg}
% ═══════════════════════════════════════════════════════════════

En ren radiovåg på t.ex. 14,200\,MHz bär ingen information alls --
den svänger bara fram och tillbaka med konstant amplitud och frekvens.
För att skicka en röst, ett morse-tecken eller data måste vi
\emph{modulera} vågen -- förändra något hos den i takt med
informationen vi vill sända.

\begin{storybox}[{\emoji{🎸} Gitarristen och förstärkaren}]
Tänk dig en gitarrist som spelar genom en förstärkare. Gitarren
skapar svaga lufttryckvariationer -- musiken -- men inte tillräckligt
starka för att höras i en stor lokal. Förstärkaren tar ett starkt
elektriskt signal (elnätet) och låter gitarrmusiken \emph{styra}
hur stark den är. Resultatet är exakt samma musik, fast mycket starkare.

Modulering fungerar på samma sätt. Radiovågen (bärfrekvensen) är
``förstärkaren'' -- stark nog att resa tusentals kilometer.
Rösten, morsetecknet eller datan är ``gitarren'' -- informationen
som ska styra bärvågen. Modulering är konsten att låta
informationen styra bärvågen.
\end{storybox}

% ────────────────────────────────────────────────────────────
\subsection{Bärfrekvensen -- stommen som allt hänger på}

Bärfrekvensen (eng. \emph{carrier frequency}) är den grundläggande
radiovåg som sändaren producerar utan modulering. Den bestämmer
vilket band och vilken frekvens du sänder på -- det är den som
mottagaren ställer in sig på.

En omodulerad bärfrekvens är en perfekt sinusvåg med konstant
amplitud, frekvens och fas. Det är dessa tre egenskaper vi
kan förändra för att modulera -- och de ger oss de tre grundläggande
moduleringstyperna: AM, FM och fasskiftsmodulering.

% ────────────────────────────────────────────────────────────
\subsection{Amplitudmodulering (AM)}

Vid amplitudmodulering (AM) håller vi frekvensen konstant men
låter signalamplituden variera i takt med det vi vill sända.
När rösten är stark är bärvågen stark; när det är tyst återgår
den till grundnivån.

\begin{storybox}[{\emoji{💡} Dimmer och glödlampa}]
En vanlig glödlampa lyser med konstant ljus -- det är bärvågen.
Nu kopplar du in en dimmer och vrider den i takt med musik.
Lampan blinkar i musikens rytm. Från långt håll kan någon se
lampan och ``höra'' musiken genom att titta på hur ljuset varierar.

Det är AM. Bärfrekvensen (lampan) är alltid på, men amplituden
(ljusstyrkan) varierar med informationen (musiken).
\end{storybox}

\begin{bluebox}[title={\emoji{🔑} Amplitudmodulering (AM)}]
Vid AM varierar \textbf{amplituden} hos bärvågen i takt med
modulationssignalen. Frekvensen är konstant.

\medskip
Matematiskt: $u(t) = [U_c + m \cdot u_m(t)] \cdot \cos(2\pi f_c t)$

där $m$ är modulationsindex ($0 \leq m \leq 1$, dvs 0--100\,\%).

\medskip
\textbf{Sidband:} AM skapar ett övre och ett undre sidband runt
bärfrekvensen. En talsignal upp till 3\,kHz ger ett AM-signal
med bandbredd $2 \times 3 = 6$\,kHz.

\medskip
\textbf{Nackdelar med AM:}
\begin{itemize}
  \item Bärvågen tar 2/3 av effekten -- men bär ingen information
  \item Känslig för amplitudstörningar (brus, atmosfäriska störningar)
  \item Dålig spektrumeffektivitet
\end{itemize}

\textbf{Används idag:} Mellanvågssändningar (AM-radio), luftfart (VHF AM).
\end{bluebox}

% ────────────────────────────────────────────────────────────
\subsection{Frekvensmodulering (FM)}

Vid frekvensmodulering (FM) håller vi amplituden konstant men
låter \emph{frekvensen} avvika från bärfrekvensen i takt med
signalen. Stark signal ger stor frekvensavvikelse; svag signal
ger liten avvikelse; tystnad ger ingen avvikelse.

\begin{storybox}[{\emoji{🎵} Cykelns växlar}]
Tänk dig en cyklist som håller konstant pedaltramp (amplituden
är konstant) men byter växlar beroende på backen. I uppförsbacke
trampar hon långsammare (lägre frekvens), i nedförsbacke snabbare
(högre frekvens). Informationen -- backen -- kodas i hur fort
hon trampar, inte hur hårt.

FM fungerar likadant. Bärvågen ``trampar'' alltid lika hårt
(konstant amplitud), men hur fort den svänger (frekvensen) varierar
med signalen.
\end{storybox}

\begin{bluebox}[title={\emoji{🔑} Frekvensmodulering (FM)}]
Vid FM varierar \textbf{frekvensen} hos bärvågen i takt med
modulationssignalen. Amplituden är konstant.

\medskip
\textbf{Frekvensdeviation} ($\Delta f$): Hur mycket frekvensen
maximalt avviker från bärfrekvensen.

\medskip
\textbf{Modulationsindex:} $\beta = \dfrac{\Delta f}{f_m}$
där $f_m$ är modulationssignalens frekvens.

\medskip
\textbf{Bandbredd (Carsons regel):} $B \approx 2(\Delta f + f_m)$

\medskip
\textbf{Fördelar med FM:}
\begin{itemize}
  \item Mycket bättre brusegenskaper än AM
  \item Konstant sändareffekt -- enklare slutsteg
  \item Bra ljudkvalitet
\end{itemize}

\textbf{Nackdelar:} Bredare bandbredd än AM och SSB.

\textbf{Används:} VHF/UHF amatörradio, FM-radio (88--108\,MHz),
repeaters.
\end{bluebox}

% ── Figur: AM vs FM ─────────────────────────────────────────
\begin{figure}[htbp]
\centering
\begin{tikzpicture}[font=\small, >=stealth]
  \fill[gray!6, rounded corners=4pt] (-0.4,-0.4) rectangle (13.4,9.8);

  % ── Modulationssignal (gemensam) ──────────────────────────
  \begin{scope}[yshift=8.2cm]
    \node[font=\small\bfseries, anchor=west] at (0,0.3) {Modulationssignal (röst/data)};
    \draw[->, gray!50] (1.0,-0.6) -- (12.8,-0.6) node[right,font=\scriptsize,gray]{$t$};
    \draw[->, gray!50] (1.0,-0.6) -- (1.0,0.7) node[above,font=\scriptsize,gray]{$u$};
    \draw[blue!70!black, thick, domain=1.1:12.6, samples=120]
      plot (\x, {0.45*sin((\x-1.1)*50)});
  \end{scope}

  % ── AM-signal ─────────────────────────────────────────────
  \begin{scope}[yshift=4.6cm]
    \node[font=\small\bfseries, anchor=west] at (0,0.9) {AM -- amplituden varierar};
    \draw[->, gray!50] (1.0,-0.9) -- (12.8,-0.9) node[right,font=\scriptsize,gray]{$t$};
    \draw[->, gray!50] (1.0,-0.9) -- (1.0,1.1) node[above,font=\scriptsize,gray]{$u$};
    \draw[red!70!black, thick, domain=1.1:12.6, samples=300]
      plot (\x, {(0.6 + 0.45*sin((\x-1.1)*50)) * sin((\x-1.1)*350)});
    % Envelope
    \draw[red!40, dashed, domain=1.1:12.6, samples=120]
      plot (\x, {0.6 + 0.45*sin((\x-1.1)*50)});
    \draw[red!40, dashed, domain=1.1:12.6, samples=120]
      plot (\x, {-(0.6 + 0.45*sin((\x-1.1)*50))});
  \end{scope}

  % ── FM-signal ─────────────────────────────────────────────
  \begin{scope}[yshift=1.0cm]
    \node[font=\small\bfseries, anchor=west] at (0,0.9) {FM -- frekvensen varierar};
    \draw[->, gray!50] (1.0,-0.9) -- (12.8,-0.9) node[right,font=\scriptsize,gray]{$t$};
    \draw[->, gray!50] (1.0,-0.9) -- (1.0,1.1) node[above,font=\scriptsize,gray]{$u$};
    \draw[green!50!black, thick, domain=1.1:12.6, samples=400]
      plot (\x, {0.7*sin((\x-1.1)*350 + 80*sin((\x-1.1)*50))});
  \end{scope}

\end{tikzpicture}
\caption{Jämförelse mellan AM och FM. Vid AM varierar amplituden
(de röda streckade linjerna visar \emph{envelopen}) medan
frekvensen är konstant. Vid FM är amplituden konstant men
frekvensen varierar -- täta svängningar = positiv deviation,
glesa svängningar = negativ deviation.}
\label{fig:am_fm}
\end{figure}

% ────────────────────────────────────────────────────────────
\subsection{Enkelsidebandssändning (SSB)}

SSB (Single Sideband) är det sändningsslag som dominerar på
kortvåg (HF) bland radioamatörer världen över. För att förstå
varför, behöver vi förstå vad som är ``onödigt'' med AM.

Vid AM skapas tre komponenter: bärvågen + övre sidobandet + undre
sidobandet. Bärvågen bär ingen information och tar halva effekten.
De båda sidobanden bär \emph{exakt samma} information -- det räcker
alltså med ett av dem.

SSB tar detta till sin logiska slutpunkt: vi filtrerar bort bärvågen
och ett av sidobanden. Kvar är bara ett sidband med all information.

\begin{bluebox}[title={\emoji{🔑} SSB -- Enkelsidebandssändning}]
\textbf{USB} (Upper Sideband) = övre sidobandet används.\\
\textbf{LSB} (Lower Sideband) = undre sidobandet används.

\medskip
\textbf{Konvention inom amatörradio:}
\begin{itemize}
  \item LSB används under 10\,MHz (40\,m, 80\,m, 160\,m)
  \item USB används över 10\,MHz (20\,m, 15\,m, 10\,m, VHF/UHF)
\end{itemize}

\medskip
\textbf{Fördelar med SSB jämfört med AM:}
\begin{itemize}
  \item Halva bandbredden (ca 2,4\,kHz istället för 6\,kHz)
  \item All sändareffekt går till informationen -- ingen bärvåg
  \item Bättre S/N-förhållande vid mottagning
  \item Möjliggör fler stationer på samma band
\end{itemize}

\textbf{Nackdelar:} Mottagaren måste ha BFO (beat frequency
oscillator) för att återskapa bärvågen och göra talet förståeligt.
Utan BFO låter SSB-tal som ``ankor''.
\end{bluebox}

% ────────────────────────────────────────────────────────────
\subsection{Digitala sändningsslag -- CW, RTTY, FT8}

Utöver analoga moduleringsformer använder radioamatörer en rad
digitala sändningsslag där information kodas som diskreta symboler
snarare än kontinuerliga signaler.

\begin{bluebox}[title={\emoji{🔑} Digitala sändningsslag}]
\textbf{CW (Continuous Wave)} -- Morsetelegrafi.
Bärvågen kopplas på och av i morse-mönster. Tekniskt sett
är det ON/OFF-modulering (OOK). CW är exceptionellt
trångbandigt (ca 100\,Hz) och kan höras långt under brusnivå
av en tränad operatör.

\medskip
\textbf{RTTY (Radio Teletype)} -- Frekvensskiftsmodulering (FSK)
mellan två frekvenser (``mark'' och ``space''). Klassisk
digitalteknik, fortfarande i bruk.

\medskip
\textbf{FT8} -- Modernt digitalt sändningsslag (2017) som
klarar extremt svaga signaler. Används globalt för DX-trafik
med minimal effekt. Bandbredd ca 50\,Hz per QSO.

\medskip
\textbf{PSK31} -- Fasskiftsmodulering med 31\,Hz bandbredd.
Populärt för textkonversationer på HF.
\end{bluebox}

\subsubsection*{Övningsfrågor -- 2.2}

\begin{qblock}{2.}
\qprov\quad Förklara skillnaden mellan AM och FM med egna ord.

a) Vad varierar vid AM? Vad är konstant?\\
b) Vad varierar vid FM? Vad är konstant?\\
c) Varför är FM mer brusresistent än AM?
\end{qblock}

\begin{qblock}{3.}
\qprov\quad En SSB-sändare sänder med 100\,W PEP på 14,225\,MHz.

a) Vilket sidband används -- USB eller LSB? Varför?\\
b) Ungefär hur bred är signalen i Hz?\\
c) Hur mycket effekt ``slösar'' SSB jämfört med AM vid samma
   sändareffekt?
\end{qblock}

\begin{qblock}{F.}
\qfallan\quad ``FM är alltid bättre än AM -- den är mer brusresistent
och bör därför användas på alla band, inklusive kortvåg.''

Varför är detta inte sant i praktiken?
\end{qblock}

\medskip
\noindent\rule{\textwidth}{0.4pt}
{\small\itshape Försök själv innan du läser svaret.}
\medskip

\begin{worked}[title={Lösning -- Fallgropsfråga 2.2\,F}]
FM:s brusresistens är verklig, men den har ett pris: \textbf{bandbredd}.
En FM-signal kräver typiskt 10--15\,kHz bandbredd (vid normal
frekvensdeviation), jämfört med ca 2,4\,kHz för SSB.

På kortvåg (HF) är spektrumet trångt och delas av tusentals
stationer över hela världen. Att använda FM där skulle innebära
att varje station tar 5--6 gånger mer plats i bandet än SSB.

Dessutom försvinner FM:s brusresistens-fördel vid svaga signaler
under ett visst tröskelvärde (``FM-tröskeleffekten'') -- under
tröskeln försämras FM dramatiskt, medan SSB degraderar mer
gradvist.

\textbf{Slutsatsen:} FM används på VHF/UHF (144\,MHz uppåt) där
spektrumet är rikligare och lokala avstånd dominerar. SSB används
på HF för långdistansförbindelser.
\end{worked}

% ═══════════════════════════════════════════════════════════════
\section{Sändaren}
% ═══════════════════════════════════════════════════════════════

En sändare tar ett ljud (eller data) och förvandlar det till
en radiovåg på rätt frekvens med rätt modulering och tillräcklig
effekt. Det låter enkelt -- men vägen dit innehåller flera
avgörande steg.

% ────────────────────────────────────────────────────────────
\subsection{Oscillatorn -- hjärtat som skapar frekvensen}

Allt börjar med oscillatorn: ett kretskort som producerar en
stabil sinusvåg på exakt rätt frekvens. Utan en stabil oscillator
hamnar du på fel frekvens, sprider ut din signal och stör andra.

\begin{storybox}[{\emoji{⌚} Klockan som aldrig stannar}]
En klocka behöver ett stabilt pendelsystem -- kvarts, pendel
eller atomröst -- för att hålla rätt tid. En sekund för snabbt
per dag och efter ett år är klockan sex minuter fel.

En radiooscillator är precis samma sak, fast för frekvens
istället för tid. Oscillatorn är klockans ``pendel'' -- den
bestämmer exakt hur snabbt signalen svänger. En kvarts-oscillator
(XO) kan hålla frekvensen stabil till en del per miljon (1 ppm)
eller bättre.
\end{storybox}

\begin{bluebox}[title={\emoji{🔑} Oscillatortyper}]
\textbf{LC-oscillator:} Använder en svängningskrets (spole + kondensator)
för att bestämma frekvensen. Enkel men relativt instabil --
temperatur och spänningsvariationer påverkar frekvensen.

\medskip
\textbf{Kristalloscillator (XO):} Använder en kvartskristalls
mekaniska resonans. Extremt stabil ($\pm$10--100\,ppm). Låst
till en fast frekvens.

\medskip
\textbf{VFO (Variable Frequency Oscillator):} LC-oscillator vars
frekvens kan varieras -- det som gör att du kan vrida på
frekvensratten. Stabilitetsproblem vid kall start och temperaturändringar.

\medskip
\textbf{PLL (Phase Locked Loop):} Se avsnitt 2.3.3 -- den moderna
lösningen som kombinerar kristallstabilitet med full frekvensvariabilitet.
\end{bluebox}

% ────────────────────────────────────────────────────────────
\subsection{Förstärkarsteg}

Oscillatorn producerar en svag signal -- kanske milliwatt. För
att nå ut med 100\,W eller mer behöver signalen förstärkas i
flera steg.

\begin{bluebox}[title={\emoji{🔑} Förstärkarklasser}]
Förstärkarsteg klassificeras efter hur stor del av sinusvågens
cykel transistorn är i ledning:

\begin{center}
\begin{tabular}{L{1.5cm}L{3.5cm}L{2.5cm}L{3.5cm}}
\rowcolor{greydark}
\tblhdr{Klass} & \tblhdr{Ledningsvinkeln} & \tblhdr{Verkningsgrad} & \tblhdr{Användning} \\
A & 360° (hela cykeln) & ca 25--50\% & Linjär förstärkning, bra linearitet \\
\rowcolor{greylight}
B & 180° (halv cykel) & ca 78\% & Push-pull-förstärkare \\
AB & 180--360° & ca 50--78\% & Vanligast i slutsteg \\
\rowcolor{greylight}
C & $<$180° & $>$78\% & RF-slutsteg, ej linjär \\
\end{tabular}
\end{center}

\smallskip
{\small\itshape Klass A ger bäst linearitet (viktigt för SSB)
men sämst verkningsgrad. Klass C ger bäst verkningsgrad men
kan bara användas med CW och FM -- inte AM eller SSB.}
\end{bluebox}

% ────────────────────────────────────────────────────────────
\subsection{PLL och frekvenssyntes}

Moderna transceivrar kan ställa in frekvens med 10\,Hz noggrannhet
och ändå hålla kristallstabilitet. Det är PLL:en som gör detta möjligt.

\begin{storybox}[{\emoji{🔒} Hunden i kopplet}]
Föreställ dig en hund (VFO) som springer fritt -- snabb och rörlig
men opålitlig. Nu kopplar du ett elastiskt koppel från hunden till
en stabil stolpe (kristallreferensen). Hunden kan fortfarande röra
sig, men dras alltid tillbaka mot rätt position om den börjar dra iväg.

PLL:en fungerar likadant. VFO:n (``hunden'') kan variera i frekvens,
men fasdetektorn (``kopplet'') jämför ständigt med kristallreferensen
och korrigerar om VFO:n börjar driva.
\end{storybox}

\begin{bluebox}[title={\emoji{🔑} PLL -- Phase Locked Loop}]
En PLL består av fyra block:

\begin{enumerate}
  \item \textbf{VCO} (Voltage Controlled Oscillator) -- oscillatorn
        vars frekvens styrs av en spänning
  \item \textbf{Frekvensdelare} -- delar ned VCO-frekvensen med
        ett programmerbart tal $N$
  \item \textbf{Fasdetektor} -- jämför den delade frekvensen med
        kristallreferensen
  \item \textbf{Lågpassfilter} -- slättar ut korrektionssignalen
        till VCO
\end{enumerate}

Systemet låser sig när VCO-frekvens $= N \times f_{\text{ref}}$.
Genom att ändra $N$ digitalt kan vi välja exakt vilken frekvens vi vill --
med kristallens stabilitet.
\end{bluebox}

% ── Figur: PLL-blockschema ───────────────────────────────────
\begin{figure}[htbp]
\centering
\begin{tikzpicture}[font=\small, >=stealth,
  block/.style={rectangle, draw=gray!60, fill=white, 
                rounded corners=3pt, minimum width=2.2cm,
                minimum height=0.9cm, align=center}]

  \fill[gray!6, rounded corners=4pt] (-0.4,-1.6) rectangle (13.0,2.8);

  % Block
  \node[block] (xtal) at (1.2,1.8) {Kristall\\referens\\$f_{\text{ref}}$};
  \node[block] (pd)   at (4.0,1.8) {Fas-\\detektor};
  \node[block] (lpf)  at (6.8,1.8) {LP-\\filter};
  \node[block] (vco)  at (9.6,1.8) {VCO};
  \node[block] (div)  at (6.8,0.0) {Delare\\$\div N$};

  % Pilar framåt
  \draw[->] (xtal) -- (pd);
  \draw[->] (pd)   -- (lpf);
  \draw[->] (lpf)  -- (vco);
  \draw[->] (vco)  -- ++(1.6,0) node[right, font=\small\bfseries] {Ut: $N \cdot f_{\text{ref}}$};

  % Återkoppling
  \draw[->] (vco) -- ++(0,-1.8) -| (div);
  \draw[->] (div) -- (pd |- div) -- (pd);

  % N-styrning
  \draw[->, dashed] (6.8,-1.3) node[below, font=\scriptsize] {Frekvensval (digitalt, $N$)} -- (div);

\end{tikzpicture}
\caption{PLL-blockschema. Fasdetektorn jämför kristallreferensen
med den delade VCO-signalen och genererar ett korrektionsspänning
som håller VCO:n låst på $f = N \times f_{\text{ref}}$. Genom att
ändra $N$ digitalt väljer vi önskad frekvens.}
\label{fig:pll}
\end{figure}

% ────────────────────────────────────────────────────────────
\subsection{Sändarskedjans blockschema}

Nu kan vi sätta ihop hela sändarkedjan -- från mikrofon till antenn.

% ── Figur: Sändarblockschema ─────────────────────────────────
\begin{figure}[htbp]
\centering
\begin{tikzpicture}[font=\small, >=stealth,
  block/.style={rectangle, draw=gray!60, fill=white,
                rounded corners=3pt, minimum width=2.0cm,
                minimum height=0.9cm, align=center},
  sblock/.style={rectangle, draw=blue!40, fill=blue!5,
                rounded corners=3pt, minimum width=2.0cm,
                minimum height=0.9cm, align=center}]

  \fill[gray!6, rounded corners=4pt] (-0.3,-0.5) rectangle (13.2,3.8);

  % Mikrofon
  \node[block] (mic)  at (0.8,1.8)  {Mikrofon};
  % Modulering
  \node[sblock] (mod) at (3.0,1.8)  {Modulering\\(AM/FM/SSB)};
  % Oscillator/PLL
  \node[block] (osc)  at (3.0,0.4)  {Oscillator\\PLL};
  % Drivsteg
  \node[sblock] (drv) at (5.8,1.8)  {Drivsteg};
  % Slutsteg
  \node[sblock] (pa)  at (8.4,1.8)  {Slutsteg\\(PA)};
  % LPF
  \node[block] (lpf)  at (10.8,1.8) {LP-filter\\$\pi$-filter};
  % Antenn
  \node[block] (ant)  at (12.8,1.8) {Ant.};

  % Pilar
  \draw[->] (mic) -- (mod);
  \draw[->] (osc) -- (mod);
  \draw[->] (mod) -- (drv);
  \draw[->] (drv) -- (pa);
  \draw[->] (pa)  -- (lpf);
  \draw[->] (lpf) -- (ant);

  % Effektmärkning
  \node[font=\scriptsize, gray] at (5.8,2.55)  {mW};
  \node[font=\scriptsize, gray] at (8.4,2.55)  {W};
  \node[font=\scriptsize, gray] at (10.8,2.55) {100\,W};

  % Förklaring
  \node[font=\scriptsize, blue!60] at (3.0,3.3) {Skapar signalen};
  \node[font=\scriptsize, blue!60] at (7.1,3.3) {Förstärker signalen};
  \node[font=\scriptsize, gray!70] at (11.6,3.3) {Rensar/sänder};

\end{tikzpicture}
\caption{Sändarskedjans blockschema. Oscillatorn genererar bärfrekvensen,
moduleringen adderar informationen, drivsteg och slutsteg förstärker
till sändningseffekt, och $\pi$-filtret dämpar övertoner innan
signalen når antennen.}
\label{fig:tx_blockschema}
\end{figure}

\subsubsection*{Övningsfrågor -- 2.3}

\begin{qblock}{4.}
\qprov\quad En SSB-sändare ska sända på 14,225\,MHz med 100\,W PEP.

a) Varför behöver slutsteget vara klass AB och inte klass C?\\
b) Varför sitter ett lågpassfilter ($\pi$-filter) efter slutsteget?\\
c) Vad händer om PLL:en tappar lock?
\end{qblock}

\begin{qblock}{F.}
\qfallan\quad ``En klass C-förstärkare har bättre verkningsgrad,
så den borde ge mer uteffekt för samma ineffekt. Därför är den
bäst för SSB-sändning.''

Vad är felet med detta resonemang?
\end{qblock}

\medskip
\noindent\rule{\textwidth}{0.4pt}
{\small\itshape Försök själv innan du läser svaret.}
\medskip

\begin{worked}[title={Lösning -- Fallgropsfråga 2.3\,F}]
Verkningsgradsresonemanget är korrekt -- klass C är mer effektiv.
Men slutsatsen är fel av ett fundamentalt skäl: \textbf{linearitet}.

SSB är ett linjärt sändningsslag -- amplituden hos den utsända
signalen måste vara proportionell mot ingångssignalen vid varje
tidpunkt. Röst är en komplex signal vars amplitud varierar
kontinuerligt.

Klass C-förstärkaren är inte linjär -- den klipper bort delar av
sinusvågen (leder bara under $<$180°). Den kan förstärka CW (bara
på/av) och FM (konstant amplitud) utan problem, eftersom dessa
sändningsslag inte kräver amplitudlinearitet.

Men för SSB och AM skulle klass C producera kraftig distortion --
signalen skulle bli oigenkännlig och dessutom sprida
intermodulationsprodukter över ett brett frekvensområde.

\textbf{Slutsats:} Bra verkningsgrad är viktigt, men linearitet
är ett absolut krav för SSB. Klass AB är kompromissen som ger
acceptabel linearitet med rimlig verkningsgrad.
\end{worked}

% ═══════════════════════════════════════════════════════════════
\section{Mottagaren}
% ═══════════════════════════════════════════════════════════════

Mottagarens uppgift är i princip spegelbilden av sändarens: ta
en radiovåg ur luften och omvandla den tillbaka till ljud eller data.
Men i praktiken är mottagartekniken mer utmanande -- signalen
som anländer kan vara miljarder gånger svagare än sändareffekten,
och mottagaren måste skilja ut just den station du vill höra
bland hundratals andra.

% ────────────────────────────────────────────────────────────
\subsection{Selektivitet -- att hitta en röst i ett kakofoni}

\begin{storybox}[{\emoji{🎭} Cocktailpartyeffekten}]
Du är på en fest med tjugo samtal som pågår samtidigt. Ändå kan
du koncentrera dig på precis det samtal du är intresserad av och
``filtrera bort'' resten. Hjärnan gör detta automatiskt baserat
på röst, riktning och sammanhang.

En radiomottagare behöver göra samma sak -- men med elektronik.
Selektiviteten är mottagarens förmåga att ``lyssna'' på bara en
frekvens och blockera allt annat.
\end{storybox}

\begin{bluebox}[title={\emoji{🔑} Selektivitet och bandbredd}]
\textbf{Selektivitet} är mottagarens förmåga att skilja ut en
signal bland angränsande signaler.

\medskip
Selektiviteten bestäms av bandpassfiltrets \textbf{bandbredd}:
\begin{itemize}
  \item CW: ca 200--500\,Hz
  \item SSB: ca 2,4\,kHz
  \item AM: ca 6--9\,kHz
  \item FM (smalbands): ca 10--15\,kHz
\end{itemize}

\textbf{Spegelfrekvens:} En viktig begränsning i superheterodyn-
mottagaren -- en signal på $f_{\text{LO}} + f_{\text{IF}}$ tas emot
precis lika väl som den önskade signalen på $f_{\text{LO}} - f_{\text{IF}}$.
Preselektion (filter före blandaren) är avgörande för att dämpa
spegelfrekvensen.
\end{bluebox}

% ────────────────────────────────────────────────────────────
\subsection{Superheterodynprincipen}

Nästan alla moderna radiomottagare -- från bilradion till avancerade
SDR-system -- bygger på superheterodynprincipen, uppfunnen av
Edwin Armstrong på 1910-talet.

Principens genialitet: istället för att försöka filtrera och
förstärka på den mottagna frekvensen (som kan vara vad som helst
från 1\,MHz till 1\,GHz), omvandlar vi alltid signalen till
\emph{en fast mellanfrekvens} (IF) där vi kan bygga optimala filter.

\begin{storybox}[{\emoji{🔭} Teleskopet som vrider sig}]
Ett astronomiskt teleskop kan riktas mot vilken stjärna som helst
på himlen. Men oavsett vart det riktas, fokuseras alltid ljuset
till exakt samma punkt i okularet. Variationen (vilken stjärna)
hanteras av hur teleskopet riktas -- analysen (hur ljuset delas
upp) sker alltid på samma plats.

Superheterodynmottagaren gör samma sak med frekvenser. LO-oscillatorn
``riktar'' mottagaren mot önskad frekvens, men signalen analyseras
(filtreras, förstärks) alltid på IF -- den fasta ``okuläret''.
\end{storybox}

\begin{bluebox}[title={\emoji{🔑} Superheterodynprincipen}]
Grundidén: blanda den mottagna signalen $f_{\text{RF}}$ med en
lokaloscillator $f_{\text{LO}}$ för att producera en
\textbf{mellanfrekvens} (IF):

\[ f_{\text{IF}} = f_{\text{LO}} - f_{\text{RF}} \quad
   \text{(eller } f_{\text{RF}} - f_{\text{LO}}\text{)} \]

\medskip
\textbf{Vanliga IF-frekvenser:}
\begin{itemize}
  \item Kortvågsmottagare: 455\,kHz, 9\,MHz, 10,7\,MHz
  \item FM-radio: 10,7\,MHz
  \item Moderna transceivrar: ofta dubbel-/trippelkonversion
\end{itemize}

\medskip
\textbf{Fördelar:}
\begin{itemize}
  \item Skarpa bandpassfilter byggs en gång (på IF)
  \item Förstärkning sker på stabil frekvens
  \item Lätt att byta frekvens (ändra bara LO)
\end{itemize}
\end{bluebox}

% ── Figur: Superheterodyn-blockschema ───────────────────────
\begin{figure}[htbp]
\centering
\begin{tikzpicture}[font=\small, >=stealth,
  block/.style={rectangle, draw=gray!60, fill=white,
                rounded corners=3pt, minimum width=2.0cm,
                minimum height=0.9cm, align=center},
  sblock/.style={rectangle, draw=green!50!black, fill=green!5,
                rounded corners=3pt, minimum width=2.0cm,
                minimum height=0.9cm, align=center}]

  \fill[gray!6, rounded corners=4pt] (-0.3,-1.5) rectangle (13.2,3.2);

  % Block
  \node[block]  (ant) at (0.8,1.8)  {Ant.};
  \node[sblock] (pre) at (2.6,1.8)  {Pre-\\selektor};
  \node[sblock] (mix) at (4.8,1.8)  {Blandare};
  \node[sblock] (if1) at (7.0,1.8)  {IF-filter\\IF-förstärk};
  \node[sblock] (det) at (9.2,1.8)  {Detektor\\(demod)};
  \node[block]  (af)  at (11.2,1.8) {NF-förstärk};
  \node[block]  (spk) at (12.9,1.8) {Högtal.};

  % LO
  \node[block] (lo) at (4.8,0.2) {LO / PLL\\$f_{\text{LO}}$};

  % Pilar
  \draw[->] (ant) -- (pre);
  \draw[->] (pre) -- (mix);
  \draw[->] (mix) -- (if1);
  \draw[->] (if1) -- (det);
  \draw[->] (det) -- (af);
  \draw[->] (af)  -- (spk);
  \draw[->] (lo)  -- (mix);

  % Frekvensmarkering
  \node[font=\scriptsize, gray] at (1.7,2.55) {$f_{\text{RF}}$};
  \node[font=\scriptsize, gray] at (5.9,2.55) {$f_{\text{IF}}$};
  \node[font=\scriptsize, gray] at (10.2,2.55) {NF};

\end{tikzpicture}
\caption{Superheterodynmottagarens blockschema. Preselektor-filtret
dämpar spegelfrekvenser och starka störare. Blandaren omvandlar
$f_{\text{RF}}$ till $f_{\text{IF}}$. IF-steget ger skarp selektion
och förstärkning. Detektorn demodulerar och återger NF-signalen.}
\label{fig:rx_blockschema}
\end{figure}

% ────────────────────────────────────────────────────────────
\subsection{Känslighet och brus}

En mottagares känslighet avgör hur svaga signaler den kan ta emot.
Men det finns en absolut gräns: det termiska bruset.

\begin{bluebox}[title={\emoji{🔑} Mottagarkänslighet och brus}]
\textbf{Termiskt brus} uppstår i varje elektrisk komponent på
grund av elektronernas slumpmässiga rörelse. Det är proportionellt
mot temperaturen och bandbredden:

\[ P_{\text{brus}} = k \cdot T \cdot B \]

där $k = 1{,}38 \times 10^{-23}$\,J/K (Boltzmanns konstant),
$T$ = temperatur (Kelvin), $B$ = bandbredd (Hz).

\medskip
\textbf{Brusfaktor (NF)} anger hur mycket mottagaren själv
bidrar till bruset utöver det termiska bruset:
\begin{itemize}
  \item Typisk kortvågsmottagare: NF ca 10--20\,dB
  \item LNA (lågbrusförstärkare) för VHF/UHF: NF ca 0,5--2\,dB
\end{itemize}

\medskip
\textbf{MDS (Minimum Discernible Signal):} Den svagaste signal
mottagaren kan detektera -- typiskt $-$130 till $-$140\,dBm
för en bra HF-mottagare.
\end{bluebox}

% ────────────────────────────────────────────────────────────
\subsection{Mottagarskedjans blockschema}

Vi såg blockschemat i 2.4.2 ovan. Låt oss titta på varje block
mer detaljerat:

\begin{bluebox}[title={\emoji{🔑} Mottagarblocken i detalj}]
\textbf{1. Preselektor / RF-filter:}
Dämpar signaler utanför det önskade bandet -- särskilt viktig
för att undertrycka spegelfrekvenser.

\medskip
\textbf{2. RF-förstärkare (LNA):}
Förstärker den svaga antennsignalen med minimal brusbidrag.
Placeras före blandaren för bäst brusfaktor för hela kedjan.

\medskip
\textbf{3. Blandare:}
Multiplicerar RF-signalen med LO-signalen. Producerar
summa- och differensfrekvenser -- vi väljer ut differensen ($f_{\text{IF}}$).

\medskip
\textbf{4. IF-filter:}
Det viktigaste selektivitets-elementet. Ofta ett keramikfilter
eller kristallfilter med hög Q-värde. Bestämmer mottagarens
bandbredd och grannkanalsdämpning.

\medskip
\textbf{5. Detektor / Demodulator:}
Återvinner informationen ur IF-signalen. Olika för olika
moduleringsslag:
\begin{itemize}
  \item AM: Diod-detektor (envelop-detektor)
  \item FM: Diskriminator eller kvadraturdetektor
  \item SSB/CW: Produktdetektor med BFO
\end{itemize}

\medskip
\textbf{6. AGC (Automatic Gain Control):}
Justerar automatiskt förstärkningen i RF- och IF-stegen för
att hålla konstant ljudnivå oavsett signalstyrka.
\end{bluebox}

\subsubsection*{Övningsfrågor -- 2.4}

\begin{qblock}{5.}
\qprov\quad En superheterodynmottagare har IF = 455\,kHz och
är inställd på 7,050\,MHz.

a) Vilken frekvens sätter LO:n?\\
b) Vad är spegelfrekvensen?\\
c) Varför är spegelfrekvensen ett problem och hur bekämpas den?
\end{qblock}

\begin{qblock}{6.}
\qprov\quad Förklara varför en SSB-mottagare behöver en BFO
(beat frequency oscillator) men en AM-mottagare inte behöver det.
\end{qblock}

\begin{qblock}{F.}
\qfallan\quad ``En mottagare med hög förstärkning är alltid
bättre -- mer förstärkning ger bättre möjlighet att höra svaga
signaler.''

Varför är detta resonemang ofullständigt?
\end{qblock}

\medskip
\noindent\rule{\textwidth}{0.4pt}
{\small\itshape Försök själv innan du läser svaret.}
\medskip

\begin{worked}[title={Lösning -- Fallgropsfråga 2.4\,F}]
Förstärkning hjälper -- men bara om bruset inte förstärks lika
mycket. Det avgörande är \textbf{signal-till-brus-förhållandet (S/N)},
inte signalnivån i sig.

Om mottagarens första förstärkarsteg har hög brusfaktor, förstärks
bruset lika mycket som signalen. Mer förstärkning ger inte bättre
hörbarhet -- bara louder noise.

Det som verkligen avgör mottagarens prestanda för svaga signaler är:
\begin{enumerate}
  \item Brusfaktorn hos \emph{första} förstärkarsteget (dominerar hela kedjans S/N)
  \item Bandbredden (smalare = mindre brus = bättre S/N)
  \item Antennens effektivitet
\end{enumerate}

Överdrivet hög förstärkning skapar dessutom ett annat problem:
starka stationer kan överstira mottagaren och orsaka
intermodulationsdistortion (IMD) -- falska signaler som verkar
verkliga.

\textbf{Slutsats:} En bra mottagare balanserar förstärkning,
brusfaktor och dynamiskt omfång -- inte maximerar förstärkningen.
\end{worked}

% ═══════════════════════════════════════════════════════════════
\section{Sändtagaren -- transceivern}
% ═══════════════════════════════════════════════════════════════

I praktiken kombineras sändare och mottagare till en enda enhet --
\textbf{transceivern} (transmitter + receiver). Det är vad du
ser framför dig om du tittar på en modern amatörradio.

\subsection{Hur sändare och mottagare delar komponenter}

Den stora fördelen med transceivern är att sändaren och mottagaren
kan dela flera dyra och komplexa block:

\begin{bluebox}[title={\emoji{🔑} Delade komponenter i en transceiver}]
\begin{itemize}
  \item \textbf{PLL/VFO:} Samma frekvensreferens styr både sändarens
        bärfrekvens och mottagarens LO. Du vrider bara en ratt.
  \item \textbf{Bandpassfilter:} Samma filter kan användas för
        preselektion i mottagarläge och undertryckning av övertoner
        i sändarläge.
  \item \textbf{Kristallfilter (IF-filter):} Används i mottagarkedjan
        och ibland även för att forma SSB-signalen i sändarkedjan.
  \item \textbf{Display och styrelektronik:}
        Gemensam mikroprocessor styr hela transceivern.
\end{itemize}

En modern transceiver är alltså mycket mer än en sammankopplad
sändare och mottagare -- den är ett integrerat system där
komponentdelning ger bättre prestanda till lägre kostnad.
\end{bluebox}

\subsection{Duplexer och VOX}

\begin{bluebox}[title={\emoji{🔑} TX/RX-växling}]
\textbf{T/R-relä (Transmit/Receive):} Kopplar antennen till antingen
sändaren eller mottagaren. Vid sändning skyddas mottagarens känsliga
ingångssteg från sändarens höga effekt.

\medskip
\textbf{VOX (Voice Operated Transmit):} Mikrofonsignalens nivå
styr automatiskt T/R-reläet -- du behöver inte trycka på PTT-knappen
(Push To Talk). Praktiskt men kräver noggrann justering för att
undvika oavsiktlig sändning.

\medskip
\textbf{Full duplex:} Möjlighet att sända och ta emot samtidigt
på \emph{olika} frekvenser -- används vid satellitförbindelser
och repeaters. Kräver antingen två antenner eller ett duplexfilter
(``duplexer'') som ger tillräcklig dämpning mellan TX och RX.
\end{bluebox}

\subsubsection*{Övningsfrågor -- 2.5}

\begin{qblock}{7.}
\qprov\quad Förklara varför ett T/R-relä behövs i en transceiver
som arbetar simplext (sänder och tar emot på samma frekvens
men inte samtidigt).

Vad händer om T/R-reläet är trasigt och fastnar i TX-läge
när du försöker ta emot?
\end{qblock}

% ═══════════════════════════════════════════════════════════════
% KAPITEL 2 AVSLUTNING
% ═══════════════════════════════════════════════════════════════
\section*{Signalen lever}
\addcontentsline{toc}{section}{Signalen lever}

Kanske kändes det tungt på vägen -- det ska det göra. Det betyder
att du faktiskt tänkte. Du har nu följt en signal hela vägen från
idé till radiovåg och tillbaka igen.

Tänk på det så här: när du en dag trycker på PTT och hör en röst
svara från andra sidan Atlanten, vet du exakt vad som händer inne
i lådan framför dig. Oscillatorn håller frekvensen, PLL:en ser
till att den inte driver, modulatorn packar din röst i bärvågen,
slutsteget pumpar upp effekten och $\pi$-filtret ser till att du
inte stör dina grannar på bandet. I mottagaren på andra sidan
gör superheterodyn-kedjan jobbet baklänges.

Det är inte magi -- det är teknik du nu förstår.

\medskip
\noindent
Och det roliga är att det svåraste nu ligger bakom dig -- det som
kommer härnäst är där signalen möter världen.

\begin{bluebox}[title={}]
\textbf{Signalen lever. Nu ska vi ge den vingar.}

\smallskip
I nästa kapitel -- Antenner -- tar vi steget ut ur lådan och upp
i luften. Du kommer att märka att allt du lärt dig om resonans,
impedans och SWR plötsligt får ett ansikte. Kapitel 1 och 2
förräntar sig.
\end{bluebox}

% ═══════════════════════════════════════════════════════════
%  FACIT – KAPITEL 2
% ═══════════════════════════════════════════════════════════
\chapter*{Facit – Kapitel 2}
\addcontentsline{toc}{chapter}{Facit – Kapitel 2}
\markboth{Facit – Kapitel 2}{}

% ── Facit 2.1 ─────────────────────────────────────────────
\subsection*{2.1 – Vad är en radiovåg?}

\begin{facitblock}{1}
\qprov\quad En radioamatör sänder på 7,05\,MHz.
a) Beräkna våglängden.
b) Vilket band tillhör frekvensen?
c) Hur lång är en halvvågsdipol för detta band (ungefär)?
\tcblower
a) $\lambda = \dfrac{300}{7{,}05} \approx \textbf{42{,}6\,m}$

b) Frekvensen ligger inom \textbf{40\,m-bandet} (7,000--7,200\,MHz).

c) En halvvågsdipol är $\lambda/2 \approx \textbf{21{,}3\,m}$ lång.
I praktiken förkortas den något (ca 5\,\%) på grund av
ändeffekter -- typiskt ca 20\,m.
\end{facitblock}

% ── Facit 2.2 ─────────────────────────────────────────────
\subsection*{2.2 – Modulering}

\begin{facitblock}{2}
\qprov\quad Förklara skillnaden mellan AM och FM med egna ord.
a) Vad varierar vid AM? Vad är konstant?
b) Vad varierar vid FM? Vad är konstant?
c) Varför är FM mer brusresistent än AM?
\tcblower
a) Vid \textbf{AM} varierar \textbf{amplituden} i takt med
modulationssignalen. Frekvensen är konstant.

b) Vid \textbf{FM} varierar \textbf{frekvensen} (deviationen)
i takt med modulationssignalen. Amplituden är konstant.

c) Brus påverkar främst amplituden hos en signal. Eftersom FM
bär informationen i frekvensen -- inte amplituden -- påverkas
FM-signalen mycket mindre av amplitudbrus. Mottagaren kan dessutom
begränsa amplituden utan att förlora information, vilket
undertrycker bruset ytterligare (FM-begränsareffekten).
\end{facitblock}

\begin{facitblock}{3}
\qprov\quad En SSB-sändare sänder med 100\,W PEP på 14,225\,MHz.
a) Vilket sidband används -- USB eller LSB? Varför?
b) Ungefär hur bred är signalen i Hz?
c) Hur mycket effekt ``slösar'' SSB jämfört med AM?
\tcblower
a) \textbf{USB} (Upper Sideband) används. Konventionen inom
amatörradio är USB över 10\,MHz -- 20\,m-bandet ligger på
14\,MHz.

b) En SSB-talsignal är ca \textbf{2,4\,kHz} bred (300\,Hz--2\,700\,Hz).

c) SSB slösar \textbf{ingen} effekt på bärvåg eller spegelsidband.
Vid AM går ca 2/3 av sändareffekten till bärvågen och det andra
sidobandet -- de bär ingen information. SSB skickar all effekt
i det nyttiga sidobandet.
\end{facitblock}

% ── Facit 2.3 ─────────────────────────────────────────────
\subsection*{2.3 – Sändaren}

\begin{facitblock}{4}
\qprov\quad En SSB-sändare ska sända på 14,225\,MHz med 100\,W PEP.
a) Varför behöver slutsteget vara klass AB och inte klass C?
b) Varför sitter ett lågpassfilter ($\pi$-filter) efter slutsteget?
c) Vad händer om PLL:en tappar lock?
\tcblower
a) SSB kräver \textbf{linjär förstärkning} -- amplituden hos
utsignalen måste vara proportionell mot insignalen vid varje
tidpunkt. Klass C är inte linjär och skulle ge kraftig distortion.
Klass AB ger tillräcklig linearitet med rimlig verkningsgrad.

b) Varje förstärkarsteg skapar \textbf{harmoniska övertoner}
(2$f$, 3$f$, \ldots) som kan störa andra frekvensband.
$\pi$-filtret dämpar dessa kraftigt innan signalen når antennen.
Det är ett lagkrav att hålla harmoniska under en viss nivå.

c) Om PLL:en tappar lock börjar VCO:n \textbf{driva fritt} --
sändaren hamnar på fel frekvens och stör potentiellt andra
stationer. Moderna transceivrar larmar och stänger av sändningen
automatiskt.
\end{facitblock}

% ── Facit 2.4 ─────────────────────────────────────────────
\subsection*{2.4 – Mottagaren}

\begin{facitblock}{5}
\qprov\quad En superheterodynmottagare har IF = 455\,kHz och
är inställd på 7,050\,MHz.
a) Vilken frekvens sätter LO:n?
b) Vad är spegelfrekvensen?
c) Varför är spegelfrekvensen ett problem och hur bekämpas den?
\tcblower
a) $f_{\text{LO}} = f_{\text{RF}} + f_{\text{IF}}
= 7{,}050 + 0{,}455 = \textbf{7{,}505\,MHz}$

b) Spegelfrekvensen ligger $2 \times f_{\text{IF}}$ bort:
$f_{\text{spegel}} = f_{\text{LO}} + f_{\text{IF}}
= 7{,}505 + 0{,}455 = \textbf{7{,}960\,MHz}$

c) En signal på spegelfrekvensen blandas ned till samma IF som
den önskade -- mottagaren kan inte skilja dem åt. Det bekämpas
med ett \textbf{preselektor-filter} före blandaren som dämpar
signaler utanför det önskade bandet.
\end{facitblock}

\begin{facitblock}{6}
\qprov\quad Förklara varför en SSB-mottagare behöver en BFO
men en AM-mottagare inte behöver det.
\tcblower
En AM-signal innehåller en \textbf{bärvåg} som detektorn använder
som referens -- en enkel diod-detektor räcker.

En SSB-signal saknar bärvåg (den är bortfiltrerad i sändaren).
Utan bärvåg har detektorn ingen referens och talet låter
obegripligt. \textbf{BFO:n} injicerar en lokal bärvåg i
mottagaren så att produktdetektorn kan rekonstruera det
ursprungliga talet.
\end{facitblock}

% ── Facit 2.5 ─────────────────────────────────────────────
\subsection*{2.5 – Sändtagaren}

\begin{facitblock}{7}
\qprov\quad Förklara varför ett T/R-relä behövs i en transceiver
som arbetar simplext. Vad händer om det fastnar i TX-läge
vid mottagning?
\tcblower
Under sändning producerar slutsteget upp till 100\,W eller mer.
Om denna effekt nådde mottagarens ingångssteg direkt skulle det
\textbf{omedelbart förstöra} de känsliga {transistorerna} -- de är
dimensionerade för nano- till mikrowatt-nivå.

T/R-reläet kopplar bort mottagaren från antennen under sändning
och återkopplar den när sändningen upphör.

Om reläet fastnar i TX-läge vid mottagning: antennen är kopplad
till sändarens slutsteg istället för mottagaren --
\textbf{du hör ingenting} trots att signaler finns på bandet.
\end{facitblock}


\chapter{Antenner}
% Kapitel 3

\chapter{Vågutbredning}
% Kapitel 4

\chapter{Mätinstrument}
% Kapitel 5

\chapter{Störningar}
% Kapitel 6

\chapter{Regler och bestämmelser}
% Kapitel 7

\chapter{Trafikmetoder}
% Kapitel 8

\chapter{Säkerhet}
% Kapitel 9

\chapter{Praktisk trafik}
% Kapitel 10

\end{document}